% Generated human-review packet template.  Source records, Lean statements,
% and raw-source-to-Spec screening fields are substituted by
% scripts/review_dashboard_packet.py.
% Do not hand-edit generated packets; annotate the PDF or reconcile notes into
% the paper-local human-review log.
\documentclass[10pt]{article}
\usepackage[margin=0.43in,footskip=0.2in]{geometry}
\usepackage{fontspec}
% Use a Unicode-complete main face as well as a Unicode-complete code face.
% Reviewer reasons and source labels can contain exact Greek symbols outside
% verbatim blocks; silently dropping one would corrupt the review packet.
\setmainfont{DejaVu Serif}
% The broad DejaVu Sans Unicode coverage preserves Lean's mathematical glyphs
% (including U+2216 set difference and superscript identifiers) in verbatim
% review blocks.  Its proportional metrics are preferable to silently missing
% exact semantic text from a narrower monospace font.
\setmonofont{DejaVu Sans}
\usepackage{hyperref}
\usepackage{amssymb}
\usepackage{fvextra}
\usepackage{enumitem}
\setlength{\parindent}{0pt}
\setlength{\parskip}{0.15em}
\linespread{0.92}
\hypersetup{colorlinks=true,linkcolor=blue,urlcolor=blue,breaklinks=true}
\DefineVerbatimEnvironment{ReviewVerbatim}{Verbatim}{breaklines=true,breakanywhere=true,fontsize=\scriptsize}
\newcommand{\reviewerbox}{%
  \par\noindent\fbox{\parbox{0.96\linewidth}{\rule{0pt}{0.38in}}}\par
}
\newcommand{\reviewmatch}[1]{%
  \par\noindent\textbf{Reviewer check:} $\square$\; #1\par
  \noindent\emph{If left unchecked, explain the discrepancy or uncertainty below.}\par
}

\begin{document}
\fontsize{9}{10}\selectfont

\begin{center}
{\LARGE Human review packet: LBG22StrategicRanking}\\[0.35em]
\small Generated from byte-pinned source excerpts, the expanded PaperInterface
specifications, and hash-bound raw-source-to-Spec screening records.
\end{center}

\noindent\textbf{Paper:} LBG22StrategicRanking\\
\textbf{Paper id:} \texttt{LBG22StrategicRanking}\\
\textbf{Source version:} \parbox[t]{0.76\linewidth}{\raggedright AISTATS 2022 / PMLR 151, pp. 2489–2518; retained active TeX members of arXiv 2109.08240}\\
\textbf{Source record:} \texttt{audit/paper\_statement\_map.json}\\
\textbf{Rows:} 18\\
\textbf{Generated:} 2026-09-09





\section*{How to use this packet}
For each source claim, compare the exact source excerpts with the expanded
PaperInterface specification. Mark up this PDF or fill the blank reviewer box in the TeX source;
return the annotated file for reconciliation into the paper-local human-review
log.

\section*{Open the interactive dashboard}
The interactive dashboard is an optional alternative to using this PDF; it
presents the same review surface rather than an additional review requirement.
After forking and cloning (or pulling) AppliedModelingLib, run the following from the
repository root. The cache command is needed only the first time in a new clone.
\begin{ReviewVerbatim}
cd <your-repository-checkout>
lake exe cache get
python3 scripts/review_dashboard.py --paper LBG22StrategicRanking --serve
\end{ReviewVerbatim}
Open \url{http://127.0.0.1:8765/} in a browser and keep that terminal running
while reviewing. The dashboard records saved annotations in this paper's local
reviewer-owned review record.

\section*{Contents}
\begin{itemize}[leftmargin=1.2em]
\item \hyperlink{governing-declaration-section-5ef5ef0364b6939c}{Governing Lean declarations (26; supporting context)}
\item \hyperlink{paper-prerequisite-section-5ef5ef0364b6939c}{Paper-specific semantic prerequisites (38; not paper claims)}
\begin{itemize}[leftmargin=1.2em]
\item \hyperlink{paper-prerequisite-73413973c7594479}{LBG22\allowbreak{}Strategic\allowbreak{}Ranking.\allowbreak{}source\allowbreak{}Environment\allowbreak{}Skill}
\item \hyperlink{paper-prerequisite-5b3be494cde30a37}{LBG22\allowbreak{}Strategic\allowbreak{}Ranking.\allowbreak{}source\allowbreak{}Two\allowbreak{}Group\allowbreak{}Measure}
\item \hyperlink{paper-prerequisite-b8e06f39e94d0049}{LBG22\allowbreak{}Strategic\allowbreak{}Ranking.\allowbreak{}Environment\allowbreak{}Equilibrium}
\item \hyperlink{paper-prerequisite-d056dd7c8060653c}{LBG22\allowbreak{}Strategic\allowbreak{}Ranking.\allowbreak{}Finite\allowbreak{}Admission\allowbreak{}Policy}
\item \hyperlink{paper-prerequisite-6ed45e46d01eadeb}{LBG22\allowbreak{}Strategic\allowbreak{}Ranking.\allowbreak{}counterfactual\allowbreak{}Tie\allowbreak{}Broken\allowbreak{}Rank}
\item \hyperlink{paper-prerequisite-b7c71e53177b76c6}{LBG22\allowbreak{}Strategic\allowbreak{}Ranking.\allowbreak{}tie\allowbreak{}Broken\allowbreak{}Rank}
\item \hyperlink{paper-prerequisite-59d181ae71b7b024}{LBG22\allowbreak{}Strategic\allowbreak{}Ranking.\allowbreak{}Source\allowbreak{}Multidimensional\allowbreak{}Best\allowbreak{}Response\allowbreak{}At}
\item \hyperlink{paper-prerequisite-a42f440ddc254fa7}{LBG22\allowbreak{}Strategic\allowbreak{}Ranking.\allowbreak{}finite\allowbreak{}Lower\allowbreak{}Rank\allowbreak{}Band}
\item \hyperlink{paper-prerequisite-e57e55071c75ce91}{LBG22\allowbreak{}Strategic\allowbreak{}Ranking.\allowbreak{}source\allowbreak{}Multidimensional\allowbreak{}Coefficient}
\item \hyperlink{paper-prerequisite-16d135799a0f0491}{LBG22\allowbreak{}Strategic\allowbreak{}Ranking.\allowbreak{}source\allowbreak{}Multidimensional\allowbreak{}Population}
\item \hyperlink{paper-prerequisite-b201854e4e4dd891}{LBG22\allowbreak{}Strategic\allowbreak{}Ranking.\allowbreak{}Multidimensional\allowbreak{}Equilibrium}
\item \hyperlink{paper-prerequisite-8e74a377842bef95}{LBG22\allowbreak{}Strategic\allowbreak{}Ranking.\allowbreak{}Multidimensional\allowbreak{}Primitives}
\item \hyperlink{paper-prerequisite-b504d72e1e0de867}{LBG22\allowbreak{}Strategic\allowbreak{}Ranking.\allowbreak{}Multitask\allowbreak{}Equilibrium}
\item \hyperlink{paper-prerequisite-3c93745e73ce0a63}{LBG22\allowbreak{}Strategic\allowbreak{}Ranking.\allowbreak{}Multitask\allowbreak{}Primitives}
\item \hyperlink{paper-prerequisite-519e321cd0957049}{LBG22\allowbreak{}Strategic\allowbreak{}Ranking.\allowbreak{}Source\allowbreak{}Finite\allowbreak{}Best\allowbreak{}Response\allowbreak{}At}
\item \hyperlink{paper-prerequisite-2a1b692f4e3136c6}{LBG22\allowbreak{}Strategic\allowbreak{}Ranking.\allowbreak{}Ranking\allowbreak{}Equilibrium}
\item \hyperlink{paper-prerequisite-9bbcc12f4a1736bf}{LBG22\allowbreak{}Strategic\allowbreak{}Ranking.\allowbreak{}Ranking\allowbreak{}Primitives}
\item \hyperlink{paper-prerequisite-2dc31b75c00d0d1b}{LBG22\allowbreak{}Strategic\allowbreak{}Ranking.\allowbreak{}Score\allowbreak{}Priority\allowbreak{}Best\allowbreak{}Response\allowbreak{}At}
\item \hyperlink{paper-prerequisite-f396601c268691e5}{LBG22\allowbreak{}Strategic\allowbreak{}Ranking.\allowbreak{}Ranking\allowbreak{}Score\allowbreak{}Priority\allowbreak{}Equilibrium}
\item \hyperlink{paper-prerequisite-38434fa78df931b7}{LBG22\allowbreak{}Strategic\allowbreak{}Ranking.\allowbreak{}Score\allowbreak{}Priority\allowbreak{}Above\allowbreak{}Cutoff}
\item \hyperlink{paper-prerequisite-f88aac58507e44ca}{LBG22\allowbreak{}Strategic\allowbreak{}Ranking.\allowbreak{}source\allowbreak{}Two\allowbreak{}Level\allowbreak{}Cutoff}
\item \hyperlink{paper-prerequisite-30d97cd43eb79017}{LBG22\allowbreak{}Strategic\allowbreak{}Ranking.\allowbreak{}source\allowbreak{}Two\allowbreak{}Level\allowbreak{}Reward}
\item \hyperlink{paper-prerequisite-c64406972a4c963f}{LBG22\allowbreak{}Strategic\allowbreak{}Ranking.\allowbreak{}Score\allowbreak{}Priority\allowbreak{}Hard\allowbreak{}Budget\allowbreak{}Best\allowbreak{}Response\allowbreak{}At}
\item \hyperlink{paper-prerequisite-fe9ea55a34d7ea9c}{LBG22\allowbreak{}Strategic\allowbreak{}Ranking.\allowbreak{}Score\allowbreak{}Priority\allowbreak{}Multidimensional\allowbreak{}Best\allowbreak{}Response\allowbreak{}At}
\item \hyperlink{paper-prerequisite-b0c4f18b812a40f1}{LBG22\allowbreak{}Strategic\allowbreak{}Ranking.\allowbreak{}Source\allowbreak{}Vector\allowbreak{}Best\allowbreak{}Response\allowbreak{}At}
\item \hyperlink{paper-prerequisite-6f1ee9b200124096}{LBG22\allowbreak{}Strategic\allowbreak{}Ranking.\allowbreak{}source\allowbreak{}Actual\allowbreak{}Disadvantaged\allowbreak{}Access}
\item \hyperlink{paper-prerequisite-22c0ae11547a174a}{LBG22\allowbreak{}Strategic\allowbreak{}Ranking.\allowbreak{}source\allowbreak{}Actual\allowbreak{}Group\allowbreak{}Welfare}
\item \hyperlink{paper-prerequisite-d3e90f3908818ca5}{LBG22\allowbreak{}Strategic\allowbreak{}Ranking.\allowbreak{}source\allowbreak{}Actual\allowbreak{}Group\allowbreak{}Welfare\allowbreak{}Gap}
\item \hyperlink{paper-prerequisite-c194f1bd51c4a4cc}{LBG22\allowbreak{}Strategic\allowbreak{}Ranking.\allowbreak{}source\allowbreak{}Actual\allowbreak{}Multitask\allowbreak{}School\allowbreak{}Utility}
\item \hyperlink{paper-prerequisite-9d0a441d993fe3f2}{LBG22\allowbreak{}Strategic\allowbreak{}Ranking.\allowbreak{}source\allowbreak{}Actual\allowbreak{}Two\allowbreak{}Level\allowbreak{}School\allowbreak{}Utility}
\item \hyperlink{paper-prerequisite-e76c8574ddaca0dc}{LBG22\allowbreak{}Strategic\allowbreak{}Ranking.\allowbreak{}source\allowbreak{}Deterministic\allowbreak{}Cutoff}
\item \hyperlink{paper-prerequisite-1aae1cdeecf21463}{LBG22\allowbreak{}Strategic\allowbreak{}Ranking.\allowbreak{}source\allowbreak{}Environment\allowbreak{}CDF}
\item \hyperlink{paper-prerequisite-b308d7da226a5fe1}{LBG22\allowbreak{}Strategic\allowbreak{}Ranking.\allowbreak{}source\allowbreak{}Environment\allowbreak{}Rank}
\item \hyperlink{paper-prerequisite-7fb5250e464d0736}{LBG22\allowbreak{}Strategic\allowbreak{}Ranking.\allowbreak{}source\allowbreak{}Finite\allowbreak{}Applicant\allowbreak{}Welfare}
\item \hyperlink{paper-prerequisite-25a6dc17f260546b}{LBG22\allowbreak{}Strategic\allowbreak{}Ranking.\allowbreak{}source\allowbreak{}Finite\allowbreak{}School\allowbreak{}Utility}
\item \hyperlink{paper-prerequisite-b2f8283910c81320}{LBG22\allowbreak{}Strategic\allowbreak{}Ranking.\allowbreak{}source\allowbreak{}Group\allowbreak{}Threshold}
\item \hyperlink{paper-prerequisite-366855b73d9ea1c5}{LBG22\allowbreak{}Strategic\allowbreak{}Ranking.\allowbreak{}source\allowbreak{}Pure\allowbreak{}Randomization\allowbreak{}Reward}
\item \hyperlink{paper-prerequisite-7224ad4c7b0b8551}{LBG22\allowbreak{}Strategic\allowbreak{}Ranking.\allowbreak{}source\allowbreak{}Societal\allowbreak{}Utility}
\end{itemize}
\item \hyperlink{source-section-f10b5f68e4acf3dc}{Source claims (18)}
\begin{itemize}[leftmargin=1.2em]
\item \hyperlink{source-claim-a38dac6d1a17c16d}{Rank\allowbreak{}Preservation\allowbreak{}Spec}
\item \hyperlink{source-claim-6f58e02df112a54c}{Second\allowbreak{}Price\allowbreak{}Effort\allowbreak{}Spec}
\item \hyperlink{source-claim-6f73613040867488}{Applicant\allowbreak{}Welfare\allowbreak{}Spec}
\item \hyperlink{source-claim-1ba4d8c13750bf79}{Private\allowbreak{}Utility\allowbreak{}Spec}
\item \hyperlink{source-claim-5cd4797b5f4721de}{Three\allowbreak{}Level\allowbreak{}Improvement\allowbreak{}Spec}
\item \hyperlink{source-claim-f1c09a02beccebed}{Societal\allowbreak{}Utility\allowbreak{}Spec}
\item \hyperlink{source-claim-d14371f504c1d7ee}{Environment\allowbreak{}Rank\allowbreak{}Preservation\allowbreak{}Spec}
\item \hyperlink{source-claim-77aa097e6083e865}{Environment\allowbreak{}Welfare\allowbreak{}Spec}
\item \hyperlink{source-claim-7af287b5fa9fa319}{Welfare\allowbreak{}Gap\allowbreak{}Derivative\allowbreak{}Spec}
\item \hyperlink{source-claim-dcc0514ddd510442}{Access\allowbreak{}Spec}
\item \hyperlink{source-claim-6cc5784dc863c5dc}{Effort\allowbreak{}Comparative\allowbreak{}Statics\allowbreak{}Spec}
\item \hyperlink{source-claim-ed834a4cf517d7fc}{Multidimensional\allowbreak{}Rank\allowbreak{}Spec}
\item \hyperlink{source-claim-ef5c828d0d65c93b}{Weighted\allowbreak{}Private\allowbreak{}Utility\allowbreak{}Spec}
\item \hyperlink{source-claim-53bb19f6829b5622}{Equal\allowbreak{}Score\allowbreak{}Rewards\allowbreak{}Spec}
\item \hyperlink{source-claim-1693f763578ff49e}{Null\allowbreak{}Deviations\allowbreak{}Spec}
\item \hyperlink{source-claim-12dd3b600e58e716}{Cost\allowbreak{}Gap\allowbreak{}Spec}
\item \hyperlink{source-claim-e24c6aaefba531d0}{Production\allowbreak{}Gap\allowbreak{}Spec}
\item \hyperlink{source-claim-4fa39a045f045b81}{Pre\allowbreak{}Effort\allowbreak{}Skill\allowbreak{}Regularity\allowbreak{}Spec}
\end{itemize}
\end{itemize}

\clearpage
\hypertarget{governing-declaration-section-5ef5ef0364b6939c}{}
\section*{Governing Lean declarations}
These exact graph-bound declaration bodies are retained by the displayed specifications or reviewed prerequisites. They are supporting context and do not add source-result rows or semantic judgments.
\hypertarget{governing-declaration-c5b2c07abbff7224}{}
\subsection*{\small\ttfamily\raggedright LBG22\allowbreak{}Strategic\allowbreak{}Ranking.\allowbreak{}effort\allowbreak{}Interval\allowbreak{}Inverse}
\textbf{Declaration location:} paper-local
\paragraph{Lean declaration body}
\begin{ReviewVerbatim}
type:
(ℝ → ℝ) → ℝ → ℝ → ℝ

value:
fun (f : ℝ → ℝ) (effortMax a : ℝ) => Function.invFunOn f (Set.Icc 0 effortMax) a
\end{ReviewVerbatim}


\hypertarget{governing-declaration-9e8db797d20dba91}{}
\subsection*{\small\ttfamily\raggedright LBG22\allowbreak{}Strategic\allowbreak{}Ranking.\allowbreak{}score\allowbreak{}Priority\allowbreak{}Rank}
\textbf{Declaration location:} paper-local
\paragraph{Lean declaration body}
\begin{ReviewVerbatim}
type:
{α : Type u_1} → [inst : MeasurableSpace α] → MeasureTheory.Measure α → (α → ℝ) → (α → ℝ) → ℝ → ℝ → ℝ

value:
fun {α : Type u_1} [MeasurableSpace α] (μ : MeasureTheory.Measure α) (score tie : α → ℝ) (v priority : ℝ) =>
  (μ {y : α | score y < v ∨ score y = v ∧ tie y ≤ priority}).toReal
\end{ReviewVerbatim}


\hypertarget{governing-declaration-29f9270f1ab8f757}{}
\subsection*{\small\ttfamily\raggedright LBG22\allowbreak{}Strategic\allowbreak{}Ranking.\allowbreak{}score\allowbreak{}Priority\allowbreak{}Band}
\textbf{Declaration location:} paper-local
\paragraph{Lean declaration body}
\begin{ReviewVerbatim}
type:
{α : Type u_1} →
  [inst : MeasurableSpace α] →
    MeasureTheory.Measure α → (α → ℝ) → (α → ℝ) → {n : ℕ} → (Fin (n + 1) → ℝ) → ℝ → ℝ → Fin (n + 1)

value:
fun {α : Type u_1} [MeasurableSpace α] (μ : MeasureTheory.Measure α) (score tie : α → ℝ) {n : ℕ}
    (cutoff : Fin (n + 1) → ℝ) (v priority : ℝ) =>
  Finset.univ.sup fun (i : Fin (n + 1)) =>
    if LBG22StrategicRanking.ScorePriorityAboveCutoff μ score tie (cutoff i) v priority then i else 0
\end{ReviewVerbatim}

\paragraph{Direct retained dependencies}
\begin{itemize}\raggedright
\item \hyperlink{paper-prerequisite-38434fa78df931b7}{LBG22\allowbreak{}Strategic\allowbreak{}Ranking.\allowbreak{}Score\allowbreak{}Priority\allowbreak{}Above\allowbreak{}Cutoff}
\end{itemize}
\hypertarget{governing-declaration-8658f19612e370c1}{}
\subsection*{\small\ttfamily\raggedright LBG22\allowbreak{}Strategic\allowbreak{}Ranking.\allowbreak{}source\allowbreak{}Clamped\allowbreak{}Skill}
\textbf{Declaration location:} paper-local
\paragraph{Lean declaration body}
\begin{ReviewVerbatim}
type:
(ℝ → ℝ) → ℝ → ℝ

value:
fun (skill : ℝ → ℝ) (t : ℝ) => skill (max 0 (min t 1))
\end{ReviewVerbatim}


\hypertarget{governing-declaration-e17c2a81c5471276}{}
\subsection*{\small\ttfamily\raggedright LBG22\allowbreak{}Strategic\allowbreak{}Ranking.\allowbreak{}source\allowbreak{}Combined\allowbreak{}Skill}
\textbf{Declaration location:} paper-local
\paragraph{Lean declaration body}
\begin{ReviewVerbatim}
type:
{ι : Type u_1} → [Fintype ι] → [Nonempty ι] → (ι → ℝ) → ℝ

value:
fun {ι : Type u_1} [Fintype ι] [Nonempty ι] (coefficient : ι → ℝ) => Finset.univ.sup' Finset.univ_nonempty coefficient
\end{ReviewVerbatim}


\hypertarget{governing-declaration-acc9172ca20a4c54}{}
\subsection*{\small\ttfamily\raggedright LBG22\allowbreak{}Strategic\allowbreak{}Ranking.\allowbreak{}source\allowbreak{}Effort\allowbreak{}At\allowbreak{}Score}
\textbf{Declaration location:} paper-local
\paragraph{Lean declaration body}
\begin{ReviewVerbatim}
type:
(ℝ → ℝ) → ℝ → ℝ → ℝ

value:
fun (production : ℝ → ℝ) (effortMax z : ℝ) =>
  LBG22StrategicRanking.effortIntervalInverse production effortMax (max z (production 0))
\end{ReviewVerbatim}

\paragraph{Direct retained dependencies}
\begin{itemize}\raggedright
\item \hyperlink{governing-declaration-c5b2c07abbff7224}{LBG22\allowbreak{}Strategic\allowbreak{}Ranking.\allowbreak{}effort\allowbreak{}Interval\allowbreak{}Inverse}
\end{itemize}
\hypertarget{governing-declaration-e466c69a7b86c3d0}{}
\subsection*{\small\ttfamily\raggedright LBG22\allowbreak{}Strategic\allowbreak{}Ranking.\allowbreak{}source\allowbreak{}Cost\allowbreak{}At\allowbreak{}Score}
\textbf{Declaration location:} paper-local
\paragraph{Lean declaration body}
\begin{ReviewVerbatim}
type:
(ℝ → ℝ) → (ℝ → ℝ) → ℝ → ℝ → ℝ

value:
fun (cost production : ℝ → ℝ) (effortMax z : ℝ) =>
  cost (LBG22StrategicRanking.sourceEffortAtScore production effortMax z)
\end{ReviewVerbatim}

\paragraph{Direct retained dependencies}
\begin{itemize}\raggedright
\item \hyperlink{governing-declaration-acc9172ca20a4c54}{LBG22\allowbreak{}Strategic\allowbreak{}Ranking.\allowbreak{}source\allowbreak{}Effort\allowbreak{}At\allowbreak{}Score}
\end{itemize}
\hypertarget{governing-declaration-e25b538a94797d96}{}
\subsection*{\small\ttfamily\raggedright LBG22\allowbreak{}Strategic\allowbreak{}Ranking.\allowbreak{}source\allowbreak{}Boundary\allowbreak{}Step\allowbreak{}Effort}
\textbf{Declaration location:} paper-local
\paragraph{Lean declaration body}
\begin{ReviewVerbatim}
type:
(ℝ → ℝ) → (ℝ → ℝ) → ℝ → ℝ → ℝ → ℝ → ℝ

value:
fun (cost production : ℝ → ℝ) (effortMax previousScore cutoffSkill rewardIncrement : ℝ) =>
  LBG22StrategicRanking.effortIntervalInverse cost effortMax
    (LBG22StrategicRanking.sourceCostAtScore cost production effortMax (previousScore / cutoffSkill) + rewardIncrement)
\end{ReviewVerbatim}

\paragraph{Direct retained dependencies}
\begin{itemize}\raggedright
\item \hyperlink{governing-declaration-c5b2c07abbff7224}{LBG22\allowbreak{}Strategic\allowbreak{}Ranking.\allowbreak{}effort\allowbreak{}Interval\allowbreak{}Inverse}
\item \hyperlink{governing-declaration-e466c69a7b86c3d0}{LBG22\allowbreak{}Strategic\allowbreak{}Ranking.\allowbreak{}source\allowbreak{}Cost\allowbreak{}At\allowbreak{}Score}
\end{itemize}
\hypertarget{governing-declaration-c452b1857fa41573}{}
\subsection*{\small\ttfamily\raggedright LBG22\allowbreak{}Strategic\allowbreak{}Ranking.\allowbreak{}source\allowbreak{}Boundary\allowbreak{}Step\allowbreak{}Score}
\textbf{Declaration location:} paper-local
\paragraph{Lean declaration body}
\begin{ReviewVerbatim}
type:
(ℝ → ℝ) → (ℝ → ℝ) → ℝ → ℝ → ℝ → ℝ → ℝ

value:
fun (cost production : ℝ → ℝ) (effortMax previousScore cutoffSkill rewardIncrement : ℝ) =>
  production
      (LBG22StrategicRanking.sourceBoundaryStepEffort cost production effortMax previousScore cutoffSkill
        rewardIncrement) *
    cutoffSkill
\end{ReviewVerbatim}

\paragraph{Direct retained dependencies}
\begin{itemize}\raggedright
\item \hyperlink{governing-declaration-e25b538a94797d96}{LBG22\allowbreak{}Strategic\allowbreak{}Ranking.\allowbreak{}source\allowbreak{}Boundary\allowbreak{}Step\allowbreak{}Effort}
\end{itemize}
\hypertarget{governing-declaration-efba6aee5c5b1874}{}
\subsection*{\small\ttfamily\raggedright LBG22\allowbreak{}Strategic\allowbreak{}Ranking.\allowbreak{}source\allowbreak{}Multitask\allowbreak{}Cost}
\textbf{Declaration location:} paper-local
\paragraph{Lean declaration body}
\begin{ReviewVerbatim}
type:
(ℝ → ℝ) → ℝ → ℝ → ℝ → ℝ

value:
fun (cost : ℝ → ℝ) (B eM eU : ℝ) => cost eM - max 0 (B - (eM + eU)) ^ 2
\end{ReviewVerbatim}


\hypertarget{governing-declaration-1166d98d13859ee1}{}
\subsection*{\small\ttfamily\raggedright LBG22\allowbreak{}Strategic\allowbreak{}Ranking.\allowbreak{}source\allowbreak{}Recursive\allowbreak{}Band\allowbreak{}Score}
\textbf{Declaration location:} paper-local
\paragraph{Lean declaration body}
\begin{ReviewVerbatim}
type:
(ℝ → ℝ) → (ℝ → ℝ) → ℝ → (ℕ → ℝ) → (ℕ → ℝ) → ℕ → ℝ

value:
fun (cost production : ℝ → ℝ) (effortMax : ℝ) (cutoffSkill reward : ℕ → ℝ) (a : ℕ) =>
  Nat.brecOn a (LBG22StrategicRanking.sourceRecursiveBandScore._f cost production effortMax cutoffSkill reward)
\end{ReviewVerbatim}

\paragraph{Direct retained dependencies}
\begin{itemize}\raggedright
\item \hyperlink{governing-declaration-c452b1857fa41573}{LBG22\allowbreak{}Strategic\allowbreak{}Ranking.\allowbreak{}source\allowbreak{}Boundary\allowbreak{}Step\allowbreak{}Score}
\end{itemize}
\hypertarget{governing-declaration-8928c46fd984c29b}{}
\subsection*{\small\ttfamily\raggedright LBG22\allowbreak{}Strategic\allowbreak{}Ranking.\allowbreak{}source\allowbreak{}Finite\allowbreak{}Rank\allowbreak{}Effort}
\textbf{Declaration location:} paper-local
\paragraph{Lean declaration body}
\begin{ReviewVerbatim}
type:
(ℝ → ℝ) → (ℝ → ℝ) → (ℝ → ℝ) → ℝ → ℕ → (ℕ → ℝ) → (ℕ → ℝ) → ℝ → ℝ

value:
fun (cost production skill : ℝ → ℝ) (effortMax : ℝ) (n : ℕ) (cutoff reward : ℕ → ℝ) (t : ℝ) =>
  LBG22StrategicRanking.sourceEffortAtScore production effortMax
    (LBG22StrategicRanking.sourceRecursiveBandScore cost production effortMax (fun (k : ℕ) => skill (cutoff k)) reward
        ↑(LBG22StrategicRanking.finiteLowerRankBand (fun (i : Fin (n + 1)) => cutoff ↑i) t) /
      LBG22StrategicRanking.sourceClampedSkill skill t)
\end{ReviewVerbatim}

\paragraph{Direct retained dependencies}
\begin{itemize}\raggedright
\item \hyperlink{paper-prerequisite-a42f440ddc254fa7}{LBG22\allowbreak{}Strategic\allowbreak{}Ranking.\allowbreak{}finite\allowbreak{}Lower\allowbreak{}Rank\allowbreak{}Band}
\item \hyperlink{governing-declaration-8658f19612e370c1}{LBG22\allowbreak{}Strategic\allowbreak{}Ranking.\allowbreak{}source\allowbreak{}Clamped\allowbreak{}Skill}
\item \hyperlink{governing-declaration-acc9172ca20a4c54}{LBG22\allowbreak{}Strategic\allowbreak{}Ranking.\allowbreak{}source\allowbreak{}Effort\allowbreak{}At\allowbreak{}Score}
\item \hyperlink{governing-declaration-1166d98d13859ee1}{LBG22\allowbreak{}Strategic\allowbreak{}Ranking.\allowbreak{}source\allowbreak{}Recursive\allowbreak{}Band\allowbreak{}Score}
\end{itemize}
\hypertarget{governing-declaration-a63079aa3a21a0d6}{}
\subsection*{\small\ttfamily\raggedright LBG22\allowbreak{}Strategic\allowbreak{}Ranking.\allowbreak{}source\allowbreak{}Finite\allowbreak{}Baseline\allowbreak{}Effort}
\textbf{Declaration location:} paper-local
\paragraph{Lean declaration body}
\begin{ReviewVerbatim}
type:
(ℝ → ℝ) → (ℝ → ℝ) → (ℝ → ℝ) → ℝ → ℝ → ℕ → (ℕ → ℝ) → (ℕ → ℝ) → ℝ → ℝ

value:
fun (cost production skill : ℝ → ℝ) (baseline effortMax : ℝ) (n : ℕ) (cutoff reward : ℕ → ℝ) (t : ℝ) =>
  baseline +
    LBG22StrategicRanking.sourceFiniteRankEffort (fun (e : ℝ) => cost (baseline + e))
      (fun (e : ℝ) => production (baseline + e)) skill effortMax n cutoff reward t
\end{ReviewVerbatim}

\paragraph{Direct retained dependencies}
\begin{itemize}\raggedright
\item \hyperlink{governing-declaration-8928c46fd984c29b}{LBG22\allowbreak{}Strategic\allowbreak{}Ranking.\allowbreak{}source\allowbreak{}Finite\allowbreak{}Rank\allowbreak{}Effort}
\end{itemize}
\hypertarget{governing-declaration-75ee87c7ce9350ef}{}
\subsection*{\small\ttfamily\raggedright LBG22\allowbreak{}Strategic\allowbreak{}Ranking.\allowbreak{}source\allowbreak{}Finite\allowbreak{}Rank\allowbreak{}Score}
\textbf{Declaration location:} paper-local
\paragraph{Lean declaration body}
\begin{ReviewVerbatim}
type:
(ℝ → ℝ) → (ℝ → ℝ) → (ℝ → ℝ) → ℝ → ℕ → (ℕ → ℝ) → (ℕ → ℝ) → ℝ → ℝ

value:
fun (cost production skill : ℝ → ℝ) (effortMax : ℝ) (n : ℕ) (cutoff reward : ℕ → ℝ) (t : ℝ) =>
  max
    (LBG22StrategicRanking.sourceRecursiveBandScore cost production effortMax (fun (k : ℕ) => skill (cutoff k)) reward
      ↑(LBG22StrategicRanking.finiteLowerRankBand (fun (i : Fin (n + 1)) => cutoff ↑i) t))
    (production 0 * LBG22StrategicRanking.sourceClampedSkill skill t)
\end{ReviewVerbatim}

\paragraph{Direct retained dependencies}
\begin{itemize}\raggedright
\item \hyperlink{paper-prerequisite-a42f440ddc254fa7}{LBG22\allowbreak{}Strategic\allowbreak{}Ranking.\allowbreak{}finite\allowbreak{}Lower\allowbreak{}Rank\allowbreak{}Band}
\item \hyperlink{governing-declaration-8658f19612e370c1}{LBG22\allowbreak{}Strategic\allowbreak{}Ranking.\allowbreak{}source\allowbreak{}Clamped\allowbreak{}Skill}
\item \hyperlink{governing-declaration-1166d98d13859ee1}{LBG22\allowbreak{}Strategic\allowbreak{}Ranking.\allowbreak{}source\allowbreak{}Recursive\allowbreak{}Band\allowbreak{}Score}
\end{itemize}
\hypertarget{governing-declaration-927fd67abc7df008}{}
\subsection*{\small\ttfamily\raggedright LBG22\allowbreak{}Strategic\allowbreak{}Ranking.\allowbreak{}source\allowbreak{}Finite\allowbreak{}Baseline\allowbreak{}Score}
\textbf{Declaration location:} paper-local
\paragraph{Lean declaration body}
\begin{ReviewVerbatim}
type:
(ℝ → ℝ) → (ℝ → ℝ) → (ℝ → ℝ) → ℝ → ℝ → ℕ → (ℕ → ℝ) → (ℕ → ℝ) → ℝ → ℝ

value:
fun (cost production skill : ℝ → ℝ) (baseline effortMax : ℝ) (n : ℕ) (cutoff reward : ℕ → ℝ) (a : ℝ) =>
  LBG22StrategicRanking.sourceFiniteRankScore (fun (e : ℝ) => cost (baseline + e))
    (fun (e : ℝ) => production (baseline + e)) skill effortMax n cutoff reward a
\end{ReviewVerbatim}

\paragraph{Direct retained dependencies}
\begin{itemize}\raggedright
\item \hyperlink{governing-declaration-75ee87c7ce9350ef}{LBG22\allowbreak{}Strategic\allowbreak{}Ranking.\allowbreak{}source\allowbreak{}Finite\allowbreak{}Rank\allowbreak{}Score}
\end{itemize}
\hypertarget{governing-declaration-2190069be417355f}{}
\subsection*{\small\ttfamily\raggedright LBG22\allowbreak{}Strategic\allowbreak{}Ranking.\allowbreak{}tie\allowbreak{}Broken\allowbreak{}Lower\allowbreak{}Contour}
\textbf{Declaration location:} paper-local
\paragraph{Lean declaration body}
\begin{ReviewVerbatim}
type:
{α : Type u_1} → (α → ℝ) → (α → ℝ) → α → α → Prop

value:
fun {α : Type u_1} (score tie : α → ℝ) (x a : α) => {y : α | score y < score x ∨ score y = score x ∧ tie y ≤ tie x} a
\end{ReviewVerbatim}


\hypertarget{governing-declaration-bfa45ba26bb8d395}{}
\subsection*{\small\ttfamily\raggedright LBG22\allowbreak{}Strategic\allowbreak{}Ranking.\allowbreak{}Source\allowbreak{}Hard\allowbreak{}Budget\allowbreak{}Best\allowbreak{}Response\allowbreak{}At}
\textbf{Declaration location:} paper-local
\paragraph{Lean declaration body}
\begin{ReviewVerbatim}
type:
{α : Type u_1} →
  [inst : MeasurableSpace α] →
    (μ : MeasureTheory.Measure α) →
      [MeasureTheory.IsFiniteMeasure μ] → (ℝ → ℝ) → (ℝ → ℝ) → (α → ℝ) → (α → ℝ) → (α → ℝ) → ℝ → ℝ → ℝ → α → ℝ → ℝ → Prop

value:
fun {α : Type u_1} [MeasurableSpace α] (μ : MeasureTheory.Measure α) [MeasureTheory.IsFiniteMeasure μ]
    (cost production : ℝ → ℝ) (skill score tie : α → ℝ) (B rho c : ℝ) (x : α) (eM eU : ℝ) =>
  0 ≤ eM ∧
    0 ≤ eU ∧
      eM + eU = B ∧
        score x = production eM * skill x ∧
          ∀ (dM dU : ℝ),
            0 ≤ dM →
              0 ≤ dU →
                dM + dU = B →
                  LBG22StrategicRanking.sourceTwoLevelReward rho c
                        ↑(LBG22StrategicRanking.finiteLowerRankBand
                            (fun (i : Fin 2) => LBG22StrategicRanking.sourceTwoLevelCutoff c ↑i)
                            (LBG22StrategicRanking.counterfactualTieBrokenRank μ score tie x
                              (production dM * skill x))) -
                      LBG22StrategicRanking.sourceMultitaskCost cost B dM dU ≤
                    LBG22StrategicRanking.sourceTwoLevelReward rho c
                        ↑(LBG22StrategicRanking.finiteLowerRankBand
                            (fun (i : Fin 2) => LBG22StrategicRanking.sourceTwoLevelCutoff c ↑i)
                            (LBG22StrategicRanking.tieBrokenRank μ score tie x)) -
                      LBG22StrategicRanking.sourceMultitaskCost cost B eM eU
\end{ReviewVerbatim}

\paragraph{Direct retained dependencies}
\begin{itemize}\raggedright
\item \hyperlink{paper-prerequisite-6ed45e46d01eadeb}{LBG22\allowbreak{}Strategic\allowbreak{}Ranking.\allowbreak{}counterfactual\allowbreak{}Tie\allowbreak{}Broken\allowbreak{}Rank}
\item \hyperlink{paper-prerequisite-a42f440ddc254fa7}{LBG22\allowbreak{}Strategic\allowbreak{}Ranking.\allowbreak{}finite\allowbreak{}Lower\allowbreak{}Rank\allowbreak{}Band}
\item \hyperlink{governing-declaration-efba6aee5c5b1874}{LBG22\allowbreak{}Strategic\allowbreak{}Ranking.\allowbreak{}source\allowbreak{}Multitask\allowbreak{}Cost}
\item \hyperlink{paper-prerequisite-f88aac58507e44ca}{LBG22\allowbreak{}Strategic\allowbreak{}Ranking.\allowbreak{}source\allowbreak{}Two\allowbreak{}Level\allowbreak{}Cutoff}
\item \hyperlink{paper-prerequisite-30d97cd43eb79017}{LBG22\allowbreak{}Strategic\allowbreak{}Ranking.\allowbreak{}source\allowbreak{}Two\allowbreak{}Level\allowbreak{}Reward}
\item \hyperlink{paper-prerequisite-b7c71e53177b76c6}{LBG22\allowbreak{}Strategic\allowbreak{}Ranking.\allowbreak{}tie\allowbreak{}Broken\allowbreak{}Rank}
\end{itemize}
\hypertarget{governing-declaration-418baddf0ad9374c}{}
\subsection*{\small\ttfamily\raggedright LBG22\allowbreak{}Strategic\allowbreak{}Ranking.\allowbreak{}Source\allowbreak{}Population\allowbreak{}Finite\allowbreak{}Best\allowbreak{}Response\allowbreak{}At}
\textbf{Declaration location:} paper-local
\paragraph{Lean declaration body}
\begin{ReviewVerbatim}
type:
{α : Type u_1} →
  [inst : MeasurableSpace α] →
    (μ : MeasureTheory.Measure α) →
      [MeasureTheory.IsFiniteMeasure μ] →
        (ℝ → ℝ) → (ℝ → ℝ) → (α → ℝ) → (α → ℝ) → (α → ℝ) → {n : ℕ} → (Fin (n + 1) → ℝ) → (ℕ → ℝ) → α → ℝ → Prop

value:
fun {α : Type u_1} [MeasurableSpace α] (μ : MeasureTheory.Measure α) [MeasureTheory.IsFiniteMeasure μ]
    (cost production : ℝ → ℝ) (skill score tie : α → ℝ) {n : ℕ} (cutoff : Fin (n + 1) → ℝ) (reward : ℕ → ℝ) (x : α)
    (effort : ℝ) =>
  0 ≤ effort ∧
    score x = production effort * skill x ∧
      ∀ (d : ℝ),
        0 ≤ d →
          reward
                ↑(LBG22StrategicRanking.finiteLowerRankBand cutoff
                    (LBG22StrategicRanking.counterfactualTieBrokenRank μ score tie x (production d * skill x))) -
              cost d ≤
            reward
                ↑(LBG22StrategicRanking.finiteLowerRankBand cutoff
                    (LBG22StrategicRanking.tieBrokenRank μ score tie x)) -
              cost effort
\end{ReviewVerbatim}

\paragraph{Direct retained dependencies}
\begin{itemize}\raggedright
\item \hyperlink{paper-prerequisite-6ed45e46d01eadeb}{LBG22\allowbreak{}Strategic\allowbreak{}Ranking.\allowbreak{}counterfactual\allowbreak{}Tie\allowbreak{}Broken\allowbreak{}Rank}
\item \hyperlink{paper-prerequisite-a42f440ddc254fa7}{LBG22\allowbreak{}Strategic\allowbreak{}Ranking.\allowbreak{}finite\allowbreak{}Lower\allowbreak{}Rank\allowbreak{}Band}
\item \hyperlink{paper-prerequisite-b7c71e53177b76c6}{LBG22\allowbreak{}Strategic\allowbreak{}Ranking.\allowbreak{}tie\allowbreak{}Broken\allowbreak{}Rank}
\end{itemize}
\hypertarget{governing-declaration-678cf48d549cb620}{}
\subsection*{\small\ttfamily\raggedright LBG22\allowbreak{}Strategic\allowbreak{}Ranking.\allowbreak{}unit\allowbreak{}Rank\allowbreak{}Measure}
\textbf{Declaration location:} paper-local
\paragraph{Lean declaration body}
\begin{ReviewVerbatim}
type:
MeasureTheory.Measure ℝ

value:
MeasureTheory.volume.restrict (Set.Ioc 0 1)
\end{ReviewVerbatim}


\hypertarget{governing-declaration-8d3290d5e2a0343e}{}
\subsection*{\small\ttfamily\raggedright LBG22\allowbreak{}Strategic\allowbreak{}Ranking.\allowbreak{}independent\allowbreak{}Skill\allowbreak{}Population}
\textbf{Declaration location:} paper-local
\paragraph{Lean declaration body}
\begin{ReviewVerbatim}
type:
{α : Type u_1} → [inst : MeasurableSpace α] → MeasureTheory.Measure α → MeasureTheory.Measure ((ℝ × ℝ) × α)

value:
fun {α : Type u_1} [MeasurableSpace α] (μ : MeasureTheory.Measure α) =>
  (LBG22StrategicRanking.unitRankMeasure.prod LBG22StrategicRanking.unitRankMeasure).prod μ
\end{ReviewVerbatim}

\paragraph{Direct retained dependencies}
\begin{itemize}\raggedright
\item \hyperlink{governing-declaration-678cf48d549cb620}{LBG22\allowbreak{}Strategic\allowbreak{}Ranking.\allowbreak{}unit\allowbreak{}Rank\allowbreak{}Measure}
\end{itemize}
\hypertarget{governing-declaration-9631d67a19f5df1f}{}
\subsection*{\small\ttfamily\raggedright LBG22\allowbreak{}Strategic\allowbreak{}Ranking.\allowbreak{}source\allowbreak{}Multidimensional\allowbreak{}Pre\allowbreak{}Rank}
\textbf{Declaration location:} paper-local
\paragraph{Lean declaration body}
\begin{ReviewVerbatim}
type:
{ι : Type u_1} → [Fintype ι] → [Nonempty ι] → (ι → ℝ) → (ι → ℝ → ℝ) → (ι → ℝ) → ℝ

value:
fun {ι : Type u_1} [Fintype ι] [Nonempty ι] (weight : ι → ℝ) (skill : ι → ℝ → ℝ) (rank : ι → ℝ) =>
  ↑(ProbabilityTheory.cdf
        (MeasureTheory.Measure.map
          (fun (x : ι → ℝ) =>
            LBG22StrategicRanking.sourceCombinedSkill
              (LBG22StrategicRanking.sourceMultidimensionalCoefficient weight skill x))
          (LBG22StrategicRanking.sourceMultidimensionalPopulation ι)))
    (LBG22StrategicRanking.sourceCombinedSkill
      (LBG22StrategicRanking.sourceMultidimensionalCoefficient weight skill rank))
\end{ReviewVerbatim}

\paragraph{Direct retained dependencies}
\begin{itemize}\raggedright
\item \hyperlink{governing-declaration-e17c2a81c5471276}{LBG22\allowbreak{}Strategic\allowbreak{}Ranking.\allowbreak{}source\allowbreak{}Combined\allowbreak{}Skill}
\item \hyperlink{paper-prerequisite-e57e55071c75ce91}{LBG22\allowbreak{}Strategic\allowbreak{}Ranking.\allowbreak{}source\allowbreak{}Multidimensional\allowbreak{}Coefficient}
\item \hyperlink{paper-prerequisite-16d135799a0f0491}{LBG22\allowbreak{}Strategic\allowbreak{}Ranking.\allowbreak{}source\allowbreak{}Multidimensional\allowbreak{}Population}
\end{itemize}
\hypertarget{governing-declaration-f1a76532081bf3d3}{}
\subsection*{\small\ttfamily\raggedright LBG22\allowbreak{}Strategic\allowbreak{}Ranking.\allowbreak{}source\allowbreak{}Multitask\allowbreak{}Admission\allowbreak{}Event}
\textbf{Declaration location:} paper-local
\paragraph{Lean declaration body}
\begin{ReviewVerbatim}
type:
{α : Type u_1} →
  [inst : MeasurableSpace α] →
    (μ : MeasureTheory.Measure α) →
      [MeasureTheory.IsProbabilityMeasure μ] → (ℝ × α → ℝ) → (ℝ × α → ℝ) → ℝ → ℝ → (ℝ × ℝ) × α → Prop

value:
fun {α : Type u_1} [MeasurableSpace α] (μ : MeasureTheory.Measure α) [MeasureTheory.IsProbabilityMeasure μ]
    (score tie : ℝ × α → ℝ) (rho c : ℝ) (a : (ℝ × ℝ) × α) =>
  {z : (ℝ × ℝ) × α |
      c < LBG22StrategicRanking.tieBrokenRank (LBG22StrategicRanking.unitRankMeasure.prod μ) score tie (z.1.1, z.2) ∧
        z.1.2 ∈ Set.Ioc 0 (rho / (1 - c))}
    a
\end{ReviewVerbatim}

\paragraph{Direct retained dependencies}
\begin{itemize}\raggedright
\item \hyperlink{paper-prerequisite-b7c71e53177b76c6}{LBG22\allowbreak{}Strategic\allowbreak{}Ranking.\allowbreak{}tie\allowbreak{}Broken\allowbreak{}Rank}
\item \hyperlink{governing-declaration-678cf48d549cb620}{LBG22\allowbreak{}Strategic\allowbreak{}Ranking.\allowbreak{}unit\allowbreak{}Rank\allowbreak{}Measure}
\end{itemize}
\hypertarget{governing-declaration-15b79809af308d40}{}
\subsection*{\small\ttfamily\raggedright LBG22\allowbreak{}Strategic\allowbreak{}Ranking.\allowbreak{}source\allowbreak{}Skill\allowbreak{}CDF}
\textbf{Declaration location:} paper-local
\paragraph{Lean declaration body}
\begin{ReviewVerbatim}
type:
(ℝ → ℝ) → ℝ → ℝ

value:
fun (skill : ℝ → ℝ) (a : ℝ) =>
  ↑(ProbabilityTheory.cdf
        (MeasureTheory.Measure.map (LBG22StrategicRanking.sourceClampedSkill skill)
          LBG22StrategicRanking.unitRankMeasure))
    a
\end{ReviewVerbatim}

\paragraph{Direct retained dependencies}
\begin{itemize}\raggedright
\item \hyperlink{governing-declaration-8658f19612e370c1}{LBG22\allowbreak{}Strategic\allowbreak{}Ranking.\allowbreak{}source\allowbreak{}Clamped\allowbreak{}Skill}
\item \hyperlink{governing-declaration-678cf48d549cb620}{LBG22\allowbreak{}Strategic\allowbreak{}Ranking.\allowbreak{}unit\allowbreak{}Rank\allowbreak{}Measure}
\end{itemize}
\hypertarget{governing-declaration-97f935133f47ad33}{}
\subsection*{\small\ttfamily\raggedright LBG22\allowbreak{}Strategic\allowbreak{}Ranking.\allowbreak{}source\allowbreak{}Disadvantaged\allowbreak{}Rank}
\textbf{Declaration location:} paper-local
\paragraph{Lean declaration body}
\begin{ReviewVerbatim}
type:
(ℝ → ℝ) → ℝ → ℝ → ℝ → ℝ

value:
fun (skill : ℝ → ℝ) (psiA psiB t : ℝ) => LBG22StrategicRanking.sourceEnvironmentRank skill psiA psiB (false, t)
\end{ReviewVerbatim}

\paragraph{Direct retained dependencies}
\begin{itemize}\raggedright
\item \hyperlink{paper-prerequisite-b308d7da226a5fe1}{LBG22\allowbreak{}Strategic\allowbreak{}Ranking.\allowbreak{}source\allowbreak{}Environment\allowbreak{}Rank}
\end{itemize}
\hypertarget{governing-declaration-b9a91eff0bdd1f65}{}
\subsection*{\small\ttfamily\raggedright LBG22\allowbreak{}Strategic\allowbreak{}Ranking.\allowbreak{}source\allowbreak{}Disadvantaged\allowbreak{}Threshold}
\textbf{Declaration location:} paper-local
\paragraph{Lean declaration body}
\begin{ReviewVerbatim}
type:
(ℝ → ℝ) → ℝ → ℝ → ℝ → ℝ

value:
fun (skill : ℝ → ℝ) (psiA psiB c : ℝ) =>
  LBG22StrategicRanking.sourceSkillCDF (LBG22StrategicRanking.sourceDisadvantagedRank skill psiA psiB) c
\end{ReviewVerbatim}

\paragraph{Direct retained dependencies}
\begin{itemize}\raggedright
\item \hyperlink{governing-declaration-97f935133f47ad33}{LBG22\allowbreak{}Strategic\allowbreak{}Ranking.\allowbreak{}source\allowbreak{}Disadvantaged\allowbreak{}Rank}
\item \hyperlink{governing-declaration-15b79809af308d40}{LBG22\allowbreak{}Strategic\allowbreak{}Ranking.\allowbreak{}source\allowbreak{}Skill\allowbreak{}CDF}
\end{itemize}
\hypertarget{governing-declaration-1f8a1d5f3261f80f}{}
\subsection*{\small\ttfamily\raggedright LBG22\allowbreak{}Strategic\allowbreak{}Ranking.\allowbreak{}upper\allowbreak{}Tail\allowbreak{}Rank}
\textbf{Declaration location:} paper-local
\paragraph{Lean declaration body}
\begin{ReviewVerbatim}
type:
ℝ → ℝ

value:
fun (x : ℝ) => 1 - Real.exp (-x)
\end{ReviewVerbatim}


\clearpage
\hypertarget{paper-prerequisite-section-5ef5ef0364b6939c}{}
\section*{Paper-specific semantic prerequisites}
\hypertarget{paper-prerequisite-73413973c7594479}{}
\subsection*{\small\ttfamily\raggedright LBG22\allowbreak{}Strategic\allowbreak{}Ranking.\allowbreak{}source\allowbreak{}Environment\allowbreak{}Skill}
\noindent\textbf{Paper-local declaration:}\par
{\footnotesize\ttfamily\raggedright LBG22\allowbreak{}Strategic\allowbreak{}Ranking.\allowbreak{}source\allowbreak{}Environment\allowbreak{}Skill\par}
\textbf{Source locator:} {\footnotesize\raggedright cited publication, lines 210-\allowbreak{}218\par}
\paragraph{Verbatim paper-source connection}
\begin{ReviewVerbatim}
\paragraph{Model and equilibria characterization} We now denote each applicant's latent skill rank as $\theta_\true \in [0,1]$. In addition to the latent skill, each applicant has an (unobserved) environmental factor $\psi \in \Psi$ that represents how favorable their environment is for attaining a higher score. Because a favorable environment results in a higher rate of return for effort, we model the environment as a multiplicative factor in the score (see~e.g., \citet{calsamiglia09decent}). Formally, the post-effort score is a function of the latent skill, the environmental factor, and the effort level:
\begin{equation*}\label{eq:env-score}
v = \psi \cdot  g(e)\cdot f(\theta_\true),
\end{equation*}
where $g, f$ are as defined in Section~\ref{sec:model}. %We also require $f$ to be strictly monotone.

%We model structural inequalities as heterogeneous returns to effort. In contrast to results in Section~\ref{sec:equillemmas}, the ranks of applicants are no longer preserved when there are heterogeneous returns to effort.

 We assume there are two groups of applicants, $\A$ and $\B$, and the distribution of skill is the same in both groups. Group $\A$ has a more favorable environment factor, that is,  $\Psi = \{\psi_\A, \psi_\B\}$ and $\psi_\A > \psi_\B$. Thus we will also refer to $\B$ as the ``disadvantaged group". To simplify our presentation, we assume each group is half of the total applicant population, though the results in this section generalize.% to the setting where groups are not equal-sized.

[Next verbatim source excerpt]

We begin by characterizing the equilibrium ranking under the designer's policy $\lambda$. The equilibrium effort levels and post-effort ranks are as defined in Definition~\ref{def:equi}, except they are now group-dependent, that is, we have $e(\theta_\true, \psi)$ and $\theta_\post(\theta_\true, \psi)$. Because of the differences in $\psi$, the post-effort ranking $\theta_\post$ is now group-dependent, and in general is not equal to $\theta_\true$. To apply Proposition~\ref{lem:rankpreservedindex} as before, we construct an ``environment-scaled pre-effort rank'' (denoted $\theta_\pre$).
\end{ReviewVerbatim}

% report-context-presentation-sha256: 90f9116e292dd71fddfd6b445ca0b060a481d026855bdb9725f2833c2cf62a31
\paragraph{Settled source reading}
\begin{sloppypar}
Equilibrium means that almost every applicant best responds against every feasible deviation;\allowbreak{} uniqueness is up to a null set.\allowbreak{}
\end{sloppypar}
\paragraph{Settled source reading}
\begin{sloppypar}
Applicants know the fixed priority used at score-\allowbreak{}boundary ties.\allowbreak{} This clarifies the stated model, not a corrected admission policy or a new economic assumption.\allowbreak{}
\end{sloppypar}

\paragraph{Lean-expanded paper semantic target}
\begin{ReviewVerbatim}
type:
(ℝ → ℝ) → ℝ → ℝ → Bool × ℝ → ℝ

value:
fun (skill : ℝ → ℝ) (psiA psiB : ℝ) (x : Bool × ℝ) =>
  (if x.1 = true then psiA else psiB) * LBG22StrategicRanking.sourceClampedSkill skill x.2
\end{ReviewVerbatim}

\paragraph{Direct retained dependencies}
\begin{itemize}\raggedright
\item \hyperlink{governing-declaration-8658f19612e370c1}{LBG22\allowbreak{}Strategic\allowbreak{}Ranking.\allowbreak{}source\allowbreak{}Clamped\allowbreak{}Skill}
\end{itemize}
\paragraph{Recorded source-to-declaration screening}
\textbf{Verdict:} \texttt{matches}
\\
{\raggedright
\textbf{Reason:} The target multiplies the common skill quantile by psiA for group A and psiB for group B.\allowbreak{} Together with the displayed clamping to the unit rank support, this is the source environment-\allowbreak{}scaled skill formula.\allowbreak{}
\par}
\reviewmatch{Matches source input}
\noindent\textbf{Reviewer annotation}\par
\reviewerbox
\clearpage
\hypertarget{paper-prerequisite-5b3be494cde30a37}{}
\subsection*{\small\ttfamily\raggedright LBG22\allowbreak{}Strategic\allowbreak{}Ranking.\allowbreak{}source\allowbreak{}Two\allowbreak{}Group\allowbreak{}Measure}
\noindent\textbf{Paper-local declaration:}\par
{\footnotesize\ttfamily\raggedright LBG22\allowbreak{}Strategic\allowbreak{}Ranking.\allowbreak{}source\allowbreak{}Two\allowbreak{}Group\allowbreak{}Measure\par}
\textbf{Source locator:} {\footnotesize\raggedright cited publication, lines 210-\allowbreak{}218\par}
\paragraph{Verbatim paper-source connection}
\begin{ReviewVerbatim}
\paragraph{Model and equilibria characterization} We now denote each applicant's latent skill rank as $\theta_\true \in [0,1]$. In addition to the latent skill, each applicant has an (unobserved) environmental factor $\psi \in \Psi$ that represents how favorable their environment is for attaining a higher score. Because a favorable environment results in a higher rate of return for effort, we model the environment as a multiplicative factor in the score (see~e.g., \citet{calsamiglia09decent}). Formally, the post-effort score is a function of the latent skill, the environmental factor, and the effort level:
\begin{equation*}\label{eq:env-score}
v = \psi \cdot  g(e)\cdot f(\theta_\true),
\end{equation*}
where $g, f$ are as defined in Section~\ref{sec:model}. %We also require $f$ to be strictly monotone.

%We model structural inequalities as heterogeneous returns to effort. In contrast to results in Section~\ref{sec:equillemmas}, the ranks of applicants are no longer preserved when there are heterogeneous returns to effort.

 We assume there are two groups of applicants, $\A$ and $\B$, and the distribution of skill is the same in both groups. Group $\A$ has a more favorable environment factor, that is,  $\Psi = \{\psi_\A, \psi_\B\}$ and $\psi_\A > \psi_\B$. Thus we will also refer to $\B$ as the ``disadvantaged group". To simplify our presentation, we assume each group is half of the total applicant population, though the results in this section generalize.% to the setting where groups are not equal-sized.

[Next verbatim source excerpt]

We begin by characterizing the equilibrium ranking under the designer's policy $\lambda$. The equilibrium effort levels and post-effort ranks are as defined in Definition~\ref{def:equi}, except they are now group-dependent, that is, we have $e(\theta_\true, \psi)$ and $\theta_\post(\theta_\true, \psi)$. Because of the differences in $\psi$, the post-effort ranking $\theta_\post$ is now group-dependent, and in general is not equal to $\theta_\true$. To apply Proposition~\ref{lem:rankpreservedindex} as before, we construct an ``environment-scaled pre-effort rank'' (denoted $\theta_\pre$).
\end{ReviewVerbatim}

% report-context-presentation-sha256: 90f9116e292dd71fddfd6b445ca0b060a481d026855bdb9725f2833c2cf62a31
\paragraph{Settled source reading}
\begin{sloppypar}
Equilibrium means that almost every applicant best responds against every feasible deviation;\allowbreak{} uniqueness is up to a null set.\allowbreak{}
\end{sloppypar}
\paragraph{Settled source reading}
\begin{sloppypar}
Applicants know the fixed priority used at score-\allowbreak{}boundary ties.\allowbreak{} This clarifies the stated model, not a corrected admission policy or a new economic assumption.\allowbreak{}
\end{sloppypar}

\paragraph{Lean-expanded paper semantic target}
\begin{ReviewVerbatim}
type:
MeasureTheory.Measure (Bool × ℝ)

value:
(1 / 2) • MeasureTheory.Measure.map (fun (t : ℝ) => (true, t)) LBG22StrategicRanking.unitRankMeasure +
  (1 / 2) • MeasureTheory.Measure.map (fun (t : ℝ) => (false, t)) LBG22StrategicRanking.unitRankMeasure
\end{ReviewVerbatim}

\paragraph{Direct retained dependencies}
\begin{itemize}\raggedright
\item \hyperlink{governing-declaration-678cf48d549cb620}{LBG22\allowbreak{}Strategic\allowbreak{}Ranking.\allowbreak{}unit\allowbreak{}Rank\allowbreak{}Measure}
\end{itemize}
\paragraph{Recorded source-to-declaration screening}
\textbf{Verdict:} \texttt{matches}
\\
{\raggedright
\textbf{Reason:} The target assigns half the population to each Boolean group and gives both groups the same uniform latent skill rank.\allowbreak{} This matches the source equal-\allowbreak{}sized two-\allowbreak{}group model with a common skill distribution.\allowbreak{}
\par}
\reviewmatch{Matches source input}
\noindent\textbf{Reviewer annotation}\par
\reviewerbox
\clearpage
\hypertarget{paper-prerequisite-b8e06f39e94d0049}{}
\subsection*{\small\ttfamily\raggedright LBG22\allowbreak{}Strategic\allowbreak{}Ranking.\allowbreak{}Environment\allowbreak{}Equilibrium}
\noindent\textbf{Paper-local declaration:}\par
{\footnotesize\ttfamily\raggedright LBG22\allowbreak{}Strategic\allowbreak{}Ranking.\allowbreak{}Environment\allowbreak{}Equilibrium\par}
\textbf{Source locator:} {\footnotesize\raggedright cited publication, lines 210-\allowbreak{}218\par}
\paragraph{Verbatim paper-source connection}
\begin{ReviewVerbatim}
\paragraph{Model and equilibria characterization} We now denote each applicant's latent skill rank as $\theta_\true \in [0,1]$. In addition to the latent skill, each applicant has an (unobserved) environmental factor $\psi \in \Psi$ that represents how favorable their environment is for attaining a higher score. Because a favorable environment results in a higher rate of return for effort, we model the environment as a multiplicative factor in the score (see~e.g., \citet{calsamiglia09decent}). Formally, the post-effort score is a function of the latent skill, the environmental factor, and the effort level:
\begin{equation*}\label{eq:env-score}
v = \psi \cdot  g(e)\cdot f(\theta_\true),
\end{equation*}
where $g, f$ are as defined in Section~\ref{sec:model}. %We also require $f$ to be strictly monotone.

%We model structural inequalities as heterogeneous returns to effort. In contrast to results in Section~\ref{sec:equillemmas}, the ranks of applicants are no longer preserved when there are heterogeneous returns to effort.

 We assume there are two groups of applicants, $\A$ and $\B$, and the distribution of skill is the same in both groups. Group $\A$ has a more favorable environment factor, that is,  $\Psi = \{\psi_\A, \psi_\B\}$ and $\psi_\A > \psi_\B$. Thus we will also refer to $\B$ as the ``disadvantaged group". To simplify our presentation, we assume each group is half of the total applicant population, though the results in this section generalize.% to the setting where groups are not equal-sized.

[Next verbatim source excerpt]

We begin by characterizing the equilibrium ranking under the designer's policy $\lambda$. The equilibrium effort levels and post-effort ranks are as defined in Definition~\ref{def:equi}, except they are now group-dependent, that is, we have $e(\theta_\true, \psi)$ and $\theta_\post(\theta_\true, \psi)$. Because of the differences in $\psi$, the post-effort ranking $\theta_\post$ is now group-dependent, and in general is not equal to $\theta_\true$. To apply Proposition~\ref{lem:rankpreservedindex} as before, we construct an ``environment-scaled pre-effort rank'' (denoted $\theta_\pre$).
\end{ReviewVerbatim}

% report-context-presentation-sha256: 90f9116e292dd71fddfd6b445ca0b060a481d026855bdb9725f2833c2cf62a31
\paragraph{Settled source reading}
\begin{sloppypar}
Equilibrium means that almost every applicant best responds against every feasible deviation;\allowbreak{} uniqueness is up to a null set.\allowbreak{}
\end{sloppypar}
\paragraph{Settled source reading}
\begin{sloppypar}
Applicants know the fixed priority used at score-\allowbreak{}boundary ties.\allowbreak{} This clarifies the stated model, not a corrected admission policy or a new economic assumption.\allowbreak{}
\end{sloppypar}

\paragraph{Lean-expanded paper semantic target}
\begin{ReviewVerbatim}
type:
(ℝ → ℝ) → (ℝ → ℝ) → (ℝ → ℝ) → ℝ → ℝ → (Bool × ℝ → ℝ) → ℕ → (ℕ → ℝ) → (ℕ → ℝ) → (Bool × ℝ → ℝ) → (Bool × ℝ → ℝ) → Prop

value:
fun (cost production skill : ℝ → ℝ) (psiA psiB : ℝ) (tie : Bool × ℝ → ℝ) (n : ℕ) (cutoff reward : ℕ → ℝ)
    (effort score : Bool × ℝ → ℝ) =>
  Measurable score ∧
    ∀ᵐ (x : Bool × ℝ) ∂LBG22StrategicRanking.sourceTwoGroupMeasure,
      LBG22StrategicRanking.SourcePopulationFiniteBestResponseAt LBG22StrategicRanking.sourceTwoGroupMeasure cost
        production (LBG22StrategicRanking.sourceEnvironmentSkill skill psiA psiB) score tie
        (fun (i : Fin (n + 1)) => cutoff ↑i) reward x (effort x)
\end{ReviewVerbatim}

\paragraph{Direct retained dependencies}
\begin{itemize}\raggedright
\item \hyperlink{governing-declaration-418baddf0ad9374c}{LBG22\allowbreak{}Strategic\allowbreak{}Ranking.\allowbreak{}Source\allowbreak{}Population\allowbreak{}Finite\allowbreak{}Best\allowbreak{}Response\allowbreak{}At}
\item \hyperlink{paper-prerequisite-73413973c7594479}{LBG22\allowbreak{}Strategic\allowbreak{}Ranking.\allowbreak{}source\allowbreak{}Environment\allowbreak{}Skill}
\item \hyperlink{paper-prerequisite-5b3be494cde30a37}{LBG22\allowbreak{}Strategic\allowbreak{}Ranking.\allowbreak{}source\allowbreak{}Two\allowbreak{}Group\allowbreak{}Measure}
\end{itemize}
\paragraph{Recorded source-to-declaration screening}
\textbf{Verdict:} \texttt{matches}
\\
{\raggedright
\textbf{Reason:} The target models the two equal-\allowbreak{}mass groups with common uniform skill rank, group multiplier psiA or psiB, measurable realized score, and one full-\allowbreak{}measure applicant set that best responds against every deviation.\allowbreak{} The displayed population, skill, finite-\allowbreak{}reward, counterfactual-\allowbreak{}rank, and fixed-\allowbreak{}priority support matches the group equilibrium source under the approved continuum and priority conventions.\allowbreak{}
\par}
\reviewmatch{Matches source input}
\noindent\textbf{Reviewer annotation}\par
\reviewerbox
\clearpage
\hypertarget{paper-prerequisite-d056dd7c8060653c}{}
\subsection*{\small\ttfamily\raggedright LBG22\allowbreak{}Strategic\allowbreak{}Ranking.\allowbreak{}Finite\allowbreak{}Admission\allowbreak{}Policy}
\noindent\textbf{Paper-local declaration:}\par
{\footnotesize\ttfamily\raggedright LBG22\allowbreak{}Strategic\allowbreak{}Ranking.\allowbreak{}Finite\allowbreak{}Admission\allowbreak{}Policy\par}
\textbf{Source locator:} {\footnotesize\raggedright cited publication, lines 1039-\allowbreak{}1044\par}
\paragraph{Verbatim paper-source connection}
\begin{ReviewVerbatim}
\parbold{Ranking designer (school)} A single \textit{school} is admitting applicants, based on their ranking. In particular, the school can choose a ranking reward function $\lambda : [0, 1] \mapsto [0, 1]$, such that an applicant with post-effort rank $\theta_\post$ is admitted with probability $\lambda(\theta_\post)$. We assume that $\lambda$ is non-decreasing and that the school has a constraint on the overall probability, such that in expectation it admits a number of applicants equal to {a} capacity constraint  $\rho \in (0,1)$, i.e., \lledit{$\E_{\theta_\post}[\lambda(\theta_\post)] = \rho $}.\footnote{\lledit{We use $\E_{\theta_\post}[\cdot]$ to denote an integral over $\theta_\post$ (and $\E[\cdot]$ to denote an integral over $\omega$)} with respect to the Lebesgue measure; informally, this can be thought of as averaging over the applicant population.} We may also refer to $\lambda$, informally, as the admission policy.



For simplicity, we further assume that $\lambda$ is a step-function with $K$ distinct levels $\ell_0 < \dots < \ell_{K-1}$, and $K-1$ cut-points parameterized by $c_1< \dots < c_{K-1}$ (with $c_0 = 0$, $c_K = 1$). %$\ell_0, \dots, \ell_{K-1}$.
In other words, we have $\lambda(\theta) = \ell_k$, for all $\theta \in \psi_k \triangleq [c_{k}, c_{k+1})$. Thus, applicants in the same post-rank interval $\theta \in \psi_k$ receive the same reward.
\end{ReviewVerbatim}

% report-context-presentation-sha256: 827a19afe0274b1faef13e344ddbcfe40d9841f2ff34b92ac9d53c0f83e29f56
\paragraph{Settled source reading}
\begin{sloppypar}
Applicants know the fixed priority used at score-\allowbreak{}boundary ties.\allowbreak{} This clarifies the stated model, not a corrected admission policy or a new economic assumption.\allowbreak{}
\end{sloppypar}

\paragraph{Lean-expanded paper semantic target}
\begin{ReviewVerbatim}
type:
ℕ → (ℕ → ℝ) → (ℕ → ℝ) → ℝ → Prop

value:
fun (n : ℕ) (cutoff reward : ℕ → ℝ) (rho : ℝ) =>
  StrictMonoOn cutoff (Set.Icc 0 (n + 1)) ∧
    cutoff 0 = 0 ∧
      cutoff (n + 1) = 1 ∧
        StrictMonoOn reward (Set.Icc 0 n) ∧
          0 ≤ reward 0 ∧
            reward n ≤ 1 ∧ rho ∈ Set.Ioo 0 1 ∧ ∑ i : Fin (n + 1), (cutoff (↑i + 1) - cutoff ↑i) * reward ↑i = rho
\end{ReviewVerbatim}


\paragraph{Recorded source-to-declaration screening}
\textbf{Verdict:} \texttt{matches}
\\
{\raggedright
\textbf{Reason:} The target requires strictly increasing cutoffs from 0 to 1, strictly increasing rewards in [0,1], capacity rho in (0,1), and the exact band-\allowbreak{}length-\allowbreak{}weighted capacity equation.\allowbreak{} These binders, domains, and formulas match the source finite step policy with K=n+1 reward levels.\allowbreak{}
\par}
\reviewmatch{Matches source input}
\noindent\textbf{Reviewer annotation}\par
\reviewerbox
\clearpage
\hypertarget{paper-prerequisite-6ed45e46d01eadeb}{}
\subsection*{\small\ttfamily\raggedright LBG22\allowbreak{}Strategic\allowbreak{}Ranking.\allowbreak{}counterfactual\allowbreak{}Tie\allowbreak{}Broken\allowbreak{}Rank}
\noindent\textbf{Paper-local declaration:}\par
{\footnotesize\ttfamily\raggedright LBG22\allowbreak{}Strategic\allowbreak{}Ranking.\allowbreak{}counterfactual\allowbreak{}Tie\allowbreak{}Broken\allowbreak{}Rank\par}
\textbf{Source locator:} {\footnotesize\raggedright cited publication, lines 1026-\allowbreak{}1030\par}
\paragraph{Verbatim paper-source connection}
\begin{ReviewVerbatim}
{\textbf{Tie Breaking.}
Given a choice of strategies $\theta_\pre \mapsto e(\theta_\pre)$, let $F$ be the CDF of the scores $v(e(\theta_\pre),\theta_\pre)$.
Since atoms for the distribution of $v$ would lead to ties for ranks defined as $F(v)$, we will use the labels of the applicants to break ties with the help of a (publicly announced) tie-breaking function $\Gamma(\omega)$, defined, e.g., via a collision free hash of the applicant names. Here  we require  that $\Gamma$  is a measurable function from $[0,1]$ to $[0,1]$ that maps different applicant labels to different values.
We then use $\Gamma$ to resolve the atoms of $F$, leading to a ranking function $v\mapsto \gamma(\omega, v)$ which is equal to $F(v)$ except when $v$ is an atom of the score distribution, {in which case it takes values in the ``gap interval'' $[F_-(v),  F(v)]$, where $F_-(v)$ is the left limit of $F$ at $v$.}  We  construct $\gamma$ in such a way that it
gives the uniform distribution for $\theta_\post(\omega)$ if we set  $\theta_\post(\omega)=\gamma(\omega,v(e(\theta_\pre(\omega)),\theta_\pre(\omega)))$.\footnote{See Remark~\ref{rem:gamma} for details on this construction.}

[Next verbatim source excerpt]

\begin{remark}[Construction of the $\gamma$ map]\label{rem:gamma}
	Formally,  if $v_0$ is an atom of the score distribution and $\Omega_{v_0}$ is the set of tied applicants with score $v_0$,
	then the discontinuity of $F$ {at  $v_0$} has height equal to the measure of $\Omega_{v_0}$, and so $\gamma$ for those applicants can be filled by using the CDF of $\Gamma(\omega)$ restricted to
	%the applicants $\omega$ who are tied
	$\Omega_{v_0}$; this gives a distribution for $\gamma(\omega, v_0)$ that is uniform over
	the gap interval $[ F_-(v_0),  F(v_0)]$ when restricted to
	%\omega\in
	$\Omega_{v_0}$, and hence leads to the claimed uniform distribution of $\theta_\post(\omega)$.
	Finally, we define $\gamma(\omega, v)$ for applicants  $\omega\notin\Omega_{v_0}$ by ``slotting them in'' in such a way that for a pair $\omega,\omega'$ with exactly one member in $\Omega_{v_0}$ and $\Gamma(\omega)<\Gamma(\omega')$ we have that
	$\gamma(v_0,\omega)\leq\gamma(v_0,\omega')$.
	Since the image of $\Omega_{v_0}$ is by construction dense in the gap interval, this uniquely determines $\gamma(v_0,\omega)$ for all applicants $\omega$.

	Note that for a given applicant $\omega$, the map $v\mapsto \gamma(\omega,v)$ is not necessarily 1-1; indeed, if $F$ is constant on an interval $I$, all $v\in I$ lead to the same rank.  But this only effects regions where the distribution of $v$ has no mass, and thus will not cause any issues.
\end{remark}
\end{ReviewVerbatim}



\paragraph{Lean-expanded paper semantic target}
\begin{ReviewVerbatim}
type:
{α : Type u_1} → [inst : MeasurableSpace α] → MeasureTheory.Measure α → (α → ℝ) → (α → ℝ) → α → ℝ → ℝ

value:
fun {α : Type u_1} [MeasurableSpace α] (μ : MeasureTheory.Measure α) (score tie : α → ℝ) (x : α) (newScore : ℝ) =>
  (μ {y : α | score y < newScore ∨ score y = newScore ∧ tie y ≤ tie x}).toReal
\end{ReviewVerbatim}


\paragraph{Recorded source-to-declaration screening}
\textbf{Verdict:} \texttt{matches}
\\
{\raggedright
\textbf{Reason:} The target changes only the selected applicant's proposed score while holding the population score function and public priority fixed, then measures lower scores plus equal-\allowbreak{}score applicants weakly preceding that priority.\allowbreak{} This is the source counterfactual gamma construction.\allowbreak{}
\par}
\reviewmatch{Matches source input}
\noindent\textbf{Reviewer annotation}\par
\reviewerbox
\clearpage
\hypertarget{paper-prerequisite-b7c71e53177b76c6}{}
\subsection*{\small\ttfamily\raggedright LBG22\allowbreak{}Strategic\allowbreak{}Ranking.\allowbreak{}tie\allowbreak{}Broken\allowbreak{}Rank}
\noindent\textbf{Paper-local declaration:}\par
{\footnotesize\ttfamily\raggedright LBG22\allowbreak{}Strategic\allowbreak{}Ranking.\allowbreak{}tie\allowbreak{}Broken\allowbreak{}Rank\par}
\textbf{Source locator:} {\footnotesize\raggedright cited publication, lines 1026-\allowbreak{}1030\par}
\paragraph{Verbatim paper-source connection}
\begin{ReviewVerbatim}
{\textbf{Tie Breaking.}
Given a choice of strategies $\theta_\pre \mapsto e(\theta_\pre)$, let $F$ be the CDF of the scores $v(e(\theta_\pre),\theta_\pre)$.
Since atoms for the distribution of $v$ would lead to ties for ranks defined as $F(v)$, we will use the labels of the applicants to break ties with the help of a (publicly announced) tie-breaking function $\Gamma(\omega)$, defined, e.g., via a collision free hash of the applicant names. Here  we require  that $\Gamma$  is a measurable function from $[0,1]$ to $[0,1]$ that maps different applicant labels to different values.
We then use $\Gamma$ to resolve the atoms of $F$, leading to a ranking function $v\mapsto \gamma(\omega, v)$ which is equal to $F(v)$ except when $v$ is an atom of the score distribution, {in which case it takes values in the ``gap interval'' $[F_-(v),  F(v)]$, where $F_-(v)$ is the left limit of $F$ at $v$.}  We  construct $\gamma$ in such a way that it
gives the uniform distribution for $\theta_\post(\omega)$ if we set  $\theta_\post(\omega)=\gamma(\omega,v(e(\theta_\pre(\omega)),\theta_\pre(\omega)))$.\footnote{See Remark~\ref{rem:gamma} for details on this construction.}

[Next verbatim source excerpt]

\begin{remark}[Construction of the $\gamma$ map]\label{rem:gamma}
	Formally,  if $v_0$ is an atom of the score distribution and $\Omega_{v_0}$ is the set of tied applicants with score $v_0$,
	then the discontinuity of $F$ {at  $v_0$} has height equal to the measure of $\Omega_{v_0}$, and so $\gamma$ for those applicants can be filled by using the CDF of $\Gamma(\omega)$ restricted to
	%the applicants $\omega$ who are tied
	$\Omega_{v_0}$; this gives a distribution for $\gamma(\omega, v_0)$ that is uniform over
	the gap interval $[ F_-(v_0),  F(v_0)]$ when restricted to
	%\omega\in
	$\Omega_{v_0}$, and hence leads to the claimed uniform distribution of $\theta_\post(\omega)$.
	Finally, we define $\gamma(\omega, v)$ for applicants  $\omega\notin\Omega_{v_0}$ by ``slotting them in'' in such a way that for a pair $\omega,\omega'$ with exactly one member in $\Omega_{v_0}$ and $\Gamma(\omega)<\Gamma(\omega')$ we have that
	$\gamma(v_0,\omega)\leq\gamma(v_0,\omega')$.
	Since the image of $\Omega_{v_0}$ is by construction dense in the gap interval, this uniquely determines $\gamma(v_0,\omega)$ for all applicants $\omega$.

	Note that for a given applicant $\omega$, the map $v\mapsto \gamma(\omega,v)$ is not necessarily 1-1; indeed, if $F$ is constant on an interval $I$, all $v\in I$ lead to the same rank.  But this only effects regions where the distribution of $v$ has no mass, and thus will not cause any issues.
\end{remark}
\end{ReviewVerbatim}



\paragraph{Lean-expanded paper semantic target}
\begin{ReviewVerbatim}
type:
{α : Type u_1} → [inst : MeasurableSpace α] → MeasureTheory.Measure α → (α → ℝ) → (α → ℝ) → α → ℝ

value:
fun {α : Type u_1} [MeasurableSpace α] (μ : MeasureTheory.Measure α) (score tie : α → ℝ) (x : α) =>
  (μ (LBG22StrategicRanking.tieBrokenLowerContour score tie x)).toReal
\end{ReviewVerbatim}

\paragraph{Direct retained dependencies}
\begin{itemize}\raggedright
\item \hyperlink{governing-declaration-2190069be417355f}{LBG22\allowbreak{}Strategic\allowbreak{}Ranking.\allowbreak{}tie\allowbreak{}Broken\allowbreak{}Lower\allowbreak{}Contour}
\end{itemize}
\paragraph{Recorded source-to-declaration screening}
\textbf{Verdict:} \texttt{matches}
\\
{\raggedright
\textbf{Reason:} The target measures the lower score contour together with equal-\allowbreak{}score applicants whose public priority is weakly below the selected applicant's.\allowbreak{} This is the source tie-\allowbreak{}broken post-\allowbreak{}effort rank gamma;\allowbreak{} the displayed support and approved record fix the same priority across actual and counterfactual scores.\allowbreak{}
\par}
\reviewmatch{Matches source input}
\noindent\textbf{Reviewer annotation}\par
\reviewerbox
\clearpage
\hypertarget{paper-prerequisite-59d181ae71b7b024}{}
\subsection*{\small\ttfamily\raggedright LBG22\allowbreak{}Strategic\allowbreak{}Ranking.\allowbreak{}Source\allowbreak{}Multidimensional\allowbreak{}Best\allowbreak{}Response\allowbreak{}At}
\noindent\textbf{Paper-local declaration:}\par
{\footnotesize\ttfamily\raggedright LBG22\allowbreak{}Strategic\allowbreak{}Ranking.\allowbreak{}Source\allowbreak{}Multidimensional\allowbreak{}Best\allowbreak{}Response\allowbreak{}At\par}
\textbf{Source locator:} {\footnotesize\raggedright cited publication, lines 1760-\allowbreak{}1764\par}
\paragraph{Verbatim paper-source connection}
\begin{ReviewVerbatim}
We show this result by characterizing the equilibrium of the following simplified setting, with further assumptions on the effort cost function $\cost$ and the effort transfer function $g$; the cost function $\cost$ is assumed to be a function of the sum of the efforts exerted:
\[p(e^1, \cdots, e^m) \triangleq p\left(\sum_{i=1}^m \e{i}\right),\] where $p$ is convex and increasing.
This assumption says that the effort exerted for any skill is entirely exchangeable, for example, two hours spent on studying math is as costly as two hours spent on studying chemistry.
%
Effort transfer function $g$ is assumed to be linear, that is, there are constant returns to effort.

[Next verbatim source excerpt]

\parbold{Model.} The model is similar to our base model \Cref{sec:model}; each applicant now has $m$ latent skill levels with respective ranks $\theta_\pre^i \in [0,1]$ for $i\in \{1, \cdots, m\}$. Each rank is drawn independently from the uniform distribution over $[0,1]$, and the skill $i$ of applicant with rank $\theta_\pre^i$ is $\f{i}(\theta_\pre^i)$. Each applicant now chooses effort levels $\{ \e{i} \}$, at cost $p^m( \e{1}, \cdots, \e{m})$, resulting in post-effort scores $\{\score{i}\}$,
\begin{equation*}
\score{i} = g(\e{i})\cdot \f{i}(\theta_\pre^i),
\end{equation*}
where $g$ is concave, increasing, as before, and $\f{i}$ is a continuous, strictly increasing quantile function (${\f{i}}^{-1}$ is the CDF function of the scores on dimension $i$). As before, the school observes the post-effort scores for each applicant, now for each dimension $i$, and designs a non-decreasing function $\lambda : [0,1] \to [0,1]$ denoting the admissions probability $\lambda(\theta_\post)$ for applicant with post-effort rank $\theta_\post$.

%Suppose the school observes $v_i$.

How does the school construct post-effort rank $\theta_\post$? In general, each applicant's \textit{combined} post-effort score may be any function of the scores on each dimension, $\{\score{i}\}$. Here, we assume the following linear score function. The school announces weights $\alpha = (\alpha_1, \cdots, \alpha_m) \in \Delta_{m-1}$ (denoting the $m$-dimensional simplex); the weights $\alpha$ represent the relative emphasis placed on each skill for the admissions decision. Then, each applicant is ranked according to their combined score, $\combscore := \sum_{i=1}^m \alpha_i\score{i}$, resulting in their combined post-effort rank $\theta_\post^\alpha$, and is admitted with probability $\lambda(\theta_\post^\alpha)$.\footnote{The linear combination of skill is similar to the ``linear mechanism'' in \citet{kleinberg18investeffort}, which studied how reward design incentivizes strategic agents to exert effort on different skill dimensions. Compared to their work, where efforts are connected to skills via an effort graph, we consider a simplified setting where each effort maps to one skill, and study how reward design affects equilibrium rankings, in presence of competition.}

%\lnote{\textbf{weighted score or weighted rank?} alternatively, the school combines the post-effort ranks for each skill. i.e. an applicant with \emph{weighted} post-effort rank $\theta_\post^\alpha := \sum_{i=1}^m \alpha_i\theta_\post^i$ is admitted with probability $\lambda(\theta_\post^\alpha)$. This case seems interesting as well. Not clear in which case it's possible to reduce to the single-dimensional case?}

Putting things together, the applicant's individual welfare is:
\begin{equation*}
\iwelf(\{ \e{i} \}_{i=1}^m, \lambda(\theta_\post^\alpha)) = \lambda(\theta_\post^\alpha) - p^m( \e{1}, \cdots, \e{m})
\end{equation*}

[Next verbatim source excerpt]

\parbold{Ranking designer (school)} A single \textit{school} is admitting applicants, based on their ranking. In particular, the school can choose a ranking reward function $\lambda : [0, 1] \mapsto [0, 1]$, such that an applicant with post-effort rank $\theta_\post$ is admitted with probability $\lambda(\theta_\post)$. We assume that $\lambda$ is non-decreasing and that the school has a constraint on the overall probability, such that in expectation it admits a number of applicants equal to {a} capacity constraint  $\rho \in (0,1)$, i.e., \lledit{$\E_{\theta_\post}[\lambda(\theta_\post)] = \rho $}.\footnote{\lledit{We use $\E_{\theta_\post}[\cdot]$ to denote an integral over $\theta_\post$ (and $\E[\cdot]$ to denote an integral over $\omega$)} with respect to the Lebesgue measure; informally, this can be thought of as averaging over the applicant population.} We may also refer to $\lambda$, informally, as the admission policy.



For simplicity, we further assume that $\lambda$ is a step-function with $K$ distinct levels $\ell_0 < \dots < \ell_{K-1}$, and $K-1$ cut-points parameterized by $c_1< \dots < c_{K-1}$ (with $c_0 = 0$, $c_K = 1$). %$\ell_0, \dots, \ell_{K-1}$.
In other words, we have $\lambda(\theta) = \ell_k$, for all $\theta \in \psi_k \triangleq [c_{k}, c_{k+1})$. Thus, applicants in the same post-rank interval $\theta \in \psi_k$ receive the same reward.

[Next verbatim source excerpt]

\begin{definition} [Equilibrium]\label{def:equi}
{Given a}
%Let the
tie-breaking function $\Gamma(\omega)$
%be a pre-announced ordering over index $\omega$, to break ties in post-effort value. Then, let $\gamma(\omega, v)$ be the 1-1 ranking function defined as the CDF of post-effort values $v$ except where there are ties in $v$, in which case the ordering $\Gamma$ is used to produce unique ranks such that the range of $\gamma$ is uniform on $[0, 1]$.\footnote{More formally, note that each discontinuity in the CDF of post-effort values $v$ has a height equal to the measure of the applicants who are tied at that value, and so $\gamma$ for those applicants can be filled in using the CDF of $\Gamma(\omega)$ restricted to the applicants $\omega$ who are tied. Thus, the range of $\gamma$ is uniform on $[0,1]$.
%
{and}
%	Then, given
	a ranking probability function $\lambda$, an \textit{equilibrium} is a set of effort levels and post-effort ranking rewards for each applicant, $\{e(\theta_\pre),  \lambda(\theta_\post(\theta_\pre))\}$ such that, for all
	$\omega$, %$\theta_\pre$,
	\begin{align*}
		e(\theta_\pre(\omega)) &\in \argmax_{e} \iwelf\left(e, \lambda\left(\gamma\left(\omega, v\left(e,\theta_\pre(\omega)\right)\right)\right)\right) %\triangleq \argmax_{e} \left[ \lambda(\gamma(v(e,\theta_\pre)) - \cost(e) \right].
		\\
		\theta_\post(\theta_\pre(\omega)) &= \gamma(\omega, v(e(\theta_\pre(\omega)),\theta_\pre(\omega))),
	\end{align*}
	{where $\gamma(\omega,v)$ is the ranking induced by the CDF of the scores resulting from effort levels $\{e(\theta_\pre)\}$ and the tie-breaking function $\Gamma$.\footnote{See Remark~\ref{rem:gamma-effort} on $\gamma$ and the set of efforts.%Formally, $\gamma$ depends on the set of efforts. Given a fixed set and $\gamma$, for each applicant $\omega$ the first condition considers the counter-factual ranking of $\omega$ with different post-effort values but using the same ranking function. As defined, $\gamma$ yields a uniform distribution of ranks with such measure $0$ changes. In Appendix \Cref{lem:deviationsmeasure0}, we further prove that two effort sets equal up to sets of measure $0$ induce the same ranking function $\gamma$, and so the condition is consistent.%, and so we set it at the $\gamma$ induced by the effort levels $\{e(\theta_\pre)\}$.%, and so it does not vary.
	}
	}
	%In particular, we denote $\theta_\post(\theta_\pre) := \gamma(v(e(\theta_\pre),\theta_\pre))$.
	%\lledit{and for any $\theta_\pre, \theta'_\pre$, we have $\lambda(\theta_\post(\theta_\pre)) > \lambda(\theta_\post(\theta'_\pre)) \implies v(e(\theta_\pre),\theta_\pre) > v(e(\theta_\pre'),\theta_\pre')$.}
\end{definition}
\end{ReviewVerbatim}

% report-context-presentation-sha256: 90f9116e292dd71fddfd6b445ca0b060a481d026855bdb9725f2833c2cf62a31
\paragraph{Settled source reading}
\begin{sloppypar}
Equilibrium means that almost every applicant best responds against every feasible deviation;\allowbreak{} uniqueness is up to a null set.\allowbreak{}
\end{sloppypar}
\paragraph{Settled source reading}
\begin{sloppypar}
Applicants know the fixed priority used at score-\allowbreak{}boundary ties.\allowbreak{} This clarifies the stated model, not a corrected admission policy or a new economic assumption.\allowbreak{}
\end{sloppypar}

\paragraph{Lean-expanded paper semantic target}
\begin{ReviewVerbatim}
type:
{α : Type u_1} →
  {ι : Type u_2} →
    [inst : MeasurableSpace α] →
      [Fintype ι] →
        (μ : MeasureTheory.Measure α) →
          [MeasureTheory.IsFiniteMeasure μ] →
            (ℝ → ℝ) → (α → ι → ℝ) → (α → ℝ) → (α → ℝ) → ℝ → (ℝ → ℝ) → α → (ι → ℝ) → Prop

value:
fun {α : Type u_1} {ι : Type u_2} [MeasurableSpace α] [Fintype ι] (μ : MeasureTheory.Measure α)
    [MeasureTheory.IsFiniteMeasure μ] (cost : ℝ → ℝ) (coefficient : α → ι → ℝ) (score tie : α → ℝ) (h : ℝ)
    (reward : ℝ → ℝ) (x : α) (effort : ι → ℝ) =>
  (∀ (i : ι), 0 ≤ effort i) ∧
    score x = h * ∑ i : ι, effort i * coefficient x i ∧
      ∀ (d : ι → ℝ),
        (∀ (i : ι), 0 ≤ d i) →
          reward
                (LBG22StrategicRanking.counterfactualTieBrokenRank μ score tie x (h * ∑ i : ι, d i * coefficient x i)) -
              cost (∑ i : ι, d i) ≤
            reward (LBG22StrategicRanking.tieBrokenRank μ score tie x) - cost (∑ i : ι, effort i)
\end{ReviewVerbatim}

\paragraph{Direct retained dependencies}
\begin{itemize}\raggedright
\item \hyperlink{paper-prerequisite-6ed45e46d01eadeb}{LBG22\allowbreak{}Strategic\allowbreak{}Ranking.\allowbreak{}counterfactual\allowbreak{}Tie\allowbreak{}Broken\allowbreak{}Rank}
\item \hyperlink{paper-prerequisite-b7c71e53177b76c6}{LBG22\allowbreak{}Strategic\allowbreak{}Ranking.\allowbreak{}tie\allowbreak{}Broken\allowbreak{}Rank}
\end{itemize}
\paragraph{Recorded source-to-declaration screening}
\textbf{Verdict:} \texttt{matches}
\\
{\raggedright
\textbf{Reason:} The target specializes multidimensional welfare to linear production h, coefficient-\allowbreak{}weighted scores, and exchangeable cost of total effort, and maximizes over every nonnegative vector deviation under the fixed ranking law.\allowbreak{} This matches the source exchangeable-\allowbreak{}cost model.\allowbreak{}
\par}
\reviewmatch{Matches source input}
\noindent\textbf{Reviewer annotation}\par
\reviewerbox
\clearpage
\hypertarget{paper-prerequisite-a42f440ddc254fa7}{}
\subsection*{\small\ttfamily\raggedright LBG22\allowbreak{}Strategic\allowbreak{}Ranking.\allowbreak{}finite\allowbreak{}Lower\allowbreak{}Rank\allowbreak{}Band}
\noindent\textbf{Paper-local declaration:}\par
{\footnotesize\ttfamily\raggedright LBG22\allowbreak{}Strategic\allowbreak{}Ranking.\allowbreak{}finite\allowbreak{}Lower\allowbreak{}Rank\allowbreak{}Band\par}
\textbf{Source locator:} {\footnotesize\raggedright cited publication, lines 1039-\allowbreak{}1044\par}
\paragraph{Verbatim paper-source connection}
\begin{ReviewVerbatim}
\parbold{Ranking designer (school)} A single \textit{school} is admitting applicants, based on their ranking. In particular, the school can choose a ranking reward function $\lambda : [0, 1] \mapsto [0, 1]$, such that an applicant with post-effort rank $\theta_\post$ is admitted with probability $\lambda(\theta_\post)$. We assume that $\lambda$ is non-decreasing and that the school has a constraint on the overall probability, such that in expectation it admits a number of applicants equal to {a} capacity constraint  $\rho \in (0,1)$, i.e., \lledit{$\E_{\theta_\post}[\lambda(\theta_\post)] = \rho $}.\footnote{\lledit{We use $\E_{\theta_\post}[\cdot]$ to denote an integral over $\theta_\post$ (and $\E[\cdot]$ to denote an integral over $\omega$)} with respect to the Lebesgue measure; informally, this can be thought of as averaging over the applicant population.} We may also refer to $\lambda$, informally, as the admission policy.



For simplicity, we further assume that $\lambda$ is a step-function with $K$ distinct levels $\ell_0 < \dots < \ell_{K-1}$, and $K-1$ cut-points parameterized by $c_1< \dots < c_{K-1}$ (with $c_0 = 0$, $c_K = 1$). %$\ell_0, \dots, \ell_{K-1}$.
In other words, we have $\lambda(\theta) = \ell_k$, for all $\theta \in \psi_k \triangleq [c_{k}, c_{k+1})$. Thus, applicants in the same post-rank interval $\theta \in \psi_k$ receive the same reward.
\end{ReviewVerbatim}

% report-context-presentation-sha256: 827a19afe0274b1faef13e344ddbcfe40d9841f2ff34b92ac9d53c0f83e29f56
\paragraph{Settled source reading}
\begin{sloppypar}
Applicants know the fixed priority used at score-\allowbreak{}boundary ties.\allowbreak{} This clarifies the stated model, not a corrected admission policy or a new economic assumption.\allowbreak{}
\end{sloppypar}

\paragraph{Lean-expanded paper semantic target}
\begin{ReviewVerbatim}
type:
{n : ℕ} → (Fin (n + 1) → ℝ) → ℝ → Fin (n + 1)

value:
fun {n : ℕ} (cutoff : Fin (n + 1) → ℝ) (rank : ℝ) =>
  Finset.univ.sup fun (i : Fin (n + 1)) => if cutoff i < rank then i else 0
\end{ReviewVerbatim}


\paragraph{Recorded source-to-declaration screening}
\textbf{Verdict:} \texttt{matches}
\\
{\raggedright
\textbf{Reason:} The target selects the greatest finite cutoff strictly below the strict-\allowbreak{}percentile rank.\allowbreak{} Under the identity-\allowbreak{}bound approved priority record, this is the permitted strict-\allowbreak{}percentile band representation used only through all-\allowbreak{}deviation equivalence to the source's actual score-\allowbreak{}priority boundary;\allowbreak{} all displayed uses respect that scope.\allowbreak{}
\par}
\reviewmatch{Matches source input}
\noindent\textbf{Reviewer annotation}\par
\reviewerbox
\clearpage
\hypertarget{paper-prerequisite-e57e55071c75ce91}{}
\subsection*{\small\ttfamily\raggedright LBG22\allowbreak{}Strategic\allowbreak{}Ranking.\allowbreak{}source\allowbreak{}Multidimensional\allowbreak{}Coefficient}
\noindent\textbf{Paper-local declaration:}\par
{\footnotesize\ttfamily\raggedright LBG22\allowbreak{}Strategic\allowbreak{}Ranking.\allowbreak{}source\allowbreak{}Multidimensional\allowbreak{}Coefficient\par}
\textbf{Source locator:} {\footnotesize\raggedright cited publication, lines 1724-\allowbreak{}1739\par}
\paragraph{Verbatim paper-source connection}
\begin{ReviewVerbatim}
\parbold{Model.} The model is similar to our base model \Cref{sec:model}; each applicant now has $m$ latent skill levels with respective ranks $\theta_\pre^i \in [0,1]$ for $i\in \{1, \cdots, m\}$. Each rank is drawn independently from the uniform distribution over $[0,1]$, and the skill $i$ of applicant with rank $\theta_\pre^i$ is $\f{i}(\theta_\pre^i)$. Each applicant now chooses effort levels $\{ \e{i} \}$, at cost $p^m( \e{1}, \cdots, \e{m})$, resulting in post-effort scores $\{\score{i}\}$,
\begin{equation*}
\score{i} = g(\e{i})\cdot \f{i}(\theta_\pre^i),
\end{equation*}
where $g$ is concave, increasing, as before, and $\f{i}$ is a continuous, strictly increasing quantile function (${\f{i}}^{-1}$ is the CDF function of the scores on dimension $i$). As before, the school observes the post-effort scores for each applicant, now for each dimension $i$, and designs a non-decreasing function $\lambda : [0,1] \to [0,1]$ denoting the admissions probability $\lambda(\theta_\post)$ for applicant with post-effort rank $\theta_\post$.

%Suppose the school observes $v_i$.

How does the school construct post-effort rank $\theta_\post$? In general, each applicant's \textit{combined} post-effort score may be any function of the scores on each dimension, $\{\score{i}\}$. Here, we assume the following linear score function. The school announces weights $\alpha = (\alpha_1, \cdots, \alpha_m) \in \Delta_{m-1}$ (denoting the $m$-dimensional simplex); the weights $\alpha$ represent the relative emphasis placed on each skill for the admissions decision. Then, each applicant is ranked according to their combined score, $\combscore := \sum_{i=1}^m \alpha_i\score{i}$, resulting in their combined post-effort rank $\theta_\post^\alpha$, and is admitted with probability $\lambda(\theta_\post^\alpha)$.\footnote{The linear combination of skill is similar to the ``linear mechanism'' in \citet{kleinberg18investeffort}, which studied how reward design incentivizes strategic agents to exert effort on different skill dimensions. Compared to their work, where efforts are connected to skills via an effort graph, we consider a simplified setting where each effort maps to one skill, and study how reward design affects equilibrium rankings, in presence of competition.}

%\lnote{\textbf{weighted score or weighted rank?} alternatively, the school combines the post-effort ranks for each skill. i.e. an applicant with \emph{weighted} post-effort rank $\theta_\post^\alpha := \sum_{i=1}^m \alpha_i\theta_\post^i$ is admitted with probability $\lambda(\theta_\post^\alpha)$. This case seems interesting as well. Not clear in which case it's possible to reduce to the single-dimensional case?}

Putting things together, the applicant's individual welfare is:
\begin{equation*}
\iwelf(\{ \e{i} \}_{i=1}^m, \lambda(\theta_\post^\alpha)) = \lambda(\theta_\post^\alpha) - p^m( \e{1}, \cdots, \e{m})
\end{equation*}

[Next verbatim source excerpt]

\begin{definition} [Equilibrium]\label{def:equi}
{Given a}
%Let the
tie-breaking function $\Gamma(\omega)$
%be a pre-announced ordering over index $\omega$, to break ties in post-effort value. Then, let $\gamma(\omega, v)$ be the 1-1 ranking function defined as the CDF of post-effort values $v$ except where there are ties in $v$, in which case the ordering $\Gamma$ is used to produce unique ranks such that the range of $\gamma$ is uniform on $[0, 1]$.\footnote{More formally, note that each discontinuity in the CDF of post-effort values $v$ has a height equal to the measure of the applicants who are tied at that value, and so $\gamma$ for those applicants can be filled in using the CDF of $\Gamma(\omega)$ restricted to the applicants $\omega$ who are tied. Thus, the range of $\gamma$ is uniform on $[0,1]$.
%
{and}
%	Then, given
	a ranking probability function $\lambda$, an \textit{equilibrium} is a set of effort levels and post-effort ranking rewards for each applicant, $\{e(\theta_\pre),  \lambda(\theta_\post(\theta_\pre))\}$ such that, for all
	$\omega$, %$\theta_\pre$,
	\begin{align*}
		e(\theta_\pre(\omega)) &\in \argmax_{e} \iwelf\left(e, \lambda\left(\gamma\left(\omega, v\left(e,\theta_\pre(\omega)\right)\right)\right)\right) %\triangleq \argmax_{e} \left[ \lambda(\gamma(v(e,\theta_\pre)) - \cost(e) \right].
		\\
		\theta_\post(\theta_\pre(\omega)) &= \gamma(\omega, v(e(\theta_\pre(\omega)),\theta_\pre(\omega))),
	\end{align*}
	{where $\gamma(\omega,v)$ is the ranking induced by the CDF of the scores resulting from effort levels $\{e(\theta_\pre)\}$ and the tie-breaking function $\Gamma$.\footnote{See Remark~\ref{rem:gamma-effort} on $\gamma$ and the set of efforts.%Formally, $\gamma$ depends on the set of efforts. Given a fixed set and $\gamma$, for each applicant $\omega$ the first condition considers the counter-factual ranking of $\omega$ with different post-effort values but using the same ranking function. As defined, $\gamma$ yields a uniform distribution of ranks with such measure $0$ changes. In Appendix \Cref{lem:deviationsmeasure0}, we further prove that two effort sets equal up to sets of measure $0$ induce the same ranking function $\gamma$, and so the condition is consistent.%, and so we set it at the $\gamma$ induced by the effort levels $\{e(\theta_\pre)\}$.%, and so it does not vary.
	}
	}
	%In particular, we denote $\theta_\post(\theta_\pre) := \gamma(v(e(\theta_\pre),\theta_\pre))$.
	%\lledit{and for any $\theta_\pre, \theta'_\pre$, we have $\lambda(\theta_\post(\theta_\pre)) > \lambda(\theta_\post(\theta'_\pre)) \implies v(e(\theta_\pre),\theta_\pre) > v(e(\theta_\pre'),\theta_\pre')$.}
\end{definition}

[Next verbatim source excerpt]

\parbold{Applicants}
There is a unit mass of \textit{applicants}, indexed by an observed index $\omega \in [0,1]$ distributed uniformly.\footnote{The index $\omega$ should be interpreted as each applicant's ``name,'' uncorrelated with skill, used solely for tie-breaking.}  Each applicant has a latent (unobserved) skill level represented by some measurable function of $\omega$.
We assume that {the distribution of the skills} has no atoms and that the CDF of this distribution is strictly increasing.
Using the CDF to map the skill of an applicant to a rank in~$[0,1]$, each applicant gets an (unobserved) rank $\theta_\pre =\theta_\pre(\omega)$ which by our assumption on the CDF is again uniformly distributed in $[0, 1]$ (the higher the rank the better).  With this setup, the skill of an applicant with rank $\theta_\pre$ can be written as $f(\theta_\pre)$ where $f$ is a strictly increasing, continuous function.  {It will be notationally convenient to label applicants by their rank $\theta_\pre$, though the reader should note that in contrast to $\omega$, $\theta_\pre(\omega)$ is assumed to be unobservable.}

%3. Each applicant exerts some effort e, at cost \cost(e).


Each applicant chooses an \textit{effort} level $e \ge 0$, the result of which is an observed, post-effort \textit{score}, {$v=v(e,\theta_\pre) = {g(e)}\cdot{f(\theta_\pre)}$}. \lledit{In other words, we assume the post-effort score is the product of two components, associated with the effort and the pre-effort skill respectively.}
{We assume} that the effort transfer function $g: [0, \infty) \mapsto [0, \infty)$ is a continuous, concave, strictly increasing function, representing that marginal effort improves one's score but has diminishing returns.
{%Since applicants can be labelled by their skill level $f(\theta_\pre)$, or equivalently, by their rank $\theta_{\pre}$,
{T}he strategies of the applicants can {then} be described by a function $\theta_\pre \mapsto e(\theta_\pre)$.
%\cb{$\omega\mapsto e(\omega)$.}%, leading to a distribution over post-effort scores $v=v(\theta_\pre))$ via the uniform distribution of $\theta_\pre$ in $[0,1]$.
}
  Each {applicant} is then ranked according to their score $v$%(where we break ties uniformly at random, if they appear)
, resulting in a \textit{post-effort rank} $\theta_\post$.
Note that the ranking $\theta_\post$ is slightly less trivial than a ranking of the skills, since the scores might have ties, which have to be resolved.

{\textbf{Tie Breaking.}
Given a choice of strategies $\theta_\pre \mapsto e(\theta_\pre)$, let $F$ be the CDF of the scores $v(e(\theta_\pre),\theta_\pre)$.
Since atoms for the distribution of $v$ would lead to ties for ranks defined as $F(v)$, we will use the labels of the applicants to break ties with the help of a (publicly announced) tie-breaking function $\Gamma(\omega)$, defined, e.g., via a collision free hash of the applicant names. Here  we require  that $\Gamma$  is a measurable function from $[0,1]$ to $[0,1]$ that maps different applicant labels to different values.
We then use $\Gamma$ to resolve the atoms of $F$, leading to a ranking function $v\mapsto \gamma(\omega, v)$ which is equal to $F(v)$ except when $v$ is an atom of the score distribution, {in which case it takes values in the ``gap interval'' $[F_-(v),  F(v)]$, where $F_-(v)$ is the left limit of $F$ at $v$.}  We  construct $\gamma$ in such a way that it
gives the uniform distribution for $\theta_\post(\omega)$ if we set  $\theta_\post(\omega)=\gamma(\omega,v(e(\theta_\pre(\omega)),\theta_\pre(\omega)))$.\footnote{See Remark~\ref{rem:gamma} for details on this construction.}
\end{ReviewVerbatim}

% report-context-presentation-sha256: 90f9116e292dd71fddfd6b445ca0b060a481d026855bdb9725f2833c2cf62a31
\paragraph{Settled source reading}
\begin{sloppypar}
Equilibrium means that almost every applicant best responds against every feasible deviation;\allowbreak{} uniqueness is up to a null set.\allowbreak{}
\end{sloppypar}
\paragraph{Settled source reading}
\begin{sloppypar}
Applicants know the fixed priority used at score-\allowbreak{}boundary ties.\allowbreak{} This clarifies the stated model, not a corrected admission policy or a new economic assumption.\allowbreak{}
\end{sloppypar}

\paragraph{Lean-expanded paper semantic target}
\begin{ReviewVerbatim}
type:
{ι : Type u_1} → (ι → ℝ) → (ι → ℝ → ℝ) → (ι → ℝ) → ι → ℝ

value:
fun {ι : Type u_1} (weight : ι → ℝ) (skill : ι → ℝ → ℝ) (rank : ι → ℝ) (i : ι) =>
  weight i * LBG22StrategicRanking.sourceClampedSkill (skill i) (rank i)
\end{ReviewVerbatim}

\paragraph{Direct retained dependencies}
\begin{itemize}\raggedright
\item \hyperlink{governing-declaration-8658f19612e370c1}{LBG22\allowbreak{}Strategic\allowbreak{}Ranking.\allowbreak{}source\allowbreak{}Clamped\allowbreak{}Skill}
\end{itemize}
\paragraph{Recorded source-to-declaration screening}
\textbf{Verdict:} \texttt{matches}
\\
{\raggedright
\textbf{Reason:} The target defines coordinate i's pre-\allowbreak{}effort score coefficient as alpha\_\allowbreak{}i times f\_\allowbreak{}i(theta\_\allowbreak{}i), with ranks restricted through the displayed clamped skill.\allowbreak{} This is the source weighted multidimensional coefficient.\allowbreak{}
\par}
\reviewmatch{Matches source input}
\noindent\textbf{Reviewer annotation}\par
\reviewerbox
\clearpage
\hypertarget{paper-prerequisite-16d135799a0f0491}{}
\subsection*{\small\ttfamily\raggedright LBG22\allowbreak{}Strategic\allowbreak{}Ranking.\allowbreak{}source\allowbreak{}Multidimensional\allowbreak{}Population}
\noindent\textbf{Paper-local declaration:}\par
{\footnotesize\ttfamily\raggedright LBG22\allowbreak{}Strategic\allowbreak{}Ranking.\allowbreak{}source\allowbreak{}Multidimensional\allowbreak{}Population\par}
\textbf{Source locator:} {\footnotesize\raggedright cited publication, lines 1724-\allowbreak{}1739\par}
\paragraph{Verbatim paper-source connection}
\begin{ReviewVerbatim}
\parbold{Model.} The model is similar to our base model \Cref{sec:model}; each applicant now has $m$ latent skill levels with respective ranks $\theta_\pre^i \in [0,1]$ for $i\in \{1, \cdots, m\}$. Each rank is drawn independently from the uniform distribution over $[0,1]$, and the skill $i$ of applicant with rank $\theta_\pre^i$ is $\f{i}(\theta_\pre^i)$. Each applicant now chooses effort levels $\{ \e{i} \}$, at cost $p^m( \e{1}, \cdots, \e{m})$, resulting in post-effort scores $\{\score{i}\}$,
\begin{equation*}
\score{i} = g(\e{i})\cdot \f{i}(\theta_\pre^i),
\end{equation*}
where $g$ is concave, increasing, as before, and $\f{i}$ is a continuous, strictly increasing quantile function (${\f{i}}^{-1}$ is the CDF function of the scores on dimension $i$). As before, the school observes the post-effort scores for each applicant, now for each dimension $i$, and designs a non-decreasing function $\lambda : [0,1] \to [0,1]$ denoting the admissions probability $\lambda(\theta_\post)$ for applicant with post-effort rank $\theta_\post$.

%Suppose the school observes $v_i$.

How does the school construct post-effort rank $\theta_\post$? In general, each applicant's \textit{combined} post-effort score may be any function of the scores on each dimension, $\{\score{i}\}$. Here, we assume the following linear score function. The school announces weights $\alpha = (\alpha_1, \cdots, \alpha_m) \in \Delta_{m-1}$ (denoting the $m$-dimensional simplex); the weights $\alpha$ represent the relative emphasis placed on each skill for the admissions decision. Then, each applicant is ranked according to their combined score, $\combscore := \sum_{i=1}^m \alpha_i\score{i}$, resulting in their combined post-effort rank $\theta_\post^\alpha$, and is admitted with probability $\lambda(\theta_\post^\alpha)$.\footnote{The linear combination of skill is similar to the ``linear mechanism'' in \citet{kleinberg18investeffort}, which studied how reward design incentivizes strategic agents to exert effort on different skill dimensions. Compared to their work, where efforts are connected to skills via an effort graph, we consider a simplified setting where each effort maps to one skill, and study how reward design affects equilibrium rankings, in presence of competition.}

%\lnote{\textbf{weighted score or weighted rank?} alternatively, the school combines the post-effort ranks for each skill. i.e. an applicant with \emph{weighted} post-effort rank $\theta_\post^\alpha := \sum_{i=1}^m \alpha_i\theta_\post^i$ is admitted with probability $\lambda(\theta_\post^\alpha)$. This case seems interesting as well. Not clear in which case it's possible to reduce to the single-dimensional case?}

Putting things together, the applicant's individual welfare is:
\begin{equation*}
\iwelf(\{ \e{i} \}_{i=1}^m, \lambda(\theta_\post^\alpha)) = \lambda(\theta_\post^\alpha) - p^m( \e{1}, \cdots, \e{m})
\end{equation*}

[Next verbatim source excerpt]

\begin{definition} [Equilibrium]\label{def:equi}
{Given a}
%Let the
tie-breaking function $\Gamma(\omega)$
%be a pre-announced ordering over index $\omega$, to break ties in post-effort value. Then, let $\gamma(\omega, v)$ be the 1-1 ranking function defined as the CDF of post-effort values $v$ except where there are ties in $v$, in which case the ordering $\Gamma$ is used to produce unique ranks such that the range of $\gamma$ is uniform on $[0, 1]$.\footnote{More formally, note that each discontinuity in the CDF of post-effort values $v$ has a height equal to the measure of the applicants who are tied at that value, and so $\gamma$ for those applicants can be filled in using the CDF of $\Gamma(\omega)$ restricted to the applicants $\omega$ who are tied. Thus, the range of $\gamma$ is uniform on $[0,1]$.
%
{and}
%	Then, given
	a ranking probability function $\lambda$, an \textit{equilibrium} is a set of effort levels and post-effort ranking rewards for each applicant, $\{e(\theta_\pre),  \lambda(\theta_\post(\theta_\pre))\}$ such that, for all
	$\omega$, %$\theta_\pre$,
	\begin{align*}
		e(\theta_\pre(\omega)) &\in \argmax_{e} \iwelf\left(e, \lambda\left(\gamma\left(\omega, v\left(e,\theta_\pre(\omega)\right)\right)\right)\right) %\triangleq \argmax_{e} \left[ \lambda(\gamma(v(e,\theta_\pre)) - \cost(e) \right].
		\\
		\theta_\post(\theta_\pre(\omega)) &= \gamma(\omega, v(e(\theta_\pre(\omega)),\theta_\pre(\omega))),
	\end{align*}
	{where $\gamma(\omega,v)$ is the ranking induced by the CDF of the scores resulting from effort levels $\{e(\theta_\pre)\}$ and the tie-breaking function $\Gamma$.\footnote{See Remark~\ref{rem:gamma-effort} on $\gamma$ and the set of efforts.%Formally, $\gamma$ depends on the set of efforts. Given a fixed set and $\gamma$, for each applicant $\omega$ the first condition considers the counter-factual ranking of $\omega$ with different post-effort values but using the same ranking function. As defined, $\gamma$ yields a uniform distribution of ranks with such measure $0$ changes. In Appendix \Cref{lem:deviationsmeasure0}, we further prove that two effort sets equal up to sets of measure $0$ induce the same ranking function $\gamma$, and so the condition is consistent.%, and so we set it at the $\gamma$ induced by the effort levels $\{e(\theta_\pre)\}$.%, and so it does not vary.
	}
	}
	%In particular, we denote $\theta_\post(\theta_\pre) := \gamma(v(e(\theta_\pre),\theta_\pre))$.
	%\lledit{and for any $\theta_\pre, \theta'_\pre$, we have $\lambda(\theta_\post(\theta_\pre)) > \lambda(\theta_\post(\theta'_\pre)) \implies v(e(\theta_\pre),\theta_\pre) > v(e(\theta_\pre'),\theta_\pre')$.}
\end{definition}

[Next verbatim source excerpt]

\parbold{Applicants}
There is a unit mass of \textit{applicants}, indexed by an observed index $\omega \in [0,1]$ distributed uniformly.\footnote{The index $\omega$ should be interpreted as each applicant's ``name,'' uncorrelated with skill, used solely for tie-breaking.}  Each applicant has a latent (unobserved) skill level represented by some measurable function of $\omega$.
We assume that {the distribution of the skills} has no atoms and that the CDF of this distribution is strictly increasing.
Using the CDF to map the skill of an applicant to a rank in~$[0,1]$, each applicant gets an (unobserved) rank $\theta_\pre =\theta_\pre(\omega)$ which by our assumption on the CDF is again uniformly distributed in $[0, 1]$ (the higher the rank the better).  With this setup, the skill of an applicant with rank $\theta_\pre$ can be written as $f(\theta_\pre)$ where $f$ is a strictly increasing, continuous function.  {It will be notationally convenient to label applicants by their rank $\theta_\pre$, though the reader should note that in contrast to $\omega$, $\theta_\pre(\omega)$ is assumed to be unobservable.}

%3. Each applicant exerts some effort e, at cost \cost(e).


Each applicant chooses an \textit{effort} level $e \ge 0$, the result of which is an observed, post-effort \textit{score}, {$v=v(e,\theta_\pre) = {g(e)}\cdot{f(\theta_\pre)}$}. \lledit{In other words, we assume the post-effort score is the product of two components, associated with the effort and the pre-effort skill respectively.}
{We assume} that the effort transfer function $g: [0, \infty) \mapsto [0, \infty)$ is a continuous, concave, strictly increasing function, representing that marginal effort improves one's score but has diminishing returns.
{%Since applicants can be labelled by their skill level $f(\theta_\pre)$, or equivalently, by their rank $\theta_{\pre}$,
{T}he strategies of the applicants can {then} be described by a function $\theta_\pre \mapsto e(\theta_\pre)$.
%\cb{$\omega\mapsto e(\omega)$.}%, leading to a distribution over post-effort scores $v=v(\theta_\pre))$ via the uniform distribution of $\theta_\pre$ in $[0,1]$.
}
  Each {applicant} is then ranked according to their score $v$%(where we break ties uniformly at random, if they appear)
, resulting in a \textit{post-effort rank} $\theta_\post$.
Note that the ranking $\theta_\post$ is slightly less trivial than a ranking of the skills, since the scores might have ties, which have to be resolved.

{\textbf{Tie Breaking.}
Given a choice of strategies $\theta_\pre \mapsto e(\theta_\pre)$, let $F$ be the CDF of the scores $v(e(\theta_\pre),\theta_\pre)$.
Since atoms for the distribution of $v$ would lead to ties for ranks defined as $F(v)$, we will use the labels of the applicants to break ties with the help of a (publicly announced) tie-breaking function $\Gamma(\omega)$, defined, e.g., via a collision free hash of the applicant names. Here  we require  that $\Gamma$  is a measurable function from $[0,1]$ to $[0,1]$ that maps different applicant labels to different values.
We then use $\Gamma$ to resolve the atoms of $F$, leading to a ranking function $v\mapsto \gamma(\omega, v)$ which is equal to $F(v)$ except when $v$ is an atom of the score distribution, {in which case it takes values in the ``gap interval'' $[F_-(v),  F(v)]$, where $F_-(v)$ is the left limit of $F$ at $v$.}  We  construct $\gamma$ in such a way that it
gives the uniform distribution for $\theta_\post(\omega)$ if we set  $\theta_\post(\omega)=\gamma(\omega,v(e(\theta_\pre(\omega)),\theta_\pre(\omega)))$.\footnote{See Remark~\ref{rem:gamma} for details on this construction.}
\end{ReviewVerbatim}

% report-context-presentation-sha256: 90f9116e292dd71fddfd6b445ca0b060a481d026855bdb9725f2833c2cf62a31
\paragraph{Settled source reading}
\begin{sloppypar}
Equilibrium means that almost every applicant best responds against every feasible deviation;\allowbreak{} uniqueness is up to a null set.\allowbreak{}
\end{sloppypar}
\paragraph{Settled source reading}
\begin{sloppypar}
Applicants know the fixed priority used at score-\allowbreak{}boundary ties.\allowbreak{} This clarifies the stated model, not a corrected admission policy or a new economic assumption.\allowbreak{}
\end{sloppypar}

\paragraph{Lean-expanded paper semantic target}
\begin{ReviewVerbatim}
type:
(ι : Type u_1) → [Fintype ι] → MeasureTheory.Measure (ι → ℝ)

value:
fun (ι : Type u_1) [Fintype ι] => MeasureTheory.Measure.pi fun (x : ι) => LBG22StrategicRanking.unitRankMeasure
\end{ReviewVerbatim}

\paragraph{Direct retained dependencies}
\begin{itemize}\raggedright
\item \hyperlink{governing-declaration-678cf48d549cb620}{LBG22\allowbreak{}Strategic\allowbreak{}Ranking.\allowbreak{}unit\allowbreak{}Rank\allowbreak{}Measure}
\end{itemize}
\paragraph{Recorded source-to-declaration screening}
\textbf{Verdict:} \texttt{matches}
\\
{\raggedright
\textbf{Reason:} The target is the finite product of unit-\allowbreak{}rank measures, so every coordinate rank is independently uniform on the unit interval as required by the source multidimensional model.\allowbreak{}
\par}
\reviewmatch{Matches source input}
\noindent\textbf{Reviewer annotation}\par
\reviewerbox
\clearpage
\hypertarget{paper-prerequisite-b201854e4e4dd891}{}
\subsection*{\small\ttfamily\raggedright LBG22\allowbreak{}Strategic\allowbreak{}Ranking.\allowbreak{}Multidimensional\allowbreak{}Equilibrium}
\noindent\textbf{Paper-local declaration:}\par
{\footnotesize\ttfamily\raggedright LBG22\allowbreak{}Strategic\allowbreak{}Ranking.\allowbreak{}Multidimensional\allowbreak{}Equilibrium\par}
\textbf{Source locator:} {\footnotesize\raggedright cited publication, lines 1760-\allowbreak{}1764\par}
\paragraph{Verbatim paper-source connection}
\begin{ReviewVerbatim}
We show this result by characterizing the equilibrium of the following simplified setting, with further assumptions on the effort cost function $\cost$ and the effort transfer function $g$; the cost function $\cost$ is assumed to be a function of the sum of the efforts exerted:
\[p(e^1, \cdots, e^m) \triangleq p\left(\sum_{i=1}^m \e{i}\right),\] where $p$ is convex and increasing.
This assumption says that the effort exerted for any skill is entirely exchangeable, for example, two hours spent on studying math is as costly as two hours spent on studying chemistry.
%
Effort transfer function $g$ is assumed to be linear, that is, there are constant returns to effort.

[Next verbatim source excerpt]

\parbold{Model.} The model is similar to our base model \Cref{sec:model}; each applicant now has $m$ latent skill levels with respective ranks $\theta_\pre^i \in [0,1]$ for $i\in \{1, \cdots, m\}$. Each rank is drawn independently from the uniform distribution over $[0,1]$, and the skill $i$ of applicant with rank $\theta_\pre^i$ is $\f{i}(\theta_\pre^i)$. Each applicant now chooses effort levels $\{ \e{i} \}$, at cost $p^m( \e{1}, \cdots, \e{m})$, resulting in post-effort scores $\{\score{i}\}$,
\begin{equation*}
\score{i} = g(\e{i})\cdot \f{i}(\theta_\pre^i),
\end{equation*}
where $g$ is concave, increasing, as before, and $\f{i}$ is a continuous, strictly increasing quantile function (${\f{i}}^{-1}$ is the CDF function of the scores on dimension $i$). As before, the school observes the post-effort scores for each applicant, now for each dimension $i$, and designs a non-decreasing function $\lambda : [0,1] \to [0,1]$ denoting the admissions probability $\lambda(\theta_\post)$ for applicant with post-effort rank $\theta_\post$.

%Suppose the school observes $v_i$.

How does the school construct post-effort rank $\theta_\post$? In general, each applicant's \textit{combined} post-effort score may be any function of the scores on each dimension, $\{\score{i}\}$. Here, we assume the following linear score function. The school announces weights $\alpha = (\alpha_1, \cdots, \alpha_m) \in \Delta_{m-1}$ (denoting the $m$-dimensional simplex); the weights $\alpha$ represent the relative emphasis placed on each skill for the admissions decision. Then, each applicant is ranked according to their combined score, $\combscore := \sum_{i=1}^m \alpha_i\score{i}$, resulting in their combined post-effort rank $\theta_\post^\alpha$, and is admitted with probability $\lambda(\theta_\post^\alpha)$.\footnote{The linear combination of skill is similar to the ``linear mechanism'' in \citet{kleinberg18investeffort}, which studied how reward design incentivizes strategic agents to exert effort on different skill dimensions. Compared to their work, where efforts are connected to skills via an effort graph, we consider a simplified setting where each effort maps to one skill, and study how reward design affects equilibrium rankings, in presence of competition.}

%\lnote{\textbf{weighted score or weighted rank?} alternatively, the school combines the post-effort ranks for each skill. i.e. an applicant with \emph{weighted} post-effort rank $\theta_\post^\alpha := \sum_{i=1}^m \alpha_i\theta_\post^i$ is admitted with probability $\lambda(\theta_\post^\alpha)$. This case seems interesting as well. Not clear in which case it's possible to reduce to the single-dimensional case?}

Putting things together, the applicant's individual welfare is:
\begin{equation*}
\iwelf(\{ \e{i} \}_{i=1}^m, \lambda(\theta_\post^\alpha)) = \lambda(\theta_\post^\alpha) - p^m( \e{1}, \cdots, \e{m})
\end{equation*}

[Next verbatim source excerpt]

\parbold{Ranking designer (school)} A single \textit{school} is admitting applicants, based on their ranking. In particular, the school can choose a ranking reward function $\lambda : [0, 1] \mapsto [0, 1]$, such that an applicant with post-effort rank $\theta_\post$ is admitted with probability $\lambda(\theta_\post)$. We assume that $\lambda$ is non-decreasing and that the school has a constraint on the overall probability, such that in expectation it admits a number of applicants equal to {a} capacity constraint  $\rho \in (0,1)$, i.e., \lledit{$\E_{\theta_\post}[\lambda(\theta_\post)] = \rho $}.\footnote{\lledit{We use $\E_{\theta_\post}[\cdot]$ to denote an integral over $\theta_\post$ (and $\E[\cdot]$ to denote an integral over $\omega$)} with respect to the Lebesgue measure; informally, this can be thought of as averaging over the applicant population.} We may also refer to $\lambda$, informally, as the admission policy.



For simplicity, we further assume that $\lambda$ is a step-function with $K$ distinct levels $\ell_0 < \dots < \ell_{K-1}$, and $K-1$ cut-points parameterized by $c_1< \dots < c_{K-1}$ (with $c_0 = 0$, $c_K = 1$). %$\ell_0, \dots, \ell_{K-1}$.
In other words, we have $\lambda(\theta) = \ell_k$, for all $\theta \in \psi_k \triangleq [c_{k}, c_{k+1})$. Thus, applicants in the same post-rank interval $\theta \in \psi_k$ receive the same reward.

[Next verbatim source excerpt]

\begin{definition} [Equilibrium]\label{def:equi}
{Given a}
%Let the
tie-breaking function $\Gamma(\omega)$
%be a pre-announced ordering over index $\omega$, to break ties in post-effort value. Then, let $\gamma(\omega, v)$ be the 1-1 ranking function defined as the CDF of post-effort values $v$ except where there are ties in $v$, in which case the ordering $\Gamma$ is used to produce unique ranks such that the range of $\gamma$ is uniform on $[0, 1]$.\footnote{More formally, note that each discontinuity in the CDF of post-effort values $v$ has a height equal to the measure of the applicants who are tied at that value, and so $\gamma$ for those applicants can be filled in using the CDF of $\Gamma(\omega)$ restricted to the applicants $\omega$ who are tied. Thus, the range of $\gamma$ is uniform on $[0,1]$.
%
{and}
%	Then, given
	a ranking probability function $\lambda$, an \textit{equilibrium} is a set of effort levels and post-effort ranking rewards for each applicant, $\{e(\theta_\pre),  \lambda(\theta_\post(\theta_\pre))\}$ such that, for all
	$\omega$, %$\theta_\pre$,
	\begin{align*}
		e(\theta_\pre(\omega)) &\in \argmax_{e} \iwelf\left(e, \lambda\left(\gamma\left(\omega, v\left(e,\theta_\pre(\omega)\right)\right)\right)\right) %\triangleq \argmax_{e} \left[ \lambda(\gamma(v(e,\theta_\pre)) - \cost(e) \right].
		\\
		\theta_\post(\theta_\pre(\omega)) &= \gamma(\omega, v(e(\theta_\pre(\omega)),\theta_\pre(\omega))),
	\end{align*}
	{where $\gamma(\omega,v)$ is the ranking induced by the CDF of the scores resulting from effort levels $\{e(\theta_\pre)\}$ and the tie-breaking function $\Gamma$.\footnote{See Remark~\ref{rem:gamma-effort} on $\gamma$ and the set of efforts.%Formally, $\gamma$ depends on the set of efforts. Given a fixed set and $\gamma$, for each applicant $\omega$ the first condition considers the counter-factual ranking of $\omega$ with different post-effort values but using the same ranking function. As defined, $\gamma$ yields a uniform distribution of ranks with such measure $0$ changes. In Appendix \Cref{lem:deviationsmeasure0}, we further prove that two effort sets equal up to sets of measure $0$ induce the same ranking function $\gamma$, and so the condition is consistent.%, and so we set it at the $\gamma$ induced by the effort levels $\{e(\theta_\pre)\}$.%, and so it does not vary.
	}
	}
	%In particular, we denote $\theta_\post(\theta_\pre) := \gamma(v(e(\theta_\pre),\theta_\pre))$.
	%\lledit{and for any $\theta_\pre, \theta'_\pre$, we have $\lambda(\theta_\post(\theta_\pre)) > \lambda(\theta_\post(\theta'_\pre)) \implies v(e(\theta_\pre),\theta_\pre) > v(e(\theta_\pre'),\theta_\pre')$.}
\end{definition}
\end{ReviewVerbatim}

% report-context-presentation-sha256: 90f9116e292dd71fddfd6b445ca0b060a481d026855bdb9725f2833c2cf62a31
\paragraph{Settled source reading}
\begin{sloppypar}
Equilibrium means that almost every applicant best responds against every feasible deviation;\allowbreak{} uniqueness is up to a null set.\allowbreak{}
\end{sloppypar}
\paragraph{Settled source reading}
\begin{sloppypar}
Applicants know the fixed priority used at score-\allowbreak{}boundary ties.\allowbreak{} This clarifies the stated model, not a corrected admission policy or a new economic assumption.\allowbreak{}
\end{sloppypar}

\paragraph{Lean-expanded paper semantic target}
\begin{ReviewVerbatim}
type:
(m : ℕ) →
  ℕ →
    (ℝ → ℝ) →
      ℝ →
        (Fin (m + 1) → ℝ) →
          (Fin (m + 1) → ℝ → ℝ) →
            (ℕ → ℝ) →
              (ℕ → ℝ) → ((Fin (m + 1) → ℝ) → Fin (m + 1) → ℝ) → ((Fin (m + 1) → ℝ) → ℝ) → ((Fin (m + 1) → ℝ) → ℝ) → Prop

value:
fun (m n : ℕ) (cost : ℝ → ℝ) (h : ℝ) (weight : Fin (m + 1) → ℝ) (skill : Fin (m + 1) → ℝ → ℝ) (cutoff reward : ℕ → ℝ)
    (effort : (Fin (m + 1) → ℝ) → Fin (m + 1) → ℝ) (score tie : (Fin (m + 1) → ℝ) → ℝ) =>
  Measurable score ∧
    ∀ᵐ (x : Fin (m + 1) → ℝ) ∂LBG22StrategicRanking.sourceMultidimensionalPopulation (Fin (m + 1)),
      LBG22StrategicRanking.SourceMultidimensionalBestResponseAt
        (LBG22StrategicRanking.sourceMultidimensionalPopulation (Fin (m + 1))) cost
        (LBG22StrategicRanking.sourceMultidimensionalCoefficient weight skill) score tie h
        (fun (r : ℝ) => reward ↑(LBG22StrategicRanking.finiteLowerRankBand (fun (i : Fin (n + 1)) => cutoff ↑i) r)) x
        (effort x)
\end{ReviewVerbatim}

\paragraph{Direct retained dependencies}
\begin{itemize}\raggedright
\item \hyperlink{paper-prerequisite-59d181ae71b7b024}{LBG22\allowbreak{}Strategic\allowbreak{}Ranking.\allowbreak{}Source\allowbreak{}Multidimensional\allowbreak{}Best\allowbreak{}Response\allowbreak{}At}
\item \hyperlink{paper-prerequisite-a42f440ddc254fa7}{LBG22\allowbreak{}Strategic\allowbreak{}Ranking.\allowbreak{}finite\allowbreak{}Lower\allowbreak{}Rank\allowbreak{}Band}
\item \hyperlink{paper-prerequisite-e57e55071c75ce91}{LBG22\allowbreak{}Strategic\allowbreak{}Ranking.\allowbreak{}source\allowbreak{}Multidimensional\allowbreak{}Coefficient}
\item \hyperlink{paper-prerequisite-16d135799a0f0491}{LBG22\allowbreak{}Strategic\allowbreak{}Ranking.\allowbreak{}source\allowbreak{}Multidimensional\allowbreak{}Population}
\end{itemize}
\paragraph{Recorded source-to-declaration screening}
\textbf{Verdict:} \texttt{matches}
\\
{\raggedright
\textbf{Reason:} The target states a measurable-\allowbreak{}score equilibrium in which almost every independently sampled rank vector maximizes finite-\allowbreak{}band reward minus exchangeable cost over every nonnegative vector deviation.\allowbreak{} The displayed coefficient, product population, combined score, rank, and policy support matches the source exchangeable-\allowbreak{}cost linear-\allowbreak{}production model under the approved a.\allowbreak{}e.\allowbreak{} and fixed-\allowbreak{}priority conventions.\allowbreak{}
\par}
\reviewmatch{Matches source input}
\noindent\textbf{Reviewer annotation}\par
\reviewerbox
\clearpage
\hypertarget{paper-prerequisite-8e74a377842bef95}{}
\subsection*{\small\ttfamily\raggedright LBG22\allowbreak{}Strategic\allowbreak{}Ranking.\allowbreak{}Multidimensional\allowbreak{}Primitives}
\noindent\textbf{Paper-local declaration:}\par
{\footnotesize\ttfamily\raggedright LBG22\allowbreak{}Strategic\allowbreak{}Ranking.\allowbreak{}Multidimensional\allowbreak{}Primitives\par}
\textbf{Source locator:} {\footnotesize\raggedright cited publication, lines 1760-\allowbreak{}1764\par}
\paragraph{Verbatim paper-source connection}
\begin{ReviewVerbatim}
We show this result by characterizing the equilibrium of the following simplified setting, with further assumptions on the effort cost function $\cost$ and the effort transfer function $g$; the cost function $\cost$ is assumed to be a function of the sum of the efforts exerted:
\[p(e^1, \cdots, e^m) \triangleq p\left(\sum_{i=1}^m \e{i}\right),\] where $p$ is convex and increasing.
This assumption says that the effort exerted for any skill is entirely exchangeable, for example, two hours spent on studying math is as costly as two hours spent on studying chemistry.
%
Effort transfer function $g$ is assumed to be linear, that is, there are constant returns to effort.

[Next verbatim source excerpt]

\parbold{Model.} The model is similar to our base model \Cref{sec:model}; each applicant now has $m$ latent skill levels with respective ranks $\theta_\pre^i \in [0,1]$ for $i\in \{1, \cdots, m\}$. Each rank is drawn independently from the uniform distribution over $[0,1]$, and the skill $i$ of applicant with rank $\theta_\pre^i$ is $\f{i}(\theta_\pre^i)$. Each applicant now chooses effort levels $\{ \e{i} \}$, at cost $p^m( \e{1}, \cdots, \e{m})$, resulting in post-effort scores $\{\score{i}\}$,
\begin{equation*}
\score{i} = g(\e{i})\cdot \f{i}(\theta_\pre^i),
\end{equation*}
where $g$ is concave, increasing, as before, and $\f{i}$ is a continuous, strictly increasing quantile function (${\f{i}}^{-1}$ is the CDF function of the scores on dimension $i$). As before, the school observes the post-effort scores for each applicant, now for each dimension $i$, and designs a non-decreasing function $\lambda : [0,1] \to [0,1]$ denoting the admissions probability $\lambda(\theta_\post)$ for applicant with post-effort rank $\theta_\post$.

%Suppose the school observes $v_i$.

How does the school construct post-effort rank $\theta_\post$? In general, each applicant's \textit{combined} post-effort score may be any function of the scores on each dimension, $\{\score{i}\}$. Here, we assume the following linear score function. The school announces weights $\alpha = (\alpha_1, \cdots, \alpha_m) \in \Delta_{m-1}$ (denoting the $m$-dimensional simplex); the weights $\alpha$ represent the relative emphasis placed on each skill for the admissions decision. Then, each applicant is ranked according to their combined score, $\combscore := \sum_{i=1}^m \alpha_i\score{i}$, resulting in their combined post-effort rank $\theta_\post^\alpha$, and is admitted with probability $\lambda(\theta_\post^\alpha)$.\footnote{The linear combination of skill is similar to the ``linear mechanism'' in \citet{kleinberg18investeffort}, which studied how reward design incentivizes strategic agents to exert effort on different skill dimensions. Compared to their work, where efforts are connected to skills via an effort graph, we consider a simplified setting where each effort maps to one skill, and study how reward design affects equilibrium rankings, in presence of competition.}

%\lnote{\textbf{weighted score or weighted rank?} alternatively, the school combines the post-effort ranks for each skill. i.e. an applicant with \emph{weighted} post-effort rank $\theta_\post^\alpha := \sum_{i=1}^m \alpha_i\theta_\post^i$ is admitted with probability $\lambda(\theta_\post^\alpha)$. This case seems interesting as well. Not clear in which case it's possible to reduce to the single-dimensional case?}

Putting things together, the applicant's individual welfare is:
\begin{equation*}
\iwelf(\{ \e{i} \}_{i=1}^m, \lambda(\theta_\post^\alpha)) = \lambda(\theta_\post^\alpha) - p^m( \e{1}, \cdots, \e{m})
\end{equation*}

[Next verbatim source excerpt]

\parbold{Ranking designer (school)} A single \textit{school} is admitting applicants, based on their ranking. In particular, the school can choose a ranking reward function $\lambda : [0, 1] \mapsto [0, 1]$, such that an applicant with post-effort rank $\theta_\post$ is admitted with probability $\lambda(\theta_\post)$. We assume that $\lambda$ is non-decreasing and that the school has a constraint on the overall probability, such that in expectation it admits a number of applicants equal to {a} capacity constraint  $\rho \in (0,1)$, i.e., \lledit{$\E_{\theta_\post}[\lambda(\theta_\post)] = \rho $}.\footnote{\lledit{We use $\E_{\theta_\post}[\cdot]$ to denote an integral over $\theta_\post$ (and $\E[\cdot]$ to denote an integral over $\omega$)} with respect to the Lebesgue measure; informally, this can be thought of as averaging over the applicant population.} We may also refer to $\lambda$, informally, as the admission policy.



For simplicity, we further assume that $\lambda$ is a step-function with $K$ distinct levels $\ell_0 < \dots < \ell_{K-1}$, and $K-1$ cut-points parameterized by $c_1< \dots < c_{K-1}$ (with $c_0 = 0$, $c_K = 1$). %$\ell_0, \dots, \ell_{K-1}$.
In other words, we have $\lambda(\theta) = \ell_k$, for all $\theta \in \psi_k \triangleq [c_{k}, c_{k+1})$. Thus, applicants in the same post-rank interval $\theta \in \psi_k$ receive the same reward.

[Next verbatim source excerpt]

\begin{definition} [Equilibrium]\label{def:equi}
{Given a}
%Let the
tie-breaking function $\Gamma(\omega)$
%be a pre-announced ordering over index $\omega$, to break ties in post-effort value. Then, let $\gamma(\omega, v)$ be the 1-1 ranking function defined as the CDF of post-effort values $v$ except where there are ties in $v$, in which case the ordering $\Gamma$ is used to produce unique ranks such that the range of $\gamma$ is uniform on $[0, 1]$.\footnote{More formally, note that each discontinuity in the CDF of post-effort values $v$ has a height equal to the measure of the applicants who are tied at that value, and so $\gamma$ for those applicants can be filled in using the CDF of $\Gamma(\omega)$ restricted to the applicants $\omega$ who are tied. Thus, the range of $\gamma$ is uniform on $[0,1]$.
%
{and}
%	Then, given
	a ranking probability function $\lambda$, an \textit{equilibrium} is a set of effort levels and post-effort ranking rewards for each applicant, $\{e(\theta_\pre),  \lambda(\theta_\post(\theta_\pre))\}$ such that, for all
	$\omega$, %$\theta_\pre$,
	\begin{align*}
		e(\theta_\pre(\omega)) &\in \argmax_{e} \iwelf\left(e, \lambda\left(\gamma\left(\omega, v\left(e,\theta_\pre(\omega)\right)\right)\right)\right) %\triangleq \argmax_{e} \left[ \lambda(\gamma(v(e,\theta_\pre)) - \cost(e) \right].
		\\
		\theta_\post(\theta_\pre(\omega)) &= \gamma(\omega, v(e(\theta_\pre(\omega)),\theta_\pre(\omega))),
	\end{align*}
	{where $\gamma(\omega,v)$ is the ranking induced by the CDF of the scores resulting from effort levels $\{e(\theta_\pre)\}$ and the tie-breaking function $\Gamma$.\footnote{See Remark~\ref{rem:gamma-effort} on $\gamma$ and the set of efforts.%Formally, $\gamma$ depends on the set of efforts. Given a fixed set and $\gamma$, for each applicant $\omega$ the first condition considers the counter-factual ranking of $\omega$ with different post-effort values but using the same ranking function. As defined, $\gamma$ yields a uniform distribution of ranks with such measure $0$ changes. In Appendix \Cref{lem:deviationsmeasure0}, we further prove that two effort sets equal up to sets of measure $0$ induce the same ranking function $\gamma$, and so the condition is consistent.%, and so we set it at the $\gamma$ induced by the effort levels $\{e(\theta_\pre)\}$.%, and so it does not vary.
	}
	}
	%In particular, we denote $\theta_\post(\theta_\pre) := \gamma(v(e(\theta_\pre),\theta_\pre))$.
	%\lledit{and for any $\theta_\pre, \theta'_\pre$, we have $\lambda(\theta_\post(\theta_\pre)) > \lambda(\theta_\post(\theta'_\pre)) \implies v(e(\theta_\pre),\theta_\pre) > v(e(\theta_\pre'),\theta_\pre')$.}
\end{definition}
\end{ReviewVerbatim}

% report-context-presentation-sha256: 90f9116e292dd71fddfd6b445ca0b060a481d026855bdb9725f2833c2cf62a31
\paragraph{Settled source reading}
\begin{sloppypar}
Equilibrium means that almost every applicant best responds against every feasible deviation;\allowbreak{} uniqueness is up to a null set.\allowbreak{}
\end{sloppypar}
\paragraph{Settled source reading}
\begin{sloppypar}
Applicants know the fixed priority used at score-\allowbreak{}boundary ties.\allowbreak{} This clarifies the stated model, not a corrected admission policy or a new economic assumption.\allowbreak{}
\end{sloppypar}

\paragraph{Lean-expanded paper semantic target}
\begin{ReviewVerbatim}
type:
(m : ℕ) → (ℝ → ℝ) → ℝ → (Fin (m + 1) → ℝ) → (Fin (m + 1) → ℝ → ℝ) → Prop

value:
fun (m : ℕ) (cost : ℝ → ℝ) (h : ℝ) (weight : Fin (m + 1) → ℝ) (skill : Fin (m + 1) → ℝ → ℝ) =>
  0 < h ∧
    (∀ (i : Fin (m + 1)), 0 ≤ weight i) ∧
      ∑ i : Fin (m + 1), weight i = 1 ∧
        ContinuousOn cost (Set.Ici 0) ∧
          StrictMonoOn cost (Set.Ici 0) ∧
            ConvexOn ℝ (Set.Ici 0) cost ∧
              (∀ (i : Fin (m + 1)), ContinuousOn (skill i) (Set.Icc 0 1)) ∧
                (∀ (i : Fin (m + 1)), StrictMonoOn (skill i) (Set.Icc 0 1)) ∧ ∀ (i : Fin (m + 1)), 0 ≤ skill i 0
\end{ReviewVerbatim}


\paragraph{Recorded source-to-declaration screening}
\textbf{Verdict:} \texttt{matches}
\\
{\raggedright
\textbf{Reason:} The target imposes positive linear production coefficient, simplex weights, continuous increasing convex total-\allowbreak{}effort cost, and continuous strictly increasing nonnegative skill quantiles.\allowbreak{} This is exactly the source's linear-\allowbreak{}production, exchangeable-\allowbreak{}cost multidimensional primitive restriction;\allowbreak{} allowing zero simplex coordinates is a broader domain that includes every positive-\allowbreak{}weight source instance.\allowbreak{}
\par}
\reviewmatch{Matches source input}
\noindent\textbf{Reviewer annotation}\par
\reviewerbox
\clearpage
\hypertarget{paper-prerequisite-b504d72e1e0de867}{}
\subsection*{\small\ttfamily\raggedright LBG22\allowbreak{}Strategic\allowbreak{}Ranking.\allowbreak{}Multitask\allowbreak{}Equilibrium}
\noindent\textbf{Paper-local declaration:}\par
{\footnotesize\ttfamily\raggedright LBG22\allowbreak{}Strategic\allowbreak{}Ranking.\allowbreak{}Multitask\allowbreak{}Equilibrium\par}
\textbf{Source locator:} {\footnotesize\raggedright cited publication, lines 1799-\allowbreak{}1806\par}
\paragraph{Verbatim paper-source connection}
\begin{ReviewVerbatim}
In our simplified setting, there are two skills $M$ and $U$: $M$ has a measurable score $v^M$ and $U$ has an unmeasurable score $v^U$. For example, $M$ could be scholastic achievement as measured by SAT scores, and $U$ could be ``creativity", a personal quality that is valued by the school but is not directly measurable.
%
Since the school cannot observe $v^U$, its admission policy is based on $\theta_\post^M$ only, that is, $\alpha_U = 0$ and $\alpha_M = 1$.

We further assume that each applicant has a fixed effort budget of $B > 0$, and is intrinsically motivated to exert effort in the unmeasurable skill $U$; in fact, they will always exert effort $e^U = B - e^M$. Formally this corresponds to the effort cost function \[\cost^m(e_M, e_U) = \cost(e^M) - (\max(0,B-(e^M+e^U)))^2\] where $\cost$ is convex and increasing. We can write the applicant's individual welfare as:
\begin{equation*}
W(e^M, e^U, \lambda(\theta_\post^M)) = \lambda(\theta_\post^M) - \cost^m(e_M, e_U).
\end{equation*}
\end{ReviewVerbatim}

% report-context-presentation-sha256: ce5679112671da8414975164981e5740bcad9608aa7765c30f1069706683e7b1
\paragraph{Settled source reading}
\begin{sloppypar}
Equilibrium means that almost every applicant best responds against every feasible deviation;\allowbreak{} uniqueness is up to a null set.\allowbreak{}
\end{sloppypar}
\paragraph{Settled source reading}
\begin{sloppypar}
The multitask model has nonnegative effort on both tasks summing to the stated budget.\allowbreak{} This is distinct from the unrestricted total effort in the multidimensional ranking model.\allowbreak{}
\end{sloppypar}
\paragraph{Settled source reading}
\begin{sloppypar}
Applicants know the fixed priority used at score-\allowbreak{}boundary ties.\allowbreak{} This clarifies the stated model, not a corrected admission policy or a new economic assumption.\allowbreak{}
\end{sloppypar}

\paragraph{Lean-expanded paper semantic target}
\begin{ReviewVerbatim}
type:
(ℝ → ℝ) → (ℝ → ℝ) → (ℝ → ℝ) → ℝ → ℝ → ℝ → (ℝ × ℝ → ℝ) → (ℝ × ℝ → ℝ) → (ℝ × ℝ → ℝ) → (ℝ × ℝ → ℝ) → Prop

value:
fun (cost production skillM : ℝ → ℝ) (B rho c : ℝ) (effortM effortU score tie : ℝ × ℝ → ℝ) =>
  Measurable score ∧
    ∀ᵐ (x : ℝ × ℝ) ∂LBG22StrategicRanking.unitRankMeasure.prod LBG22StrategicRanking.unitRankMeasure,
      LBG22StrategicRanking.SourceHardBudgetBestResponseAt
        (LBG22StrategicRanking.unitRankMeasure.prod LBG22StrategicRanking.unitRankMeasure) cost production
        (fun (x : ℝ × ℝ) => skillM x.1) score tie B rho c x (effortM x) (effortU x)
\end{ReviewVerbatim}

\paragraph{Direct retained dependencies}
\begin{itemize}\raggedright
\item \hyperlink{governing-declaration-bfa45ba26bb8d395}{LBG22\allowbreak{}Strategic\allowbreak{}Ranking.\allowbreak{}Source\allowbreak{}Hard\allowbreak{}Budget\allowbreak{}Best\allowbreak{}Response\allowbreak{}At}
\item \hyperlink{governing-declaration-678cf48d549cb620}{LBG22\allowbreak{}Strategic\allowbreak{}Ranking.\allowbreak{}unit\allowbreak{}Rank\allowbreak{}Measure}
\end{itemize}
\paragraph{Recorded source-to-declaration screening}
\textbf{Verdict:} \texttt{matches}
\\
{\raggedright
\textbf{Reason:} The target uses the independent two-\allowbreak{}rank population and requires almost every applicant to choose nonnegative measurable and unmeasurable efforts exhausting B while maximizing admission reward minus measurable-\allowbreak{}effort cost against all hard-\allowbreak{}budget deviations.\allowbreak{} This matches the approved hard-\allowbreak{}equality-\allowbreak{}budget reading and fixed-\allowbreak{}priority source model.\allowbreak{}
\par}
\reviewmatch{Matches source input}
\noindent\textbf{Reviewer annotation}\par
\reviewerbox
\clearpage
\hypertarget{paper-prerequisite-3c93745e73ce0a63}{}
\subsection*{\small\ttfamily\raggedright LBG22\allowbreak{}Strategic\allowbreak{}Ranking.\allowbreak{}Multitask\allowbreak{}Primitives}
\noindent\textbf{Paper-local declaration:}\par
{\footnotesize\ttfamily\raggedright LBG22\allowbreak{}Strategic\allowbreak{}Ranking.\allowbreak{}Multitask\allowbreak{}Primitives\par}
\textbf{Source locator:} {\footnotesize\raggedright cited publication, lines 1799-\allowbreak{}1806\par}
\paragraph{Verbatim paper-source connection}
\begin{ReviewVerbatim}
In our simplified setting, there are two skills $M$ and $U$: $M$ has a measurable score $v^M$ and $U$ has an unmeasurable score $v^U$. For example, $M$ could be scholastic achievement as measured by SAT scores, and $U$ could be ``creativity", a personal quality that is valued by the school but is not directly measurable.
%
Since the school cannot observe $v^U$, its admission policy is based on $\theta_\post^M$ only, that is, $\alpha_U = 0$ and $\alpha_M = 1$.

We further assume that each applicant has a fixed effort budget of $B > 0$, and is intrinsically motivated to exert effort in the unmeasurable skill $U$; in fact, they will always exert effort $e^U = B - e^M$. Formally this corresponds to the effort cost function \[\cost^m(e_M, e_U) = \cost(e^M) - (\max(0,B-(e^M+e^U)))^2\] where $\cost$ is convex and increasing. We can write the applicant's individual welfare as:
\begin{equation*}
W(e^M, e^U, \lambda(\theta_\post^M)) = \lambda(\theta_\post^M) - \cost^m(e_M, e_U).
\end{equation*}
\end{ReviewVerbatim}

% report-context-presentation-sha256: ce5679112671da8414975164981e5740bcad9608aa7765c30f1069706683e7b1
\paragraph{Settled source reading}
\begin{sloppypar}
Equilibrium means that almost every applicant best responds against every feasible deviation;\allowbreak{} uniqueness is up to a null set.\allowbreak{}
\end{sloppypar}
\paragraph{Settled source reading}
\begin{sloppypar}
The multitask model has nonnegative effort on both tasks summing to the stated budget.\allowbreak{} This is distinct from the unrestricted total effort in the multidimensional ranking model.\allowbreak{}
\end{sloppypar}
\paragraph{Settled source reading}
\begin{sloppypar}
Applicants know the fixed priority used at score-\allowbreak{}boundary ties.\allowbreak{} This clarifies the stated model, not a corrected admission policy or a new economic assumption.\allowbreak{}
\end{sloppypar}

\paragraph{Lean-expanded paper semantic target}
\begin{ReviewVerbatim}
type:
(ℝ → ℝ) → (ℝ → ℝ) → (ℝ → ℝ) → (ℝ → ℝ) → ℝ → ℝ → Prop

value:
fun (cost production skillM skillU : ℝ → ℝ) (B rho : ℝ) =>
  0 < B ∧
    rho ∈ Set.Ioo 0 1 ∧
      ContinuousOn cost (Set.Ici 0) ∧
        MonotoneOn cost (Set.Ici 0) ∧
          ConvexOn ℝ (Set.Ici 0) cost ∧
            (∀ e ∈ Set.Ici 0, 0 ≤ cost e) ∧
              ContinuousOn production (Set.Ici 0) ∧
                StrictMonoOn production (Set.Ici 0) ∧
                  ConcaveOn ℝ (Set.Ici 0) production ∧
                    0 ≤ production 0 ∧
                      ContinuousOn skillM (Set.Icc 0 1) ∧
                        StrictMonoOn skillM (Set.Icc 0 1) ∧
                          0 ≤ skillM 0 ∧
                            ContinuousOn skillU (Set.Icc 0 1) ∧ StrictMonoOn skillU (Set.Icc 0 1) ∧ 0 ≤ skillU 0
\end{ReviewVerbatim}


\paragraph{Recorded source-to-declaration screening}
\textbf{Verdict:} \texttt{matches}
\\
{\raggedright
\textbf{Reason:} The target gives B>0, capacity in (0,1), convex increasing nonnegative cost, concave strictly increasing nonnegative production, and regular nonnegative skill quantiles for both tasks.\allowbreak{} These are the source primitive conditions, with no added frontier or optimizer assumption.\allowbreak{}
\par}
\reviewmatch{Matches source input}
\noindent\textbf{Reviewer annotation}\par
\reviewerbox
\clearpage
\hypertarget{paper-prerequisite-519e321cd0957049}{}
\subsection*{\small\ttfamily\raggedright LBG22\allowbreak{}Strategic\allowbreak{}Ranking.\allowbreak{}Source\allowbreak{}Finite\allowbreak{}Best\allowbreak{}Response\allowbreak{}At}
\noindent\textbf{Paper-local declaration:}\par
{\footnotesize\ttfamily\raggedright LBG22\allowbreak{}Strategic\allowbreak{}Ranking.\allowbreak{}Source\allowbreak{}Finite\allowbreak{}Best\allowbreak{}Response\allowbreak{}At\par}
\textbf{Source locator:} {\footnotesize\raggedright cited publication, lines 1061-\allowbreak{}1069\par}
\paragraph{Verbatim paper-source connection}
\begin{ReviewVerbatim}
\parbold{Individual applicant welfare and equilibrium} Given the designer's function $\lambda$ and the effort levels of other applicants, each applicant chooses effort $e$ to maximize their individual welfare, % $W$,
\begin{equation*}
	\iwelf(e, \lambda(\theta_\post)) = \lambda(\theta_\post) - \cost(e),
\end{equation*}
where the effort cost function $\cost$ is non-negative, continuous, and strictly convex {on $[0,\infty)$, with $e_0:=\argmin_e \cost(e)$ and $p(e_0) = 0$.}
%\lnote{also want to assume that $\cost$ attains its minimum at $0$?}
%Let $e_0:=\argmin_e \cost(e)$ and $p(e_0) = 0$.

Applicants are assumed not to personally benefit from increasing their score $v$, except through the corresponding increase in their rank and reward. While the definition of $\cost$ does not preclude the applicant receiving any intrinsic benefit from exerting effort, we assume that the \emph{net} benefit from effort is non-positive.

[Next verbatim source excerpt]

\begin{definition} [Equilibrium]\label{def:equi}
{Given a}
%Let the
tie-breaking function $\Gamma(\omega)$
%be a pre-announced ordering over index $\omega$, to break ties in post-effort value. Then, let $\gamma(\omega, v)$ be the 1-1 ranking function defined as the CDF of post-effort values $v$ except where there are ties in $v$, in which case the ordering $\Gamma$ is used to produce unique ranks such that the range of $\gamma$ is uniform on $[0, 1]$.\footnote{More formally, note that each discontinuity in the CDF of post-effort values $v$ has a height equal to the measure of the applicants who are tied at that value, and so $\gamma$ for those applicants can be filled in using the CDF of $\Gamma(\omega)$ restricted to the applicants $\omega$ who are tied. Thus, the range of $\gamma$ is uniform on $[0,1]$.
%
{and}
%	Then, given
	a ranking probability function $\lambda$, an \textit{equilibrium} is a set of effort levels and post-effort ranking rewards for each applicant, $\{e(\theta_\pre),  \lambda(\theta_\post(\theta_\pre))\}$ such that, for all
	$\omega$, %$\theta_\pre$,
	\begin{align*}
		e(\theta_\pre(\omega)) &\in \argmax_{e} \iwelf\left(e, \lambda\left(\gamma\left(\omega, v\left(e,\theta_\pre(\omega)\right)\right)\right)\right) %\triangleq \argmax_{e} \left[ \lambda(\gamma(v(e,\theta_\pre)) - \cost(e) \right].
		\\
		\theta_\post(\theta_\pre(\omega)) &= \gamma(\omega, v(e(\theta_\pre(\omega)),\theta_\pre(\omega))),
	\end{align*}
	{where $\gamma(\omega,v)$ is the ranking induced by the CDF of the scores resulting from effort levels $\{e(\theta_\pre)\}$ and the tie-breaking function $\Gamma$.\footnote{See Remark~\ref{rem:gamma-effort} on $\gamma$ and the set of efforts.%Formally, $\gamma$ depends on the set of efforts. Given a fixed set and $\gamma$, for each applicant $\omega$ the first condition considers the counter-factual ranking of $\omega$ with different post-effort values but using the same ranking function. As defined, $\gamma$ yields a uniform distribution of ranks with such measure $0$ changes. In Appendix \Cref{lem:deviationsmeasure0}, we further prove that two effort sets equal up to sets of measure $0$ induce the same ranking function $\gamma$, and so the condition is consistent.%, and so we set it at the $\gamma$ induced by the effort levels $\{e(\theta_\pre)\}$.%, and so it does not vary.
	}
	}
	%In particular, we denote $\theta_\post(\theta_\pre) := \gamma(v(e(\theta_\pre),\theta_\pre))$.
	%\lledit{and for any $\theta_\pre, \theta'_\pre$, we have $\lambda(\theta_\post(\theta_\pre)) > \lambda(\theta_\post(\theta'_\pre)) \implies v(e(\theta_\pre),\theta_\pre) > v(e(\theta_\pre'),\theta_\pre')$.}
\end{definition}
\end{ReviewVerbatim}

% report-context-presentation-sha256: 90f9116e292dd71fddfd6b445ca0b060a481d026855bdb9725f2833c2cf62a31
\paragraph{Settled source reading}
\begin{sloppypar}
Equilibrium means that almost every applicant best responds against every feasible deviation;\allowbreak{} uniqueness is up to a null set.\allowbreak{}
\end{sloppypar}
\paragraph{Settled source reading}
\begin{sloppypar}
Applicants know the fixed priority used at score-\allowbreak{}boundary ties.\allowbreak{} This clarifies the stated model, not a corrected admission policy or a new economic assumption.\allowbreak{}
\end{sloppypar}

\paragraph{Lean-expanded paper semantic target}
\begin{ReviewVerbatim}
type:
(ℝ → ℝ) → (ℝ → ℝ) → (ℝ → ℝ) → (ℝ → ℝ) → (ℝ → ℝ) → {n : ℕ} → (Fin (n + 1) → ℝ) → (ℕ → ℝ) → ℝ → ℝ → Prop

value:
fun (cost production skill score tie : ℝ → ℝ) {n : ℕ} (cutoff : Fin (n + 1) → ℝ) (reward : ℕ → ℝ) (t effort : ℝ) =>
  0 ≤ effort ∧
    score t = production effort * skill t ∧
      ∀ (d : ℝ),
        0 ≤ d →
          reward
                ↑(LBG22StrategicRanking.finiteLowerRankBand cutoff
                    (LBG22StrategicRanking.counterfactualTieBrokenRank LBG22StrategicRanking.unitRankMeasure score tie t
                      (production d * skill t))) -
              cost d ≤
            reward
                ↑(LBG22StrategicRanking.finiteLowerRankBand cutoff
                    (LBG22StrategicRanking.tieBrokenRank LBG22StrategicRanking.unitRankMeasure score tie t)) -
              cost effort
\end{ReviewVerbatim}

\paragraph{Direct retained dependencies}
\begin{itemize}\raggedright
\item \hyperlink{paper-prerequisite-6ed45e46d01eadeb}{LBG22\allowbreak{}Strategic\allowbreak{}Ranking.\allowbreak{}counterfactual\allowbreak{}Tie\allowbreak{}Broken\allowbreak{}Rank}
\item \hyperlink{paper-prerequisite-a42f440ddc254fa7}{LBG22\allowbreak{}Strategic\allowbreak{}Ranking.\allowbreak{}finite\allowbreak{}Lower\allowbreak{}Rank\allowbreak{}Band}
\item \hyperlink{paper-prerequisite-b7c71e53177b76c6}{LBG22\allowbreak{}Strategic\allowbreak{}Ranking.\allowbreak{}tie\allowbreak{}Broken\allowbreak{}Rank}
\item \hyperlink{governing-declaration-678cf48d549cb620}{LBG22\allowbreak{}Strategic\allowbreak{}Ranking.\allowbreak{}unit\allowbreak{}Rank\allowbreak{}Measure}
\end{itemize}
\paragraph{Recorded source-to-declaration screening}
\textbf{Verdict:} \texttt{matches}
\\
{\raggedright
\textbf{Reason:} The target gives the scalar nonnegative effort, production-\allowbreak{}times-\allowbreak{}skill score identity, and utility maximization against every nonnegative deviation using the fixed counterfactual rank and finite reward band.\allowbreak{} The approved context establishes its strict-\allowbreak{}percentile form as equivalent to actual score-\allowbreak{}priority admission for all deviations.\allowbreak{}
\par}
\reviewmatch{Matches source input}
\noindent\textbf{Reviewer annotation}\par
\reviewerbox
\clearpage
\hypertarget{paper-prerequisite-2a1b692f4e3136c6}{}
\subsection*{\small\ttfamily\raggedright LBG22\allowbreak{}Strategic\allowbreak{}Ranking.\allowbreak{}Ranking\allowbreak{}Equilibrium}
\noindent\textbf{Paper-local declaration:}\par
{\footnotesize\ttfamily\raggedright LBG22\allowbreak{}Strategic\allowbreak{}Ranking.\allowbreak{}Ranking\allowbreak{}Equilibrium\par}
\textbf{Source locator:} {\footnotesize\raggedright cited publication, lines 1077-\allowbreak{}1097\par}
\paragraph{Verbatim paper-source connection}
\begin{ReviewVerbatim}
\begin{definition} [Equilibrium]\label{def:equi}
{Given a}
%Let the
tie-breaking function $\Gamma(\omega)$
%be a pre-announced ordering over index $\omega$, to break ties in post-effort value. Then, let $\gamma(\omega, v)$ be the 1-1 ranking function defined as the CDF of post-effort values $v$ except where there are ties in $v$, in which case the ordering $\Gamma$ is used to produce unique ranks such that the range of $\gamma$ is uniform on $[0, 1]$.\footnote{More formally, note that each discontinuity in the CDF of post-effort values $v$ has a height equal to the measure of the applicants who are tied at that value, and so $\gamma$ for those applicants can be filled in using the CDF of $\Gamma(\omega)$ restricted to the applicants $\omega$ who are tied. Thus, the range of $\gamma$ is uniform on $[0,1]$.
%
{and}
%	Then, given
	a ranking probability function $\lambda$, an \textit{equilibrium} is a set of effort levels and post-effort ranking rewards for each applicant, $\{e(\theta_\pre),  \lambda(\theta_\post(\theta_\pre))\}$ such that, for all
	$\omega$, %$\theta_\pre$,
	\begin{align*}
		e(\theta_\pre(\omega)) &\in \argmax_{e} \iwelf\left(e, \lambda\left(\gamma\left(\omega, v\left(e,\theta_\pre(\omega)\right)\right)\right)\right) %\triangleq \argmax_{e} \left[ \lambda(\gamma(v(e,\theta_\pre)) - \cost(e) \right].
		\\
		\theta_\post(\theta_\pre(\omega)) &= \gamma(\omega, v(e(\theta_\pre(\omega)),\theta_\pre(\omega))),
	\end{align*}
	{where $\gamma(\omega,v)$ is the ranking induced by the CDF of the scores resulting from effort levels $\{e(\theta_\pre)\}$ and the tie-breaking function $\Gamma$.\footnote{See Remark~\ref{rem:gamma-effort} on $\gamma$ and the set of efforts.%Formally, $\gamma$ depends on the set of efforts. Given a fixed set and $\gamma$, for each applicant $\omega$ the first condition considers the counter-factual ranking of $\omega$ with different post-effort values but using the same ranking function. As defined, $\gamma$ yields a uniform distribution of ranks with such measure $0$ changes. In Appendix \Cref{lem:deviationsmeasure0}, we further prove that two effort sets equal up to sets of measure $0$ induce the same ranking function $\gamma$, and so the condition is consistent.%, and so we set it at the $\gamma$ induced by the effort levels $\{e(\theta_\pre)\}$.%, and so it does not vary.
	}
	}
	%In particular, we denote $\theta_\post(\theta_\pre) := \gamma(v(e(\theta_\pre),\theta_\pre))$.
	%\lledit{and for any $\theta_\pre, \theta'_\pre$, we have $\lambda(\theta_\post(\theta_\pre)) > \lambda(\theta_\post(\theta'_\pre)) \implies v(e(\theta_\pre),\theta_\pre) > v(e(\theta_\pre'),\theta_\pre')$.}
\end{definition}
\end{ReviewVerbatim}

% report-context-presentation-sha256: 90f9116e292dd71fddfd6b445ca0b060a481d026855bdb9725f2833c2cf62a31
\paragraph{Settled source reading}
\begin{sloppypar}
Equilibrium means that almost every applicant best responds against every feasible deviation;\allowbreak{} uniqueness is up to a null set.\allowbreak{}
\end{sloppypar}
\paragraph{Settled source reading}
\begin{sloppypar}
Applicants know the fixed priority used at score-\allowbreak{}boundary ties.\allowbreak{} This clarifies the stated model, not a corrected admission policy or a new economic assumption.\allowbreak{}
\end{sloppypar}

\paragraph{Lean-expanded paper semantic target}
\begin{ReviewVerbatim}
type:
(ℝ → ℝ) → (ℝ → ℝ) → (ℝ → ℝ) → (ℝ → ℝ) → ℕ → (ℕ → ℝ) → (ℕ → ℝ) → (ℝ → ℝ) → (ℝ → ℝ) → Prop

value:
fun (cost production skill tie : ℝ → ℝ) (n : ℕ) (cutoff reward : ℕ → ℝ) (effort score : ℝ → ℝ) =>
  Measurable score ∧
    ∀ᵐ (t : ℝ) ∂LBG22StrategicRanking.unitRankMeasure,
      LBG22StrategicRanking.SourceFiniteBestResponseAt cost production skill score tie
        (fun (i : Fin (n + 1)) => cutoff ↑i) reward t (effort t)
\end{ReviewVerbatim}

\paragraph{Direct retained dependencies}
\begin{itemize}\raggedright
\item \hyperlink{paper-prerequisite-519e321cd0957049}{LBG22\allowbreak{}Strategic\allowbreak{}Ranking.\allowbreak{}Source\allowbreak{}Finite\allowbreak{}Best\allowbreak{}Response\allowbreak{}At}
\item \hyperlink{governing-declaration-678cf48d549cb620}{LBG22\allowbreak{}Strategic\allowbreak{}Ranking.\allowbreak{}unit\allowbreak{}Rank\allowbreak{}Measure}
\end{itemize}
\paragraph{Recorded source-to-declaration screening}
\textbf{Verdict:} \texttt{matches}
\\
{\raggedright
\textbf{Reason:} The target requires a measurable realized score and one full-\allowbreak{}measure set of applicants whose nonnegative effort maximizes finite-\allowbreak{}policy reward minus cost against every nonnegative deviation under a fixed score law and priority.\allowbreak{} This is the source equilibrium definition under the displayed approved strong quantifier order and strict-\allowbreak{}percentile representation bridge.\allowbreak{}
\par}
\reviewmatch{Matches source input}
\noindent\textbf{Reviewer annotation}\par
\reviewerbox
\clearpage
\hypertarget{paper-prerequisite-9bbcc12f4a1736bf}{}
\subsection*{\small\ttfamily\raggedright LBG22\allowbreak{}Strategic\allowbreak{}Ranking.\allowbreak{}Ranking\allowbreak{}Primitives}
\noindent\textbf{Paper-local declaration:}\par
{\footnotesize\ttfamily\raggedright LBG22\allowbreak{}Strategic\allowbreak{}Ranking.\allowbreak{}Ranking\allowbreak{}Primitives\par}
\textbf{Source locator:} {\footnotesize\raggedright cited publication, lines 1008-\allowbreak{}1024\par}
\paragraph{Verbatim paper-source connection}
\begin{ReviewVerbatim}
\parbold{Applicants}
There is a unit mass of \textit{applicants}, indexed by an observed index $\omega \in [0,1]$ distributed uniformly.\footnote{The index $\omega$ should be interpreted as each applicant's ``name,'' uncorrelated with skill, used solely for tie-breaking.}  Each applicant has a latent (unobserved) skill level represented by some measurable function of $\omega$.
We assume that {the distribution of the skills} has no atoms and that the CDF of this distribution is strictly increasing.
Using the CDF to map the skill of an applicant to a rank in~$[0,1]$, each applicant gets an (unobserved) rank $\theta_\pre =\theta_\pre(\omega)$ which by our assumption on the CDF is again uniformly distributed in $[0, 1]$ (the higher the rank the better).  With this setup, the skill of an applicant with rank $\theta_\pre$ can be written as $f(\theta_\pre)$ where $f$ is a strictly increasing, continuous function.  {It will be notationally convenient to label applicants by their rank $\theta_\pre$, though the reader should note that in contrast to $\omega$, $\theta_\pre(\omega)$ is assumed to be unobservable.}

%3. Each applicant exerts some effort e, at cost \cost(e).


Each applicant chooses an \textit{effort} level $e \ge 0$, the result of which is an observed, post-effort \textit{score}, {$v=v(e,\theta_\pre) = {g(e)}\cdot{f(\theta_\pre)}$}. \lledit{In other words, we assume the post-effort score is the product of two components, associated with the effort and the pre-effort skill respectively.}
{We assume} that the effort transfer function $g: [0, \infty) \mapsto [0, \infty)$ is a continuous, concave, strictly increasing function, representing that marginal effort improves one's score but has diminishing returns.
{%Since applicants can be labelled by their skill level $f(\theta_\pre)$, or equivalently, by their rank $\theta_{\pre}$,
{T}he strategies of the applicants can {then} be described by a function $\theta_\pre \mapsto e(\theta_\pre)$.
%\cb{$\omega\mapsto e(\omega)$.}%, leading to a distribution over post-effort scores $v=v(\theta_\pre))$ via the uniform distribution of $\theta_\pre$ in $[0,1]$.
}
  Each {applicant} is then ranked according to their score $v$%(where we break ties uniformly at random, if they appear)
, resulting in a \textit{post-effort rank} $\theta_\post$.
Note that the ranking $\theta_\post$ is slightly less trivial than a ranking of the skills, since the scores might have ties, which have to be resolved.

[Next verbatim source excerpt]

\parbold{Individual applicant welfare and equilibrium} Given the designer's function $\lambda$ and the effort levels of other applicants, each applicant chooses effort $e$ to maximize their individual welfare, % $W$,
\begin{equation*}
	\iwelf(e, \lambda(\theta_\post)) = \lambda(\theta_\post) - \cost(e),
\end{equation*}
where the effort cost function $\cost$ is non-negative, continuous, and strictly convex {on $[0,\infty)$, with $e_0:=\argmin_e \cost(e)$ and $p(e_0) = 0$.}
%\lnote{also want to assume that $\cost$ attains its minimum at $0$?}
%Let $e_0:=\argmin_e \cost(e)$ and $p(e_0) = 0$.

Applicants are assumed not to personally benefit from increasing their score $v$, except through the corresponding increase in their rank and reward. While the definition of $\cost$ does not preclude the applicant receiving any intrinsic benefit from exerting effort, we assume that the \emph{net} benefit from effort is non-positive.
\end{ReviewVerbatim}



\paragraph{Lean-expanded paper semantic target}
\begin{ReviewVerbatim}
type:
(ℝ → ℝ) → (ℝ → ℝ) → (ℝ → ℝ) → ℝ → Prop

value:
fun (cost production skill : ℝ → ℝ) (baseline : ℝ) =>
  0 ≤ baseline ∧
    ContinuousOn cost (Set.Ici 0) ∧
      StrictConvexOn ℝ (Set.Ici 0) cost ∧
        (∀ e ∈ Set.Ici 0, 0 ≤ cost e) ∧
          cost baseline = 0 ∧
            ContinuousOn production (Set.Ici 0) ∧
              StrictMonoOn production (Set.Ici 0) ∧
                ConcaveOn ℝ (Set.Ici 0) production ∧
                  0 ≤ production 0 ∧ ContinuousOn skill (Set.Icc 0 1) ∧ StrictMonoOn skill (Set.Icc 0 1) ∧ 0 ≤ skill 0
\end{ReviewVerbatim}


\paragraph{Recorded source-to-declaration screening}
\textbf{Verdict:} \texttt{matches}
\\
{\raggedright
\textbf{Reason:} The target reproduces the source scalar primitives:\allowbreak{} a nonnegative baseline where nonnegative strictly convex cost attains zero, continuous concave strictly increasing nonnegative production, and a continuous strictly increasing nonnegative skill quantile on [0,1].\allowbreak{} The displayed support confirms all uses remain on these source domains.\allowbreak{}
\par}
\reviewmatch{Matches source input}
\noindent\textbf{Reviewer annotation}\par
\reviewerbox
\clearpage
\hypertarget{paper-prerequisite-2dc31b75c00d0d1b}{}
\subsection*{\small\ttfamily\raggedright LBG22\allowbreak{}Strategic\allowbreak{}Ranking.\allowbreak{}Score\allowbreak{}Priority\allowbreak{}Best\allowbreak{}Response\allowbreak{}At}
\noindent\textbf{Paper-local declaration:}\par
{\footnotesize\ttfamily\raggedright LBG22\allowbreak{}Strategic\allowbreak{}Ranking.\allowbreak{}Score\allowbreak{}Priority\allowbreak{}Best\allowbreak{}Response\allowbreak{}At\par}
\textbf{Source locator:} {\footnotesize\raggedright cited publication, lines 1061-\allowbreak{}1069\par}
\paragraph{Verbatim paper-source connection}
\begin{ReviewVerbatim}
\parbold{Individual applicant welfare and equilibrium} Given the designer's function $\lambda$ and the effort levels of other applicants, each applicant chooses effort $e$ to maximize their individual welfare, % $W$,
\begin{equation*}
	\iwelf(e, \lambda(\theta_\post)) = \lambda(\theta_\post) - \cost(e),
\end{equation*}
where the effort cost function $\cost$ is non-negative, continuous, and strictly convex {on $[0,\infty)$, with $e_0:=\argmin_e \cost(e)$ and $p(e_0) = 0$.}
%\lnote{also want to assume that $\cost$ attains its minimum at $0$?}
%Let $e_0:=\argmin_e \cost(e)$ and $p(e_0) = 0$.

Applicants are assumed not to personally benefit from increasing their score $v$, except through the corresponding increase in their rank and reward. While the definition of $\cost$ does not preclude the applicant receiving any intrinsic benefit from exerting effort, we assume that the \emph{net} benefit from effort is non-positive.

[Next verbatim source excerpt]

\begin{definition} [Equilibrium]\label{def:equi}
{Given a}
%Let the
tie-breaking function $\Gamma(\omega)$
%be a pre-announced ordering over index $\omega$, to break ties in post-effort value. Then, let $\gamma(\omega, v)$ be the 1-1 ranking function defined as the CDF of post-effort values $v$ except where there are ties in $v$, in which case the ordering $\Gamma$ is used to produce unique ranks such that the range of $\gamma$ is uniform on $[0, 1]$.\footnote{More formally, note that each discontinuity in the CDF of post-effort values $v$ has a height equal to the measure of the applicants who are tied at that value, and so $\gamma$ for those applicants can be filled in using the CDF of $\Gamma(\omega)$ restricted to the applicants $\omega$ who are tied. Thus, the range of $\gamma$ is uniform on $[0,1]$.
%
{and}
%	Then, given
	a ranking probability function $\lambda$, an \textit{equilibrium} is a set of effort levels and post-effort ranking rewards for each applicant, $\{e(\theta_\pre),  \lambda(\theta_\post(\theta_\pre))\}$ such that, for all
	$\omega$, %$\theta_\pre$,
	\begin{align*}
		e(\theta_\pre(\omega)) &\in \argmax_{e} \iwelf\left(e, \lambda\left(\gamma\left(\omega, v\left(e,\theta_\pre(\omega)\right)\right)\right)\right) %\triangleq \argmax_{e} \left[ \lambda(\gamma(v(e,\theta_\pre)) - \cost(e) \right].
		\\
		\theta_\post(\theta_\pre(\omega)) &= \gamma(\omega, v(e(\theta_\pre(\omega)),\theta_\pre(\omega))),
	\end{align*}
	{where $\gamma(\omega,v)$ is the ranking induced by the CDF of the scores resulting from effort levels $\{e(\theta_\pre)\}$ and the tie-breaking function $\Gamma$.\footnote{See Remark~\ref{rem:gamma-effort} on $\gamma$ and the set of efforts.%Formally, $\gamma$ depends on the set of efforts. Given a fixed set and $\gamma$, for each applicant $\omega$ the first condition considers the counter-factual ranking of $\omega$ with different post-effort values but using the same ranking function. As defined, $\gamma$ yields a uniform distribution of ranks with such measure $0$ changes. In Appendix \Cref{lem:deviationsmeasure0}, we further prove that two effort sets equal up to sets of measure $0$ induce the same ranking function $\gamma$, and so the condition is consistent.%, and so we set it at the $\gamma$ induced by the effort levels $\{e(\theta_\pre)\}$.%, and so it does not vary.
	}
	}
	%In particular, we denote $\theta_\post(\theta_\pre) := \gamma(v(e(\theta_\pre),\theta_\pre))$.
	%\lledit{and for any $\theta_\pre, \theta'_\pre$, we have $\lambda(\theta_\post(\theta_\pre)) > \lambda(\theta_\post(\theta'_\pre)) \implies v(e(\theta_\pre),\theta_\pre) > v(e(\theta_\pre'),\theta_\pre')$.}
\end{definition}
\end{ReviewVerbatim}

% report-context-presentation-sha256: 90f9116e292dd71fddfd6b445ca0b060a481d026855bdb9725f2833c2cf62a31
\paragraph{Settled source reading}
\begin{sloppypar}
Equilibrium means that almost every applicant best responds against every feasible deviation;\allowbreak{} uniqueness is up to a null set.\allowbreak{}
\end{sloppypar}
\paragraph{Settled source reading}
\begin{sloppypar}
Applicants know the fixed priority used at score-\allowbreak{}boundary ties.\allowbreak{} This clarifies the stated model, not a corrected admission policy or a new economic assumption.\allowbreak{}
\end{sloppypar}

\paragraph{Lean-expanded paper semantic target}
\begin{ReviewVerbatim}
type:
{α : Type u_1} →
  [inst : MeasurableSpace α] →
    MeasureTheory.Measure α →
      (ℝ → ℝ) → (ℝ → ℝ) → (α → ℝ) → (α → ℝ) → (α → ℝ) → {n : ℕ} → (Fin (n + 1) → ℝ) → (ℕ → ℝ) → α → ℝ → Prop

value:
fun {α : Type u_1} [MeasurableSpace α] (μ : MeasureTheory.Measure α) (cost production : ℝ → ℝ) (skill score tie : α → ℝ)
    {n : ℕ} (cutoff : Fin (n + 1) → ℝ) (reward : ℕ → ℝ) (x : α) (e : ℝ) =>
  0 ≤ e ∧
    score x = production e * skill x ∧
      ∀ (d : ℝ),
        0 ≤ d →
          reward ↑(LBG22StrategicRanking.scorePriorityBand μ score tie cutoff (production d * skill x) (tie x)) -
              cost d ≤
            reward ↑(LBG22StrategicRanking.scorePriorityBand μ score tie cutoff (score x) (tie x)) - cost e
\end{ReviewVerbatim}

\paragraph{Direct retained dependencies}
\begin{itemize}\raggedright
\item \hyperlink{governing-declaration-29f9270f1ab8f757}{LBG22\allowbreak{}Strategic\allowbreak{}Ranking.\allowbreak{}score\allowbreak{}Priority\allowbreak{}Band}
\end{itemize}
\paragraph{Recorded source-to-declaration screening}
\textbf{Verdict:} \texttt{matches}
\\
{\raggedright
\textbf{Reason:} The target requires nonnegative effort, the source production-\allowbreak{}times-\allowbreak{}skill score identity, and reward minus cost at least as large as for every nonnegative deviation, with the score population and priority held fixed.\allowbreak{} This matches the individual welfare and approved all-\allowbreak{}deviation best-\allowbreak{}response reading.\allowbreak{}
\par}
\reviewmatch{Matches source input}
\noindent\textbf{Reviewer annotation}\par
\reviewerbox
\clearpage
\hypertarget{paper-prerequisite-f396601c268691e5}{}
\subsection*{\small\ttfamily\raggedright LBG22\allowbreak{}Strategic\allowbreak{}Ranking.\allowbreak{}Ranking\allowbreak{}Score\allowbreak{}Priority\allowbreak{}Equilibrium}
\noindent\textbf{Paper-local declaration:}\par
{\footnotesize\ttfamily\raggedright LBG22\allowbreak{}Strategic\allowbreak{}Ranking.\allowbreak{}Ranking\allowbreak{}Score\allowbreak{}Priority\allowbreak{}Equilibrium\par}
\textbf{Source locator:} {\footnotesize\raggedright cited publication, lines 1077-\allowbreak{}1097\par}
\paragraph{Verbatim paper-source connection}
\begin{ReviewVerbatim}
\begin{definition} [Equilibrium]\label{def:equi}
{Given a}
%Let the
tie-breaking function $\Gamma(\omega)$
%be a pre-announced ordering over index $\omega$, to break ties in post-effort value. Then, let $\gamma(\omega, v)$ be the 1-1 ranking function defined as the CDF of post-effort values $v$ except where there are ties in $v$, in which case the ordering $\Gamma$ is used to produce unique ranks such that the range of $\gamma$ is uniform on $[0, 1]$.\footnote{More formally, note that each discontinuity in the CDF of post-effort values $v$ has a height equal to the measure of the applicants who are tied at that value, and so $\gamma$ for those applicants can be filled in using the CDF of $\Gamma(\omega)$ restricted to the applicants $\omega$ who are tied. Thus, the range of $\gamma$ is uniform on $[0,1]$.
%
{and}
%	Then, given
	a ranking probability function $\lambda$, an \textit{equilibrium} is a set of effort levels and post-effort ranking rewards for each applicant, $\{e(\theta_\pre),  \lambda(\theta_\post(\theta_\pre))\}$ such that, for all
	$\omega$, %$\theta_\pre$,
	\begin{align*}
		e(\theta_\pre(\omega)) &\in \argmax_{e} \iwelf\left(e, \lambda\left(\gamma\left(\omega, v\left(e,\theta_\pre(\omega)\right)\right)\right)\right) %\triangleq \argmax_{e} \left[ \lambda(\gamma(v(e,\theta_\pre)) - \cost(e) \right].
		\\
		\theta_\post(\theta_\pre(\omega)) &= \gamma(\omega, v(e(\theta_\pre(\omega)),\theta_\pre(\omega))),
	\end{align*}
	{where $\gamma(\omega,v)$ is the ranking induced by the CDF of the scores resulting from effort levels $\{e(\theta_\pre)\}$ and the tie-breaking function $\Gamma$.\footnote{See Remark~\ref{rem:gamma-effort} on $\gamma$ and the set of efforts.%Formally, $\gamma$ depends on the set of efforts. Given a fixed set and $\gamma$, for each applicant $\omega$ the first condition considers the counter-factual ranking of $\omega$ with different post-effort values but using the same ranking function. As defined, $\gamma$ yields a uniform distribution of ranks with such measure $0$ changes. In Appendix \Cref{lem:deviationsmeasure0}, we further prove that two effort sets equal up to sets of measure $0$ induce the same ranking function $\gamma$, and so the condition is consistent.%, and so we set it at the $\gamma$ induced by the effort levels $\{e(\theta_\pre)\}$.%, and so it does not vary.
	}
	}
	%In particular, we denote $\theta_\post(\theta_\pre) := \gamma(v(e(\theta_\pre),\theta_\pre))$.
	%\lledit{and for any $\theta_\pre, \theta'_\pre$, we have $\lambda(\theta_\post(\theta_\pre)) > \lambda(\theta_\post(\theta'_\pre)) \implies v(e(\theta_\pre),\theta_\pre) > v(e(\theta_\pre'),\theta_\pre')$.}
\end{definition}
\end{ReviewVerbatim}

% report-context-presentation-sha256: 90f9116e292dd71fddfd6b445ca0b060a481d026855bdb9725f2833c2cf62a31
\paragraph{Settled source reading}
\begin{sloppypar}
Equilibrium means that almost every applicant best responds against every feasible deviation;\allowbreak{} uniqueness is up to a null set.\allowbreak{}
\end{sloppypar}
\paragraph{Settled source reading}
\begin{sloppypar}
Applicants know the fixed priority used at score-\allowbreak{}boundary ties.\allowbreak{} This clarifies the stated model, not a corrected admission policy or a new economic assumption.\allowbreak{}
\end{sloppypar}

\paragraph{Lean-expanded paper semantic target}
\begin{ReviewVerbatim}
type:
(ℝ → ℝ) → (ℝ → ℝ) → (ℝ → ℝ) → (ℝ → ℝ) → ℕ → (ℕ → ℝ) → (ℕ → ℝ) → (ℝ → ℝ) → (ℝ → ℝ) → Prop

value:
fun (cost production skill tie : ℝ → ℝ) (n : ℕ) (cutoff reward : ℕ → ℝ) (effort score : ℝ → ℝ) =>
  Measurable score ∧
    ∀ᵐ (t : ℝ) ∂LBG22StrategicRanking.unitRankMeasure,
      LBG22StrategicRanking.ScorePriorityBestResponseAt LBG22StrategicRanking.unitRankMeasure cost production skill
        score tie (fun (i : Fin (n + 1)) => cutoff ↑i) reward t (effort t)
\end{ReviewVerbatim}

\paragraph{Direct retained dependencies}
\begin{itemize}\raggedright
\item \hyperlink{paper-prerequisite-2dc31b75c00d0d1b}{LBG22\allowbreak{}Strategic\allowbreak{}Ranking.\allowbreak{}Score\allowbreak{}Priority\allowbreak{}Best\allowbreak{}Response\allowbreak{}At}
\item \hyperlink{governing-declaration-678cf48d549cb620}{LBG22\allowbreak{}Strategic\allowbreak{}Ranking.\allowbreak{}unit\allowbreak{}Rank\allowbreak{}Measure}
\end{itemize}
\paragraph{Recorded source-to-declaration screening}
\textbf{Verdict:} \texttt{matches}
\\
{\raggedright
\textbf{Reason:} The target is the source equilibrium expressed directly with actual score-\allowbreak{}priority boundary admission:\allowbreak{} measurable score and almost-\allowbreak{}everywhere best response against all deviations using one fixed priority key.\allowbreak{} It matches the approved equilibrium and tie-\allowbreak{}boundary conventions.\allowbreak{}
\par}
\reviewmatch{Matches source input}
\noindent\textbf{Reviewer annotation}\par
\reviewerbox
\clearpage
\hypertarget{paper-prerequisite-38434fa78df931b7}{}
\subsection*{\small\ttfamily\raggedright LBG22\allowbreak{}Strategic\allowbreak{}Ranking.\allowbreak{}Score\allowbreak{}Priority\allowbreak{}Above\allowbreak{}Cutoff}
\noindent\textbf{Paper-local declaration:}\par
{\footnotesize\ttfamily\raggedright LBG22\allowbreak{}Strategic\allowbreak{}Ranking.\allowbreak{}Score\allowbreak{}Priority\allowbreak{}Above\allowbreak{}Cutoff\par}
\textbf{Source locator:} {\footnotesize\raggedright cited publication, lines 1026-\allowbreak{}1030\par}
\paragraph{Verbatim paper-source connection}
\begin{ReviewVerbatim}
{\textbf{Tie Breaking.}
Given a choice of strategies $\theta_\pre \mapsto e(\theta_\pre)$, let $F$ be the CDF of the scores $v(e(\theta_\pre),\theta_\pre)$.
Since atoms for the distribution of $v$ would lead to ties for ranks defined as $F(v)$, we will use the labels of the applicants to break ties with the help of a (publicly announced) tie-breaking function $\Gamma(\omega)$, defined, e.g., via a collision free hash of the applicant names. Here  we require  that $\Gamma$  is a measurable function from $[0,1]$ to $[0,1]$ that maps different applicant labels to different values.
We then use $\Gamma$ to resolve the atoms of $F$, leading to a ranking function $v\mapsto \gamma(\omega, v)$ which is equal to $F(v)$ except when $v$ is an atom of the score distribution, {in which case it takes values in the ``gap interval'' $[F_-(v),  F(v)]$, where $F_-(v)$ is the left limit of $F$ at $v$.}  We  construct $\gamma$ in such a way that it
gives the uniform distribution for $\theta_\post(\omega)$ if we set  $\theta_\post(\omega)=\gamma(\omega,v(e(\theta_\pre(\omega)),\theta_\pre(\omega)))$.\footnote{See Remark~\ref{rem:gamma} for details on this construction.}

[Next verbatim source excerpt]

\begin{remark}[Construction of the $\gamma$ map]\label{rem:gamma}
	Formally,  if $v_0$ is an atom of the score distribution and $\Omega_{v_0}$ is the set of tied applicants with score $v_0$,
	then the discontinuity of $F$ {at  $v_0$} has height equal to the measure of $\Omega_{v_0}$, and so $\gamma$ for those applicants can be filled by using the CDF of $\Gamma(\omega)$ restricted to
	%the applicants $\omega$ who are tied
	$\Omega_{v_0}$; this gives a distribution for $\gamma(\omega, v_0)$ that is uniform over
	the gap interval $[ F_-(v_0),  F(v_0)]$ when restricted to
	%\omega\in
	$\Omega_{v_0}$, and hence leads to the claimed uniform distribution of $\theta_\post(\omega)$.
	Finally, we define $\gamma(\omega, v)$ for applicants  $\omega\notin\Omega_{v_0}$ by ``slotting them in'' in such a way that for a pair $\omega,\omega'$ with exactly one member in $\Omega_{v_0}$ and $\Gamma(\omega)<\Gamma(\omega')$ we have that
	$\gamma(v_0,\omega)\leq\gamma(v_0,\omega')$.
	Since the image of $\Omega_{v_0}$ is by construction dense in the gap interval, this uniquely determines $\gamma(v_0,\omega)$ for all applicants $\omega$.

	Note that for a given applicant $\omega$, the map $v\mapsto \gamma(\omega,v)$ is not necessarily 1-1; indeed, if $F$ is constant on an interval $I$, all $v\in I$ lead to the same rank.  But this only effects regions where the distribution of $v$ has no mass, and thus will not cause any issues.
\end{remark}
\end{ReviewVerbatim}



\paragraph{Lean-expanded paper semantic target}
\begin{ReviewVerbatim}
type:
{α : Type u_1} → [inst : MeasurableSpace α] → MeasureTheory.Measure α → (α → ℝ) → (α → ℝ) → ℝ → ℝ → ℝ → Prop

value:
fun {α : Type u_1} [MeasurableSpace α] (μ : MeasureTheory.Measure α) (score tie : α → ℝ) (cutoff v priority : ℝ) =>
  ∀ (w q : ℝ), v < w ∨ v = w ∧ priority < q → cutoff < LBG22StrategicRanking.scorePriorityRank μ score tie w q
\end{ReviewVerbatim}

\paragraph{Direct retained dependencies}
\begin{itemize}\raggedright
\item \hyperlink{governing-declaration-9e8db797d20dba91}{LBG22\allowbreak{}Strategic\allowbreak{}Ranking.\allowbreak{}score\allowbreak{}Priority\allowbreak{}Rank}
\end{itemize}
\paragraph{Recorded source-to-declaration screening}
\textbf{Verdict:} \texttt{matches}
\\
{\raggedright
\textbf{Reason:} For a proposed score and priority, the target says every lexicographically larger score-\allowbreak{}priority pair has rank strictly above the cutoff.\allowbreak{} With the displayed lower-\allowbreak{}contour rank, this is precisely membership above the actual score-\allowbreak{}priority boundary in the source tie-\allowbreak{}breaking construction.\allowbreak{}
\par}
\reviewmatch{Matches source input}
\noindent\textbf{Reviewer annotation}\par
\reviewerbox
\clearpage
\hypertarget{paper-prerequisite-f88aac58507e44ca}{}
\subsection*{\small\ttfamily\raggedright LBG22\allowbreak{}Strategic\allowbreak{}Ranking.\allowbreak{}source\allowbreak{}Two\allowbreak{}Level\allowbreak{}Cutoff}
\noindent\textbf{Paper-local declaration:}\par
{\footnotesize\ttfamily\raggedright LBG22\allowbreak{}Strategic\allowbreak{}Ranking.\allowbreak{}source\allowbreak{}Two\allowbreak{}Level\allowbreak{}Cutoff\par}
\textbf{Source locator:} {\footnotesize\raggedright cited publication, lines 102-\allowbreak{}103\par}
\paragraph{Verbatim paper-source connection}
\begin{ReviewVerbatim}
\begin{definition}[Two-level policy]\label{assump:2-level-fixed-cap} In our baseline two-level function parameterized by cut-off $c \in (0,1-\rho]$, each applicant with post-effort rank  $\theta_\post \geq c$ is admitted with probability $\ell_1 =\frac{\rho}{1-c} \in (0,1]$. Others are rejected, $\ell_0 = 0$.
\end{definition}
\end{ReviewVerbatim}

% report-context-presentation-sha256: 827a19afe0274b1faef13e344ddbcfe40d9841f2ff34b92ac9d53c0f83e29f56
\paragraph{Settled source reading}
\begin{sloppypar}
Applicants know the fixed priority used at score-\allowbreak{}boundary ties.\allowbreak{} This clarifies the stated model, not a corrected admission policy or a new economic assumption.\allowbreak{}
\end{sloppypar}

\paragraph{Lean-expanded paper semantic target}
\begin{ReviewVerbatim}
type:
ℝ → ℕ → ℝ

value:
fun (c : ℝ) (k : ℕ) => if k = 0 then 0 else if k = 1 then c else 1
\end{ReviewVerbatim}


\paragraph{Recorded source-to-declaration screening}
\textbf{Verdict:} \texttt{matches}
\\
{\raggedright
\textbf{Reason:} The target encodes two-\allowbreak{}level policy cutoffs 0, c, and 1.\allowbreak{} With the displayed band representation and fixed-\allowbreak{}priority convention, this is the source partition at cutoff c.\allowbreak{}
\par}
\reviewmatch{Matches source input}
\noindent\textbf{Reviewer annotation}\par
\reviewerbox
\clearpage
\hypertarget{paper-prerequisite-30d97cd43eb79017}{}
\subsection*{\small\ttfamily\raggedright LBG22\allowbreak{}Strategic\allowbreak{}Ranking.\allowbreak{}source\allowbreak{}Two\allowbreak{}Level\allowbreak{}Reward}
\noindent\textbf{Paper-local declaration:}\par
{\footnotesize\ttfamily\raggedright LBG22\allowbreak{}Strategic\allowbreak{}Ranking.\allowbreak{}source\allowbreak{}Two\allowbreak{}Level\allowbreak{}Reward\par}
\textbf{Source locator:} {\footnotesize\raggedright cited publication, lines 102-\allowbreak{}103\par}
\paragraph{Verbatim paper-source connection}
\begin{ReviewVerbatim}
\begin{definition}[Two-level policy]\label{assump:2-level-fixed-cap} In our baseline two-level function parameterized by cut-off $c \in (0,1-\rho]$, each applicant with post-effort rank  $\theta_\post \geq c$ is admitted with probability $\ell_1 =\frac{\rho}{1-c} \in (0,1]$. Others are rejected, $\ell_0 = 0$.
\end{definition}
\end{ReviewVerbatim}

% report-context-presentation-sha256: 827a19afe0274b1faef13e344ddbcfe40d9841f2ff34b92ac9d53c0f83e29f56
\paragraph{Settled source reading}
\begin{sloppypar}
Applicants know the fixed priority used at score-\allowbreak{}boundary ties.\allowbreak{} This clarifies the stated model, not a corrected admission policy or a new economic assumption.\allowbreak{}
\end{sloppypar}

\paragraph{Lean-expanded paper semantic target}
\begin{ReviewVerbatim}
type:
ℝ → ℝ → ℕ → ℝ

value:
fun (rho c : ℝ) (k : ℕ) => if k = 0 then 0 else rho / (1 - c)
\end{ReviewVerbatim}


\paragraph{Recorded source-to-declaration screening}
\textbf{Verdict:} \texttt{matches}
\\
{\raggedright
\textbf{Reason:} The target assigns reward zero below the two-\allowbreak{}level cutoff and rho/\allowbreak{}(1-\allowbreak{}c) above it, exactly the source capacity-\allowbreak{}preserving two-\allowbreak{}level reward formula.\allowbreak{}
\par}
\reviewmatch{Matches source input}
\noindent\textbf{Reviewer annotation}\par
\reviewerbox
\clearpage
\hypertarget{paper-prerequisite-c64406972a4c963f}{}
\subsection*{\small\ttfamily\raggedright LBG22\allowbreak{}Strategic\allowbreak{}Ranking.\allowbreak{}Score\allowbreak{}Priority\allowbreak{}Hard\allowbreak{}Budget\allowbreak{}Best\allowbreak{}Response\allowbreak{}At}
\noindent\textbf{Paper-local declaration:}\par
{\footnotesize\ttfamily\raggedright LBG22\allowbreak{}Strategic\allowbreak{}Ranking.\allowbreak{}Score\allowbreak{}Priority\allowbreak{}Hard\allowbreak{}Budget\allowbreak{}Best\allowbreak{}Response\allowbreak{}At\par}
\textbf{Source locator:} {\footnotesize\raggedright cited publication, lines 1799-\allowbreak{}1806\par}
\paragraph{Verbatim paper-source connection}
\begin{ReviewVerbatim}
In our simplified setting, there are two skills $M$ and $U$: $M$ has a measurable score $v^M$ and $U$ has an unmeasurable score $v^U$. For example, $M$ could be scholastic achievement as measured by SAT scores, and $U$ could be ``creativity", a personal quality that is valued by the school but is not directly measurable.
%
Since the school cannot observe $v^U$, its admission policy is based on $\theta_\post^M$ only, that is, $\alpha_U = 0$ and $\alpha_M = 1$.

We further assume that each applicant has a fixed effort budget of $B > 0$, and is intrinsically motivated to exert effort in the unmeasurable skill $U$; in fact, they will always exert effort $e^U = B - e^M$. Formally this corresponds to the effort cost function \[\cost^m(e_M, e_U) = \cost(e^M) - (\max(0,B-(e^M+e^U)))^2\] where $\cost$ is convex and increasing. We can write the applicant's individual welfare as:
\begin{equation*}
W(e^M, e^U, \lambda(\theta_\post^M)) = \lambda(\theta_\post^M) - \cost^m(e_M, e_U).
\end{equation*}
\end{ReviewVerbatim}

% report-context-presentation-sha256: ce5679112671da8414975164981e5740bcad9608aa7765c30f1069706683e7b1
\paragraph{Settled source reading}
\begin{sloppypar}
Equilibrium means that almost every applicant best responds against every feasible deviation;\allowbreak{} uniqueness is up to a null set.\allowbreak{}
\end{sloppypar}
\paragraph{Settled source reading}
\begin{sloppypar}
The multitask model has nonnegative effort on both tasks summing to the stated budget.\allowbreak{} This is distinct from the unrestricted total effort in the multidimensional ranking model.\allowbreak{}
\end{sloppypar}
\paragraph{Settled source reading}
\begin{sloppypar}
Applicants know the fixed priority used at score-\allowbreak{}boundary ties.\allowbreak{} This clarifies the stated model, not a corrected admission policy or a new economic assumption.\allowbreak{}
\end{sloppypar}

\paragraph{Lean-expanded paper semantic target}
\begin{ReviewVerbatim}
type:
{α : Type u_1} →
  [inst : MeasurableSpace α] →
    (μ : MeasureTheory.Measure α) →
      [MeasureTheory.IsFiniteMeasure μ] → (ℝ → ℝ) → (ℝ → ℝ) → (α → ℝ) → (α → ℝ) → (α → ℝ) → ℝ → ℝ → ℝ → α → ℝ → ℝ → Prop

value:
fun {α : Type u_1} [MeasurableSpace α] (μ : MeasureTheory.Measure α) [MeasureTheory.IsFiniteMeasure μ]
    (cost production : ℝ → ℝ) (skill score tie : α → ℝ) (B rho c : ℝ) (x : α) (eM eU : ℝ) =>
  0 ≤ eM ∧
    0 ≤ eU ∧
      eM + eU = B ∧
        score x = production eM * skill x ∧
          ∀ (dM dU : ℝ),
            0 ≤ dM →
              0 ≤ dU →
                dM + dU = B →
                  LBG22StrategicRanking.sourceTwoLevelReward rho c
                        ↑(LBG22StrategicRanking.scorePriorityBand μ score tie
                            (fun (i : Fin 2) => LBG22StrategicRanking.sourceTwoLevelCutoff c ↑i)
                            (production dM * skill x) (tie x)) -
                      LBG22StrategicRanking.sourceMultitaskCost cost B dM dU ≤
                    LBG22StrategicRanking.sourceTwoLevelReward rho c
                        ↑(LBG22StrategicRanking.scorePriorityBand μ score tie
                            (fun (i : Fin 2) => LBG22StrategicRanking.sourceTwoLevelCutoff c ↑i) (score x) (tie x)) -
                      LBG22StrategicRanking.sourceMultitaskCost cost B eM eU
\end{ReviewVerbatim}

\paragraph{Direct retained dependencies}
\begin{itemize}\raggedright
\item \hyperlink{governing-declaration-29f9270f1ab8f757}{LBG22\allowbreak{}Strategic\allowbreak{}Ranking.\allowbreak{}score\allowbreak{}Priority\allowbreak{}Band}
\item \hyperlink{governing-declaration-efba6aee5c5b1874}{LBG22\allowbreak{}Strategic\allowbreak{}Ranking.\allowbreak{}source\allowbreak{}Multitask\allowbreak{}Cost}
\item \hyperlink{paper-prerequisite-f88aac58507e44ca}{LBG22\allowbreak{}Strategic\allowbreak{}Ranking.\allowbreak{}source\allowbreak{}Two\allowbreak{}Level\allowbreak{}Cutoff}
\item \hyperlink{paper-prerequisite-30d97cd43eb79017}{LBG22\allowbreak{}Strategic\allowbreak{}Ranking.\allowbreak{}source\allowbreak{}Two\allowbreak{}Level\allowbreak{}Reward}
\end{itemize}
\paragraph{Recorded source-to-declaration screening}
\textbf{Verdict:} \texttt{matches}
\\
{\raggedright
\textbf{Reason:} The target restricts actual and deviating task efforts to nonnegative pairs summing to B, uses only measurable-\allowbreak{}skill score for admission, and compares reward minus the displayed multitask cost.\allowbreak{} On the hard-\allowbreak{}budget feasible set that cost is p(eM), exactly matching the approved source model.\allowbreak{}
\par}
\reviewmatch{Matches source input}
\noindent\textbf{Reviewer annotation}\par
\reviewerbox
\clearpage
\hypertarget{paper-prerequisite-fe9ea55a34d7ea9c}{}
\subsection*{\small\ttfamily\raggedright LBG22\allowbreak{}Strategic\allowbreak{}Ranking.\allowbreak{}Score\allowbreak{}Priority\allowbreak{}Multidimensional\allowbreak{}Best\allowbreak{}Response\allowbreak{}At}
\noindent\textbf{Paper-local declaration:}\par
{\footnotesize\ttfamily\raggedright LBG22\allowbreak{}Strategic\allowbreak{}Ranking.\allowbreak{}Score\allowbreak{}Priority\allowbreak{}Multidimensional\allowbreak{}Best\allowbreak{}Response\allowbreak{}At\par}
\textbf{Source locator:} {\footnotesize\raggedright cited publication, lines 1724-\allowbreak{}1739\par}
\paragraph{Verbatim paper-source connection}
\begin{ReviewVerbatim}
\parbold{Model.} The model is similar to our base model \Cref{sec:model}; each applicant now has $m$ latent skill levels with respective ranks $\theta_\pre^i \in [0,1]$ for $i\in \{1, \cdots, m\}$. Each rank is drawn independently from the uniform distribution over $[0,1]$, and the skill $i$ of applicant with rank $\theta_\pre^i$ is $\f{i}(\theta_\pre^i)$. Each applicant now chooses effort levels $\{ \e{i} \}$, at cost $p^m( \e{1}, \cdots, \e{m})$, resulting in post-effort scores $\{\score{i}\}$,
\begin{equation*}
\score{i} = g(\e{i})\cdot \f{i}(\theta_\pre^i),
\end{equation*}
where $g$ is concave, increasing, as before, and $\f{i}$ is a continuous, strictly increasing quantile function (${\f{i}}^{-1}$ is the CDF function of the scores on dimension $i$). As before, the school observes the post-effort scores for each applicant, now for each dimension $i$, and designs a non-decreasing function $\lambda : [0,1] \to [0,1]$ denoting the admissions probability $\lambda(\theta_\post)$ for applicant with post-effort rank $\theta_\post$.

%Suppose the school observes $v_i$.

How does the school construct post-effort rank $\theta_\post$? In general, each applicant's \textit{combined} post-effort score may be any function of the scores on each dimension, $\{\score{i}\}$. Here, we assume the following linear score function. The school announces weights $\alpha = (\alpha_1, \cdots, \alpha_m) \in \Delta_{m-1}$ (denoting the $m$-dimensional simplex); the weights $\alpha$ represent the relative emphasis placed on each skill for the admissions decision. Then, each applicant is ranked according to their combined score, $\combscore := \sum_{i=1}^m \alpha_i\score{i}$, resulting in their combined post-effort rank $\theta_\post^\alpha$, and is admitted with probability $\lambda(\theta_\post^\alpha)$.\footnote{The linear combination of skill is similar to the ``linear mechanism'' in \citet{kleinberg18investeffort}, which studied how reward design incentivizes strategic agents to exert effort on different skill dimensions. Compared to their work, where efforts are connected to skills via an effort graph, we consider a simplified setting where each effort maps to one skill, and study how reward design affects equilibrium rankings, in presence of competition.}

%\lnote{\textbf{weighted score or weighted rank?} alternatively, the school combines the post-effort ranks for each skill. i.e. an applicant with \emph{weighted} post-effort rank $\theta_\post^\alpha := \sum_{i=1}^m \alpha_i\theta_\post^i$ is admitted with probability $\lambda(\theta_\post^\alpha)$. This case seems interesting as well. Not clear in which case it's possible to reduce to the single-dimensional case?}

Putting things together, the applicant's individual welfare is:
\begin{equation*}
\iwelf(\{ \e{i} \}_{i=1}^m, \lambda(\theta_\post^\alpha)) = \lambda(\theta_\post^\alpha) - p^m( \e{1}, \cdots, \e{m})
\end{equation*}

[Next verbatim source excerpt]

\begin{definition} [Equilibrium]\label{def:equi}
{Given a}
%Let the
tie-breaking function $\Gamma(\omega)$
%be a pre-announced ordering over index $\omega$, to break ties in post-effort value. Then, let $\gamma(\omega, v)$ be the 1-1 ranking function defined as the CDF of post-effort values $v$ except where there are ties in $v$, in which case the ordering $\Gamma$ is used to produce unique ranks such that the range of $\gamma$ is uniform on $[0, 1]$.\footnote{More formally, note that each discontinuity in the CDF of post-effort values $v$ has a height equal to the measure of the applicants who are tied at that value, and so $\gamma$ for those applicants can be filled in using the CDF of $\Gamma(\omega)$ restricted to the applicants $\omega$ who are tied. Thus, the range of $\gamma$ is uniform on $[0,1]$.
%
{and}
%	Then, given
	a ranking probability function $\lambda$, an \textit{equilibrium} is a set of effort levels and post-effort ranking rewards for each applicant, $\{e(\theta_\pre),  \lambda(\theta_\post(\theta_\pre))\}$ such that, for all
	$\omega$, %$\theta_\pre$,
	\begin{align*}
		e(\theta_\pre(\omega)) &\in \argmax_{e} \iwelf\left(e, \lambda\left(\gamma\left(\omega, v\left(e,\theta_\pre(\omega)\right)\right)\right)\right) %\triangleq \argmax_{e} \left[ \lambda(\gamma(v(e,\theta_\pre)) - \cost(e) \right].
		\\
		\theta_\post(\theta_\pre(\omega)) &= \gamma(\omega, v(e(\theta_\pre(\omega)),\theta_\pre(\omega))),
	\end{align*}
	{where $\gamma(\omega,v)$ is the ranking induced by the CDF of the scores resulting from effort levels $\{e(\theta_\pre)\}$ and the tie-breaking function $\Gamma$.\footnote{See Remark~\ref{rem:gamma-effort} on $\gamma$ and the set of efforts.%Formally, $\gamma$ depends on the set of efforts. Given a fixed set and $\gamma$, for each applicant $\omega$ the first condition considers the counter-factual ranking of $\omega$ with different post-effort values but using the same ranking function. As defined, $\gamma$ yields a uniform distribution of ranks with such measure $0$ changes. In Appendix \Cref{lem:deviationsmeasure0}, we further prove that two effort sets equal up to sets of measure $0$ induce the same ranking function $\gamma$, and so the condition is consistent.%, and so we set it at the $\gamma$ induced by the effort levels $\{e(\theta_\pre)\}$.%, and so it does not vary.
	}
	}
	%In particular, we denote $\theta_\post(\theta_\pre) := \gamma(v(e(\theta_\pre),\theta_\pre))$.
	%\lledit{and for any $\theta_\pre, \theta'_\pre$, we have $\lambda(\theta_\post(\theta_\pre)) > \lambda(\theta_\post(\theta'_\pre)) \implies v(e(\theta_\pre),\theta_\pre) > v(e(\theta_\pre'),\theta_\pre')$.}
\end{definition}

[Next verbatim source excerpt]

\parbold{Applicants}
There is a unit mass of \textit{applicants}, indexed by an observed index $\omega \in [0,1]$ distributed uniformly.\footnote{The index $\omega$ should be interpreted as each applicant's ``name,'' uncorrelated with skill, used solely for tie-breaking.}  Each applicant has a latent (unobserved) skill level represented by some measurable function of $\omega$.
We assume that {the distribution of the skills} has no atoms and that the CDF of this distribution is strictly increasing.
Using the CDF to map the skill of an applicant to a rank in~$[0,1]$, each applicant gets an (unobserved) rank $\theta_\pre =\theta_\pre(\omega)$ which by our assumption on the CDF is again uniformly distributed in $[0, 1]$ (the higher the rank the better).  With this setup, the skill of an applicant with rank $\theta_\pre$ can be written as $f(\theta_\pre)$ where $f$ is a strictly increasing, continuous function.  {It will be notationally convenient to label applicants by their rank $\theta_\pre$, though the reader should note that in contrast to $\omega$, $\theta_\pre(\omega)$ is assumed to be unobservable.}

%3. Each applicant exerts some effort e, at cost \cost(e).


Each applicant chooses an \textit{effort} level $e \ge 0$, the result of which is an observed, post-effort \textit{score}, {$v=v(e,\theta_\pre) = {g(e)}\cdot{f(\theta_\pre)}$}. \lledit{In other words, we assume the post-effort score is the product of two components, associated with the effort and the pre-effort skill respectively.}
{We assume} that the effort transfer function $g: [0, \infty) \mapsto [0, \infty)$ is a continuous, concave, strictly increasing function, representing that marginal effort improves one's score but has diminishing returns.
{%Since applicants can be labelled by their skill level $f(\theta_\pre)$, or equivalently, by their rank $\theta_{\pre}$,
{T}he strategies of the applicants can {then} be described by a function $\theta_\pre \mapsto e(\theta_\pre)$.
%\cb{$\omega\mapsto e(\omega)$.}%, leading to a distribution over post-effort scores $v=v(\theta_\pre))$ via the uniform distribution of $\theta_\pre$ in $[0,1]$.
}
  Each {applicant} is then ranked according to their score $v$%(where we break ties uniformly at random, if they appear)
, resulting in a \textit{post-effort rank} $\theta_\post$.
Note that the ranking $\theta_\post$ is slightly less trivial than a ranking of the skills, since the scores might have ties, which have to be resolved.

{\textbf{Tie Breaking.}
Given a choice of strategies $\theta_\pre \mapsto e(\theta_\pre)$, let $F$ be the CDF of the scores $v(e(\theta_\pre),\theta_\pre)$.
Since atoms for the distribution of $v$ would lead to ties for ranks defined as $F(v)$, we will use the labels of the applicants to break ties with the help of a (publicly announced) tie-breaking function $\Gamma(\omega)$, defined, e.g., via a collision free hash of the applicant names. Here  we require  that $\Gamma$  is a measurable function from $[0,1]$ to $[0,1]$ that maps different applicant labels to different values.
We then use $\Gamma$ to resolve the atoms of $F$, leading to a ranking function $v\mapsto \gamma(\omega, v)$ which is equal to $F(v)$ except when $v$ is an atom of the score distribution, {in which case it takes values in the ``gap interval'' $[F_-(v),  F(v)]$, where $F_-(v)$ is the left limit of $F$ at $v$.}  We  construct $\gamma$ in such a way that it
gives the uniform distribution for $\theta_\post(\omega)$ if we set  $\theta_\post(\omega)=\gamma(\omega,v(e(\theta_\pre(\omega)),\theta_\pre(\omega)))$.\footnote{See Remark~\ref{rem:gamma} for details on this construction.}
\end{ReviewVerbatim}

% report-context-presentation-sha256: 90f9116e292dd71fddfd6b445ca0b060a481d026855bdb9725f2833c2cf62a31
\paragraph{Settled source reading}
\begin{sloppypar}
Equilibrium means that almost every applicant best responds against every feasible deviation;\allowbreak{} uniqueness is up to a null set.\allowbreak{}
\end{sloppypar}
\paragraph{Settled source reading}
\begin{sloppypar}
Applicants know the fixed priority used at score-\allowbreak{}boundary ties.\allowbreak{} This clarifies the stated model, not a corrected admission policy or a new economic assumption.\allowbreak{}
\end{sloppypar}

\paragraph{Lean-expanded paper semantic target}
\begin{ReviewVerbatim}
type:
{α : Type u_1} →
  {ι : Type u_2} →
    [inst : MeasurableSpace α] →
      [Fintype ι] →
        (μ : MeasureTheory.Measure α) →
          [MeasureTheory.IsFiniteMeasure μ] →
            (ℝ → ℝ) → (α → ι → ℝ) → (α → ℝ) → (α → ℝ) → ℝ → {n : ℕ} → (Fin (n + 1) → ℝ) → (ℕ → ℝ) → α → (ι → ℝ) → Prop

value:
fun {α : Type u_1} {ι : Type u_2} [MeasurableSpace α] [Fintype ι] (μ : MeasureTheory.Measure α)
    [MeasureTheory.IsFiniteMeasure μ] (cost : ℝ → ℝ) (coefficient : α → ι → ℝ) (score tie : α → ℝ) (h : ℝ) {n : ℕ}
    (cutoff : Fin (n + 1) → ℝ) (reward : ℕ → ℝ) (x : α) (effort : ι → ℝ) =>
  (∀ (i : ι), 0 ≤ effort i) ∧
    score x = h * ∑ i : ι, effort i * coefficient x i ∧
      ∀ (d : ι → ℝ),
        (∀ (i : ι), 0 ≤ d i) →
          reward
                ↑(LBG22StrategicRanking.scorePriorityBand μ score tie cutoff (h * ∑ i : ι, d i * coefficient x i)
                    (tie x)) -
              cost (∑ i : ι, d i) ≤
            reward ↑(LBG22StrategicRanking.scorePriorityBand μ score tie cutoff (score x) (tie x)) -
              cost (∑ i : ι, effort i)
\end{ReviewVerbatim}

\paragraph{Direct retained dependencies}
\begin{itemize}\raggedright
\item \hyperlink{governing-declaration-29f9270f1ab8f757}{LBG22\allowbreak{}Strategic\allowbreak{}Ranking.\allowbreak{}score\allowbreak{}Priority\allowbreak{}Band}
\end{itemize}
\paragraph{Recorded source-to-declaration screening}
\textbf{Verdict:} \texttt{matches}
\\
{\raggedright
\textbf{Reason:} The target quantifies over all nonnegative effort vectors, forms the linear combined score h times the weighted skill-\allowbreak{}effort sum, charges cost of total effort, and compares finite-\allowbreak{}policy reward under fixed score priority.\allowbreak{} This is the source multidimensional best-\allowbreak{}response problem in its exchangeable linear-\allowbreak{}production specialization.\allowbreak{}
\par}
\reviewmatch{Matches source input}
\noindent\textbf{Reviewer annotation}\par
\reviewerbox
\clearpage
\hypertarget{paper-prerequisite-b0c4f18b812a40f1}{}
\subsection*{\small\ttfamily\raggedright LBG22\allowbreak{}Strategic\allowbreak{}Ranking.\allowbreak{}Source\allowbreak{}Vector\allowbreak{}Best\allowbreak{}Response\allowbreak{}At}
\noindent\textbf{Paper-local declaration:}\par
{\footnotesize\ttfamily\raggedright LBG22\allowbreak{}Strategic\allowbreak{}Ranking.\allowbreak{}Source\allowbreak{}Vector\allowbreak{}Best\allowbreak{}Response\allowbreak{}At\par}
\textbf{Source locator:} {\footnotesize\raggedright cited publication, lines 1724-\allowbreak{}1739\par}
\paragraph{Verbatim paper-source connection}
\begin{ReviewVerbatim}
\parbold{Model.} The model is similar to our base model \Cref{sec:model}; each applicant now has $m$ latent skill levels with respective ranks $\theta_\pre^i \in [0,1]$ for $i\in \{1, \cdots, m\}$. Each rank is drawn independently from the uniform distribution over $[0,1]$, and the skill $i$ of applicant with rank $\theta_\pre^i$ is $\f{i}(\theta_\pre^i)$. Each applicant now chooses effort levels $\{ \e{i} \}$, at cost $p^m( \e{1}, \cdots, \e{m})$, resulting in post-effort scores $\{\score{i}\}$,
\begin{equation*}
\score{i} = g(\e{i})\cdot \f{i}(\theta_\pre^i),
\end{equation*}
where $g$ is concave, increasing, as before, and $\f{i}$ is a continuous, strictly increasing quantile function (${\f{i}}^{-1}$ is the CDF function of the scores on dimension $i$). As before, the school observes the post-effort scores for each applicant, now for each dimension $i$, and designs a non-decreasing function $\lambda : [0,1] \to [0,1]$ denoting the admissions probability $\lambda(\theta_\post)$ for applicant with post-effort rank $\theta_\post$.

%Suppose the school observes $v_i$.

How does the school construct post-effort rank $\theta_\post$? In general, each applicant's \textit{combined} post-effort score may be any function of the scores on each dimension, $\{\score{i}\}$. Here, we assume the following linear score function. The school announces weights $\alpha = (\alpha_1, \cdots, \alpha_m) \in \Delta_{m-1}$ (denoting the $m$-dimensional simplex); the weights $\alpha$ represent the relative emphasis placed on each skill for the admissions decision. Then, each applicant is ranked according to their combined score, $\combscore := \sum_{i=1}^m \alpha_i\score{i}$, resulting in their combined post-effort rank $\theta_\post^\alpha$, and is admitted with probability $\lambda(\theta_\post^\alpha)$.\footnote{The linear combination of skill is similar to the ``linear mechanism'' in \citet{kleinberg18investeffort}, which studied how reward design incentivizes strategic agents to exert effort on different skill dimensions. Compared to their work, where efforts are connected to skills via an effort graph, we consider a simplified setting where each effort maps to one skill, and study how reward design affects equilibrium rankings, in presence of competition.}

%\lnote{\textbf{weighted score or weighted rank?} alternatively, the school combines the post-effort ranks for each skill. i.e. an applicant with \emph{weighted} post-effort rank $\theta_\post^\alpha := \sum_{i=1}^m \alpha_i\theta_\post^i$ is admitted with probability $\lambda(\theta_\post^\alpha)$. This case seems interesting as well. Not clear in which case it's possible to reduce to the single-dimensional case?}

Putting things together, the applicant's individual welfare is:
\begin{equation*}
\iwelf(\{ \e{i} \}_{i=1}^m, \lambda(\theta_\post^\alpha)) = \lambda(\theta_\post^\alpha) - p^m( \e{1}, \cdots, \e{m})
\end{equation*}

[Next verbatim source excerpt]

\begin{definition} [Equilibrium]\label{def:equi}
{Given a}
%Let the
tie-breaking function $\Gamma(\omega)$
%be a pre-announced ordering over index $\omega$, to break ties in post-effort value. Then, let $\gamma(\omega, v)$ be the 1-1 ranking function defined as the CDF of post-effort values $v$ except where there are ties in $v$, in which case the ordering $\Gamma$ is used to produce unique ranks such that the range of $\gamma$ is uniform on $[0, 1]$.\footnote{More formally, note that each discontinuity in the CDF of post-effort values $v$ has a height equal to the measure of the applicants who are tied at that value, and so $\gamma$ for those applicants can be filled in using the CDF of $\Gamma(\omega)$ restricted to the applicants $\omega$ who are tied. Thus, the range of $\gamma$ is uniform on $[0,1]$.
%
{and}
%	Then, given
	a ranking probability function $\lambda$, an \textit{equilibrium} is a set of effort levels and post-effort ranking rewards for each applicant, $\{e(\theta_\pre),  \lambda(\theta_\post(\theta_\pre))\}$ such that, for all
	$\omega$, %$\theta_\pre$,
	\begin{align*}
		e(\theta_\pre(\omega)) &\in \argmax_{e} \iwelf\left(e, \lambda\left(\gamma\left(\omega, v\left(e,\theta_\pre(\omega)\right)\right)\right)\right) %\triangleq \argmax_{e} \left[ \lambda(\gamma(v(e,\theta_\pre)) - \cost(e) \right].
		\\
		\theta_\post(\theta_\pre(\omega)) &= \gamma(\omega, v(e(\theta_\pre(\omega)),\theta_\pre(\omega))),
	\end{align*}
	{where $\gamma(\omega,v)$ is the ranking induced by the CDF of the scores resulting from effort levels $\{e(\theta_\pre)\}$ and the tie-breaking function $\Gamma$.\footnote{See Remark~\ref{rem:gamma-effort} on $\gamma$ and the set of efforts.%Formally, $\gamma$ depends on the set of efforts. Given a fixed set and $\gamma$, for each applicant $\omega$ the first condition considers the counter-factual ranking of $\omega$ with different post-effort values but using the same ranking function. As defined, $\gamma$ yields a uniform distribution of ranks with such measure $0$ changes. In Appendix \Cref{lem:deviationsmeasure0}, we further prove that two effort sets equal up to sets of measure $0$ induce the same ranking function $\gamma$, and so the condition is consistent.%, and so we set it at the $\gamma$ induced by the effort levels $\{e(\theta_\pre)\}$.%, and so it does not vary.
	}
	}
	%In particular, we denote $\theta_\post(\theta_\pre) := \gamma(v(e(\theta_\pre),\theta_\pre))$.
	%\lledit{and for any $\theta_\pre, \theta'_\pre$, we have $\lambda(\theta_\post(\theta_\pre)) > \lambda(\theta_\post(\theta'_\pre)) \implies v(e(\theta_\pre),\theta_\pre) > v(e(\theta_\pre'),\theta_\pre')$.}
\end{definition}

[Next verbatim source excerpt]

\parbold{Applicants}
There is a unit mass of \textit{applicants}, indexed by an observed index $\omega \in [0,1]$ distributed uniformly.\footnote{The index $\omega$ should be interpreted as each applicant's ``name,'' uncorrelated with skill, used solely for tie-breaking.}  Each applicant has a latent (unobserved) skill level represented by some measurable function of $\omega$.
We assume that {the distribution of the skills} has no atoms and that the CDF of this distribution is strictly increasing.
Using the CDF to map the skill of an applicant to a rank in~$[0,1]$, each applicant gets an (unobserved) rank $\theta_\pre =\theta_\pre(\omega)$ which by our assumption on the CDF is again uniformly distributed in $[0, 1]$ (the higher the rank the better).  With this setup, the skill of an applicant with rank $\theta_\pre$ can be written as $f(\theta_\pre)$ where $f$ is a strictly increasing, continuous function.  {It will be notationally convenient to label applicants by their rank $\theta_\pre$, though the reader should note that in contrast to $\omega$, $\theta_\pre(\omega)$ is assumed to be unobservable.}

%3. Each applicant exerts some effort e, at cost \cost(e).


Each applicant chooses an \textit{effort} level $e \ge 0$, the result of which is an observed, post-effort \textit{score}, {$v=v(e,\theta_\pre) = {g(e)}\cdot{f(\theta_\pre)}$}. \lledit{In other words, we assume the post-effort score is the product of two components, associated with the effort and the pre-effort skill respectively.}
{We assume} that the effort transfer function $g: [0, \infty) \mapsto [0, \infty)$ is a continuous, concave, strictly increasing function, representing that marginal effort improves one's score but has diminishing returns.
{%Since applicants can be labelled by their skill level $f(\theta_\pre)$, or equivalently, by their rank $\theta_{\pre}$,
{T}he strategies of the applicants can {then} be described by a function $\theta_\pre \mapsto e(\theta_\pre)$.
%\cb{$\omega\mapsto e(\omega)$.}%, leading to a distribution over post-effort scores $v=v(\theta_\pre))$ via the uniform distribution of $\theta_\pre$ in $[0,1]$.
}
  Each {applicant} is then ranked according to their score $v$%(where we break ties uniformly at random, if they appear)
, resulting in a \textit{post-effort rank} $\theta_\post$.
Note that the ranking $\theta_\post$ is slightly less trivial than a ranking of the skills, since the scores might have ties, which have to be resolved.

{\textbf{Tie Breaking.}
Given a choice of strategies $\theta_\pre \mapsto e(\theta_\pre)$, let $F$ be the CDF of the scores $v(e(\theta_\pre),\theta_\pre)$.
Since atoms for the distribution of $v$ would lead to ties for ranks defined as $F(v)$, we will use the labels of the applicants to break ties with the help of a (publicly announced) tie-breaking function $\Gamma(\omega)$, defined, e.g., via a collision free hash of the applicant names. Here  we require  that $\Gamma$  is a measurable function from $[0,1]$ to $[0,1]$ that maps different applicant labels to different values.
We then use $\Gamma$ to resolve the atoms of $F$, leading to a ranking function $v\mapsto \gamma(\omega, v)$ which is equal to $F(v)$ except when $v$ is an atom of the score distribution, {in which case it takes values in the ``gap interval'' $[F_-(v),  F(v)]$, where $F_-(v)$ is the left limit of $F$ at $v$.}  We  construct $\gamma$ in such a way that it
gives the uniform distribution for $\theta_\post(\omega)$ if we set  $\theta_\post(\omega)=\gamma(\omega,v(e(\theta_\pre(\omega)),\theta_\pre(\omega)))$.\footnote{See Remark~\ref{rem:gamma} for details on this construction.}
\end{ReviewVerbatim}

% report-context-presentation-sha256: 90f9116e292dd71fddfd6b445ca0b060a481d026855bdb9725f2833c2cf62a31
\paragraph{Settled source reading}
\begin{sloppypar}
Equilibrium means that almost every applicant best responds against every feasible deviation;\allowbreak{} uniqueness is up to a null set.\allowbreak{}
\end{sloppypar}
\paragraph{Settled source reading}
\begin{sloppypar}
Applicants know the fixed priority used at score-\allowbreak{}boundary ties.\allowbreak{} This clarifies the stated model, not a corrected admission policy or a new economic assumption.\allowbreak{}
\end{sloppypar}

\paragraph{Lean-expanded paper semantic target}
\begin{ReviewVerbatim}
type:
{α : Type u_1} →
  {ι : Type u_2} →
    [inst : MeasurableSpace α] →
      [Fintype ι] →
        (μ : MeasureTheory.Measure α) →
          [MeasureTheory.IsFiniteMeasure μ] →
            ((ι → ℝ) → ℝ) → (ℝ → ℝ) → (α → ι → ℝ) → (α → ℝ) → (α → ℝ) → (ℝ → ℝ) → α → (ι → ℝ) → Prop

value:
fun {α : Type u_1} {ι : Type u_2} [MeasurableSpace α] [Fintype ι] (μ : MeasureTheory.Measure α)
    [MeasureTheory.IsFiniteMeasure μ] (cost : (ι → ℝ) → ℝ) (production : ℝ → ℝ) (coefficient : α → ι → ℝ)
    (score tie : α → ℝ) (reward : ℝ → ℝ) (x : α) (effort : ι → ℝ) =>
  (∀ (i : ι), 0 ≤ effort i) ∧
    score x = ∑ i : ι, production (effort i) * coefficient x i ∧
      ∀ (d : ι → ℝ),
        (∀ (i : ι), 0 ≤ d i) →
          reward
                (LBG22StrategicRanking.counterfactualTieBrokenRank μ score tie x
                  (∑ i : ι, production (d i) * coefficient x i)) -
              cost d ≤
            reward (LBG22StrategicRanking.tieBrokenRank μ score tie x) - cost effort
\end{ReviewVerbatim}

\paragraph{Direct retained dependencies}
\begin{itemize}\raggedright
\item \hyperlink{paper-prerequisite-6ed45e46d01eadeb}{LBG22\allowbreak{}Strategic\allowbreak{}Ranking.\allowbreak{}counterfactual\allowbreak{}Tie\allowbreak{}Broken\allowbreak{}Rank}
\item \hyperlink{paper-prerequisite-b7c71e53177b76c6}{LBG22\allowbreak{}Strategic\allowbreak{}Ranking.\allowbreak{}tie\allowbreak{}Broken\allowbreak{}Rank}
\end{itemize}
\paragraph{Recorded source-to-declaration screening}
\textbf{Verdict:} \texttt{matches}
\\
{\raggedright
\textbf{Reason:} The target states the general multidimensional source problem:\allowbreak{} arbitrary vector cost, coordinatewise production and skill coefficients, summed score, and reward minus cost maximized over every nonnegative vector deviation with fixed ranking.\allowbreak{} Its broader reusable types specialize to every source instance.\allowbreak{}
\par}
\reviewmatch{Matches source input}
\noindent\textbf{Reviewer annotation}\par
\reviewerbox
\clearpage
\hypertarget{paper-prerequisite-6f1ee9b200124096}{}
\subsection*{\small\ttfamily\raggedright LBG22\allowbreak{}Strategic\allowbreak{}Ranking.\allowbreak{}source\allowbreak{}Actual\allowbreak{}Disadvantaged\allowbreak{}Access}
\noindent\textbf{Paper-local declaration:}\par
{\footnotesize\ttfamily\raggedright LBG22\allowbreak{}Strategic\allowbreak{}Ranking.\allowbreak{}source\allowbreak{}Actual\allowbreak{}Disadvantaged\allowbreak{}Access\par}
\textbf{Source locator:} {\footnotesize\raggedright cited publication, lines 269-\allowbreak{}275\par}
\paragraph{Verbatim paper-source connection}
\begin{ReviewVerbatim}
\begin{definition}[Access]
 \emph{Access} is the overall probability of admission of the disadvantaged group.
	\begin{equation*}
		\access :=  \E_{\theta_\true}[\lambda(\theta_\post(\theta_\true, \psi_\B))].
	\end{equation*}
% \nikhil{Defn is a bit confusing. What is the expectation over, and what does that conditioning mean?}
\end{definition}
\end{ReviewVerbatim}

% report-context-presentation-sha256: 827a19afe0274b1faef13e344ddbcfe40d9841f2ff34b92ac9d53c0f83e29f56
\paragraph{Settled source reading}
\begin{sloppypar}
Applicants know the fixed priority used at score-\allowbreak{}boundary ties.\allowbreak{} This clarifies the stated model, not a corrected admission policy or a new economic assumption.\allowbreak{}
\end{sloppypar}

\paragraph{Lean-expanded paper semantic target}
\begin{ReviewVerbatim}
type:
{n : ℕ} → (Bool × ℝ → ℝ) → (Bool × ℝ → ℝ) → (Fin (n + 1) → ℝ) → (ℕ → ℝ) → ℝ

value:
fun {n : ℕ} (score tie : Bool × ℝ → ℝ) (cutoff : Fin (n + 1) → ℝ) (reward : ℕ → ℝ) =>
  ∫ (t : ℝ),
    reward
      ↑(LBG22StrategicRanking.finiteLowerRankBand cutoff
          (LBG22StrategicRanking.tieBrokenRank LBG22StrategicRanking.sourceTwoGroupMeasure score tie
            (false, t))) ∂LBG22StrategicRanking.unitRankMeasure
\end{ReviewVerbatim}

\paragraph{Direct retained dependencies}
\begin{itemize}\raggedright
\item \hyperlink{paper-prerequisite-a42f440ddc254fa7}{LBG22\allowbreak{}Strategic\allowbreak{}Ranking.\allowbreak{}finite\allowbreak{}Lower\allowbreak{}Rank\allowbreak{}Band}
\item \hyperlink{paper-prerequisite-5b3be494cde30a37}{LBG22\allowbreak{}Strategic\allowbreak{}Ranking.\allowbreak{}source\allowbreak{}Two\allowbreak{}Group\allowbreak{}Measure}
\item \hyperlink{paper-prerequisite-b7c71e53177b76c6}{LBG22\allowbreak{}Strategic\allowbreak{}Ranking.\allowbreak{}tie\allowbreak{}Broken\allowbreak{}Rank}
\item \hyperlink{governing-declaration-678cf48d549cb620}{LBG22\allowbreak{}Strategic\allowbreak{}Ranking.\allowbreak{}unit\allowbreak{}Rank\allowbreak{}Measure}
\end{itemize}
\paragraph{Recorded source-to-declaration screening}
\textbf{Verdict:} \texttt{matches}
\\
{\raggedright
\textbf{Reason:} The target integrates the finite-\allowbreak{}policy admission reward of group B over its uniform latent rank using the common two-\allowbreak{}group realized-\allowbreak{}score ranking.\allowbreak{} This is exactly the source definition of disadvantaged-\allowbreak{}group access under the fixed-\allowbreak{}priority convention.\allowbreak{}
\par}
\reviewmatch{Matches source input}
\noindent\textbf{Reviewer annotation}\par
\reviewerbox
\clearpage
\hypertarget{paper-prerequisite-22c0ae11547a174a}{}
\subsection*{\small\ttfamily\raggedright LBG22\allowbreak{}Strategic\allowbreak{}Ranking.\allowbreak{}source\allowbreak{}Actual\allowbreak{}Group\allowbreak{}Welfare}
\noindent\textbf{Paper-local declaration:}\par
{\footnotesize\ttfamily\raggedright LBG22\allowbreak{}Strategic\allowbreak{}Ranking.\allowbreak{}source\allowbreak{}Actual\allowbreak{}Group\allowbreak{}Welfare\par}
\textbf{Source locator:} {\footnotesize\raggedright cited publication, lines 250-\allowbreak{}265\par}
\paragraph{Verbatim paper-source connection}
\begin{ReviewVerbatim}
\begin{definition}[Welfare gap]
	Let $\swelf^\G(\theta_\true)$ denote post-effort welfare of an applicant with latent skill ranking $\theta_\true$ from group $\G \in \{\A, \B\}$, i.e.,
	\begin{equation}
	\swelf^\G(\theta_\true):= \lambda(\theta_\post(\theta_\true, \psi_\G)) - \cost(e(\theta_\true, \psi_\G)).
	\end{equation}
	%\lnote{Using function notation for $\theta_\post$ and $e$? or some other way to express the welfare as a function of $\theta_\true$.}
	We define the \emph{%pointwise
		welfare gap} as %a function of rank $\theta_\true$ is
	\begin{equation*}
		\welfgap(\theta_\true):= \swelf^\A(\theta_\true) - \swelf^\B(\theta_\true).
	\end{equation*}
%	We define the \emph{subpopulation welfare gap} as %a function of baseline rank $\overline{\theta}_\true$ is
%	\begin{equation*}
%		\popwelfgap(\overline{\theta}_\true):=\E[\welfgap(\theta_\true) \mid \theta_\true > \overline{\theta}_\true].
%	\end{equation*}
\end{definition}
\end{ReviewVerbatim}

% report-context-presentation-sha256: 827a19afe0274b1faef13e344ddbcfe40d9841f2ff34b92ac9d53c0f83e29f56
\paragraph{Settled source reading}
\begin{sloppypar}
Applicants know the fixed priority used at score-\allowbreak{}boundary ties.\allowbreak{} This clarifies the stated model, not a corrected admission policy or a new economic assumption.\allowbreak{}
\end{sloppypar}

\paragraph{Lean-expanded paper semantic target}
\begin{ReviewVerbatim}
type:
{n : ℕ} → (ℝ → ℝ) → (Bool × ℝ → ℝ) → (Bool × ℝ → ℝ) → (Bool × ℝ → ℝ) → (Fin (n + 1) → ℝ) → (ℕ → ℝ) → Bool → ℝ → ℝ

value:
fun {n : ℕ} (cost : ℝ → ℝ) (effort score tie : Bool × ℝ → ℝ) (cutoff : Fin (n + 1) → ℝ) (reward : ℕ → ℝ) (group : Bool)
    (t : ℝ) =>
  reward
      ↑(LBG22StrategicRanking.finiteLowerRankBand cutoff
          (LBG22StrategicRanking.tieBrokenRank LBG22StrategicRanking.sourceTwoGroupMeasure score tie (group, t))) -
    cost (effort (group, t))
\end{ReviewVerbatim}

\paragraph{Direct retained dependencies}
\begin{itemize}\raggedright
\item \hyperlink{paper-prerequisite-a42f440ddc254fa7}{LBG22\allowbreak{}Strategic\allowbreak{}Ranking.\allowbreak{}finite\allowbreak{}Lower\allowbreak{}Rank\allowbreak{}Band}
\item \hyperlink{paper-prerequisite-5b3be494cde30a37}{LBG22\allowbreak{}Strategic\allowbreak{}Ranking.\allowbreak{}source\allowbreak{}Two\allowbreak{}Group\allowbreak{}Measure}
\item \hyperlink{paper-prerequisite-b7c71e53177b76c6}{LBG22\allowbreak{}Strategic\allowbreak{}Ranking.\allowbreak{}tie\allowbreak{}Broken\allowbreak{}Rank}
\end{itemize}
\paragraph{Recorded source-to-declaration screening}
\textbf{Verdict:} \texttt{matches}
\\
{\raggedright
\textbf{Reason:} For a group and latent rank, the target subtracts effort cost from the finite-\allowbreak{}policy reward at that applicant's actual common-\allowbreak{}population post-\allowbreak{}rank.\allowbreak{} This is the source pointwise group welfare formula.\allowbreak{}
\par}
\reviewmatch{Matches source input}
\noindent\textbf{Reviewer annotation}\par
\reviewerbox
\clearpage
\hypertarget{paper-prerequisite-d3e90f3908818ca5}{}
\subsection*{\small\ttfamily\raggedright LBG22\allowbreak{}Strategic\allowbreak{}Ranking.\allowbreak{}source\allowbreak{}Actual\allowbreak{}Group\allowbreak{}Welfare\allowbreak{}Gap}
\noindent\textbf{Paper-local declaration:}\par
{\footnotesize\ttfamily\raggedright LBG22\allowbreak{}Strategic\allowbreak{}Ranking.\allowbreak{}source\allowbreak{}Actual\allowbreak{}Group\allowbreak{}Welfare\allowbreak{}Gap\par}
\textbf{Source locator:} {\footnotesize\raggedright cited publication, lines 250-\allowbreak{}265\par}
\paragraph{Verbatim paper-source connection}
\begin{ReviewVerbatim}
\begin{definition}[Welfare gap]
	Let $\swelf^\G(\theta_\true)$ denote post-effort welfare of an applicant with latent skill ranking $\theta_\true$ from group $\G \in \{\A, \B\}$, i.e.,
	\begin{equation}
	\swelf^\G(\theta_\true):= \lambda(\theta_\post(\theta_\true, \psi_\G)) - \cost(e(\theta_\true, \psi_\G)).
	\end{equation}
	%\lnote{Using function notation for $\theta_\post$ and $e$? or some other way to express the welfare as a function of $\theta_\true$.}
	We define the \emph{%pointwise
		welfare gap} as %a function of rank $\theta_\true$ is
	\begin{equation*}
		\welfgap(\theta_\true):= \swelf^\A(\theta_\true) - \swelf^\B(\theta_\true).
	\end{equation*}
%	We define the \emph{subpopulation welfare gap} as %a function of baseline rank $\overline{\theta}_\true$ is
%	\begin{equation*}
%		\popwelfgap(\overline{\theta}_\true):=\E[\welfgap(\theta_\true) \mid \theta_\true > \overline{\theta}_\true].
%	\end{equation*}
\end{definition}
\end{ReviewVerbatim}

% report-context-presentation-sha256: 827a19afe0274b1faef13e344ddbcfe40d9841f2ff34b92ac9d53c0f83e29f56
\paragraph{Settled source reading}
\begin{sloppypar}
Applicants know the fixed priority used at score-\allowbreak{}boundary ties.\allowbreak{} This clarifies the stated model, not a corrected admission policy or a new economic assumption.\allowbreak{}
\end{sloppypar}

\paragraph{Lean-expanded paper semantic target}
\begin{ReviewVerbatim}
type:
{n : ℕ} → (ℝ → ℝ) → (Bool × ℝ → ℝ) → (Bool × ℝ → ℝ) → (Bool × ℝ → ℝ) → (Fin (n + 1) → ℝ) → (ℕ → ℝ) → ℝ → ℝ

value:
fun {n : ℕ} (cost : ℝ → ℝ) (effort score tie : Bool × ℝ → ℝ) (cutoff : Fin (n + 1) → ℝ) (reward : ℕ → ℝ) (t : ℝ) =>
  LBG22StrategicRanking.sourceActualGroupWelfare cost effort score tie cutoff reward true t -
    LBG22StrategicRanking.sourceActualGroupWelfare cost effort score tie cutoff reward false t
\end{ReviewVerbatim}

\paragraph{Direct retained dependencies}
\begin{itemize}\raggedright
\item \hyperlink{paper-prerequisite-22c0ae11547a174a}{LBG22\allowbreak{}Strategic\allowbreak{}Ranking.\allowbreak{}source\allowbreak{}Actual\allowbreak{}Group\allowbreak{}Welfare}
\end{itemize}
\paragraph{Recorded source-to-declaration screening}
\textbf{Verdict:} \texttt{matches}
\\
{\raggedright
\textbf{Reason:} The target defines the pointwise welfare gap as advantaged-\allowbreak{}group welfare minus disadvantaged-\allowbreak{}group welfare at the same latent rank, using the displayed source welfare function.\allowbreak{} This matches the source definition exactly.\allowbreak{}
\par}
\reviewmatch{Matches source input}
\noindent\textbf{Reviewer annotation}\par
\reviewerbox
\clearpage
\hypertarget{paper-prerequisite-c194f1bd51c4a4cc}{}
\subsection*{\small\ttfamily\raggedright LBG22\allowbreak{}Strategic\allowbreak{}Ranking.\allowbreak{}source\allowbreak{}Actual\allowbreak{}Multitask\allowbreak{}School\allowbreak{}Utility}
\noindent\textbf{Paper-local declaration:}\par
{\footnotesize\ttfamily\raggedright LBG22\allowbreak{}Strategic\allowbreak{}Ranking.\allowbreak{}source\allowbreak{}Actual\allowbreak{}Multitask\allowbreak{}School\allowbreak{}Utility\par}
\textbf{Source locator:} {\footnotesize\raggedright cited publication, lines 1809-\allowbreak{}1812\par}
\paragraph{Verbatim paper-source connection}
\begin{ReviewVerbatim}
The school's private utility is now weighted by $\beta \in (0,1)$, which quantifies the relative value the school places on the measurable skill over the unmeasurable skill.
\begin{equation*}
\Util\pri_\beta = \E[\beta\cdot v^M + (1-\beta)\cdot v^U\mid Z=1].
\end{equation*}
\end{ReviewVerbatim}

% report-context-presentation-sha256: f1061b67864489b0b3af8011b46b77ed4e156ac6d59630d74f61e31a6aa20b6f
\paragraph{Settled source reading}
\begin{sloppypar}
The multitask model has nonnegative effort on both tasks summing to the stated budget.\allowbreak{} This is distinct from the unrestricted total effort in the multidimensional ranking model.\allowbreak{}
\end{sloppypar}
\paragraph{Settled source reading}
\begin{sloppypar}
Applicants know the fixed priority used at score-\allowbreak{}boundary ties.\allowbreak{} This clarifies the stated model, not a corrected admission policy or a new economic assumption.\allowbreak{}
\end{sloppypar}

\paragraph{Lean-expanded paper semantic target}
\begin{ReviewVerbatim}
type:
{α : Type u_1} →
  [inst : MeasurableSpace α] →
    (μ : MeasureTheory.Measure α) →
      [MeasureTheory.IsProbabilityMeasure μ] →
        (ℝ → ℝ) → (ℝ → ℝ) → (α → ℝ) → (ℝ × α → ℝ) → (ℝ × α → ℝ) → (ℝ × α → ℝ) → (ℝ × α → ℝ) → ℝ → ℝ → ℝ → ℝ

value:
fun {α : Type u_1} [MeasurableSpace α] (μ : MeasureTheory.Measure α) [MeasureTheory.IsProbabilityMeasure μ]
    (production skillM : ℝ → ℝ) (skillU : α → ℝ) (effortM effortU score tie : ℝ × α → ℝ) (rho beta c : ℝ) =>
  ∫ (z : (ℝ × ℝ) × α),
    beta * (production (effortM (z.1.1, z.2)) * skillM z.1.1) +
      (1 - beta) *
        (production (effortU (z.1.1, z.2)) *
          skillU
            z.2) ∂(LBG22StrategicRanking.independentSkillPopulation
        μ)[|LBG22StrategicRanking.sourceMultitaskAdmissionEvent μ score tie rho c]
\end{ReviewVerbatim}

\paragraph{Direct retained dependencies}
\begin{itemize}\raggedright
\item \hyperlink{governing-declaration-8d3290d5e2a0343e}{LBG22\allowbreak{}Strategic\allowbreak{}Ranking.\allowbreak{}independent\allowbreak{}Skill\allowbreak{}Population}
\item \hyperlink{governing-declaration-f1a76532081bf3d3}{LBG22\allowbreak{}Strategic\allowbreak{}Ranking.\allowbreak{}source\allowbreak{}Multitask\allowbreak{}Admission\allowbreak{}Event}
\end{itemize}
\paragraph{Recorded source-to-declaration screening}
\textbf{Verdict:} \texttt{matches}
\\
{\raggedright
\textbf{Reason:} The target conditions the independent measurable/\allowbreak{}unmeasurable skill population on actual admission and integrates beta times measurable production plus (1-\allowbreak{}beta) times unmeasurable production.\allowbreak{} This is the source conditional weighted school utility, with the approved hard-\allowbreak{}budget and fixed-\allowbreak{}priority readings.\allowbreak{}
\par}
\reviewmatch{Matches source input}
\noindent\textbf{Reviewer annotation}\par
\reviewerbox
\clearpage
\hypertarget{paper-prerequisite-9d0a441d993fe3f2}{}
\subsection*{\small\ttfamily\raggedright LBG22\allowbreak{}Strategic\allowbreak{}Ranking.\allowbreak{}source\allowbreak{}Actual\allowbreak{}Two\allowbreak{}Level\allowbreak{}School\allowbreak{}Utility}
\noindent\textbf{Paper-local declaration:}\par
{\footnotesize\ttfamily\raggedright LBG22\allowbreak{}Strategic\allowbreak{}Ranking.\allowbreak{}source\allowbreak{}Actual\allowbreak{}Two\allowbreak{}Level\allowbreak{}School\allowbreak{}Utility\par}
\textbf{Source locator:} {\footnotesize\raggedright cited publication, lines 1118-\allowbreak{}1122\par}
\paragraph{Verbatim paper-source connection}
\begin{ReviewVerbatim}
Finally, the design $\lambda$ is controlled by a school, who may only draw value from those who enroll. The school's \textbf{private utility} is the expected score of admitted applicants,
which in our continuum formulation is the expectation of $v$ weighted by $\lambda(v)$: %\footnote{See Remark~\ref{rem:prob} on interpreting $\Util\pri$ as a conditional expectation.}
%\begin{equation*}
\(	\Util\pri:= \E[v\cdot \lambda(\theta_\post)]\).
%\end{equation*}
\end{ReviewVerbatim}

% report-context-presentation-sha256: 827a19afe0274b1faef13e344ddbcfe40d9841f2ff34b92ac9d53c0f83e29f56
\paragraph{Settled source reading}
\begin{sloppypar}
Applicants know the fixed priority used at score-\allowbreak{}boundary ties.\allowbreak{} This clarifies the stated model, not a corrected admission policy or a new economic assumption.\allowbreak{}
\end{sloppypar}

\paragraph{Lean-expanded paper semantic target}
\begin{ReviewVerbatim}
type:
(ℝ → ℝ) → (ℝ → ℝ) → ℝ → ℝ → ℝ

value:
fun (score tie : ℝ → ℝ) (rho c : ℝ) =>
  ∫ (t : ℝ),
    score t *
      LBG22StrategicRanking.sourceTwoLevelReward rho c
        ↑(LBG22StrategicRanking.finiteLowerRankBand (fun (i : Fin 2) => LBG22StrategicRanking.sourceTwoLevelCutoff c ↑i)
            (LBG22StrategicRanking.tieBrokenRank LBG22StrategicRanking.unitRankMeasure score tie
              t)) ∂LBG22StrategicRanking.unitRankMeasure
\end{ReviewVerbatim}

\paragraph{Direct retained dependencies}
\begin{itemize}\raggedright
\item \hyperlink{paper-prerequisite-a42f440ddc254fa7}{LBG22\allowbreak{}Strategic\allowbreak{}Ranking.\allowbreak{}finite\allowbreak{}Lower\allowbreak{}Rank\allowbreak{}Band}
\item \hyperlink{paper-prerequisite-f88aac58507e44ca}{LBG22\allowbreak{}Strategic\allowbreak{}Ranking.\allowbreak{}source\allowbreak{}Two\allowbreak{}Level\allowbreak{}Cutoff}
\item \hyperlink{paper-prerequisite-30d97cd43eb79017}{LBG22\allowbreak{}Strategic\allowbreak{}Ranking.\allowbreak{}source\allowbreak{}Two\allowbreak{}Level\allowbreak{}Reward}
\item \hyperlink{paper-prerequisite-b7c71e53177b76c6}{LBG22\allowbreak{}Strategic\allowbreak{}Ranking.\allowbreak{}tie\allowbreak{}Broken\allowbreak{}Rank}
\item \hyperlink{governing-declaration-678cf48d549cb620}{LBG22\allowbreak{}Strategic\allowbreak{}Ranking.\allowbreak{}unit\allowbreak{}Rank\allowbreak{}Measure}
\end{itemize}
\paragraph{Recorded source-to-declaration screening}
\textbf{Verdict:} \texttt{matches}
\\
{\raggedright
\textbf{Reason:} The target integrates realized score multiplied by the actual two-\allowbreak{}level admission reward over the unit applicant population.\allowbreak{} This is the source private utility E[v lambda(theta\_\allowbreak{}post)], with all ranking and band support displayed.\allowbreak{}
\par}
\reviewmatch{Matches source input}
\noindent\textbf{Reviewer annotation}\par
\reviewerbox
\clearpage
\hypertarget{paper-prerequisite-e76c8574ddaca0dc}{}
\subsection*{\small\ttfamily\raggedright LBG22\allowbreak{}Strategic\allowbreak{}Ranking.\allowbreak{}source\allowbreak{}Deterministic\allowbreak{}Cutoff}
\noindent\textbf{Paper-local declaration:}\par
{\footnotesize\ttfamily\raggedright LBG22\allowbreak{}Strategic\allowbreak{}Ranking.\allowbreak{}source\allowbreak{}Deterministic\allowbreak{}Cutoff\par}
\textbf{Source locator:} {\footnotesize\raggedright cited publication, lines 105-\allowbreak{}105\par}
\paragraph{Verbatim paper-source connection}
\begin{ReviewVerbatim}
Note that standard non-randomized admissions policies are equivalent to case where $c = 1 - \rho$: the highest scoring applicants up to the capacity constraint are accepted with probability 1 and all others are rejected. We call this case \textbf{non-randomized admissions}. The other extreme is the one-level policy, \textbf{pure randomization}, where each applicant is admitted with probability ${\rho}$, i.e., $\ell_0 =\rho$. Decreasing the cut-off $c$ can be viewed as increasing the level of randomization in the admissions policy.
\end{ReviewVerbatim}

% report-context-presentation-sha256: 827a19afe0274b1faef13e344ddbcfe40d9841f2ff34b92ac9d53c0f83e29f56
\paragraph{Settled source reading}
\begin{sloppypar}
Applicants know the fixed priority used at score-\allowbreak{}boundary ties.\allowbreak{} This clarifies the stated model, not a corrected admission policy or a new economic assumption.\allowbreak{}
\end{sloppypar}

\paragraph{Lean-expanded paper semantic target}
\begin{ReviewVerbatim}
type:
ℝ → ℝ

value:
fun (rho : ℝ) => 1 - rho
\end{ReviewVerbatim}


\paragraph{Recorded source-to-declaration screening}
\textbf{Verdict:} \texttt{matches}
\\
{\raggedright
\textbf{Reason:} The target is exactly c=1-\allowbreak{}rho, the source cutoff for accepting the highest rho share with probability one under non-\allowbreak{}randomized admissions.\allowbreak{}
\par}
\reviewmatch{Matches source input}
\noindent\textbf{Reviewer annotation}\par
\reviewerbox
\clearpage
\hypertarget{paper-prerequisite-1aae1cdeecf21463}{}
\subsection*{\small\ttfamily\raggedright LBG22\allowbreak{}Strategic\allowbreak{}Ranking.\allowbreak{}source\allowbreak{}Environment\allowbreak{}CDF}
\noindent\textbf{Paper-local declaration:}\par
{\footnotesize\ttfamily\raggedright LBG22\allowbreak{}Strategic\allowbreak{}Ranking.\allowbreak{}source\allowbreak{}Environment\allowbreak{}CDF\par}
\textbf{Source locator:} {\footnotesize\raggedright cited publication, lines 224-\allowbreak{}233\par}
\paragraph{Verbatim paper-source connection}
\begin{ReviewVerbatim}
\begin{restatable}[Equilibrium under group differences]{proposition}{propEquiEnv}\label{prop:equi-env}
Define $\theta_\pre$ as:
%	\begin{equation*}
	\(\theta_\pre := f^{-1}_\mix(f(\theta_\true)\cdot \psi)\),
%	\end{equation*}
		where $f^{-1}_\mix$	is the CDF for the environment-scaled skill, $f(\theta_\true)\cdot \psi$:
	\begin{equation*}
	f^{-1}_\mix(x) := \frac{1}{2}f^{-1}(x/\psi_\A)+ \frac{1}{2}f^{-1}(x/\psi_\B).
	\end{equation*}
\end{ReviewVerbatim}



\paragraph{Lean-expanded paper semantic target}
\begin{ReviewVerbatim}
type:
(ℝ → ℝ) → ℝ → ℝ → ℝ → ℝ

value:
fun (skill : ℝ → ℝ) (psiA psiB a : ℝ) =>
  ↑(ProbabilityTheory.cdf
        (MeasureTheory.Measure.map (LBG22StrategicRanking.sourceEnvironmentSkill skill psiA psiB)
          LBG22StrategicRanking.sourceTwoGroupMeasure))
    a
\end{ReviewVerbatim}

\paragraph{Direct retained dependencies}
\begin{itemize}\raggedright
\item \hyperlink{paper-prerequisite-73413973c7594479}{LBG22\allowbreak{}Strategic\allowbreak{}Ranking.\allowbreak{}source\allowbreak{}Environment\allowbreak{}Skill}
\item \hyperlink{paper-prerequisite-5b3be494cde30a37}{LBG22\allowbreak{}Strategic\allowbreak{}Ranking.\allowbreak{}source\allowbreak{}Two\allowbreak{}Group\allowbreak{}Measure}
\end{itemize}
\paragraph{Recorded source-to-declaration screening}
\textbf{Verdict:} \texttt{matches}
\\
{\raggedright
\textbf{Reason:} The target is the CDF of environment-\allowbreak{}scaled skill under the equal-\allowbreak{}weight two-\allowbreak{}group measure.\allowbreak{} Its displayed map, clamped skill, and group measure match the source mixture CDF f\_\allowbreak{}mix inverse.\allowbreak{}
\par}
\reviewmatch{Matches source input}
\noindent\textbf{Reviewer annotation}\par
\reviewerbox
\clearpage
\hypertarget{paper-prerequisite-b308d7da226a5fe1}{}
\subsection*{\small\ttfamily\raggedright LBG22\allowbreak{}Strategic\allowbreak{}Ranking.\allowbreak{}source\allowbreak{}Environment\allowbreak{}Rank}
\noindent\textbf{Paper-local declaration:}\par
{\footnotesize\ttfamily\raggedright LBG22\allowbreak{}Strategic\allowbreak{}Ranking.\allowbreak{}source\allowbreak{}Environment\allowbreak{}Rank\par}
\textbf{Source locator:} {\footnotesize\raggedright cited publication, lines 224-\allowbreak{}233\par}
\paragraph{Verbatim paper-source connection}
\begin{ReviewVerbatim}
\begin{restatable}[Equilibrium under group differences]{proposition}{propEquiEnv}\label{prop:equi-env}
Define $\theta_\pre$ as:
%	\begin{equation*}
	\(\theta_\pre := f^{-1}_\mix(f(\theta_\true)\cdot \psi)\),
%	\end{equation*}
		where $f^{-1}_\mix$	is the CDF for the environment-scaled skill, $f(\theta_\true)\cdot \psi$:
	\begin{equation*}
	f^{-1}_\mix(x) := \frac{1}{2}f^{-1}(x/\psi_\A)+ \frac{1}{2}f^{-1}(x/\psi_\B).
	\end{equation*}
\end{ReviewVerbatim}



\paragraph{Lean-expanded paper semantic target}
\begin{ReviewVerbatim}
type:
(ℝ → ℝ) → ℝ → ℝ → Bool × ℝ → ℝ

value:
fun (skill : ℝ → ℝ) (psiA psiB : ℝ) (x : Bool × ℝ) =>
  LBG22StrategicRanking.sourceEnvironmentCDF skill psiA psiB
    (LBG22StrategicRanking.sourceEnvironmentSkill skill psiA psiB x)
\end{ReviewVerbatim}

\paragraph{Direct retained dependencies}
\begin{itemize}\raggedright
\item \hyperlink{paper-prerequisite-1aae1cdeecf21463}{LBG22\allowbreak{}Strategic\allowbreak{}Ranking.\allowbreak{}source\allowbreak{}Environment\allowbreak{}CDF}
\item \hyperlink{paper-prerequisite-73413973c7594479}{LBG22\allowbreak{}Strategic\allowbreak{}Ranking.\allowbreak{}source\allowbreak{}Environment\allowbreak{}Skill}
\end{itemize}
\paragraph{Recorded source-to-declaration screening}
\textbf{Verdict:} \texttt{matches}
\\
{\raggedright
\textbf{Reason:} The target evaluates the environment-\allowbreak{}scaled-\allowbreak{}skill mixture CDF at the applicant's own scaled skill, which is precisely the source environment-\allowbreak{}scaled pre-\allowbreak{}effort rank.\allowbreak{}
\par}
\reviewmatch{Matches source input}
\noindent\textbf{Reviewer annotation}\par
\reviewerbox
\clearpage
\hypertarget{paper-prerequisite-7fb5250e464d0736}{}
\subsection*{\small\ttfamily\raggedright LBG22\allowbreak{}Strategic\allowbreak{}Ranking.\allowbreak{}source\allowbreak{}Finite\allowbreak{}Applicant\allowbreak{}Welfare}
\noindent\textbf{Paper-local declaration:}\par
{\footnotesize\ttfamily\raggedright LBG22\allowbreak{}Strategic\allowbreak{}Ranking.\allowbreak{}source\allowbreak{}Finite\allowbreak{}Applicant\allowbreak{}Welfare\par}
\textbf{Source locator:} {\footnotesize\raggedright cited publication, lines 1107-\allowbreak{}1113\par}
\paragraph{Verbatim paper-source connection}
\begin{ReviewVerbatim}
\parbold{Aggregate welfare and utility} We define three aggregate welfare functions of the equilibrium efforts and scores to capture the interests of different stakeholders.
The \textbf{applicant welfare} is defined as the population average of the individual applicant welfare at equilibrium:
%\begin{definition}[applicant Welfare]
%	The population applicant welfare is defined as
\begin{align*}
	\swelf :=& \E[\iwelf(e,\lambda(\theta_\post))]	= \rho - \E[\cost(e)]. %& \text{Reward function $\lambda$ constraint $\rho$}
\end{align*}
\end{ReviewVerbatim}

% report-context-presentation-sha256: 827a19afe0274b1faef13e344ddbcfe40d9841f2ff34b92ac9d53c0f83e29f56
\paragraph{Settled source reading}
\begin{sloppypar}
Applicants know the fixed priority used at score-\allowbreak{}boundary ties.\allowbreak{} This clarifies the stated model, not a corrected admission policy or a new economic assumption.\allowbreak{}
\end{sloppypar}

\paragraph{Lean-expanded paper semantic target}
\begin{ReviewVerbatim}
type:
{n : ℕ} → (ℝ → ℝ) → (ℝ → ℝ) → (ℝ → ℝ) → (ℝ → ℝ) → (Fin (n + 1) → ℝ) → (ℕ → ℝ) → ℝ

value:
fun {n : ℕ} (cost effort score tie : ℝ → ℝ) (cutoff : Fin (n + 1) → ℝ) (reward : ℕ → ℝ) =>
  ∫ (t : ℝ),
    reward
        ↑(LBG22StrategicRanking.finiteLowerRankBand cutoff
            (LBG22StrategicRanking.tieBrokenRank LBG22StrategicRanking.unitRankMeasure score tie t)) -
      cost (effort t) ∂LBG22StrategicRanking.unitRankMeasure
\end{ReviewVerbatim}

\paragraph{Direct retained dependencies}
\begin{itemize}\raggedright
\item \hyperlink{paper-prerequisite-a42f440ddc254fa7}{LBG22\allowbreak{}Strategic\allowbreak{}Ranking.\allowbreak{}finite\allowbreak{}Lower\allowbreak{}Rank\allowbreak{}Band}
\item \hyperlink{paper-prerequisite-b7c71e53177b76c6}{LBG22\allowbreak{}Strategic\allowbreak{}Ranking.\allowbreak{}tie\allowbreak{}Broken\allowbreak{}Rank}
\item \hyperlink{governing-declaration-678cf48d549cb620}{LBG22\allowbreak{}Strategic\allowbreak{}Ranking.\allowbreak{}unit\allowbreak{}Rank\allowbreak{}Measure}
\end{itemize}
\paragraph{Recorded source-to-declaration screening}
\textbf{Verdict:} \texttt{matches}
\\
{\raggedright
\textbf{Reason:} The target integrates actual finite-\allowbreak{}policy admission reward minus effort cost over the uniform applicant population.\allowbreak{} This is the source aggregate applicant welfare definition and specializes to rho minus expected cost under capacity.\allowbreak{}
\par}
\reviewmatch{Matches source input}
\noindent\textbf{Reviewer annotation}\par
\reviewerbox
\clearpage
\hypertarget{paper-prerequisite-25a6dc17f260546b}{}
\subsection*{\small\ttfamily\raggedright LBG22\allowbreak{}Strategic\allowbreak{}Ranking.\allowbreak{}source\allowbreak{}Finite\allowbreak{}School\allowbreak{}Utility}
\noindent\textbf{Paper-local declaration:}\par
{\footnotesize\ttfamily\raggedright LBG22\allowbreak{}Strategic\allowbreak{}Ranking.\allowbreak{}source\allowbreak{}Finite\allowbreak{}School\allowbreak{}Utility\par}
\textbf{Source locator:} {\footnotesize\raggedright cited publication, lines 1118-\allowbreak{}1122\par}
\paragraph{Verbatim paper-source connection}
\begin{ReviewVerbatim}
Finally, the design $\lambda$ is controlled by a school, who may only draw value from those who enroll. The school's \textbf{private utility} is the expected score of admitted applicants,
which in our continuum formulation is the expectation of $v$ weighted by $\lambda(v)$: %\footnote{See Remark~\ref{rem:prob} on interpreting $\Util\pri$ as a conditional expectation.}
%\begin{equation*}
\(	\Util\pri:= \E[v\cdot \lambda(\theta_\post)]\).
%\end{equation*}
\end{ReviewVerbatim}

% report-context-presentation-sha256: 827a19afe0274b1faef13e344ddbcfe40d9841f2ff34b92ac9d53c0f83e29f56
\paragraph{Settled source reading}
\begin{sloppypar}
Applicants know the fixed priority used at score-\allowbreak{}boundary ties.\allowbreak{} This clarifies the stated model, not a corrected admission policy or a new economic assumption.\allowbreak{}
\end{sloppypar}

\paragraph{Lean-expanded paper semantic target}
\begin{ReviewVerbatim}
type:
{n : ℕ} → (ℝ → ℝ) → (ℝ → ℝ) → (Fin (n + 1) → ℝ) → (ℕ → ℝ) → ℝ

value:
fun {n : ℕ} (score tie : ℝ → ℝ) (cutoff : Fin (n + 1) → ℝ) (reward : ℕ → ℝ) =>
  ∫ (t : ℝ),
    score t *
      reward
        ↑(LBG22StrategicRanking.finiteLowerRankBand cutoff
            (LBG22StrategicRanking.tieBrokenRank LBG22StrategicRanking.unitRankMeasure score tie
              t)) ∂LBG22StrategicRanking.unitRankMeasure
\end{ReviewVerbatim}

\paragraph{Direct retained dependencies}
\begin{itemize}\raggedright
\item \hyperlink{paper-prerequisite-a42f440ddc254fa7}{LBG22\allowbreak{}Strategic\allowbreak{}Ranking.\allowbreak{}finite\allowbreak{}Lower\allowbreak{}Rank\allowbreak{}Band}
\item \hyperlink{paper-prerequisite-b7c71e53177b76c6}{LBG22\allowbreak{}Strategic\allowbreak{}Ranking.\allowbreak{}tie\allowbreak{}Broken\allowbreak{}Rank}
\item \hyperlink{governing-declaration-678cf48d549cb620}{LBG22\allowbreak{}Strategic\allowbreak{}Ranking.\allowbreak{}unit\allowbreak{}Rank\allowbreak{}Measure}
\end{itemize}
\paragraph{Recorded source-to-declaration screening}
\textbf{Verdict:} \texttt{matches}
\\
{\raggedright
\textbf{Reason:} The target integrates realized score times actual finite-\allowbreak{}policy admission reward over the applicant population.\allowbreak{} This exactly matches the source private school utility definition.\allowbreak{}
\par}
\reviewmatch{Matches source input}
\noindent\textbf{Reviewer annotation}\par
\reviewerbox
\clearpage
\hypertarget{paper-prerequisite-b2f8283910c81320}{}
\subsection*{\small\ttfamily\raggedright LBG22\allowbreak{}Strategic\allowbreak{}Ranking.\allowbreak{}source\allowbreak{}Group\allowbreak{}Threshold}
\noindent\textbf{Paper-local declaration:}\par
{\footnotesize\ttfamily\raggedright LBG22\allowbreak{}Strategic\allowbreak{}Ranking.\allowbreak{}source\allowbreak{}Group\allowbreak{}Threshold\par}
\textbf{Source locator:} {\footnotesize\raggedright cited publication, lines 282-\allowbreak{}285\par}
\paragraph{Verbatim paper-source connection}
\begin{ReviewVerbatim}
Denote the group-specific rank threshold for group $\G$ as
	\begin{equation*}
	\thresG{c} := f^{-1}\left(\frac{f_\mix(c)}{\psi_\G}\right).
	\end{equation*}
\end{ReviewVerbatim}



\paragraph{Lean-expanded paper semantic target}
\begin{ReviewVerbatim}
type:
(ℝ → ℝ) → ℝ → ℝ → Bool → ℝ → ℝ

value:
fun (skill : ℝ → ℝ) (psiA psiB : ℝ) (group : Bool) (c : ℝ) =>
  if group = true then LBG22StrategicRanking.sourceDisadvantagedThreshold skill psiB psiA c
  else LBG22StrategicRanking.sourceDisadvantagedThreshold skill psiA psiB c
\end{ReviewVerbatim}

\paragraph{Direct retained dependencies}
\begin{itemize}\raggedright
\item \hyperlink{governing-declaration-b9a91eff0bdd1f65}{LBG22\allowbreak{}Strategic\allowbreak{}Ranking.\allowbreak{}source\allowbreak{}Disadvantaged\allowbreak{}Threshold}
\end{itemize}
\paragraph{Recorded source-to-declaration screening}
\textbf{Verdict:} \texttt{matches}
\\
{\raggedright
\textbf{Reason:} The target gives each group's latent-\allowbreak{}rank cutoff by applying the disadvantaged-\allowbreak{}threshold construction with the appropriate own and opposite environment multipliers.\allowbreak{} The displayed mixture CDF and skill support make this equivalent to f inverse of f\_\allowbreak{}mix(c)/\allowbreak{}psi\_\allowbreak{}G in the source.\allowbreak{}
\par}
\reviewmatch{Matches source input}
\noindent\textbf{Reviewer annotation}\par
\reviewerbox
\clearpage
\hypertarget{paper-prerequisite-366855b73d9ea1c5}{}
\subsection*{\small\ttfamily\raggedright LBG22\allowbreak{}Strategic\allowbreak{}Ranking.\allowbreak{}source\allowbreak{}Pure\allowbreak{}Randomization\allowbreak{}Reward}
\noindent\textbf{Paper-local declaration:}\par
{\footnotesize\ttfamily\raggedright LBG22\allowbreak{}Strategic\allowbreak{}Ranking.\allowbreak{}source\allowbreak{}Pure\allowbreak{}Randomization\allowbreak{}Reward\par}
\textbf{Source locator:} {\footnotesize\raggedright cited publication, lines 105-\allowbreak{}105\par}
\paragraph{Verbatim paper-source connection}
\begin{ReviewVerbatim}
Note that standard non-randomized admissions policies are equivalent to case where $c = 1 - \rho$: the highest scoring applicants up to the capacity constraint are accepted with probability 1 and all others are rejected. We call this case \textbf{non-randomized admissions}. The other extreme is the one-level policy, \textbf{pure randomization}, where each applicant is admitted with probability ${\rho}$, i.e., $\ell_0 =\rho$. Decreasing the cut-off $c$ can be viewed as increasing the level of randomization in the admissions policy.
\end{ReviewVerbatim}



\paragraph{Lean-expanded paper semantic target}
\begin{ReviewVerbatim}
type:
ℝ → ℕ → ℝ

value:
fun (rho : ℝ) (_level : ℕ) => rho
\end{ReviewVerbatim}


\paragraph{Recorded source-to-declaration screening}
\textbf{Verdict:} \texttt{matches}
\\
{\raggedright
\textbf{Reason:} The target returns rho at every reward level, exactly the source one-\allowbreak{}level pure-\allowbreak{}randomization policy.\allowbreak{}
\par}
\reviewmatch{Matches source input}
\noindent\textbf{Reviewer annotation}\par
\reviewerbox
\clearpage
\hypertarget{paper-prerequisite-7224ad4c7b0b8551}{}
\subsection*{\small\ttfamily\raggedright LBG22\allowbreak{}Strategic\allowbreak{}Ranking.\allowbreak{}source\allowbreak{}Societal\allowbreak{}Utility}
\noindent\textbf{Paper-local declaration:}\par
{\footnotesize\ttfamily\raggedright LBG22\allowbreak{}Strategic\allowbreak{}Ranking.\allowbreak{}source\allowbreak{}Societal\allowbreak{}Utility\par}
\textbf{Source locator:} {\footnotesize\raggedright cited publication, lines 1114-\allowbreak{}1114\par}
\paragraph{Verbatim paper-source connection}
\begin{ReviewVerbatim}
On the other hand, society derives value from the scores of applicants post-effort, leading to the following \textbf{societal utility}:$\quad\quad  	\Util\soc:= \E[ v].$
\end{ReviewVerbatim}



\paragraph{Lean-expanded paper semantic target}
\begin{ReviewVerbatim}
type:
(ℝ → ℝ) → ℝ

value:
fun (score : ℝ → ℝ) => ∫ (t : ℝ), score t ∂LBG22StrategicRanking.unitRankMeasure
\end{ReviewVerbatim}

\paragraph{Direct retained dependencies}
\begin{itemize}\raggedright
\item \hyperlink{governing-declaration-678cf48d549cb620}{LBG22\allowbreak{}Strategic\allowbreak{}Ranking.\allowbreak{}unit\allowbreak{}Rank\allowbreak{}Measure}
\end{itemize}
\paragraph{Recorded source-to-declaration screening}
\textbf{Verdict:} \texttt{matches}
\\
{\raggedright
\textbf{Reason:} The target integrates post-\allowbreak{}effort score over the unit applicant population, exactly the source societal utility E[v].\allowbreak{}
\par}
\reviewmatch{Matches source input}
\noindent\textbf{Reviewer annotation}\par
\reviewerbox

\clearpage
\hypertarget{source-claim-a38dac6d1a17c16d}{}
\hypertarget{source-section-f10b5f68e4acf3dc}{}
\section*{Source claims}
\subsection*{Review row 1}
\textbf{Source locator:} {\footnotesize\raggedright cited publication, lines 1161-\allowbreak{}1165\par}
\paragraph{Verbatim source input}
\begin{ReviewVerbatim}
\begin{restatable}[Rank preservation]{proposition}{lemassortative}
	\label{lem:rankpreservedindex}
In every equilibrium, $ \lambda(\theta_\post(\theta_\pre))=\lambda(\theta_\pre)$, up to %changes with
{sets of} measure 0. %, up to tie-breaking for applicants with identical post-effort skills.
\end{restatable}

[Next verbatim source excerpt]

\parbold{Applicants}
There is a unit mass of \textit{applicants}, indexed by an observed index $\omega \in [0,1]$ distributed uniformly.\footnote{The index $\omega$ should be interpreted as each applicant's ``name,'' uncorrelated with skill, used solely for tie-breaking.}  Each applicant has a latent (unobserved) skill level represented by some measurable function of $\omega$.
We assume that {the distribution of the skills} has no atoms and that the CDF of this distribution is strictly increasing.
Using the CDF to map the skill of an applicant to a rank in~$[0,1]$, each applicant gets an (unobserved) rank $\theta_\pre =\theta_\pre(\omega)$ which by our assumption on the CDF is again uniformly distributed in $[0, 1]$ (the higher the rank the better).  With this setup, the skill of an applicant with rank $\theta_\pre$ can be written as $f(\theta_\pre)$ where $f$ is a strictly increasing, continuous function.  {It will be notationally convenient to label applicants by their rank $\theta_\pre$, though the reader should note that in contrast to $\omega$, $\theta_\pre(\omega)$ is assumed to be unobservable.}

%3. Each applicant exerts some effort e, at cost \cost(e).


Each applicant chooses an \textit{effort} level $e \ge 0$, the result of which is an observed, post-effort \textit{score}, {$v=v(e,\theta_\pre) = {g(e)}\cdot{f(\theta_\pre)}$}. \lledit{In other words, we assume the post-effort score is the product of two components, associated with the effort and the pre-effort skill respectively.}
{We assume} that the effort transfer function $g: [0, \infty) \mapsto [0, \infty)$ is a continuous, concave, strictly increasing function, representing that marginal effort improves one's score but has diminishing returns.
{%Since applicants can be labelled by their skill level $f(\theta_\pre)$, or equivalently, by their rank $\theta_{\pre}$,
{T}he strategies of the applicants can {then} be described by a function $\theta_\pre \mapsto e(\theta_\pre)$.
%\cb{$\omega\mapsto e(\omega)$.}%, leading to a distribution over post-effort scores $v=v(\theta_\pre))$ via the uniform distribution of $\theta_\pre$ in $[0,1]$.
}
  Each {applicant} is then ranked according to their score $v$%(where we break ties uniformly at random, if they appear)
, resulting in a \textit{post-effort rank} $\theta_\post$.
Note that the ranking $\theta_\post$ is slightly less trivial than a ranking of the skills, since the scores might have ties, which have to be resolved.

{\textbf{Tie Breaking.}
Given a choice of strategies $\theta_\pre \mapsto e(\theta_\pre)$, let $F$ be the CDF of the scores $v(e(\theta_\pre),\theta_\pre)$.
Since atoms for the distribution of $v$ would lead to ties for ranks defined as $F(v)$, we will use the labels of the applicants to break ties with the help of a (publicly announced) tie-breaking function $\Gamma(\omega)$, defined, e.g., via a collision free hash of the applicant names. Here  we require  that $\Gamma$  is a measurable function from $[0,1]$ to $[0,1]$ that maps different applicant labels to different values.
We then use $\Gamma$ to resolve the atoms of $F$, leading to a ranking function $v\mapsto \gamma(\omega, v)$ which is equal to $F(v)$ except when $v$ is an atom of the score distribution, {in which case it takes values in the ``gap interval'' $[F_-(v),  F(v)]$, where $F_-(v)$ is the left limit of $F$ at $v$.}  We  construct $\gamma$ in such a way that it
gives the uniform distribution for $\theta_\post(\omega)$ if we set  $\theta_\post(\omega)=\gamma(\omega,v(e(\theta_\pre(\omega)),\theta_\pre(\omega)))$.\footnote{See Remark~\ref{rem:gamma} for details on this construction.}

[Next verbatim source excerpt]

\parbold{Ranking designer (school)} A single \textit{school} is admitting applicants, based on their ranking. In particular, the school can choose a ranking reward function $\lambda : [0, 1] \mapsto [0, 1]$, such that an applicant with post-effort rank $\theta_\post$ is admitted with probability $\lambda(\theta_\post)$. We assume that $\lambda$ is non-decreasing and that the school has a constraint on the overall probability, such that in expectation it admits a number of applicants equal to {a} capacity constraint  $\rho \in (0,1)$, i.e., \lledit{$\E_{\theta_\post}[\lambda(\theta_\post)] = \rho $}.\footnote{\lledit{We use $\E_{\theta_\post}[\cdot]$ to denote an integral over $\theta_\post$ (and $\E[\cdot]$ to denote an integral over $\omega$)} with respect to the Lebesgue measure; informally, this can be thought of as averaging over the applicant population.} We may also refer to $\lambda$, informally, as the admission policy.



For simplicity, we further assume that $\lambda$ is a step-function with $K$ distinct levels $\ell_0 < \dots < \ell_{K-1}$, and $K-1$ cut-points parameterized by $c_1< \dots < c_{K-1}$ (with $c_0 = 0$, $c_K = 1$). %$\ell_0, \dots, \ell_{K-1}$.
In other words, we have $\lambda(\theta) = \ell_k$, for all $\theta \in \psi_k \triangleq [c_{k}, c_{k+1})$. Thus, applicants in the same post-rank interval $\theta \in \psi_k$ receive the same reward.

[Next verbatim source excerpt]

\parbold{Individual applicant welfare and equilibrium} Given the designer's function $\lambda$ and the effort levels of other applicants, each applicant chooses effort $e$ to maximize their individual welfare, % $W$,
\begin{equation*}
	\iwelf(e, \lambda(\theta_\post)) = \lambda(\theta_\post) - \cost(e),
\end{equation*}
where the effort cost function $\cost$ is non-negative, continuous, and strictly convex {on $[0,\infty)$, with $e_0:=\argmin_e \cost(e)$ and $p(e_0) = 0$.}
%\lnote{also want to assume that $\cost$ attains its minimum at $0$?}
%Let $e_0:=\argmin_e \cost(e)$ and $p(e_0) = 0$.

Applicants are assumed not to personally benefit from increasing their score $v$, except through the corresponding increase in their rank and reward. While the definition of $\cost$ does not preclude the applicant receiving any intrinsic benefit from exerting effort, we assume that the \emph{net} benefit from effort is non-positive.

[Next verbatim source excerpt]

\begin{definition} [Equilibrium]\label{def:equi}
{Given a}
%Let the
tie-breaking function $\Gamma(\omega)$
%be a pre-announced ordering over index $\omega$, to break ties in post-effort value. Then, let $\gamma(\omega, v)$ be the 1-1 ranking function defined as the CDF of post-effort values $v$ except where there are ties in $v$, in which case the ordering $\Gamma$ is used to produce unique ranks such that the range of $\gamma$ is uniform on $[0, 1]$.\footnote{More formally, note that each discontinuity in the CDF of post-effort values $v$ has a height equal to the measure of the applicants who are tied at that value, and so $\gamma$ for those applicants can be filled in using the CDF of $\Gamma(\omega)$ restricted to the applicants $\omega$ who are tied. Thus, the range of $\gamma$ is uniform on $[0,1]$.
%
{and}
%	Then, given
	a ranking probability function $\lambda$, an \textit{equilibrium} is a set of effort levels and post-effort ranking rewards for each applicant, $\{e(\theta_\pre),  \lambda(\theta_\post(\theta_\pre))\}$ such that, for all
	$\omega$, %$\theta_\pre$,
	\begin{align*}
		e(\theta_\pre(\omega)) &\in \argmax_{e} \iwelf\left(e, \lambda\left(\gamma\left(\omega, v\left(e,\theta_\pre(\omega)\right)\right)\right)\right) %\triangleq \argmax_{e} \left[ \lambda(\gamma(v(e,\theta_\pre)) - \cost(e) \right].
		\\
		\theta_\post(\theta_\pre(\omega)) &= \gamma(\omega, v(e(\theta_\pre(\omega)),\theta_\pre(\omega))),
	\end{align*}
	{where $\gamma(\omega,v)$ is the ranking induced by the CDF of the scores resulting from effort levels $\{e(\theta_\pre)\}$ and the tie-breaking function $\Gamma$.\footnote{See Remark~\ref{rem:gamma-effort} on $\gamma$ and the set of efforts.%Formally, $\gamma$ depends on the set of efforts. Given a fixed set and $\gamma$, for each applicant $\omega$ the first condition considers the counter-factual ranking of $\omega$ with different post-effort values but using the same ranking function. As defined, $\gamma$ yields a uniform distribution of ranks with such measure $0$ changes. In Appendix \Cref{lem:deviationsmeasure0}, we further prove that two effort sets equal up to sets of measure $0$ induce the same ranking function $\gamma$, and so the condition is consistent.%, and so we set it at the $\gamma$ induced by the effort levels $\{e(\theta_\pre)\}$.%, and so it does not vary.
	}
	}
	%In particular, we denote $\theta_\post(\theta_\pre) := \gamma(v(e(\theta_\pre),\theta_\pre))$.
	%\lledit{and for any $\theta_\pre, \theta'_\pre$, we have $\lambda(\theta_\post(\theta_\pre)) > \lambda(\theta_\post(\theta'_\pre)) \implies v(e(\theta_\pre),\theta_\pre) > v(e(\theta_\pre'),\theta_\pre')$.}
\end{definition}
\end{ReviewVerbatim}

% report-context-presentation-sha256: 90f9116e292dd71fddfd6b445ca0b060a481d026855bdb9725f2833c2cf62a31
\paragraph{Settled source reading}
\begin{sloppypar}
Equilibrium means that almost every applicant best responds against every feasible deviation;\allowbreak{} uniqueness is up to a null set.\allowbreak{}
\end{sloppypar}
\paragraph{Settled source reading}
\begin{sloppypar}
Applicants know the fixed priority used at score-\allowbreak{}boundary ties.\allowbreak{} This clarifies the stated model, not a corrected admission policy or a new economic assumption.\allowbreak{}
\end{sloppypar}

\paragraph{Expanded PaperInterface specification}
\begin{ReviewVerbatim}
∀ (cost production skill tie effort score : ℝ → ℝ) (baseline rho : ℝ) (n : ℕ) (cutoff reward : ℕ → ℝ),
  LBG22StrategicRanking.RankingPrimitives cost production skill baseline →
    LBG22StrategicRanking.FiniteAdmissionPolicy n cutoff reward rho →
      Measurable tie →
        Set.InjOn tie (Set.Ioc 0 1) →
          LBG22StrategicRanking.RankingEquilibrium cost production skill tie n cutoff reward effort score →
            ∀ᵐ (t : ℝ) ∂LBG22StrategicRanking.unitRankMeasure,
              reward
                  ↑(LBG22StrategicRanking.finiteLowerRankBand (fun (i : Fin (n + 1)) => cutoff ↑i)
                      (LBG22StrategicRanking.tieBrokenRank LBG22StrategicRanking.unitRankMeasure score tie t)) =
                reward ↑(LBG22StrategicRanking.finiteLowerRankBand (fun (i : Fin (n + 1)) => cutoff ↑i) t)
\end{ReviewVerbatim}

\paragraph{Retained governing declarations}
\begin{itemize}\raggedright
\item \hyperlink{paper-prerequisite-d056dd7c8060653c}{LBG22\allowbreak{}Strategic\allowbreak{}Ranking.\allowbreak{}Finite\allowbreak{}Admission\allowbreak{}Policy}
\item \hyperlink{paper-prerequisite-2a1b692f4e3136c6}{LBG22\allowbreak{}Strategic\allowbreak{}Ranking.\allowbreak{}Ranking\allowbreak{}Equilibrium}
\item \hyperlink{paper-prerequisite-9bbcc12f4a1736bf}{LBG22\allowbreak{}Strategic\allowbreak{}Ranking.\allowbreak{}Ranking\allowbreak{}Primitives}
\item \hyperlink{paper-prerequisite-a42f440ddc254fa7}{LBG22\allowbreak{}Strategic\allowbreak{}Ranking.\allowbreak{}finite\allowbreak{}Lower\allowbreak{}Rank\allowbreak{}Band}
\item \hyperlink{paper-prerequisite-b7c71e53177b76c6}{LBG22\allowbreak{}Strategic\allowbreak{}Ranking.\allowbreak{}tie\allowbreak{}Broken\allowbreak{}Rank}
\item \hyperlink{governing-declaration-678cf48d549cb620}{LBG22\allowbreak{}Strategic\allowbreak{}Ranking.\allowbreak{}unit\allowbreak{}Rank\allowbreak{}Measure}
\end{itemize}
\paragraph{Lean-elaborated claim roles}\begin{ReviewVerbatim}
1. parameter: ℝ → ℝ
2. parameter: ℝ → ℝ
3. parameter: ℝ → ℝ
4. parameter: ℝ → ℝ
5. parameter: ℝ → ℝ
6. parameter: ℝ → ℝ
7. parameter: ℝ
8. parameter: ℝ
9. parameter: ℕ
10. parameter: ℕ → ℝ
11. parameter: ℕ → ℝ
12. assumption: LBG22StrategicRanking.RankingPrimitives cost production skill baseline
13. assumption: LBG22StrategicRanking.FiniteAdmissionPolicy n cutoff reward rho
14. assumption: Measurable tie
15. assumption: Set.InjOn tie (Set.Ioc 0 1)
16. assumption: LBG22StrategicRanking.RankingEquilibrium cost production skill tie n cutoff reward effort score
17. conclusion: ∀ᵐ (t : ℝ) ∂LBG22StrategicRanking.unitRankMeasure,
  reward
      ↑(LBG22StrategicRanking.finiteLowerRankBand (fun (i : Fin (n + 1)) => cutoff ↑i)
          (LBG22StrategicRanking.tieBrokenRank LBG22StrategicRanking.unitRankMeasure score tie t)) =
    reward ↑(LBG22StrategicRanking.finiteLowerRankBand (fun (i : Fin (n + 1)) => cutoff ↑i) t)
\end{ReviewVerbatim}

\paragraph{Lean proof endpoint (built separately)}\begin{ReviewVerbatim}
LBG22StrategicRanking.ProofInterface.rankPreservation
\end{ReviewVerbatim}

\paragraph{Recorded source-to-Spec screening}
\textbf{Verdict:} \texttt{matches}
\\
{\raggedright
\textbf{Reason:} In every finite-\allowbreak{}policy equilibrium, the target states almost-\allowbreak{}everywhere equality between reward at actual fixed-\allowbreak{}priority post-\allowbreak{}rank and reward at pre-\allowbreak{}effort rank.\allowbreak{} This is Proposition 2.\allowbreak{}1, with its measure-\allowbreak{}zero qualification and all-\allowbreak{}deviation equilibrium reading supplied by the approved contexts and displayed dependencies.\allowbreak{}
\par}
\reviewmatch{Matches source input}
\noindent\textbf{Reviewer annotation}\par
\reviewerbox
\clearpage
\hypertarget{source-claim-6f58e02df112a54c}{}
\section*{Review row 2}
\textbf{Source locator:} {\footnotesize\raggedright cited publication, lines 1192-\allowbreak{}1207\par}
\paragraph{Verbatim source input}
\begin{ReviewVerbatim}
\begin{restatable}[Second price effort]{theorem}{lemeffort}
	\label{lem:effort}
	There exists a equilibrium such that $\lambda(\theta_\post(\theta_\pre))=\lambda(\theta_\pre)$ and applicants with $\theta_\pre \in \psi_k\triangleq{ [c_{k}, c_{k+1})}$ exert effort $e_k(\theta_\pre)$, where %\todo{replace i with k}
	\begin{align*}
		e_k(\theta_\pre) &= \begin{cases}
			e_0 & \text{for } k = 0\\
			\max\left(g^{-1}\left(\frac{g(\tilde e_{k-1})\cdot f(c_{k})}{f(\theta_\pre)} \right), e_0\right) & \text{otherwise,}
		\end{cases}
	\end{align*}
	\label{lem:secondpricecontimuum}
	with $\tilde e_{k-1}{\geq e_{k-1}(c_k)}$ {inductively defined by} $$
	%	\frac{\ell_k - \ell_{k-1}}{\cost(\tilde e_{k-1}) - \cost(e_{k-1}(c_k))} = 1.
	{\cost(\tilde e_{k-1})=\cost(e_{k-1}(c_k))+\ell_k-\ell_{k-1}}
		$$
		The equilibrium is unique up to sets of  measure 0. % up to ties in post-effort skills and {the} admissions probability.
\end{restatable}

[Next verbatim source excerpt]

\parbold{Applicants}
There is a unit mass of \textit{applicants}, indexed by an observed index $\omega \in [0,1]$ distributed uniformly.\footnote{The index $\omega$ should be interpreted as each applicant's ``name,'' uncorrelated with skill, used solely for tie-breaking.}  Each applicant has a latent (unobserved) skill level represented by some measurable function of $\omega$.
We assume that {the distribution of the skills} has no atoms and that the CDF of this distribution is strictly increasing.
Using the CDF to map the skill of an applicant to a rank in~$[0,1]$, each applicant gets an (unobserved) rank $\theta_\pre =\theta_\pre(\omega)$ which by our assumption on the CDF is again uniformly distributed in $[0, 1]$ (the higher the rank the better).  With this setup, the skill of an applicant with rank $\theta_\pre$ can be written as $f(\theta_\pre)$ where $f$ is a strictly increasing, continuous function.  {It will be notationally convenient to label applicants by their rank $\theta_\pre$, though the reader should note that in contrast to $\omega$, $\theta_\pre(\omega)$ is assumed to be unobservable.}

%3. Each applicant exerts some effort e, at cost \cost(e).


Each applicant chooses an \textit{effort} level $e \ge 0$, the result of which is an observed, post-effort \textit{score}, {$v=v(e,\theta_\pre) = {g(e)}\cdot{f(\theta_\pre)}$}. \lledit{In other words, we assume the post-effort score is the product of two components, associated with the effort and the pre-effort skill respectively.}
{We assume} that the effort transfer function $g: [0, \infty) \mapsto [0, \infty)$ is a continuous, concave, strictly increasing function, representing that marginal effort improves one's score but has diminishing returns.
{%Since applicants can be labelled by their skill level $f(\theta_\pre)$, or equivalently, by their rank $\theta_{\pre}$,
{T}he strategies of the applicants can {then} be described by a function $\theta_\pre \mapsto e(\theta_\pre)$.
%\cb{$\omega\mapsto e(\omega)$.}%, leading to a distribution over post-effort scores $v=v(\theta_\pre))$ via the uniform distribution of $\theta_\pre$ in $[0,1]$.
}
  Each {applicant} is then ranked according to their score $v$%(where we break ties uniformly at random, if they appear)
, resulting in a \textit{post-effort rank} $\theta_\post$.
Note that the ranking $\theta_\post$ is slightly less trivial than a ranking of the skills, since the scores might have ties, which have to be resolved.

{\textbf{Tie Breaking.}
Given a choice of strategies $\theta_\pre \mapsto e(\theta_\pre)$, let $F$ be the CDF of the scores $v(e(\theta_\pre),\theta_\pre)$.
Since atoms for the distribution of $v$ would lead to ties for ranks defined as $F(v)$, we will use the labels of the applicants to break ties with the help of a (publicly announced) tie-breaking function $\Gamma(\omega)$, defined, e.g., via a collision free hash of the applicant names. Here  we require  that $\Gamma$  is a measurable function from $[0,1]$ to $[0,1]$ that maps different applicant labels to different values.
We then use $\Gamma$ to resolve the atoms of $F$, leading to a ranking function $v\mapsto \gamma(\omega, v)$ which is equal to $F(v)$ except when $v$ is an atom of the score distribution, {in which case it takes values in the ``gap interval'' $[F_-(v),  F(v)]$, where $F_-(v)$ is the left limit of $F$ at $v$.}  We  construct $\gamma$ in such a way that it
gives the uniform distribution for $\theta_\post(\omega)$ if we set  $\theta_\post(\omega)=\gamma(\omega,v(e(\theta_\pre(\omega)),\theta_\pre(\omega)))$.\footnote{See Remark~\ref{rem:gamma} for details on this construction.}

[Next verbatim source excerpt]

\parbold{Ranking designer (school)} A single \textit{school} is admitting applicants, based on their ranking. In particular, the school can choose a ranking reward function $\lambda : [0, 1] \mapsto [0, 1]$, such that an applicant with post-effort rank $\theta_\post$ is admitted with probability $\lambda(\theta_\post)$. We assume that $\lambda$ is non-decreasing and that the school has a constraint on the overall probability, such that in expectation it admits a number of applicants equal to {a} capacity constraint  $\rho \in (0,1)$, i.e., \lledit{$\E_{\theta_\post}[\lambda(\theta_\post)] = \rho $}.\footnote{\lledit{We use $\E_{\theta_\post}[\cdot]$ to denote an integral over $\theta_\post$ (and $\E[\cdot]$ to denote an integral over $\omega$)} with respect to the Lebesgue measure; informally, this can be thought of as averaging over the applicant population.} We may also refer to $\lambda$, informally, as the admission policy.



For simplicity, we further assume that $\lambda$ is a step-function with $K$ distinct levels $\ell_0 < \dots < \ell_{K-1}$, and $K-1$ cut-points parameterized by $c_1< \dots < c_{K-1}$ (with $c_0 = 0$, $c_K = 1$). %$\ell_0, \dots, \ell_{K-1}$.
In other words, we have $\lambda(\theta) = \ell_k$, for all $\theta \in \psi_k \triangleq [c_{k}, c_{k+1})$. Thus, applicants in the same post-rank interval $\theta \in \psi_k$ receive the same reward.

[Next verbatim source excerpt]

\parbold{Individual applicant welfare and equilibrium} Given the designer's function $\lambda$ and the effort levels of other applicants, each applicant chooses effort $e$ to maximize their individual welfare, % $W$,
\begin{equation*}
	\iwelf(e, \lambda(\theta_\post)) = \lambda(\theta_\post) - \cost(e),
\end{equation*}
where the effort cost function $\cost$ is non-negative, continuous, and strictly convex {on $[0,\infty)$, with $e_0:=\argmin_e \cost(e)$ and $p(e_0) = 0$.}
%\lnote{also want to assume that $\cost$ attains its minimum at $0$?}
%Let $e_0:=\argmin_e \cost(e)$ and $p(e_0) = 0$.

Applicants are assumed not to personally benefit from increasing their score $v$, except through the corresponding increase in their rank and reward. While the definition of $\cost$ does not preclude the applicant receiving any intrinsic benefit from exerting effort, we assume that the \emph{net} benefit from effort is non-positive.

[Next verbatim source excerpt]

\begin{definition} [Equilibrium]\label{def:equi}
{Given a}
%Let the
tie-breaking function $\Gamma(\omega)$
%be a pre-announced ordering over index $\omega$, to break ties in post-effort value. Then, let $\gamma(\omega, v)$ be the 1-1 ranking function defined as the CDF of post-effort values $v$ except where there are ties in $v$, in which case the ordering $\Gamma$ is used to produce unique ranks such that the range of $\gamma$ is uniform on $[0, 1]$.\footnote{More formally, note that each discontinuity in the CDF of post-effort values $v$ has a height equal to the measure of the applicants who are tied at that value, and so $\gamma$ for those applicants can be filled in using the CDF of $\Gamma(\omega)$ restricted to the applicants $\omega$ who are tied. Thus, the range of $\gamma$ is uniform on $[0,1]$.
%
{and}
%	Then, given
	a ranking probability function $\lambda$, an \textit{equilibrium} is a set of effort levels and post-effort ranking rewards for each applicant, $\{e(\theta_\pre),  \lambda(\theta_\post(\theta_\pre))\}$ such that, for all
	$\omega$, %$\theta_\pre$,
	\begin{align*}
		e(\theta_\pre(\omega)) &\in \argmax_{e} \iwelf\left(e, \lambda\left(\gamma\left(\omega, v\left(e,\theta_\pre(\omega)\right)\right)\right)\right) %\triangleq \argmax_{e} \left[ \lambda(\gamma(v(e,\theta_\pre)) - \cost(e) \right].
		\\
		\theta_\post(\theta_\pre(\omega)) &= \gamma(\omega, v(e(\theta_\pre(\omega)),\theta_\pre(\omega))),
	\end{align*}
	{where $\gamma(\omega,v)$ is the ranking induced by the CDF of the scores resulting from effort levels $\{e(\theta_\pre)\}$ and the tie-breaking function $\Gamma$.\footnote{See Remark~\ref{rem:gamma-effort} on $\gamma$ and the set of efforts.%Formally, $\gamma$ depends on the set of efforts. Given a fixed set and $\gamma$, for each applicant $\omega$ the first condition considers the counter-factual ranking of $\omega$ with different post-effort values but using the same ranking function. As defined, $\gamma$ yields a uniform distribution of ranks with such measure $0$ changes. In Appendix \Cref{lem:deviationsmeasure0}, we further prove that two effort sets equal up to sets of measure $0$ induce the same ranking function $\gamma$, and so the condition is consistent.%, and so we set it at the $\gamma$ induced by the effort levels $\{e(\theta_\pre)\}$.%, and so it does not vary.
	}
	}
	%In particular, we denote $\theta_\post(\theta_\pre) := \gamma(v(e(\theta_\pre),\theta_\pre))$.
	%\lledit{and for any $\theta_\pre, \theta'_\pre$, we have $\lambda(\theta_\post(\theta_\pre)) > \lambda(\theta_\post(\theta'_\pre)) \implies v(e(\theta_\pre),\theta_\pre) > v(e(\theta_\pre'),\theta_\pre')$.}
\end{definition}
\end{ReviewVerbatim}

% report-context-presentation-sha256: 90f9116e292dd71fddfd6b445ca0b060a481d026855bdb9725f2833c2cf62a31
\paragraph{Settled source reading}
\begin{sloppypar}
Equilibrium means that almost every applicant best responds against every feasible deviation;\allowbreak{} uniqueness is up to a null set.\allowbreak{}
\end{sloppypar}
\paragraph{Settled source reading}
\begin{sloppypar}
Applicants know the fixed priority used at score-\allowbreak{}boundary ties.\allowbreak{} This clarifies the stated model, not a corrected admission policy or a new economic assumption.\allowbreak{}
\end{sloppypar}

\paragraph{Expanded PaperInterface specification}
\begin{ReviewVerbatim}
∀ (cost production skill tie : ℝ → ℝ) (baseline rho : ℝ) (n : ℕ) (cutoff reward : ℕ → ℝ),
  LBG22StrategicRanking.RankingPrimitives cost production skill baseline →
    LBG22StrategicRanking.FiniteAdmissionPolicy n cutoff reward rho →
      Measurable tie →
        Set.InjOn tie (Set.Ioc 0 1) →
          ∃ (E : ℝ),
            0 < E ∧
              cost (baseline + E) = 1 ∧
                have effort : ℝ → ℝ :=
                  LBG22StrategicRanking.sourceFiniteBaselineEffort cost production skill baseline E n cutoff reward;
                have score : ℝ → ℝ :=
                  LBG22StrategicRanking.sourceFiniteBaselineScore cost production skill baseline E n cutoff reward;
                AEMeasurable effort LBG22StrategicRanking.unitRankMeasure ∧
                  LBG22StrategicRanking.RankingEquilibrium cost production skill tie n cutoff reward effort score ∧
                    MeasureTheory.Measure.map
                          (LBG22StrategicRanking.tieBrokenRank LBG22StrategicRanking.unitRankMeasure score tie)
                          LBG22StrategicRanking.unitRankMeasure =
                        LBG22StrategicRanking.unitRankMeasure ∧
                      (∀ᵐ (t : ℝ) ∂LBG22StrategicRanking.unitRankMeasure,
                          reward
                              ↑(LBG22StrategicRanking.finiteLowerRankBand (fun (i : Fin (n + 1)) => cutoff ↑i)
                                  (LBG22StrategicRanking.tieBrokenRank LBG22StrategicRanking.unitRankMeasure score tie
                                    t)) =
                            reward
                              ↑(LBG22StrategicRanking.finiteLowerRankBand (fun (i : Fin (n + 1)) => cutoff ↑i) t)) ∧
                        ∀ (otherEffort otherScore : ℝ → ℝ),
                          LBG22StrategicRanking.RankingEquilibrium cost production skill tie n cutoff reward otherEffort
                              otherScore →
                            otherEffort =ᵐ[LBG22StrategicRanking.unitRankMeasure] effort
\end{ReviewVerbatim}

\paragraph{Retained governing declarations}
\begin{itemize}\raggedright
\item \hyperlink{paper-prerequisite-d056dd7c8060653c}{LBG22\allowbreak{}Strategic\allowbreak{}Ranking.\allowbreak{}Finite\allowbreak{}Admission\allowbreak{}Policy}
\item \hyperlink{paper-prerequisite-2a1b692f4e3136c6}{LBG22\allowbreak{}Strategic\allowbreak{}Ranking.\allowbreak{}Ranking\allowbreak{}Equilibrium}
\item \hyperlink{paper-prerequisite-9bbcc12f4a1736bf}{LBG22\allowbreak{}Strategic\allowbreak{}Ranking.\allowbreak{}Ranking\allowbreak{}Primitives}
\item \hyperlink{paper-prerequisite-a42f440ddc254fa7}{LBG22\allowbreak{}Strategic\allowbreak{}Ranking.\allowbreak{}finite\allowbreak{}Lower\allowbreak{}Rank\allowbreak{}Band}
\item \hyperlink{governing-declaration-a63079aa3a21a0d6}{LBG22\allowbreak{}Strategic\allowbreak{}Ranking.\allowbreak{}source\allowbreak{}Finite\allowbreak{}Baseline\allowbreak{}Effort}
\item \hyperlink{governing-declaration-927fd67abc7df008}{LBG22\allowbreak{}Strategic\allowbreak{}Ranking.\allowbreak{}source\allowbreak{}Finite\allowbreak{}Baseline\allowbreak{}Score}
\item \hyperlink{paper-prerequisite-b7c71e53177b76c6}{LBG22\allowbreak{}Strategic\allowbreak{}Ranking.\allowbreak{}tie\allowbreak{}Broken\allowbreak{}Rank}
\item \hyperlink{governing-declaration-678cf48d549cb620}{LBG22\allowbreak{}Strategic\allowbreak{}Ranking.\allowbreak{}unit\allowbreak{}Rank\allowbreak{}Measure}
\end{itemize}
\paragraph{Lean-elaborated claim roles}\begin{ReviewVerbatim}
1. parameter: ℝ → ℝ
2. parameter: ℝ → ℝ
3. parameter: ℝ → ℝ
4. parameter: ℝ → ℝ
5. parameter: ℝ
6. parameter: ℝ
7. parameter: ℕ
8. parameter: ℕ → ℝ
9. parameter: ℕ → ℝ
10. assumption: LBG22StrategicRanking.RankingPrimitives cost production skill baseline
11. assumption: LBG22StrategicRanking.FiniteAdmissionPolicy n cutoff reward rho
12. assumption: Measurable tie
13. assumption: Set.InjOn tie (Set.Ioc 0 1)
14. conclusion: ∃ (E : ℝ),
  0 < E ∧
    cost (baseline + E) = 1 ∧
      have effort : ℝ → ℝ :=
        LBG22StrategicRanking.sourceFiniteBaselineEffort cost production skill baseline E n cutoff reward;
      have score : ℝ → ℝ :=
        LBG22StrategicRanking.sourceFiniteBaselineScore cost production skill baseline E n cutoff reward;
      AEMeasurable effort LBG22StrategicRanking.unitRankMeasure ∧
        LBG22StrategicRanking.RankingEquilibrium cost production skill tie n cutoff reward effort score ∧
          MeasureTheory.Measure.map
                (LBG22StrategicRanking.tieBrokenRank LBG22StrategicRanking.unitRankMeasure score tie)
                LBG22StrategicRanking.unitRankMeasure =
              LBG22StrategicRanking.unitRankMeasure ∧
            (∀ᵐ (t : ℝ) ∂LBG22StrategicRanking.unitRankMeasure,
                reward
                    ↑(LBG22StrategicRanking.finiteLowerRankBand (fun (i : Fin (n + 1)) => cutoff ↑i)
                        (LBG22StrategicRanking.tieBrokenRank LBG22StrategicRanking.unitRankMeasure score tie t)) =
                  reward ↑(LBG22StrategicRanking.finiteLowerRankBand (fun (i : Fin (n + 1)) => cutoff ↑i) t)) ∧
              ∀ (otherEffort otherScore : ℝ → ℝ),
                LBG22StrategicRanking.RankingEquilibrium cost production skill tie n cutoff reward otherEffort
                    otherScore →
                  otherEffort =ᵐ[LBG22StrategicRanking.unitRankMeasure] effort
\end{ReviewVerbatim}

\paragraph{Lean proof endpoint (built separately)}\begin{ReviewVerbatim}
LBG22StrategicRanking.ProofInterface.secondPriceEffort
\end{ReviewVerbatim}

\paragraph{Recorded source-to-Spec screening}
\textbf{Verdict:} \texttt{matches}
\\
{\raggedright
\textbf{Reason:} The target constructs the source recursive boundary effort and score through the fully displayed inverse-\allowbreak{}cost, inverse-\allowbreak{}production, boundary-\allowbreak{}step, and band-\allowbreak{}recursion dependencies;\allowbreak{} it asserts a measurable equilibrium, uniform realized ranks, almost-\allowbreak{}everywhere reward preservation, and uniqueness of equilibrium effort almost everywhere.\allowbreak{} These are the existence, formula, rank-\allowbreak{}reward, and uniqueness conclusions of Theorem 2.\allowbreak{}2 under the approved continuum and priority conventions.\allowbreak{}
\par}
\reviewmatch{Matches source input}
\noindent\textbf{Reviewer annotation}\par
\reviewerbox
\clearpage
\hypertarget{source-claim-6f73613040867488}{}
\section*{Review row 3}
\textbf{Source locator:} {\footnotesize\raggedright cited publication, lines 135-\allowbreak{}137\par}
\paragraph{Verbatim source input}
\begin{ReviewVerbatim}
\begin{restatable}[Applicant welfare]{proposition}{propStudentWelfare}\label{prop:student-welfare}
	Among all $\lambda$ with $K$ levels, for $K\geq 1$, applicant welfare $\swelf$ is maximized by the one-level policy with pure randomization. Further, in the class of two-level policies, $\swelf$ is monotonically non-increasing in $c$.
\end{restatable}

[Next verbatim source excerpt]

\parbold{Applicants}
There is a unit mass of \textit{applicants}, indexed by an observed index $\omega \in [0,1]$ distributed uniformly.\footnote{The index $\omega$ should be interpreted as each applicant's ``name,'' uncorrelated with skill, used solely for tie-breaking.}  Each applicant has a latent (unobserved) skill level represented by some measurable function of $\omega$.
We assume that {the distribution of the skills} has no atoms and that the CDF of this distribution is strictly increasing.
Using the CDF to map the skill of an applicant to a rank in~$[0,1]$, each applicant gets an (unobserved) rank $\theta_\pre =\theta_\pre(\omega)$ which by our assumption on the CDF is again uniformly distributed in $[0, 1]$ (the higher the rank the better).  With this setup, the skill of an applicant with rank $\theta_\pre$ can be written as $f(\theta_\pre)$ where $f$ is a strictly increasing, continuous function.  {It will be notationally convenient to label applicants by their rank $\theta_\pre$, though the reader should note that in contrast to $\omega$, $\theta_\pre(\omega)$ is assumed to be unobservable.}

%3. Each applicant exerts some effort e, at cost \cost(e).


Each applicant chooses an \textit{effort} level $e \ge 0$, the result of which is an observed, post-effort \textit{score}, {$v=v(e,\theta_\pre) = {g(e)}\cdot{f(\theta_\pre)}$}. \lledit{In other words, we assume the post-effort score is the product of two components, associated with the effort and the pre-effort skill respectively.}
{We assume} that the effort transfer function $g: [0, \infty) \mapsto [0, \infty)$ is a continuous, concave, strictly increasing function, representing that marginal effort improves one's score but has diminishing returns.
{%Since applicants can be labelled by their skill level $f(\theta_\pre)$, or equivalently, by their rank $\theta_{\pre}$,
{T}he strategies of the applicants can {then} be described by a function $\theta_\pre \mapsto e(\theta_\pre)$.
%\cb{$\omega\mapsto e(\omega)$.}%, leading to a distribution over post-effort scores $v=v(\theta_\pre))$ via the uniform distribution of $\theta_\pre$ in $[0,1]$.
}
  Each {applicant} is then ranked according to their score $v$%(where we break ties uniformly at random, if they appear)
, resulting in a \textit{post-effort rank} $\theta_\post$.
Note that the ranking $\theta_\post$ is slightly less trivial than a ranking of the skills, since the scores might have ties, which have to be resolved.

{\textbf{Tie Breaking.}
Given a choice of strategies $\theta_\pre \mapsto e(\theta_\pre)$, let $F$ be the CDF of the scores $v(e(\theta_\pre),\theta_\pre)$.
Since atoms for the distribution of $v$ would lead to ties for ranks defined as $F(v)$, we will use the labels of the applicants to break ties with the help of a (publicly announced) tie-breaking function $\Gamma(\omega)$, defined, e.g., via a collision free hash of the applicant names. Here  we require  that $\Gamma$  is a measurable function from $[0,1]$ to $[0,1]$ that maps different applicant labels to different values.
We then use $\Gamma$ to resolve the atoms of $F$, leading to a ranking function $v\mapsto \gamma(\omega, v)$ which is equal to $F(v)$ except when $v$ is an atom of the score distribution, {in which case it takes values in the ``gap interval'' $[F_-(v),  F(v)]$, where $F_-(v)$ is the left limit of $F$ at $v$.}  We  construct $\gamma$ in such a way that it
gives the uniform distribution for $\theta_\post(\omega)$ if we set  $\theta_\post(\omega)=\gamma(\omega,v(e(\theta_\pre(\omega)),\theta_\pre(\omega)))$.\footnote{See Remark~\ref{rem:gamma} for details on this construction.}

[Next verbatim source excerpt]

\parbold{Ranking designer (school)} A single \textit{school} is admitting applicants, based on their ranking. In particular, the school can choose a ranking reward function $\lambda : [0, 1] \mapsto [0, 1]$, such that an applicant with post-effort rank $\theta_\post$ is admitted with probability $\lambda(\theta_\post)$. We assume that $\lambda$ is non-decreasing and that the school has a constraint on the overall probability, such that in expectation it admits a number of applicants equal to {a} capacity constraint  $\rho \in (0,1)$, i.e., \lledit{$\E_{\theta_\post}[\lambda(\theta_\post)] = \rho $}.\footnote{\lledit{We use $\E_{\theta_\post}[\cdot]$ to denote an integral over $\theta_\post$ (and $\E[\cdot]$ to denote an integral over $\omega$)} with respect to the Lebesgue measure; informally, this can be thought of as averaging over the applicant population.} We may also refer to $\lambda$, informally, as the admission policy.



For simplicity, we further assume that $\lambda$ is a step-function with $K$ distinct levels $\ell_0 < \dots < \ell_{K-1}$, and $K-1$ cut-points parameterized by $c_1< \dots < c_{K-1}$ (with $c_0 = 0$, $c_K = 1$). %$\ell_0, \dots, \ell_{K-1}$.
In other words, we have $\lambda(\theta) = \ell_k$, for all $\theta \in \psi_k \triangleq [c_{k}, c_{k+1})$. Thus, applicants in the same post-rank interval $\theta \in \psi_k$ receive the same reward.

[Next verbatim source excerpt]

\parbold{Individual applicant welfare and equilibrium} Given the designer's function $\lambda$ and the effort levels of other applicants, each applicant chooses effort $e$ to maximize their individual welfare, % $W$,
\begin{equation*}
	\iwelf(e, \lambda(\theta_\post)) = \lambda(\theta_\post) - \cost(e),
\end{equation*}
where the effort cost function $\cost$ is non-negative, continuous, and strictly convex {on $[0,\infty)$, with $e_0:=\argmin_e \cost(e)$ and $p(e_0) = 0$.}
%\lnote{also want to assume that $\cost$ attains its minimum at $0$?}
%Let $e_0:=\argmin_e \cost(e)$ and $p(e_0) = 0$.

Applicants are assumed not to personally benefit from increasing their score $v$, except through the corresponding increase in their rank and reward. While the definition of $\cost$ does not preclude the applicant receiving any intrinsic benefit from exerting effort, we assume that the \emph{net} benefit from effort is non-positive.

[Next verbatim source excerpt]

\begin{definition} [Equilibrium]\label{def:equi}
{Given a}
%Let the
tie-breaking function $\Gamma(\omega)$
%be a pre-announced ordering over index $\omega$, to break ties in post-effort value. Then, let $\gamma(\omega, v)$ be the 1-1 ranking function defined as the CDF of post-effort values $v$ except where there are ties in $v$, in which case the ordering $\Gamma$ is used to produce unique ranks such that the range of $\gamma$ is uniform on $[0, 1]$.\footnote{More formally, note that each discontinuity in the CDF of post-effort values $v$ has a height equal to the measure of the applicants who are tied at that value, and so $\gamma$ for those applicants can be filled in using the CDF of $\Gamma(\omega)$ restricted to the applicants $\omega$ who are tied. Thus, the range of $\gamma$ is uniform on $[0,1]$.
%
{and}
%	Then, given
	a ranking probability function $\lambda$, an \textit{equilibrium} is a set of effort levels and post-effort ranking rewards for each applicant, $\{e(\theta_\pre),  \lambda(\theta_\post(\theta_\pre))\}$ such that, for all
	$\omega$, %$\theta_\pre$,
	\begin{align*}
		e(\theta_\pre(\omega)) &\in \argmax_{e} \iwelf\left(e, \lambda\left(\gamma\left(\omega, v\left(e,\theta_\pre(\omega)\right)\right)\right)\right) %\triangleq \argmax_{e} \left[ \lambda(\gamma(v(e,\theta_\pre)) - \cost(e) \right].
		\\
		\theta_\post(\theta_\pre(\omega)) &= \gamma(\omega, v(e(\theta_\pre(\omega)),\theta_\pre(\omega))),
	\end{align*}
	{where $\gamma(\omega,v)$ is the ranking induced by the CDF of the scores resulting from effort levels $\{e(\theta_\pre)\}$ and the tie-breaking function $\Gamma$.\footnote{See Remark~\ref{rem:gamma-effort} on $\gamma$ and the set of efforts.%Formally, $\gamma$ depends on the set of efforts. Given a fixed set and $\gamma$, for each applicant $\omega$ the first condition considers the counter-factual ranking of $\omega$ with different post-effort values but using the same ranking function. As defined, $\gamma$ yields a uniform distribution of ranks with such measure $0$ changes. In Appendix \Cref{lem:deviationsmeasure0}, we further prove that two effort sets equal up to sets of measure $0$ induce the same ranking function $\gamma$, and so the condition is consistent.%, and so we set it at the $\gamma$ induced by the effort levels $\{e(\theta_\pre)\}$.%, and so it does not vary.
	}
	}
	%In particular, we denote $\theta_\post(\theta_\pre) := \gamma(v(e(\theta_\pre),\theta_\pre))$.
	%\lledit{and for any $\theta_\pre, \theta'_\pre$, we have $\lambda(\theta_\post(\theta_\pre)) > \lambda(\theta_\post(\theta'_\pre)) \implies v(e(\theta_\pre),\theta_\pre) > v(e(\theta_\pre'),\theta_\pre')$.}
\end{definition}

[Next verbatim source excerpt]

\begin{definition}[Two-level policy]\label{assump:2-level-fixed-cap} In our baseline two-level function parameterized by cut-off $c \in (0,1-\rho]$, each applicant with post-effort rank  $\theta_\post \geq c$ is admitted with probability $\ell_1 =\frac{\rho}{1-c} \in (0,1]$. Others are rejected, $\ell_0 = 0$.
\end{definition}

Note that standard non-randomized admissions policies are equivalent to case where $c = 1 - \rho$: the highest scoring applicants up to the capacity constraint are accepted with probability 1 and all others are rejected. We call this case \textbf{non-randomized admissions}. The other extreme is the one-level policy, \textbf{pure randomization}, where each applicant is admitted with probability ${\rho}$, i.e., $\ell_0 =\rho$. Decreasing the cut-off $c$ can be viewed as increasing the level of randomization in the admissions policy.


{We now reason about how various welfare and fairness metrics vary with $\lambda$ (in the two-level policy class, just $c$). To simplify the presentation, we assume that $f, g$ and $\cost$ are differentiable and that baseline effort $e_0$ is~$0$.}

[Next verbatim source excerpt]

\parbold{Aggregate welfare and utility} We define three aggregate welfare functions of the equilibrium efforts and scores to capture the interests of different stakeholders.
The \textbf{applicant welfare} is defined as the population average of the individual applicant welfare at equilibrium:
%\begin{definition}[applicant Welfare]
%	The population applicant welfare is defined as
\begin{align*}
	\swelf :=& \E[\iwelf(e,\lambda(\theta_\post))]	= \rho - \E[\cost(e)]. %& \text{Reward function $\lambda$ constraint $\rho$}
\end{align*}
On the other hand, society derives value from the scores of applicants post-effort, leading to the following \textbf{societal utility}:$\quad\quad  	\Util\soc:= \E[ v].$

In other words, society prefers the entire applicant population to achieve higher scores, not only those who are admitted to the school, since higher test scores are correlated with labor productivity and economic growth \citep{hanushek2010high}.

Finally, the design $\lambda$ is controlled by a school, who may only draw value from those who enroll. The school's \textbf{private utility} is the expected score of admitted applicants,
which in our continuum formulation is the expectation of $v$ weighted by $\lambda(v)$: %\footnote{See Remark~\ref{rem:prob} on interpreting $\Util\pri$ as a conditional expectation.}
%\begin{equation*}
\(	\Util\pri:= \E[v\cdot \lambda(\theta_\post)]\).
%\end{equation*}
\end{ReviewVerbatim}

% report-context-presentation-sha256: 90f9116e292dd71fddfd6b445ca0b060a481d026855bdb9725f2833c2cf62a31
\paragraph{Settled source reading}
\begin{sloppypar}
Equilibrium means that almost every applicant best responds against every feasible deviation;\allowbreak{} uniqueness is up to a null set.\allowbreak{}
\end{sloppypar}
\paragraph{Settled source reading}
\begin{sloppypar}
Applicants know the fixed priority used at score-\allowbreak{}boundary ties.\allowbreak{} This clarifies the stated model, not a corrected admission policy or a new economic assumption.\allowbreak{}
\end{sloppypar}
\paragraph{Additional assumptions}
\begin{sloppypar}
Only for two-level cutoff monotonicity: x maps to log f(1-exp(-x)) is concave on (0,infinity).; Only for two-level cutoff monotonicity: for E>0 with p(E)=1 and C=p composed with the inverse of g, z maps to log C(exp z) is concave on \{z : g(0)<exp z<=g(E)\}. Pure-randomization optimality has neither additional shape premise.
\end{sloppypar}
\paragraph{Formalized review target}
The archival source above is not asserted equivalent to this formalized target.
\begin{ReviewVerbatim}
Under the Section 2 primitives, baseline effort zero, the welfare section's differentiability conventions, capacity 0<rho<1, and fixed known tie priorities, pure randomization has an equilibrium attaining applicant welfare rho and dominates every finite-level-policy equilibrium, without additional shape assumptions. For the separate two-level monotonicity clause let E>0 satisfy p(E)=1 and C(x)=p(g^{-1}(x)) on the production interval [g(0),g(E)]. If x maps to log f(1-exp(-x)) is concave for x>0 and z maps to log C(exp z) is concave on g(0)<exp z<=g(E), then for any equilibrium family with the same primitives, population, priorities, and capacity, W(d)<=W(c) whenever 0<c<=d<=1-rho. Neither g(0)=0 nor f(0)>0 is imposed. The two sufficient concavity assumptions govern only this second clause.
\end{ReviewVerbatim}

\textbf{Archival source anchor:} {\footnotesize\raggedright cited publication, lines 135-\allowbreak{}137\par}
\paragraph{Expanded PaperInterface specification}
\begin{ReviewVerbatim}
∀ (cost production skill tie : ℝ → ℝ) (rho : ℝ),
  LBG22StrategicRanking.RankingPrimitives cost production skill 0 →
    rho ∈ Set.Ioo 0 1 →
      DifferentiableOn ℝ cost (Set.Ioi 0) →
        DifferentiableOn ℝ production (Set.Ioi 0) →
          DifferentiableOn ℝ skill (Set.Ioo 0 1) →
            Measurable tie →
              Set.InjOn tie (Set.Ioc 0 1) →
                (∃ (pureEffort : ℝ → ℝ) (pureScore : ℝ → ℝ),
                    LBG22StrategicRanking.RankingEquilibrium cost production skill tie 0
                        (fun (k : ℕ) => if k = 0 then 0 else 1)
                        (LBG22StrategicRanking.sourcePureRandomizationReward rho) pureEffort pureScore ∧
                      LBG22StrategicRanking.sourceFiniteApplicantWelfare cost pureEffort pureScore tie
                            (fun (x : Fin 1) => 0) (LBG22StrategicRanking.sourcePureRandomizationReward rho) =
                          rho ∧
                        ∀ (n : ℕ) (cutoff reward : ℕ → ℝ) (effort score : ℝ → ℝ),
                          LBG22StrategicRanking.FiniteAdmissionPolicy n cutoff reward rho →
                            LBG22StrategicRanking.RankingEquilibrium cost production skill tie n cutoff reward effort
                                score →
                              LBG22StrategicRanking.sourceFiniteApplicantWelfare cost effort score tie
                                  (fun (i : Fin (n + 1)) => cutoff ↑i) reward ≤
                                LBG22StrategicRanking.sourceFiniteApplicantWelfare cost pureEffort pureScore tie
                                  (fun (x : Fin 1) => 0) (LBG22StrategicRanking.sourcePureRandomizationReward rho)) ∧
                  ∀ (E : ℝ),
                    0 < E →
                      cost E = 1 →
                        (ConcaveOn ℝ (Set.Ioi 0) fun (x : ℝ) =>
                            Real.log (skill (LBG22StrategicRanking.upperTailRank x))) →
                          (ConcaveOn ℝ {z : ℝ | production 0 < Real.exp z ∧ Real.exp z ≤ production E} fun (z : ℝ) =>
                              Real.log (cost (LBG22StrategicRanking.effortIntervalInverse production E (Real.exp z)))) →
                            ∀ (effort score : ℝ → ℝ → ℝ),
                              (∀ c ∈ Set.Ioc 0 (1 - rho),
                                  LBG22StrategicRanking.RankingEquilibrium cost production skill tie 1
                                    (LBG22StrategicRanking.sourceTwoLevelCutoff c)
                                    (LBG22StrategicRanking.sourceTwoLevelReward rho c) (effort c) (score c)) →
                                AntitoneOn
                                  (fun (c : ℝ) =>
                                    LBG22StrategicRanking.sourceFiniteApplicantWelfare cost (effort c) (score c) tie
                                      (fun (i : Fin 2) => LBG22StrategicRanking.sourceTwoLevelCutoff c ↑i)
                                      (LBG22StrategicRanking.sourceTwoLevelReward rho c))
                                  (Set.Ioc 0 (1 - rho))
\end{ReviewVerbatim}

\paragraph{Retained governing declarations}
\begin{itemize}\raggedright
\item \hyperlink{paper-prerequisite-d056dd7c8060653c}{LBG22\allowbreak{}Strategic\allowbreak{}Ranking.\allowbreak{}Finite\allowbreak{}Admission\allowbreak{}Policy}
\item \hyperlink{paper-prerequisite-2a1b692f4e3136c6}{LBG22\allowbreak{}Strategic\allowbreak{}Ranking.\allowbreak{}Ranking\allowbreak{}Equilibrium}
\item \hyperlink{paper-prerequisite-9bbcc12f4a1736bf}{LBG22\allowbreak{}Strategic\allowbreak{}Ranking.\allowbreak{}Ranking\allowbreak{}Primitives}
\item \hyperlink{governing-declaration-c5b2c07abbff7224}{LBG22\allowbreak{}Strategic\allowbreak{}Ranking.\allowbreak{}effort\allowbreak{}Interval\allowbreak{}Inverse}
\item \hyperlink{paper-prerequisite-7fb5250e464d0736}{LBG22\allowbreak{}Strategic\allowbreak{}Ranking.\allowbreak{}source\allowbreak{}Finite\allowbreak{}Applicant\allowbreak{}Welfare}
\item \hyperlink{paper-prerequisite-366855b73d9ea1c5}{LBG22\allowbreak{}Strategic\allowbreak{}Ranking.\allowbreak{}source\allowbreak{}Pure\allowbreak{}Randomization\allowbreak{}Reward}
\item \hyperlink{paper-prerequisite-f88aac58507e44ca}{LBG22\allowbreak{}Strategic\allowbreak{}Ranking.\allowbreak{}source\allowbreak{}Two\allowbreak{}Level\allowbreak{}Cutoff}
\item \hyperlink{paper-prerequisite-30d97cd43eb79017}{LBG22\allowbreak{}Strategic\allowbreak{}Ranking.\allowbreak{}source\allowbreak{}Two\allowbreak{}Level\allowbreak{}Reward}
\item \hyperlink{governing-declaration-1f8a1d5f3261f80f}{LBG22\allowbreak{}Strategic\allowbreak{}Ranking.\allowbreak{}upper\allowbreak{}Tail\allowbreak{}Rank}
\end{itemize}
\paragraph{Lean-elaborated claim roles}\begin{ReviewVerbatim}
1. parameter: ℝ → ℝ
2. parameter: ℝ → ℝ
3. parameter: ℝ → ℝ
4. parameter: ℝ → ℝ
5. parameter: ℝ
6. assumption: LBG22StrategicRanking.RankingPrimitives cost production skill 0
7. assumption: rho ∈ Set.Ioo 0 1
8. assumption: DifferentiableOn ℝ cost (Set.Ioi 0)
9. assumption: DifferentiableOn ℝ production (Set.Ioi 0)
10. assumption: DifferentiableOn ℝ skill (Set.Ioo 0 1)
11. assumption: Measurable tie
12. assumption: Set.InjOn tie (Set.Ioc 0 1)
13. conclusion: (∃ (pureEffort : ℝ → ℝ) (pureScore : ℝ → ℝ),
    LBG22StrategicRanking.RankingEquilibrium cost production skill tie 0 (fun (k : ℕ) => if k = 0 then 0 else 1)
        (LBG22StrategicRanking.sourcePureRandomizationReward rho) pureEffort pureScore ∧
      LBG22StrategicRanking.sourceFiniteApplicantWelfare cost pureEffort pureScore tie (fun (x : Fin 1) => 0)
            (LBG22StrategicRanking.sourcePureRandomizationReward rho) =
          rho ∧
        ∀ (n : ℕ) (cutoff reward : ℕ → ℝ) (effort score : ℝ → ℝ),
          LBG22StrategicRanking.FiniteAdmissionPolicy n cutoff reward rho →
            LBG22StrategicRanking.RankingEquilibrium cost production skill tie n cutoff reward effort score →
              LBG22StrategicRanking.sourceFiniteApplicantWelfare cost effort score tie
                  (fun (i : Fin (n + 1)) => cutoff ↑i) reward ≤
                LBG22StrategicRanking.sourceFiniteApplicantWelfare cost pureEffort pureScore tie (fun (x : Fin 1) => 0)
                  (LBG22StrategicRanking.sourcePureRandomizationReward rho)) ∧
  ∀ (E : ℝ),
    0 < E →
      cost E = 1 →
        (ConcaveOn ℝ (Set.Ioi 0) fun (x : ℝ) => Real.log (skill (LBG22StrategicRanking.upperTailRank x))) →
          (ConcaveOn ℝ {z : ℝ | production 0 < Real.exp z ∧ Real.exp z ≤ production E} fun (z : ℝ) =>
              Real.log (cost (LBG22StrategicRanking.effortIntervalInverse production E (Real.exp z)))) →
            ∀ (effort score : ℝ → ℝ → ℝ),
              (∀ c ∈ Set.Ioc 0 (1 - rho),
                  LBG22StrategicRanking.RankingEquilibrium cost production skill tie 1
                    (LBG22StrategicRanking.sourceTwoLevelCutoff c) (LBG22StrategicRanking.sourceTwoLevelReward rho c)
                    (effort c) (score c)) →
                AntitoneOn
                  (fun (c : ℝ) =>
                    LBG22StrategicRanking.sourceFiniteApplicantWelfare cost (effort c) (score c) tie
                      (fun (i : Fin 2) => LBG22StrategicRanking.sourceTwoLevelCutoff c ↑i)
                      (LBG22StrategicRanking.sourceTwoLevelReward rho c))
                  (Set.Ioc 0 (1 - rho))
\end{ReviewVerbatim}

\paragraph{Lean proof endpoint (built separately)}\begin{ReviewVerbatim}
LBG22StrategicRanking.ProofInterface.applicantWelfare
\end{ReviewVerbatim}

\paragraph{Recorded source-to-Spec screening}
\textbf{Verdict:} \texttt{matches\_approved\_corrected\_target}
\\
{\raggedright
\textbf{Reason:} The selected source route disclaims archival equivalence and binds defect LBG22-\allowbreak{}C1 to a complete approved target.\allowbreak{} The expanded target separates the two clauses exactly:\allowbreak{} pure randomization has an equilibrium with welfare rho and dominates every finite-\allowbreak{}policy equilibrium without added shape assumptions, while two-\allowbreak{}level cutoff antitonicity alone assumes the approved two log-\allowbreak{}concavity conditions through upperTailRank and effortIntervalInverse, with E>0 and cost E=1.\allowbreak{} The displayed equilibrium, welfare integral, finite-\allowbreak{}policy, and priority prerequisites preserve the corrected target's quantifiers and domains.\allowbreak{}
\par}
\reviewmatch{Matches formalized target}
\noindent\textbf{Reviewer annotation}\par
\reviewerbox
\clearpage
\hypertarget{source-claim-1ba4d8c13750bf79}{}
\section*{Review row 4}
\textbf{Source locator:} {\footnotesize\raggedright cited publication, lines 148-\allowbreak{}151\par}
\paragraph{Verbatim source input}
\begin{ReviewVerbatim}
\begin{restatable}[School's private utility for two-level policies]{proposition}{propPrivateUtil}% on measurable skill is maximized by deterministic policy]
\label{prop:private-util-max}
In a two-level policy, the school's private utility~$\Util\pri$ is monotonically non-decreasing in $c$, and consequently is maximized by non-randomized admissions.
\end{restatable}

[Next verbatim source excerpt]

\parbold{Applicants}
There is a unit mass of \textit{applicants}, indexed by an observed index $\omega \in [0,1]$ distributed uniformly.\footnote{The index $\omega$ should be interpreted as each applicant's ``name,'' uncorrelated with skill, used solely for tie-breaking.}  Each applicant has a latent (unobserved) skill level represented by some measurable function of $\omega$.
We assume that {the distribution of the skills} has no atoms and that the CDF of this distribution is strictly increasing.
Using the CDF to map the skill of an applicant to a rank in~$[0,1]$, each applicant gets an (unobserved) rank $\theta_\pre =\theta_\pre(\omega)$ which by our assumption on the CDF is again uniformly distributed in $[0, 1]$ (the higher the rank the better).  With this setup, the skill of an applicant with rank $\theta_\pre$ can be written as $f(\theta_\pre)$ where $f$ is a strictly increasing, continuous function.  {It will be notationally convenient to label applicants by their rank $\theta_\pre$, though the reader should note that in contrast to $\omega$, $\theta_\pre(\omega)$ is assumed to be unobservable.}

%3. Each applicant exerts some effort e, at cost \cost(e).


Each applicant chooses an \textit{effort} level $e \ge 0$, the result of which is an observed, post-effort \textit{score}, {$v=v(e,\theta_\pre) = {g(e)}\cdot{f(\theta_\pre)}$}. \lledit{In other words, we assume the post-effort score is the product of two components, associated with the effort and the pre-effort skill respectively.}
{We assume} that the effort transfer function $g: [0, \infty) \mapsto [0, \infty)$ is a continuous, concave, strictly increasing function, representing that marginal effort improves one's score but has diminishing returns.
{%Since applicants can be labelled by their skill level $f(\theta_\pre)$, or equivalently, by their rank $\theta_{\pre}$,
{T}he strategies of the applicants can {then} be described by a function $\theta_\pre \mapsto e(\theta_\pre)$.
%\cb{$\omega\mapsto e(\omega)$.}%, leading to a distribution over post-effort scores $v=v(\theta_\pre))$ via the uniform distribution of $\theta_\pre$ in $[0,1]$.
}
  Each {applicant} is then ranked according to their score $v$%(where we break ties uniformly at random, if they appear)
, resulting in a \textit{post-effort rank} $\theta_\post$.
Note that the ranking $\theta_\post$ is slightly less trivial than a ranking of the skills, since the scores might have ties, which have to be resolved.

{\textbf{Tie Breaking.}
Given a choice of strategies $\theta_\pre \mapsto e(\theta_\pre)$, let $F$ be the CDF of the scores $v(e(\theta_\pre),\theta_\pre)$.
Since atoms for the distribution of $v$ would lead to ties for ranks defined as $F(v)$, we will use the labels of the applicants to break ties with the help of a (publicly announced) tie-breaking function $\Gamma(\omega)$, defined, e.g., via a collision free hash of the applicant names. Here  we require  that $\Gamma$  is a measurable function from $[0,1]$ to $[0,1]$ that maps different applicant labels to different values.
We then use $\Gamma$ to resolve the atoms of $F$, leading to a ranking function $v\mapsto \gamma(\omega, v)$ which is equal to $F(v)$ except when $v$ is an atom of the score distribution, {in which case it takes values in the ``gap interval'' $[F_-(v),  F(v)]$, where $F_-(v)$ is the left limit of $F$ at $v$.}  We  construct $\gamma$ in such a way that it
gives the uniform distribution for $\theta_\post(\omega)$ if we set  $\theta_\post(\omega)=\gamma(\omega,v(e(\theta_\pre(\omega)),\theta_\pre(\omega)))$.\footnote{See Remark~\ref{rem:gamma} for details on this construction.}

[Next verbatim source excerpt]

\parbold{Ranking designer (school)} A single \textit{school} is admitting applicants, based on their ranking. In particular, the school can choose a ranking reward function $\lambda : [0, 1] \mapsto [0, 1]$, such that an applicant with post-effort rank $\theta_\post$ is admitted with probability $\lambda(\theta_\post)$. We assume that $\lambda$ is non-decreasing and that the school has a constraint on the overall probability, such that in expectation it admits a number of applicants equal to {a} capacity constraint  $\rho \in (0,1)$, i.e., \lledit{$\E_{\theta_\post}[\lambda(\theta_\post)] = \rho $}.\footnote{\lledit{We use $\E_{\theta_\post}[\cdot]$ to denote an integral over $\theta_\post$ (and $\E[\cdot]$ to denote an integral over $\omega$)} with respect to the Lebesgue measure; informally, this can be thought of as averaging over the applicant population.} We may also refer to $\lambda$, informally, as the admission policy.



For simplicity, we further assume that $\lambda$ is a step-function with $K$ distinct levels $\ell_0 < \dots < \ell_{K-1}$, and $K-1$ cut-points parameterized by $c_1< \dots < c_{K-1}$ (with $c_0 = 0$, $c_K = 1$). %$\ell_0, \dots, \ell_{K-1}$.
In other words, we have $\lambda(\theta) = \ell_k$, for all $\theta \in \psi_k \triangleq [c_{k}, c_{k+1})$. Thus, applicants in the same post-rank interval $\theta \in \psi_k$ receive the same reward.

[Next verbatim source excerpt]

\parbold{Individual applicant welfare and equilibrium} Given the designer's function $\lambda$ and the effort levels of other applicants, each applicant chooses effort $e$ to maximize their individual welfare, % $W$,
\begin{equation*}
	\iwelf(e, \lambda(\theta_\post)) = \lambda(\theta_\post) - \cost(e),
\end{equation*}
where the effort cost function $\cost$ is non-negative, continuous, and strictly convex {on $[0,\infty)$, with $e_0:=\argmin_e \cost(e)$ and $p(e_0) = 0$.}
%\lnote{also want to assume that $\cost$ attains its minimum at $0$?}
%Let $e_0:=\argmin_e \cost(e)$ and $p(e_0) = 0$.

Applicants are assumed not to personally benefit from increasing their score $v$, except through the corresponding increase in their rank and reward. While the definition of $\cost$ does not preclude the applicant receiving any intrinsic benefit from exerting effort, we assume that the \emph{net} benefit from effort is non-positive.

[Next verbatim source excerpt]

\begin{definition} [Equilibrium]\label{def:equi}
{Given a}
%Let the
tie-breaking function $\Gamma(\omega)$
%be a pre-announced ordering over index $\omega$, to break ties in post-effort value. Then, let $\gamma(\omega, v)$ be the 1-1 ranking function defined as the CDF of post-effort values $v$ except where there are ties in $v$, in which case the ordering $\Gamma$ is used to produce unique ranks such that the range of $\gamma$ is uniform on $[0, 1]$.\footnote{More formally, note that each discontinuity in the CDF of post-effort values $v$ has a height equal to the measure of the applicants who are tied at that value, and so $\gamma$ for those applicants can be filled in using the CDF of $\Gamma(\omega)$ restricted to the applicants $\omega$ who are tied. Thus, the range of $\gamma$ is uniform on $[0,1]$.
%
{and}
%	Then, given
	a ranking probability function $\lambda$, an \textit{equilibrium} is a set of effort levels and post-effort ranking rewards for each applicant, $\{e(\theta_\pre),  \lambda(\theta_\post(\theta_\pre))\}$ such that, for all
	$\omega$, %$\theta_\pre$,
	\begin{align*}
		e(\theta_\pre(\omega)) &\in \argmax_{e} \iwelf\left(e, \lambda\left(\gamma\left(\omega, v\left(e,\theta_\pre(\omega)\right)\right)\right)\right) %\triangleq \argmax_{e} \left[ \lambda(\gamma(v(e,\theta_\pre)) - \cost(e) \right].
		\\
		\theta_\post(\theta_\pre(\omega)) &= \gamma(\omega, v(e(\theta_\pre(\omega)),\theta_\pre(\omega))),
	\end{align*}
	{where $\gamma(\omega,v)$ is the ranking induced by the CDF of the scores resulting from effort levels $\{e(\theta_\pre)\}$ and the tie-breaking function $\Gamma$.\footnote{See Remark~\ref{rem:gamma-effort} on $\gamma$ and the set of efforts.%Formally, $\gamma$ depends on the set of efforts. Given a fixed set and $\gamma$, for each applicant $\omega$ the first condition considers the counter-factual ranking of $\omega$ with different post-effort values but using the same ranking function. As defined, $\gamma$ yields a uniform distribution of ranks with such measure $0$ changes. In Appendix \Cref{lem:deviationsmeasure0}, we further prove that two effort sets equal up to sets of measure $0$ induce the same ranking function $\gamma$, and so the condition is consistent.%, and so we set it at the $\gamma$ induced by the effort levels $\{e(\theta_\pre)\}$.%, and so it does not vary.
	}
	}
	%In particular, we denote $\theta_\post(\theta_\pre) := \gamma(v(e(\theta_\pre),\theta_\pre))$.
	%\lledit{and for any $\theta_\pre, \theta'_\pre$, we have $\lambda(\theta_\post(\theta_\pre)) > \lambda(\theta_\post(\theta'_\pre)) \implies v(e(\theta_\pre),\theta_\pre) > v(e(\theta_\pre'),\theta_\pre')$.}
\end{definition}

[Next verbatim source excerpt]

\begin{definition}[Two-level policy]\label{assump:2-level-fixed-cap} In our baseline two-level function parameterized by cut-off $c \in (0,1-\rho]$, each applicant with post-effort rank  $\theta_\post \geq c$ is admitted with probability $\ell_1 =\frac{\rho}{1-c} \in (0,1]$. Others are rejected, $\ell_0 = 0$.
\end{definition}

Note that standard non-randomized admissions policies are equivalent to case where $c = 1 - \rho$: the highest scoring applicants up to the capacity constraint are accepted with probability 1 and all others are rejected. We call this case \textbf{non-randomized admissions}. The other extreme is the one-level policy, \textbf{pure randomization}, where each applicant is admitted with probability ${\rho}$, i.e., $\ell_0 =\rho$. Decreasing the cut-off $c$ can be viewed as increasing the level of randomization in the admissions policy.


{We now reason about how various welfare and fairness metrics vary with $\lambda$ (in the two-level policy class, just $c$). To simplify the presentation, we assume that $f, g$ and $\cost$ are differentiable and that baseline effort $e_0$ is~$0$.}

[Next verbatim source excerpt]

\parbold{Aggregate welfare and utility} We define three aggregate welfare functions of the equilibrium efforts and scores to capture the interests of different stakeholders.
The \textbf{applicant welfare} is defined as the population average of the individual applicant welfare at equilibrium:
%\begin{definition}[applicant Welfare]
%	The population applicant welfare is defined as
\begin{align*}
	\swelf :=& \E[\iwelf(e,\lambda(\theta_\post))]	= \rho - \E[\cost(e)]. %& \text{Reward function $\lambda$ constraint $\rho$}
\end{align*}
On the other hand, society derives value from the scores of applicants post-effort, leading to the following \textbf{societal utility}:$\quad\quad  	\Util\soc:= \E[ v].$

In other words, society prefers the entire applicant population to achieve higher scores, not only those who are admitted to the school, since higher test scores are correlated with labor productivity and economic growth \citep{hanushek2010high}.

Finally, the design $\lambda$ is controlled by a school, who may only draw value from those who enroll. The school's \textbf{private utility} is the expected score of admitted applicants,
which in our continuum formulation is the expectation of $v$ weighted by $\lambda(v)$: %\footnote{See Remark~\ref{rem:prob} on interpreting $\Util\pri$ as a conditional expectation.}
%\begin{equation*}
\(	\Util\pri:= \E[v\cdot \lambda(\theta_\post)]\).
%\end{equation*}
\end{ReviewVerbatim}

% report-context-presentation-sha256: 90f9116e292dd71fddfd6b445ca0b060a481d026855bdb9725f2833c2cf62a31
\paragraph{Settled source reading}
\begin{sloppypar}
Equilibrium means that almost every applicant best responds against every feasible deviation;\allowbreak{} uniqueness is up to a null set.\allowbreak{}
\end{sloppypar}
\paragraph{Settled source reading}
\begin{sloppypar}
Applicants know the fixed priority used at score-\allowbreak{}boundary ties.\allowbreak{} This clarifies the stated model, not a corrected admission policy or a new economic assumption.\allowbreak{}
\end{sloppypar}

\paragraph{Expanded PaperInterface specification}
\begin{ReviewVerbatim}
∀ (cost production skill tie : ℝ → ℝ) (effort score : ℝ → ℝ → ℝ) (rho : ℝ),
  LBG22StrategicRanking.RankingPrimitives cost production skill 0 →
    rho ∈ Set.Ioo 0 1 →
      DifferentiableOn ℝ cost (Set.Ioi 0) →
        DifferentiableOn ℝ production (Set.Ioi 0) →
          DifferentiableOn ℝ skill (Set.Ioo 0 1) →
            Measurable tie →
              Set.InjOn tie (Set.Ioc 0 1) →
                (∀ c ∈ Set.Ioc 0 (1 - rho),
                    LBG22StrategicRanking.RankingEquilibrium cost production skill tie 1
                      (LBG22StrategicRanking.sourceTwoLevelCutoff c) (LBG22StrategicRanking.sourceTwoLevelReward rho c)
                      (effort c) (score c)) →
                  MonotoneOn
                      (fun (c : ℝ) => LBG22StrategicRanking.sourceActualTwoLevelSchoolUtility (score c) tie rho c)
                      (Set.Ioc 0 (1 - rho)) ∧
                    ∀ c ∈ Set.Ioc 0 (1 - rho),
                      LBG22StrategicRanking.sourceActualTwoLevelSchoolUtility (score c) tie rho c ≤
                        LBG22StrategicRanking.sourceActualTwoLevelSchoolUtility
                          (score (LBG22StrategicRanking.sourceDeterministicCutoff rho)) tie rho
                          (LBG22StrategicRanking.sourceDeterministicCutoff rho)
\end{ReviewVerbatim}

\paragraph{Retained governing declarations}
\begin{itemize}\raggedright
\item \hyperlink{paper-prerequisite-2a1b692f4e3136c6}{LBG22\allowbreak{}Strategic\allowbreak{}Ranking.\allowbreak{}Ranking\allowbreak{}Equilibrium}
\item \hyperlink{paper-prerequisite-9bbcc12f4a1736bf}{LBG22\allowbreak{}Strategic\allowbreak{}Ranking.\allowbreak{}Ranking\allowbreak{}Primitives}
\item \hyperlink{paper-prerequisite-9d0a441d993fe3f2}{LBG22\allowbreak{}Strategic\allowbreak{}Ranking.\allowbreak{}source\allowbreak{}Actual\allowbreak{}Two\allowbreak{}Level\allowbreak{}School\allowbreak{}Utility}
\item \hyperlink{paper-prerequisite-e76c8574ddaca0dc}{LBG22\allowbreak{}Strategic\allowbreak{}Ranking.\allowbreak{}source\allowbreak{}Deterministic\allowbreak{}Cutoff}
\item \hyperlink{paper-prerequisite-f88aac58507e44ca}{LBG22\allowbreak{}Strategic\allowbreak{}Ranking.\allowbreak{}source\allowbreak{}Two\allowbreak{}Level\allowbreak{}Cutoff}
\item \hyperlink{paper-prerequisite-30d97cd43eb79017}{LBG22\allowbreak{}Strategic\allowbreak{}Ranking.\allowbreak{}source\allowbreak{}Two\allowbreak{}Level\allowbreak{}Reward}
\end{itemize}
\paragraph{Lean-elaborated claim roles}\begin{ReviewVerbatim}
1. parameter: ℝ → ℝ
2. parameter: ℝ → ℝ
3. parameter: ℝ → ℝ
4. parameter: ℝ → ℝ
5. parameter: ℝ → ℝ → ℝ
6. parameter: ℝ → ℝ → ℝ
7. parameter: ℝ
8. assumption: LBG22StrategicRanking.RankingPrimitives cost production skill 0
9. assumption: rho ∈ Set.Ioo 0 1
10. assumption: DifferentiableOn ℝ cost (Set.Ioi 0)
11. assumption: DifferentiableOn ℝ production (Set.Ioi 0)
12. assumption: DifferentiableOn ℝ skill (Set.Ioo 0 1)
13. assumption: Measurable tie
14. assumption: Set.InjOn tie (Set.Ioc 0 1)
15. assumption: ∀ c ∈ Set.Ioc 0 (1 - rho),
  LBG22StrategicRanking.RankingEquilibrium cost production skill tie 1 (LBG22StrategicRanking.sourceTwoLevelCutoff c)
    (LBG22StrategicRanking.sourceTwoLevelReward rho c) (effort c) (score c)
16. conclusion: MonotoneOn (fun (c : ℝ) => LBG22StrategicRanking.sourceActualTwoLevelSchoolUtility (score c) tie rho c)
    (Set.Ioc 0 (1 - rho)) ∧
  ∀ c ∈ Set.Ioc 0 (1 - rho),
    LBG22StrategicRanking.sourceActualTwoLevelSchoolUtility (score c) tie rho c ≤
      LBG22StrategicRanking.sourceActualTwoLevelSchoolUtility
        (score (LBG22StrategicRanking.sourceDeterministicCutoff rho)) tie rho
        (LBG22StrategicRanking.sourceDeterministicCutoff rho)
\end{ReviewVerbatim}

\paragraph{Lean proof endpoint (built separately)}\begin{ReviewVerbatim}
LBG22StrategicRanking.ProofInterface.privateUtility
\end{ReviewVerbatim}

\paragraph{Recorded source-to-Spec screening}
\textbf{Verdict:} \texttt{matches}
\\
{\raggedright
\textbf{Reason:} For every equilibrium family over admissible two-\allowbreak{}level cutoffs, the target makes the displayed expected score-\allowbreak{}weighted admission utility monotone in c and bounded above by its value at c=1-\allowbreak{}rho.\allowbreak{} This matches Proposition 3.\allowbreak{}2, including the source primitives, differentiability conventions, capacity range, and approved fixed-\allowbreak{}priority equilibrium representation.\allowbreak{}
\par}
\reviewmatch{Matches source input}
\noindent\textbf{Reviewer annotation}\par
\reviewerbox
\clearpage
\hypertarget{source-claim-5cd4797b5f4721de}{}
\section*{Review row 5}
\textbf{Source locator:} {\footnotesize\raggedright cited publication, lines 161-\allowbreak{}163\par}
\paragraph{Verbatim source input}
\begin{ReviewVerbatim}
\begin{restatable}[]{proposition}{counterexPrivateUtil}\label{prop-counterex-priv-util}
	 A $3$-level policy may achieve strictly higher $\Util\pri$ than non-randomized admissions.
\end{restatable}

[Next verbatim source excerpt]

\parbold{Applicants}
There is a unit mass of \textit{applicants}, indexed by an observed index $\omega \in [0,1]$ distributed uniformly.\footnote{The index $\omega$ should be interpreted as each applicant's ``name,'' uncorrelated with skill, used solely for tie-breaking.}  Each applicant has a latent (unobserved) skill level represented by some measurable function of $\omega$.
We assume that {the distribution of the skills} has no atoms and that the CDF of this distribution is strictly increasing.
Using the CDF to map the skill of an applicant to a rank in~$[0,1]$, each applicant gets an (unobserved) rank $\theta_\pre =\theta_\pre(\omega)$ which by our assumption on the CDF is again uniformly distributed in $[0, 1]$ (the higher the rank the better).  With this setup, the skill of an applicant with rank $\theta_\pre$ can be written as $f(\theta_\pre)$ where $f$ is a strictly increasing, continuous function.  {It will be notationally convenient to label applicants by their rank $\theta_\pre$, though the reader should note that in contrast to $\omega$, $\theta_\pre(\omega)$ is assumed to be unobservable.}

%3. Each applicant exerts some effort e, at cost \cost(e).


Each applicant chooses an \textit{effort} level $e \ge 0$, the result of which is an observed, post-effort \textit{score}, {$v=v(e,\theta_\pre) = {g(e)}\cdot{f(\theta_\pre)}$}. \lledit{In other words, we assume the post-effort score is the product of two components, associated with the effort and the pre-effort skill respectively.}
{We assume} that the effort transfer function $g: [0, \infty) \mapsto [0, \infty)$ is a continuous, concave, strictly increasing function, representing that marginal effort improves one's score but has diminishing returns.
{%Since applicants can be labelled by their skill level $f(\theta_\pre)$, or equivalently, by their rank $\theta_{\pre}$,
{T}he strategies of the applicants can {then} be described by a function $\theta_\pre \mapsto e(\theta_\pre)$.
%\cb{$\omega\mapsto e(\omega)$.}%, leading to a distribution over post-effort scores $v=v(\theta_\pre))$ via the uniform distribution of $\theta_\pre$ in $[0,1]$.
}
  Each {applicant} is then ranked according to their score $v$%(where we break ties uniformly at random, if they appear)
, resulting in a \textit{post-effort rank} $\theta_\post$.
Note that the ranking $\theta_\post$ is slightly less trivial than a ranking of the skills, since the scores might have ties, which have to be resolved.

{\textbf{Tie Breaking.}
Given a choice of strategies $\theta_\pre \mapsto e(\theta_\pre)$, let $F$ be the CDF of the scores $v(e(\theta_\pre),\theta_\pre)$.
Since atoms for the distribution of $v$ would lead to ties for ranks defined as $F(v)$, we will use the labels of the applicants to break ties with the help of a (publicly announced) tie-breaking function $\Gamma(\omega)$, defined, e.g., via a collision free hash of the applicant names. Here  we require  that $\Gamma$  is a measurable function from $[0,1]$ to $[0,1]$ that maps different applicant labels to different values.
We then use $\Gamma$ to resolve the atoms of $F$, leading to a ranking function $v\mapsto \gamma(\omega, v)$ which is equal to $F(v)$ except when $v$ is an atom of the score distribution, {in which case it takes values in the ``gap interval'' $[F_-(v),  F(v)]$, where $F_-(v)$ is the left limit of $F$ at $v$.}  We  construct $\gamma$ in such a way that it
gives the uniform distribution for $\theta_\post(\omega)$ if we set  $\theta_\post(\omega)=\gamma(\omega,v(e(\theta_\pre(\omega)),\theta_\pre(\omega)))$.\footnote{See Remark~\ref{rem:gamma} for details on this construction.}

[Next verbatim source excerpt]

\parbold{Ranking designer (school)} A single \textit{school} is admitting applicants, based on their ranking. In particular, the school can choose a ranking reward function $\lambda : [0, 1] \mapsto [0, 1]$, such that an applicant with post-effort rank $\theta_\post$ is admitted with probability $\lambda(\theta_\post)$. We assume that $\lambda$ is non-decreasing and that the school has a constraint on the overall probability, such that in expectation it admits a number of applicants equal to {a} capacity constraint  $\rho \in (0,1)$, i.e., \lledit{$\E_{\theta_\post}[\lambda(\theta_\post)] = \rho $}.\footnote{\lledit{We use $\E_{\theta_\post}[\cdot]$ to denote an integral over $\theta_\post$ (and $\E[\cdot]$ to denote an integral over $\omega$)} with respect to the Lebesgue measure; informally, this can be thought of as averaging over the applicant population.} We may also refer to $\lambda$, informally, as the admission policy.



For simplicity, we further assume that $\lambda$ is a step-function with $K$ distinct levels $\ell_0 < \dots < \ell_{K-1}$, and $K-1$ cut-points parameterized by $c_1< \dots < c_{K-1}$ (with $c_0 = 0$, $c_K = 1$). %$\ell_0, \dots, \ell_{K-1}$.
In other words, we have $\lambda(\theta) = \ell_k$, for all $\theta \in \psi_k \triangleq [c_{k}, c_{k+1})$. Thus, applicants in the same post-rank interval $\theta \in \psi_k$ receive the same reward.

[Next verbatim source excerpt]

\parbold{Individual applicant welfare and equilibrium} Given the designer's function $\lambda$ and the effort levels of other applicants, each applicant chooses effort $e$ to maximize their individual welfare, % $W$,
\begin{equation*}
	\iwelf(e, \lambda(\theta_\post)) = \lambda(\theta_\post) - \cost(e),
\end{equation*}
where the effort cost function $\cost$ is non-negative, continuous, and strictly convex {on $[0,\infty)$, with $e_0:=\argmin_e \cost(e)$ and $p(e_0) = 0$.}
%\lnote{also want to assume that $\cost$ attains its minimum at $0$?}
%Let $e_0:=\argmin_e \cost(e)$ and $p(e_0) = 0$.

Applicants are assumed not to personally benefit from increasing their score $v$, except through the corresponding increase in their rank and reward. While the definition of $\cost$ does not preclude the applicant receiving any intrinsic benefit from exerting effort, we assume that the \emph{net} benefit from effort is non-positive.

[Next verbatim source excerpt]

\begin{definition} [Equilibrium]\label{def:equi}
{Given a}
%Let the
tie-breaking function $\Gamma(\omega)$
%be a pre-announced ordering over index $\omega$, to break ties in post-effort value. Then, let $\gamma(\omega, v)$ be the 1-1 ranking function defined as the CDF of post-effort values $v$ except where there are ties in $v$, in which case the ordering $\Gamma$ is used to produce unique ranks such that the range of $\gamma$ is uniform on $[0, 1]$.\footnote{More formally, note that each discontinuity in the CDF of post-effort values $v$ has a height equal to the measure of the applicants who are tied at that value, and so $\gamma$ for those applicants can be filled in using the CDF of $\Gamma(\omega)$ restricted to the applicants $\omega$ who are tied. Thus, the range of $\gamma$ is uniform on $[0,1]$.
%
{and}
%	Then, given
	a ranking probability function $\lambda$, an \textit{equilibrium} is a set of effort levels and post-effort ranking rewards for each applicant, $\{e(\theta_\pre),  \lambda(\theta_\post(\theta_\pre))\}$ such that, for all
	$\omega$, %$\theta_\pre$,
	\begin{align*}
		e(\theta_\pre(\omega)) &\in \argmax_{e} \iwelf\left(e, \lambda\left(\gamma\left(\omega, v\left(e,\theta_\pre(\omega)\right)\right)\right)\right) %\triangleq \argmax_{e} \left[ \lambda(\gamma(v(e,\theta_\pre)) - \cost(e) \right].
		\\
		\theta_\post(\theta_\pre(\omega)) &= \gamma(\omega, v(e(\theta_\pre(\omega)),\theta_\pre(\omega))),
	\end{align*}
	{where $\gamma(\omega,v)$ is the ranking induced by the CDF of the scores resulting from effort levels $\{e(\theta_\pre)\}$ and the tie-breaking function $\Gamma$.\footnote{See Remark~\ref{rem:gamma-effort} on $\gamma$ and the set of efforts.%Formally, $\gamma$ depends on the set of efforts. Given a fixed set and $\gamma$, for each applicant $\omega$ the first condition considers the counter-factual ranking of $\omega$ with different post-effort values but using the same ranking function. As defined, $\gamma$ yields a uniform distribution of ranks with such measure $0$ changes. In Appendix \Cref{lem:deviationsmeasure0}, we further prove that two effort sets equal up to sets of measure $0$ induce the same ranking function $\gamma$, and so the condition is consistent.%, and so we set it at the $\gamma$ induced by the effort levels $\{e(\theta_\pre)\}$.%, and so it does not vary.
	}
	}
	%In particular, we denote $\theta_\post(\theta_\pre) := \gamma(v(e(\theta_\pre),\theta_\pre))$.
	%\lledit{and for any $\theta_\pre, \theta'_\pre$, we have $\lambda(\theta_\post(\theta_\pre)) > \lambda(\theta_\post(\theta'_\pre)) \implies v(e(\theta_\pre),\theta_\pre) > v(e(\theta_\pre'),\theta_\pre')$.}
\end{definition}

[Next verbatim source excerpt]

\begin{definition}[Two-level policy]\label{assump:2-level-fixed-cap} In our baseline two-level function parameterized by cut-off $c \in (0,1-\rho]$, each applicant with post-effort rank  $\theta_\post \geq c$ is admitted with probability $\ell_1 =\frac{\rho}{1-c} \in (0,1]$. Others are rejected, $\ell_0 = 0$.
\end{definition}

Note that standard non-randomized admissions policies are equivalent to case where $c = 1 - \rho$: the highest scoring applicants up to the capacity constraint are accepted with probability 1 and all others are rejected. We call this case \textbf{non-randomized admissions}. The other extreme is the one-level policy, \textbf{pure randomization}, where each applicant is admitted with probability ${\rho}$, i.e., $\ell_0 =\rho$. Decreasing the cut-off $c$ can be viewed as increasing the level of randomization in the admissions policy.


{We now reason about how various welfare and fairness metrics vary with $\lambda$ (in the two-level policy class, just $c$). To simplify the presentation, we assume that $f, g$ and $\cost$ are differentiable and that baseline effort $e_0$ is~$0$.}

[Next verbatim source excerpt]

\parbold{Aggregate welfare and utility} We define three aggregate welfare functions of the equilibrium efforts and scores to capture the interests of different stakeholders.
The \textbf{applicant welfare} is defined as the population average of the individual applicant welfare at equilibrium:
%\begin{definition}[applicant Welfare]
%	The population applicant welfare is defined as
\begin{align*}
	\swelf :=& \E[\iwelf(e,\lambda(\theta_\post))]	= \rho - \E[\cost(e)]. %& \text{Reward function $\lambda$ constraint $\rho$}
\end{align*}
On the other hand, society derives value from the scores of applicants post-effort, leading to the following \textbf{societal utility}:$\quad\quad  	\Util\soc:= \E[ v].$

In other words, society prefers the entire applicant population to achieve higher scores, not only those who are admitted to the school, since higher test scores are correlated with labor productivity and economic growth \citep{hanushek2010high}.

Finally, the design $\lambda$ is controlled by a school, who may only draw value from those who enroll. The school's \textbf{private utility} is the expected score of admitted applicants,
which in our continuum formulation is the expectation of $v$ weighted by $\lambda(v)$: %\footnote{See Remark~\ref{rem:prob} on interpreting $\Util\pri$ as a conditional expectation.}
%\begin{equation*}
\(	\Util\pri:= \E[v\cdot \lambda(\theta_\post)]\).
%\end{equation*}
\end{ReviewVerbatim}

% report-context-presentation-sha256: 90f9116e292dd71fddfd6b445ca0b060a481d026855bdb9725f2833c2cf62a31
\paragraph{Settled source reading}
\begin{sloppypar}
Equilibrium means that almost every applicant best responds against every feasible deviation;\allowbreak{} uniqueness is up to a null set.\allowbreak{}
\end{sloppypar}
\paragraph{Settled source reading}
\begin{sloppypar}
Applicants know the fixed priority used at score-\allowbreak{}boundary ties.\allowbreak{} This clarifies the stated model, not a corrected admission policy or a new economic assumption.\allowbreak{}
\end{sloppypar}

\paragraph{Expanded PaperInterface specification}
\begin{ReviewVerbatim}
∀ (tie : ℝ → ℝ),
  Measurable tie →
    Set.InjOn tie (Set.Ioc 0 1) →
      ∃ (cost : ℝ → ℝ) (production : ℝ → ℝ) (skill : ℝ → ℝ) (rho : ℝ) (cutoff : ℕ → ℝ) (reward : ℕ → ℝ) (effort : ℝ → ℝ)
        (score : ℝ → ℝ) (deterministicEffort : ℝ → ℝ) (deterministicScore : ℝ → ℝ),
        LBG22StrategicRanking.RankingPrimitives cost production skill 0 ∧
          DifferentiableOn ℝ cost (Set.Ioi 0) ∧
            DifferentiableOn ℝ production (Set.Ioi 0) ∧
              DifferentiableOn ℝ skill (Set.Ioo 0 1) ∧
                LBG22StrategicRanking.FiniteAdmissionPolicy 2 cutoff reward rho ∧
                  LBG22StrategicRanking.RankingEquilibrium cost production skill tie 2 cutoff reward effort score ∧
                    LBG22StrategicRanking.RankingEquilibrium cost production skill tie 1
                        (LBG22StrategicRanking.sourceTwoLevelCutoff
                          (LBG22StrategicRanking.sourceDeterministicCutoff rho))
                        (LBG22StrategicRanking.sourceTwoLevelReward rho
                          (LBG22StrategicRanking.sourceDeterministicCutoff rho))
                        deterministicEffort deterministicScore ∧
                      LBG22StrategicRanking.sourceFiniteSchoolUtility deterministicScore tie
                          (fun (i : Fin 2) =>
                            LBG22StrategicRanking.sourceTwoLevelCutoff
                              (LBG22StrategicRanking.sourceDeterministicCutoff rho) ↑i)
                          (LBG22StrategicRanking.sourceTwoLevelReward rho
                            (LBG22StrategicRanking.sourceDeterministicCutoff rho)) <
                        LBG22StrategicRanking.sourceFiniteSchoolUtility score tie (fun (i : Fin 3) => cutoff ↑i) reward
\end{ReviewVerbatim}

\paragraph{Retained governing declarations}
\begin{itemize}\raggedright
\item \hyperlink{paper-prerequisite-d056dd7c8060653c}{LBG22\allowbreak{}Strategic\allowbreak{}Ranking.\allowbreak{}Finite\allowbreak{}Admission\allowbreak{}Policy}
\item \hyperlink{paper-prerequisite-2a1b692f4e3136c6}{LBG22\allowbreak{}Strategic\allowbreak{}Ranking.\allowbreak{}Ranking\allowbreak{}Equilibrium}
\item \hyperlink{paper-prerequisite-9bbcc12f4a1736bf}{LBG22\allowbreak{}Strategic\allowbreak{}Ranking.\allowbreak{}Ranking\allowbreak{}Primitives}
\item \hyperlink{paper-prerequisite-e76c8574ddaca0dc}{LBG22\allowbreak{}Strategic\allowbreak{}Ranking.\allowbreak{}source\allowbreak{}Deterministic\allowbreak{}Cutoff}
\item \hyperlink{paper-prerequisite-25a6dc17f260546b}{LBG22\allowbreak{}Strategic\allowbreak{}Ranking.\allowbreak{}source\allowbreak{}Finite\allowbreak{}School\allowbreak{}Utility}
\item \hyperlink{paper-prerequisite-f88aac58507e44ca}{LBG22\allowbreak{}Strategic\allowbreak{}Ranking.\allowbreak{}source\allowbreak{}Two\allowbreak{}Level\allowbreak{}Cutoff}
\item \hyperlink{paper-prerequisite-30d97cd43eb79017}{LBG22\allowbreak{}Strategic\allowbreak{}Ranking.\allowbreak{}source\allowbreak{}Two\allowbreak{}Level\allowbreak{}Reward}
\end{itemize}
\paragraph{Lean-elaborated claim roles}\begin{ReviewVerbatim}
1. parameter: ℝ → ℝ
2. assumption: Measurable tie
3. assumption: Set.InjOn tie (Set.Ioc 0 1)
4. conclusion: ∃ (cost : ℝ → ℝ) (production : ℝ → ℝ) (skill : ℝ → ℝ) (rho : ℝ) (cutoff : ℕ → ℝ) (reward : ℕ → ℝ) (effort : ℝ → ℝ)
  (score : ℝ → ℝ) (deterministicEffort : ℝ → ℝ) (deterministicScore : ℝ → ℝ),
  LBG22StrategicRanking.RankingPrimitives cost production skill 0 ∧
    DifferentiableOn ℝ cost (Set.Ioi 0) ∧
      DifferentiableOn ℝ production (Set.Ioi 0) ∧
        DifferentiableOn ℝ skill (Set.Ioo 0 1) ∧
          LBG22StrategicRanking.FiniteAdmissionPolicy 2 cutoff reward rho ∧
            LBG22StrategicRanking.RankingEquilibrium cost production skill tie 2 cutoff reward effort score ∧
              LBG22StrategicRanking.RankingEquilibrium cost production skill tie 1
                  (LBG22StrategicRanking.sourceTwoLevelCutoff (LBG22StrategicRanking.sourceDeterministicCutoff rho))
                  (LBG22StrategicRanking.sourceTwoLevelReward rho (LBG22StrategicRanking.sourceDeterministicCutoff rho))
                  deterministicEffort deterministicScore ∧
                LBG22StrategicRanking.sourceFiniteSchoolUtility deterministicScore tie
                    (fun (i : Fin 2) =>
                      LBG22StrategicRanking.sourceTwoLevelCutoff (LBG22StrategicRanking.sourceDeterministicCutoff rho)
                        ↑i)
                    (LBG22StrategicRanking.sourceTwoLevelReward rho
                      (LBG22StrategicRanking.sourceDeterministicCutoff rho)) <
                  LBG22StrategicRanking.sourceFiniteSchoolUtility score tie (fun (i : Fin 3) => cutoff ↑i) reward
\end{ReviewVerbatim}

\paragraph{Lean proof endpoint (built separately)}\begin{ReviewVerbatim}
LBG22StrategicRanking.ProofInterface.threeLevelImprovement
\end{ReviewVerbatim}

\paragraph{Recorded source-to-Spec screening}
\textbf{Verdict:} \texttt{matches}
\\
{\raggedright
\textbf{Reason:} The existential target provides a valid three-\allowbreak{}level finite policy and equilibrium together with the deterministic two-\allowbreak{}level equilibrium, and concludes a strict increase in the displayed expected score-\allowbreak{}weighted admission utility.\allowbreak{} That is exactly the may-\allowbreak{}achieve counterexample asserted by Proposition 3.\allowbreak{}3.\allowbreak{}
\par}
\reviewmatch{Matches source input}
\noindent\textbf{Reviewer annotation}\par
\reviewerbox
\clearpage
\hypertarget{source-claim-f1c09a02beccebed}{}
\section*{Review row 6}
\textbf{Source locator:} {\footnotesize\raggedright cited publication, lines 173-\allowbreak{}177\par}
\paragraph{Verbatim source input}
\begin{ReviewVerbatim}
\begin{restatable}[Societal utility maximization for two-level policies]{proposition}{propSocUtil}
	Among two-level functions, there exists a setting in which societal utility $\Util\soc$ is maximized by choosing $c \in(0, 1-\rho)$, $\ell_1 = \frac{\rho}{1-c} \in (\rho,1)$.
%There exists a 2-level model with fixed capacity (Assumption~\ref{assump:2-level-fixed-cap}) where the societal utility on the measurable skill only, $\Util\soc$, is maximized by choosing $c < 1-\rho$, $\ell_1 = \frac{\rho}{1-c} < 1$.
\label{prop:SocUtil}
\end{restatable}

[Next verbatim source excerpt]

\parbold{Applicants}
There is a unit mass of \textit{applicants}, indexed by an observed index $\omega \in [0,1]$ distributed uniformly.\footnote{The index $\omega$ should be interpreted as each applicant's ``name,'' uncorrelated with skill, used solely for tie-breaking.}  Each applicant has a latent (unobserved) skill level represented by some measurable function of $\omega$.
We assume that {the distribution of the skills} has no atoms and that the CDF of this distribution is strictly increasing.
Using the CDF to map the skill of an applicant to a rank in~$[0,1]$, each applicant gets an (unobserved) rank $\theta_\pre =\theta_\pre(\omega)$ which by our assumption on the CDF is again uniformly distributed in $[0, 1]$ (the higher the rank the better).  With this setup, the skill of an applicant with rank $\theta_\pre$ can be written as $f(\theta_\pre)$ where $f$ is a strictly increasing, continuous function.  {It will be notationally convenient to label applicants by their rank $\theta_\pre$, though the reader should note that in contrast to $\omega$, $\theta_\pre(\omega)$ is assumed to be unobservable.}

%3. Each applicant exerts some effort e, at cost \cost(e).


Each applicant chooses an \textit{effort} level $e \ge 0$, the result of which is an observed, post-effort \textit{score}, {$v=v(e,\theta_\pre) = {g(e)}\cdot{f(\theta_\pre)}$}. \lledit{In other words, we assume the post-effort score is the product of two components, associated with the effort and the pre-effort skill respectively.}
{We assume} that the effort transfer function $g: [0, \infty) \mapsto [0, \infty)$ is a continuous, concave, strictly increasing function, representing that marginal effort improves one's score but has diminishing returns.
{%Since applicants can be labelled by their skill level $f(\theta_\pre)$, or equivalently, by their rank $\theta_{\pre}$,
{T}he strategies of the applicants can {then} be described by a function $\theta_\pre \mapsto e(\theta_\pre)$.
%\cb{$\omega\mapsto e(\omega)$.}%, leading to a distribution over post-effort scores $v=v(\theta_\pre))$ via the uniform distribution of $\theta_\pre$ in $[0,1]$.
}
  Each {applicant} is then ranked according to their score $v$%(where we break ties uniformly at random, if they appear)
, resulting in a \textit{post-effort rank} $\theta_\post$.
Note that the ranking $\theta_\post$ is slightly less trivial than a ranking of the skills, since the scores might have ties, which have to be resolved.

{\textbf{Tie Breaking.}
Given a choice of strategies $\theta_\pre \mapsto e(\theta_\pre)$, let $F$ be the CDF of the scores $v(e(\theta_\pre),\theta_\pre)$.
Since atoms for the distribution of $v$ would lead to ties for ranks defined as $F(v)$, we will use the labels of the applicants to break ties with the help of a (publicly announced) tie-breaking function $\Gamma(\omega)$, defined, e.g., via a collision free hash of the applicant names. Here  we require  that $\Gamma$  is a measurable function from $[0,1]$ to $[0,1]$ that maps different applicant labels to different values.
We then use $\Gamma$ to resolve the atoms of $F$, leading to a ranking function $v\mapsto \gamma(\omega, v)$ which is equal to $F(v)$ except when $v$ is an atom of the score distribution, {in which case it takes values in the ``gap interval'' $[F_-(v),  F(v)]$, where $F_-(v)$ is the left limit of $F$ at $v$.}  We  construct $\gamma$ in such a way that it
gives the uniform distribution for $\theta_\post(\omega)$ if we set  $\theta_\post(\omega)=\gamma(\omega,v(e(\theta_\pre(\omega)),\theta_\pre(\omega)))$.\footnote{See Remark~\ref{rem:gamma} for details on this construction.}

[Next verbatim source excerpt]

\parbold{Ranking designer (school)} A single \textit{school} is admitting applicants, based on their ranking. In particular, the school can choose a ranking reward function $\lambda : [0, 1] \mapsto [0, 1]$, such that an applicant with post-effort rank $\theta_\post$ is admitted with probability $\lambda(\theta_\post)$. We assume that $\lambda$ is non-decreasing and that the school has a constraint on the overall probability, such that in expectation it admits a number of applicants equal to {a} capacity constraint  $\rho \in (0,1)$, i.e., \lledit{$\E_{\theta_\post}[\lambda(\theta_\post)] = \rho $}.\footnote{\lledit{We use $\E_{\theta_\post}[\cdot]$ to denote an integral over $\theta_\post$ (and $\E[\cdot]$ to denote an integral over $\omega$)} with respect to the Lebesgue measure; informally, this can be thought of as averaging over the applicant population.} We may also refer to $\lambda$, informally, as the admission policy.



For simplicity, we further assume that $\lambda$ is a step-function with $K$ distinct levels $\ell_0 < \dots < \ell_{K-1}$, and $K-1$ cut-points parameterized by $c_1< \dots < c_{K-1}$ (with $c_0 = 0$, $c_K = 1$). %$\ell_0, \dots, \ell_{K-1}$.
In other words, we have $\lambda(\theta) = \ell_k$, for all $\theta \in \psi_k \triangleq [c_{k}, c_{k+1})$. Thus, applicants in the same post-rank interval $\theta \in \psi_k$ receive the same reward.

[Next verbatim source excerpt]

\parbold{Individual applicant welfare and equilibrium} Given the designer's function $\lambda$ and the effort levels of other applicants, each applicant chooses effort $e$ to maximize their individual welfare, % $W$,
\begin{equation*}
	\iwelf(e, \lambda(\theta_\post)) = \lambda(\theta_\post) - \cost(e),
\end{equation*}
where the effort cost function $\cost$ is non-negative, continuous, and strictly convex {on $[0,\infty)$, with $e_0:=\argmin_e \cost(e)$ and $p(e_0) = 0$.}
%\lnote{also want to assume that $\cost$ attains its minimum at $0$?}
%Let $e_0:=\argmin_e \cost(e)$ and $p(e_0) = 0$.

Applicants are assumed not to personally benefit from increasing their score $v$, except through the corresponding increase in their rank and reward. While the definition of $\cost$ does not preclude the applicant receiving any intrinsic benefit from exerting effort, we assume that the \emph{net} benefit from effort is non-positive.

[Next verbatim source excerpt]

\begin{definition} [Equilibrium]\label{def:equi}
{Given a}
%Let the
tie-breaking function $\Gamma(\omega)$
%be a pre-announced ordering over index $\omega$, to break ties in post-effort value. Then, let $\gamma(\omega, v)$ be the 1-1 ranking function defined as the CDF of post-effort values $v$ except where there are ties in $v$, in which case the ordering $\Gamma$ is used to produce unique ranks such that the range of $\gamma$ is uniform on $[0, 1]$.\footnote{More formally, note that each discontinuity in the CDF of post-effort values $v$ has a height equal to the measure of the applicants who are tied at that value, and so $\gamma$ for those applicants can be filled in using the CDF of $\Gamma(\omega)$ restricted to the applicants $\omega$ who are tied. Thus, the range of $\gamma$ is uniform on $[0,1]$.
%
{and}
%	Then, given
	a ranking probability function $\lambda$, an \textit{equilibrium} is a set of effort levels and post-effort ranking rewards for each applicant, $\{e(\theta_\pre),  \lambda(\theta_\post(\theta_\pre))\}$ such that, for all
	$\omega$, %$\theta_\pre$,
	\begin{align*}
		e(\theta_\pre(\omega)) &\in \argmax_{e} \iwelf\left(e, \lambda\left(\gamma\left(\omega, v\left(e,\theta_\pre(\omega)\right)\right)\right)\right) %\triangleq \argmax_{e} \left[ \lambda(\gamma(v(e,\theta_\pre)) - \cost(e) \right].
		\\
		\theta_\post(\theta_\pre(\omega)) &= \gamma(\omega, v(e(\theta_\pre(\omega)),\theta_\pre(\omega))),
	\end{align*}
	{where $\gamma(\omega,v)$ is the ranking induced by the CDF of the scores resulting from effort levels $\{e(\theta_\pre)\}$ and the tie-breaking function $\Gamma$.\footnote{See Remark~\ref{rem:gamma-effort} on $\gamma$ and the set of efforts.%Formally, $\gamma$ depends on the set of efforts. Given a fixed set and $\gamma$, for each applicant $\omega$ the first condition considers the counter-factual ranking of $\omega$ with different post-effort values but using the same ranking function. As defined, $\gamma$ yields a uniform distribution of ranks with such measure $0$ changes. In Appendix \Cref{lem:deviationsmeasure0}, we further prove that two effort sets equal up to sets of measure $0$ induce the same ranking function $\gamma$, and so the condition is consistent.%, and so we set it at the $\gamma$ induced by the effort levels $\{e(\theta_\pre)\}$.%, and so it does not vary.
	}
	}
	%In particular, we denote $\theta_\post(\theta_\pre) := \gamma(v(e(\theta_\pre),\theta_\pre))$.
	%\lledit{and for any $\theta_\pre, \theta'_\pre$, we have $\lambda(\theta_\post(\theta_\pre)) > \lambda(\theta_\post(\theta'_\pre)) \implies v(e(\theta_\pre),\theta_\pre) > v(e(\theta_\pre'),\theta_\pre')$.}
\end{definition}

[Next verbatim source excerpt]

\begin{definition}[Two-level policy]\label{assump:2-level-fixed-cap} In our baseline two-level function parameterized by cut-off $c \in (0,1-\rho]$, each applicant with post-effort rank  $\theta_\post \geq c$ is admitted with probability $\ell_1 =\frac{\rho}{1-c} \in (0,1]$. Others are rejected, $\ell_0 = 0$.
\end{definition}

Note that standard non-randomized admissions policies are equivalent to case where $c = 1 - \rho$: the highest scoring applicants up to the capacity constraint are accepted with probability 1 and all others are rejected. We call this case \textbf{non-randomized admissions}. The other extreme is the one-level policy, \textbf{pure randomization}, where each applicant is admitted with probability ${\rho}$, i.e., $\ell_0 =\rho$. Decreasing the cut-off $c$ can be viewed as increasing the level of randomization in the admissions policy.


{We now reason about how various welfare and fairness metrics vary with $\lambda$ (in the two-level policy class, just $c$). To simplify the presentation, we assume that $f, g$ and $\cost$ are differentiable and that baseline effort $e_0$ is~$0$.}

[Next verbatim source excerpt]

\parbold{Aggregate welfare and utility} We define three aggregate welfare functions of the equilibrium efforts and scores to capture the interests of different stakeholders.
The \textbf{applicant welfare} is defined as the population average of the individual applicant welfare at equilibrium:
%\begin{definition}[applicant Welfare]
%	The population applicant welfare is defined as
\begin{align*}
	\swelf :=& \E[\iwelf(e,\lambda(\theta_\post))]	= \rho - \E[\cost(e)]. %& \text{Reward function $\lambda$ constraint $\rho$}
\end{align*}
On the other hand, society derives value from the scores of applicants post-effort, leading to the following \textbf{societal utility}:$\quad\quad  	\Util\soc:= \E[ v].$

In other words, society prefers the entire applicant population to achieve higher scores, not only those who are admitted to the school, since higher test scores are correlated with labor productivity and economic growth \citep{hanushek2010high}.

Finally, the design $\lambda$ is controlled by a school, who may only draw value from those who enroll. The school's \textbf{private utility} is the expected score of admitted applicants,
which in our continuum formulation is the expectation of $v$ weighted by $\lambda(v)$: %\footnote{See Remark~\ref{rem:prob} on interpreting $\Util\pri$ as a conditional expectation.}
%\begin{equation*}
\(	\Util\pri:= \E[v\cdot \lambda(\theta_\post)]\).
%\end{equation*}
\end{ReviewVerbatim}

% report-context-presentation-sha256: 90f9116e292dd71fddfd6b445ca0b060a481d026855bdb9725f2833c2cf62a31
\paragraph{Settled source reading}
\begin{sloppypar}
Equilibrium means that almost every applicant best responds against every feasible deviation;\allowbreak{} uniqueness is up to a null set.\allowbreak{}
\end{sloppypar}
\paragraph{Settled source reading}
\begin{sloppypar}
Applicants know the fixed priority used at score-\allowbreak{}boundary ties.\allowbreak{} This clarifies the stated model, not a corrected admission policy or a new economic assumption.\allowbreak{}
\end{sloppypar}

\paragraph{Expanded PaperInterface specification}
\begin{ReviewVerbatim}
∀ (tie : ℝ → ℝ),
  Measurable tie →
    Set.InjOn tie (Set.Ioc 0 1) →
      ∃ (cost : ℝ → ℝ) (production : ℝ → ℝ) (skill : ℝ → ℝ) (rho : ℝ) (cStar : ℝ) (effort : ℝ → ℝ → ℝ) (score :
        ℝ → ℝ → ℝ),
        LBG22StrategicRanking.RankingPrimitives cost production skill 0 ∧
          DifferentiableOn ℝ cost (Set.Ioi 0) ∧
            DifferentiableOn ℝ production (Set.Ioi 0) ∧
              DifferentiableOn ℝ skill (Set.Ioo 0 1) ∧
                rho ∈ Set.Ioo 0 1 ∧
                  cStar ∈ Set.Ioo 0 (1 - rho) ∧
                    rho / (1 - cStar) ∈ Set.Ioo rho 1 ∧
                      (∀ c ∈ Set.Ioc 0 (1 - rho),
                          LBG22StrategicRanking.RankingEquilibrium cost production skill tie 1
                              (LBG22StrategicRanking.sourceTwoLevelCutoff c)
                              (LBG22StrategicRanking.sourceTwoLevelReward rho c) (effort c) (score c) ∧
                            ∫ (t : ℝ),
                                LBG22StrategicRanking.sourceTwoLevelReward rho c
                                  ↑(LBG22StrategicRanking.finiteLowerRankBand
                                      (fun (i : Fin 2) => LBG22StrategicRanking.sourceTwoLevelCutoff c ↑i)
                                      (LBG22StrategicRanking.tieBrokenRank LBG22StrategicRanking.unitRankMeasure
                                        (score c) tie t)) ∂LBG22StrategicRanking.unitRankMeasure =
                              rho) ∧
                        ∀ c ∈ Set.Ioc 0 (1 - rho),
                          MeasureTheory.Integrable (score c) LBG22StrategicRanking.unitRankMeasure ∧
                            LBG22StrategicRanking.sourceSocietalUtility (score c) ≤
                              LBG22StrategicRanking.sourceSocietalUtility (score cStar)
\end{ReviewVerbatim}

\paragraph{Retained governing declarations}
\begin{itemize}\raggedright
\item \hyperlink{paper-prerequisite-2a1b692f4e3136c6}{LBG22\allowbreak{}Strategic\allowbreak{}Ranking.\allowbreak{}Ranking\allowbreak{}Equilibrium}
\item \hyperlink{paper-prerequisite-9bbcc12f4a1736bf}{LBG22\allowbreak{}Strategic\allowbreak{}Ranking.\allowbreak{}Ranking\allowbreak{}Primitives}
\item \hyperlink{paper-prerequisite-a42f440ddc254fa7}{LBG22\allowbreak{}Strategic\allowbreak{}Ranking.\allowbreak{}finite\allowbreak{}Lower\allowbreak{}Rank\allowbreak{}Band}
\item \hyperlink{paper-prerequisite-7224ad4c7b0b8551}{LBG22\allowbreak{}Strategic\allowbreak{}Ranking.\allowbreak{}source\allowbreak{}Societal\allowbreak{}Utility}
\item \hyperlink{paper-prerequisite-f88aac58507e44ca}{LBG22\allowbreak{}Strategic\allowbreak{}Ranking.\allowbreak{}source\allowbreak{}Two\allowbreak{}Level\allowbreak{}Cutoff}
\item \hyperlink{paper-prerequisite-30d97cd43eb79017}{LBG22\allowbreak{}Strategic\allowbreak{}Ranking.\allowbreak{}source\allowbreak{}Two\allowbreak{}Level\allowbreak{}Reward}
\item \hyperlink{paper-prerequisite-b7c71e53177b76c6}{LBG22\allowbreak{}Strategic\allowbreak{}Ranking.\allowbreak{}tie\allowbreak{}Broken\allowbreak{}Rank}
\item \hyperlink{governing-declaration-678cf48d549cb620}{LBG22\allowbreak{}Strategic\allowbreak{}Ranking.\allowbreak{}unit\allowbreak{}Rank\allowbreak{}Measure}
\end{itemize}
\paragraph{Lean-elaborated claim roles}\begin{ReviewVerbatim}
1. parameter: ℝ → ℝ
2. assumption: Measurable tie
3. assumption: Set.InjOn tie (Set.Ioc 0 1)
4. conclusion: ∃ (cost : ℝ → ℝ) (production : ℝ → ℝ) (skill : ℝ → ℝ) (rho : ℝ) (cStar : ℝ) (effort : ℝ → ℝ → ℝ) (score : ℝ → ℝ → ℝ),
  LBG22StrategicRanking.RankingPrimitives cost production skill 0 ∧
    DifferentiableOn ℝ cost (Set.Ioi 0) ∧
      DifferentiableOn ℝ production (Set.Ioi 0) ∧
        DifferentiableOn ℝ skill (Set.Ioo 0 1) ∧
          rho ∈ Set.Ioo 0 1 ∧
            cStar ∈ Set.Ioo 0 (1 - rho) ∧
              rho / (1 - cStar) ∈ Set.Ioo rho 1 ∧
                (∀ c ∈ Set.Ioc 0 (1 - rho),
                    LBG22StrategicRanking.RankingEquilibrium cost production skill tie 1
                        (LBG22StrategicRanking.sourceTwoLevelCutoff c)
                        (LBG22StrategicRanking.sourceTwoLevelReward rho c) (effort c) (score c) ∧
                      ∫ (t : ℝ),
                          LBG22StrategicRanking.sourceTwoLevelReward rho c
                            ↑(LBG22StrategicRanking.finiteLowerRankBand
                                (fun (i : Fin 2) => LBG22StrategicRanking.sourceTwoLevelCutoff c ↑i)
                                (LBG22StrategicRanking.tieBrokenRank LBG22StrategicRanking.unitRankMeasure (score c) tie
                                  t)) ∂LBG22StrategicRanking.unitRankMeasure =
                        rho) ∧
                  ∀ c ∈ Set.Ioc 0 (1 - rho),
                    MeasureTheory.Integrable (score c) LBG22StrategicRanking.unitRankMeasure ∧
                      LBG22StrategicRanking.sourceSocietalUtility (score c) ≤
                        LBG22StrategicRanking.sourceSocietalUtility (score cStar)
\end{ReviewVerbatim}

\paragraph{Lean proof endpoint (built separately)}\begin{ReviewVerbatim}
LBG22StrategicRanking.ProofInterface.societalUtility
\end{ReviewVerbatim}

\paragraph{Recorded source-to-Spec screening}
\textbf{Verdict:} \texttt{matches}
\\
{\raggedright
\textbf{Reason:} The existential target supplies source-\allowbreak{}admissible primitives, capacity, an interior cutoff cStar with rho/\allowbreak{}(1-\allowbreak{}cStar) strictly between rho and one, equilibria at every feasible cutoff, capacity accounting, and global maximality of expected post-\allowbreak{}effort score at cStar.\allowbreak{} This is precisely the existence claim of Proposition 3.\allowbreak{}4.\allowbreak{}
\par}
\reviewmatch{Matches source input}
\noindent\textbf{Reviewer annotation}\par
\reviewerbox
\clearpage
\hypertarget{source-claim-d14371f504c1d7ee}{}
\section*{Review row 7}
\textbf{Source locator:} {\footnotesize\raggedright cited publication, lines 224-\allowbreak{}235\par}
\paragraph{Verbatim source input}
\begin{ReviewVerbatim}
\begin{restatable}[Equilibrium under group differences]{proposition}{propEquiEnv}\label{prop:equi-env}
Define $\theta_\pre$ as:
%	\begin{equation*}
	\(\theta_\pre := f^{-1}_\mix(f(\theta_\true)\cdot \psi)\),
%	\end{equation*}
		where $f^{-1}_\mix$	is the CDF for the environment-scaled skill, $f(\theta_\true)\cdot \psi$:
	\begin{equation*}
	f^{-1}_\mix(x) := \frac{1}{2}f^{-1}(x/\psi_\A)+ \frac{1}{2}f^{-1}(x/\psi_\B).
	\end{equation*}

	Then, in every equilibrium, $\lambda(\theta_{\post}(\theta_\true, \psi))=\lambda(\theta_\pre)$.%, up to the appropriate tie-breaking.
\end{restatable}

[Next verbatim source excerpt]

\parbold{Applicants}
There is a unit mass of \textit{applicants}, indexed by an observed index $\omega \in [0,1]$ distributed uniformly.\footnote{The index $\omega$ should be interpreted as each applicant's ``name,'' uncorrelated with skill, used solely for tie-breaking.}  Each applicant has a latent (unobserved) skill level represented by some measurable function of $\omega$.
We assume that {the distribution of the skills} has no atoms and that the CDF of this distribution is strictly increasing.
Using the CDF to map the skill of an applicant to a rank in~$[0,1]$, each applicant gets an (unobserved) rank $\theta_\pre =\theta_\pre(\omega)$ which by our assumption on the CDF is again uniformly distributed in $[0, 1]$ (the higher the rank the better).  With this setup, the skill of an applicant with rank $\theta_\pre$ can be written as $f(\theta_\pre)$ where $f$ is a strictly increasing, continuous function.  {It will be notationally convenient to label applicants by their rank $\theta_\pre$, though the reader should note that in contrast to $\omega$, $\theta_\pre(\omega)$ is assumed to be unobservable.}

%3. Each applicant exerts some effort e, at cost \cost(e).


Each applicant chooses an \textit{effort} level $e \ge 0$, the result of which is an observed, post-effort \textit{score}, {$v=v(e,\theta_\pre) = {g(e)}\cdot{f(\theta_\pre)}$}. \lledit{In other words, we assume the post-effort score is the product of two components, associated with the effort and the pre-effort skill respectively.}
{We assume} that the effort transfer function $g: [0, \infty) \mapsto [0, \infty)$ is a continuous, concave, strictly increasing function, representing that marginal effort improves one's score but has diminishing returns.
{%Since applicants can be labelled by their skill level $f(\theta_\pre)$, or equivalently, by their rank $\theta_{\pre}$,
{T}he strategies of the applicants can {then} be described by a function $\theta_\pre \mapsto e(\theta_\pre)$.
%\cb{$\omega\mapsto e(\omega)$.}%, leading to a distribution over post-effort scores $v=v(\theta_\pre))$ via the uniform distribution of $\theta_\pre$ in $[0,1]$.
}
  Each {applicant} is then ranked according to their score $v$%(where we break ties uniformly at random, if they appear)
, resulting in a \textit{post-effort rank} $\theta_\post$.
Note that the ranking $\theta_\post$ is slightly less trivial than a ranking of the skills, since the scores might have ties, which have to be resolved.

{\textbf{Tie Breaking.}
Given a choice of strategies $\theta_\pre \mapsto e(\theta_\pre)$, let $F$ be the CDF of the scores $v(e(\theta_\pre),\theta_\pre)$.
Since atoms for the distribution of $v$ would lead to ties for ranks defined as $F(v)$, we will use the labels of the applicants to break ties with the help of a (publicly announced) tie-breaking function $\Gamma(\omega)$, defined, e.g., via a collision free hash of the applicant names. Here  we require  that $\Gamma$  is a measurable function from $[0,1]$ to $[0,1]$ that maps different applicant labels to different values.
We then use $\Gamma$ to resolve the atoms of $F$, leading to a ranking function $v\mapsto \gamma(\omega, v)$ which is equal to $F(v)$ except when $v$ is an atom of the score distribution, {in which case it takes values in the ``gap interval'' $[F_-(v),  F(v)]$, where $F_-(v)$ is the left limit of $F$ at $v$.}  We  construct $\gamma$ in such a way that it
gives the uniform distribution for $\theta_\post(\omega)$ if we set  $\theta_\post(\omega)=\gamma(\omega,v(e(\theta_\pre(\omega)),\theta_\pre(\omega)))$.\footnote{See Remark~\ref{rem:gamma} for details on this construction.}

[Next verbatim source excerpt]

\parbold{Ranking designer (school)} A single \textit{school} is admitting applicants, based on their ranking. In particular, the school can choose a ranking reward function $\lambda : [0, 1] \mapsto [0, 1]$, such that an applicant with post-effort rank $\theta_\post$ is admitted with probability $\lambda(\theta_\post)$. We assume that $\lambda$ is non-decreasing and that the school has a constraint on the overall probability, such that in expectation it admits a number of applicants equal to {a} capacity constraint  $\rho \in (0,1)$, i.e., \lledit{$\E_{\theta_\post}[\lambda(\theta_\post)] = \rho $}.\footnote{\lledit{We use $\E_{\theta_\post}[\cdot]$ to denote an integral over $\theta_\post$ (and $\E[\cdot]$ to denote an integral over $\omega$)} with respect to the Lebesgue measure; informally, this can be thought of as averaging over the applicant population.} We may also refer to $\lambda$, informally, as the admission policy.



For simplicity, we further assume that $\lambda$ is a step-function with $K$ distinct levels $\ell_0 < \dots < \ell_{K-1}$, and $K-1$ cut-points parameterized by $c_1< \dots < c_{K-1}$ (with $c_0 = 0$, $c_K = 1$). %$\ell_0, \dots, \ell_{K-1}$.
In other words, we have $\lambda(\theta) = \ell_k$, for all $\theta \in \psi_k \triangleq [c_{k}, c_{k+1})$. Thus, applicants in the same post-rank interval $\theta \in \psi_k$ receive the same reward.

[Next verbatim source excerpt]

\parbold{Individual applicant welfare and equilibrium} Given the designer's function $\lambda$ and the effort levels of other applicants, each applicant chooses effort $e$ to maximize their individual welfare, % $W$,
\begin{equation*}
	\iwelf(e, \lambda(\theta_\post)) = \lambda(\theta_\post) - \cost(e),
\end{equation*}
where the effort cost function $\cost$ is non-negative, continuous, and strictly convex {on $[0,\infty)$, with $e_0:=\argmin_e \cost(e)$ and $p(e_0) = 0$.}
%\lnote{also want to assume that $\cost$ attains its minimum at $0$?}
%Let $e_0:=\argmin_e \cost(e)$ and $p(e_0) = 0$.

Applicants are assumed not to personally benefit from increasing their score $v$, except through the corresponding increase in their rank and reward. While the definition of $\cost$ does not preclude the applicant receiving any intrinsic benefit from exerting effort, we assume that the \emph{net} benefit from effort is non-positive.

[Next verbatim source excerpt]

\begin{definition} [Equilibrium]\label{def:equi}
{Given a}
%Let the
tie-breaking function $\Gamma(\omega)$
%be a pre-announced ordering over index $\omega$, to break ties in post-effort value. Then, let $\gamma(\omega, v)$ be the 1-1 ranking function defined as the CDF of post-effort values $v$ except where there are ties in $v$, in which case the ordering $\Gamma$ is used to produce unique ranks such that the range of $\gamma$ is uniform on $[0, 1]$.\footnote{More formally, note that each discontinuity in the CDF of post-effort values $v$ has a height equal to the measure of the applicants who are tied at that value, and so $\gamma$ for those applicants can be filled in using the CDF of $\Gamma(\omega)$ restricted to the applicants $\omega$ who are tied. Thus, the range of $\gamma$ is uniform on $[0,1]$.
%
{and}
%	Then, given
	a ranking probability function $\lambda$, an \textit{equilibrium} is a set of effort levels and post-effort ranking rewards for each applicant, $\{e(\theta_\pre),  \lambda(\theta_\post(\theta_\pre))\}$ such that, for all
	$\omega$, %$\theta_\pre$,
	\begin{align*}
		e(\theta_\pre(\omega)) &\in \argmax_{e} \iwelf\left(e, \lambda\left(\gamma\left(\omega, v\left(e,\theta_\pre(\omega)\right)\right)\right)\right) %\triangleq \argmax_{e} \left[ \lambda(\gamma(v(e,\theta_\pre)) - \cost(e) \right].
		\\
		\theta_\post(\theta_\pre(\omega)) &= \gamma(\omega, v(e(\theta_\pre(\omega)),\theta_\pre(\omega))),
	\end{align*}
	{where $\gamma(\omega,v)$ is the ranking induced by the CDF of the scores resulting from effort levels $\{e(\theta_\pre)\}$ and the tie-breaking function $\Gamma$.\footnote{See Remark~\ref{rem:gamma-effort} on $\gamma$ and the set of efforts.%Formally, $\gamma$ depends on the set of efforts. Given a fixed set and $\gamma$, for each applicant $\omega$ the first condition considers the counter-factual ranking of $\omega$ with different post-effort values but using the same ranking function. As defined, $\gamma$ yields a uniform distribution of ranks with such measure $0$ changes. In Appendix \Cref{lem:deviationsmeasure0}, we further prove that two effort sets equal up to sets of measure $0$ induce the same ranking function $\gamma$, and so the condition is consistent.%, and so we set it at the $\gamma$ induced by the effort levels $\{e(\theta_\pre)\}$.%, and so it does not vary.
	}
	}
	%In particular, we denote $\theta_\post(\theta_\pre) := \gamma(v(e(\theta_\pre),\theta_\pre))$.
	%\lledit{and for any $\theta_\pre, \theta'_\pre$, we have $\lambda(\theta_\post(\theta_\pre)) > \lambda(\theta_\post(\theta'_\pre)) \implies v(e(\theta_\pre),\theta_\pre) > v(e(\theta_\pre'),\theta_\pre')$.}
\end{definition}

[Next verbatim source excerpt]

\paragraph{Model and equilibria characterization} We now denote each applicant's latent skill rank as $\theta_\true \in [0,1]$. In addition to the latent skill, each applicant has an (unobserved) environmental factor $\psi \in \Psi$ that represents how favorable their environment is for attaining a higher score. Because a favorable environment results in a higher rate of return for effort, we model the environment as a multiplicative factor in the score (see~e.g., \citet{calsamiglia09decent}). Formally, the post-effort score is a function of the latent skill, the environmental factor, and the effort level:
\begin{equation*}\label{eq:env-score}
v = \psi \cdot  g(e)\cdot f(\theta_\true),
\end{equation*}
where $g, f$ are as defined in Section~\ref{sec:model}. %We also require $f$ to be strictly monotone.

%We model structural inequalities as heterogeneous returns to effort. In contrast to results in Section~\ref{sec:equillemmas}, the ranks of applicants are no longer preserved when there are heterogeneous returns to effort.

 We assume there are two groups of applicants, $\A$ and $\B$, and the distribution of skill is the same in both groups. Group $\A$ has a more favorable environment factor, that is,  $\Psi = \{\psi_\A, \psi_\B\}$ and $\psi_\A > \psi_\B$. Thus we will also refer to $\B$ as the ``disadvantaged group". To simplify our presentation, we assume each group is half of the total applicant population, though the results in this section generalize.% to the setting where groups are not equal-sized.
%
We defer all proofs in this section to Appendix~\ref{app:environment}.

We begin by characterizing the equilibrium ranking under the designer's policy $\lambda$. The equilibrium effort levels and post-effort ranks are as defined in Definition~\ref{def:equi}, except they are now group-dependent, that is, we have $e(\theta_\true, \psi)$ and $\theta_\post(\theta_\true, \psi)$. Because of the differences in $\psi$, the post-effort ranking $\theta_\post$ is now group-dependent, and in general is not equal to $\theta_\true$. To apply Proposition~\ref{lem:rankpreservedindex} as before, we construct an ``environment-scaled pre-effort rank'' (denoted $\theta_\pre$).

[Next verbatim source excerpt]

\begin{restatable}[Equilibrium under group differences]{proposition}{propEquiEnv}\label{prop:equi-env}
Define $\theta_\pre$ as:
%	\begin{equation*}
	\(\theta_\pre := f^{-1}_\mix(f(\theta_\true)\cdot \psi)\),
%	\end{equation*}
		where $f^{-1}_\mix$	is the CDF for the environment-scaled skill, $f(\theta_\true)\cdot \psi$:
	\begin{equation*}
	f^{-1}_\mix(x) := \frac{1}{2}f^{-1}(x/\psi_\A)+ \frac{1}{2}f^{-1}(x/\psi_\B).
	\end{equation*}


[Next verbatim source excerpt]

\begin{definition}[Welfare gap]
	Let $\swelf^\G(\theta_\true)$ denote post-effort welfare of an applicant with latent skill ranking $\theta_\true$ from group $\G \in \{\A, \B\}$, i.e.,
	\begin{equation}
	\swelf^\G(\theta_\true):= \lambda(\theta_\post(\theta_\true, \psi_\G)) - \cost(e(\theta_\true, \psi_\G)).
	\end{equation}
	%\lnote{Using function notation for $\theta_\post$ and $e$? or some other way to express the welfare as a function of $\theta_\true$.}
	We define the \emph{%pointwise
		welfare gap} as %a function of rank $\theta_\true$ is
	\begin{equation*}
		\welfgap(\theta_\true):= \swelf^\A(\theta_\true) - \swelf^\B(\theta_\true).
	\end{equation*}
%	We define the \emph{subpopulation welfare gap} as %a function of baseline rank $\overline{\theta}_\true$ is
%	\begin{equation*}
%		\popwelfgap(\overline{\theta}_\true):=\E[\welfgap(\theta_\true) \mid \theta_\true > \overline{\theta}_\true].
%	\end{equation*}
\end{definition}

The welfare gap captures differences in admission probabilities \textit{and} in the effort needed to achieve such probabilities. Our next notion, \textit{access}, captures whether a decision policy \emph{includes} the disadvantaged group in the admitted class, regardless of effort. %This is captured by the notion of \emph{access}.

\begin{definition}[Access]
 \emph{Access} is the overall probability of admission of the disadvantaged group.
	\begin{equation*}
		\access :=  \E_{\theta_\true}[\lambda(\theta_\post(\theta_\true, \psi_\B))].
	\end{equation*}
% \nikhil{Defn is a bit confusing. What is the expectation over, and what does that conditioning mean?}
\end{definition}

[Next verbatim source excerpt]

	Denote the group-specific rank threshold for group $\G$ as
	\begin{equation*}
	\thresG{c} := f^{-1}\left(\frac{f_\mix(c)}{\psi_\G}\right).
	\end{equation*}
\end{ReviewVerbatim}

% report-context-presentation-sha256: 90f9116e292dd71fddfd6b445ca0b060a481d026855bdb9725f2833c2cf62a31
\paragraph{Settled source reading}
\begin{sloppypar}
Equilibrium means that almost every applicant best responds against every feasible deviation;\allowbreak{} uniqueness is up to a null set.\allowbreak{}
\end{sloppypar}
\paragraph{Settled source reading}
\begin{sloppypar}
Applicants know the fixed priority used at score-\allowbreak{}boundary ties.\allowbreak{} This clarifies the stated model, not a corrected admission policy or a new economic assumption.\allowbreak{}
\end{sloppypar}

\paragraph{Expanded PaperInterface specification}
\begin{ReviewVerbatim}
∀ (cost production skill : ℝ → ℝ) (effort score tie : Bool × ℝ → ℝ) (baseline rho psiA psiB : ℝ) (n : ℕ)
  (cutoff reward : ℕ → ℝ),
  LBG22StrategicRanking.RankingPrimitives cost production skill baseline →
    LBG22StrategicRanking.FiniteAdmissionPolicy n cutoff reward rho →
      0 < psiB →
        psiB < psiA →
          Measurable tie →
            Set.InjOn tie {x : Bool × ℝ | x.2 ∈ Set.Ioc 0 1} →
              LBG22StrategicRanking.EnvironmentEquilibrium cost production skill psiA psiB tie n cutoff reward effort
                  score →
                (∀ (z : ℝ),
                    LBG22StrategicRanking.sourceEnvironmentCDF skill psiA psiB z =
                      1 / 2 * LBG22StrategicRanking.sourceSkillCDF skill (z / psiA) +
                        1 / 2 * LBG22StrategicRanking.sourceSkillCDF skill (z / psiB)) ∧
                  ∀ᵐ (x : Bool × ℝ) ∂LBG22StrategicRanking.sourceTwoGroupMeasure,
                    reward
                        ↑(LBG22StrategicRanking.finiteLowerRankBand (fun (i : Fin (n + 1)) => cutoff ↑i)
                            (LBG22StrategicRanking.tieBrokenRank LBG22StrategicRanking.sourceTwoGroupMeasure score tie
                              x)) =
                      reward
                        ↑(LBG22StrategicRanking.finiteLowerRankBand (fun (i : Fin (n + 1)) => cutoff ↑i)
                            (LBG22StrategicRanking.sourceEnvironmentRank skill psiA psiB x))
\end{ReviewVerbatim}

\paragraph{Retained governing declarations}
\begin{itemize}\raggedright
\item \hyperlink{paper-prerequisite-b8e06f39e94d0049}{LBG22\allowbreak{}Strategic\allowbreak{}Ranking.\allowbreak{}Environment\allowbreak{}Equilibrium}
\item \hyperlink{paper-prerequisite-d056dd7c8060653c}{LBG22\allowbreak{}Strategic\allowbreak{}Ranking.\allowbreak{}Finite\allowbreak{}Admission\allowbreak{}Policy}
\item \hyperlink{paper-prerequisite-9bbcc12f4a1736bf}{LBG22\allowbreak{}Strategic\allowbreak{}Ranking.\allowbreak{}Ranking\allowbreak{}Primitives}
\item \hyperlink{paper-prerequisite-a42f440ddc254fa7}{LBG22\allowbreak{}Strategic\allowbreak{}Ranking.\allowbreak{}finite\allowbreak{}Lower\allowbreak{}Rank\allowbreak{}Band}
\item \hyperlink{paper-prerequisite-1aae1cdeecf21463}{LBG22\allowbreak{}Strategic\allowbreak{}Ranking.\allowbreak{}source\allowbreak{}Environment\allowbreak{}CDF}
\item \hyperlink{paper-prerequisite-b308d7da226a5fe1}{LBG22\allowbreak{}Strategic\allowbreak{}Ranking.\allowbreak{}source\allowbreak{}Environment\allowbreak{}Rank}
\item \hyperlink{governing-declaration-15b79809af308d40}{LBG22\allowbreak{}Strategic\allowbreak{}Ranking.\allowbreak{}source\allowbreak{}Skill\allowbreak{}CDF}
\item \hyperlink{paper-prerequisite-5b3be494cde30a37}{LBG22\allowbreak{}Strategic\allowbreak{}Ranking.\allowbreak{}source\allowbreak{}Two\allowbreak{}Group\allowbreak{}Measure}
\item \hyperlink{paper-prerequisite-b7c71e53177b76c6}{LBG22\allowbreak{}Strategic\allowbreak{}Ranking.\allowbreak{}tie\allowbreak{}Broken\allowbreak{}Rank}
\end{itemize}
\paragraph{Lean-elaborated claim roles}\begin{ReviewVerbatim}
1. parameter: ℝ → ℝ
2. parameter: ℝ → ℝ
3. parameter: ℝ → ℝ
4. parameter: Bool × ℝ → ℝ
5. parameter: Bool × ℝ → ℝ
6. parameter: Bool × ℝ → ℝ
7. parameter: ℝ
8. parameter: ℝ
9. parameter: ℝ
10. parameter: ℝ
11. parameter: ℕ
12. parameter: ℕ → ℝ
13. parameter: ℕ → ℝ
14. assumption: LBG22StrategicRanking.RankingPrimitives cost production skill baseline
15. assumption: LBG22StrategicRanking.FiniteAdmissionPolicy n cutoff reward rho
16. assumption: 0 < psiB
17. assumption: psiB < psiA
18. assumption: Measurable tie
19. assumption: Set.InjOn tie {x : Bool × ℝ | x.2 ∈ Set.Ioc 0 1}
20. assumption: LBG22StrategicRanking.EnvironmentEquilibrium cost production skill psiA psiB tie n cutoff reward effort score
21. conclusion: (∀ (z : ℝ),
    LBG22StrategicRanking.sourceEnvironmentCDF skill psiA psiB z =
      1 / 2 * LBG22StrategicRanking.sourceSkillCDF skill (z / psiA) +
        1 / 2 * LBG22StrategicRanking.sourceSkillCDF skill (z / psiB)) ∧
  ∀ᵐ (x : Bool × ℝ) ∂LBG22StrategicRanking.sourceTwoGroupMeasure,
    reward
        ↑(LBG22StrategicRanking.finiteLowerRankBand (fun (i : Fin (n + 1)) => cutoff ↑i)
            (LBG22StrategicRanking.tieBrokenRank LBG22StrategicRanking.sourceTwoGroupMeasure score tie x)) =
      reward
        ↑(LBG22StrategicRanking.finiteLowerRankBand (fun (i : Fin (n + 1)) => cutoff ↑i)
            (LBG22StrategicRanking.sourceEnvironmentRank skill psiA psiB x))
\end{ReviewVerbatim}

\paragraph{Lean proof endpoint (built separately)}\begin{ReviewVerbatim}
LBG22StrategicRanking.ProofInterface.environmentRankPreservation
\end{ReviewVerbatim}

\paragraph{Recorded source-to-Spec screening}
\textbf{Verdict:} \texttt{matches}
\\
{\raggedright
\textbf{Reason:} The target gives the equal-\allowbreak{}half mixture CDF formula for environment-\allowbreak{}scaled skill and, in every displayed group equilibrium, equality almost everywhere between the reward at actual post-\allowbreak{}effort rank and at the environment-\allowbreak{}scaled pre-\allowbreak{}rank.\allowbreak{} This is Proposition 4.\allowbreak{}1, with positive ordered environment multipliers, common skill distribution, finite policy, and fixed priority supplied by the displayed prerequisites.\allowbreak{}
\par}
\reviewmatch{Matches source input}
\noindent\textbf{Reviewer annotation}\par
\reviewerbox
\clearpage
\hypertarget{source-claim-77aa097e6083e865}{}
\section*{Review row 8}
\textbf{Source locator:} {\footnotesize\raggedright cited publication, lines 281-\allowbreak{}310\par}
\paragraph{Verbatim source input}
\begin{ReviewVerbatim}
\begin{restatable}[Admission and pointwise welfare gap for two-level policies]{proposition}{propEnvDiff}\label{prop:environment-diff}
	Denote the group-specific rank threshold for group $\G$ as
	\begin{equation*}
	\thresG{c} := f^{-1}\left(\frac{f_\mix(c)}{\psi_\G}\right).
	\end{equation*}
A two-level policy with $c \in (0,1-\rho]$ admits a group~$\A$ applicant with $\theta_\true \geq \thresA{c} $ with probability $\frac{\rho}{1-c}$, a group~$\B$ applicant with $\theta_\true \geq \thresB{c} $ with probability $\frac{\rho}{1-c}$, and all other applicants with probability~0. % Let $\effortA{\theta_\true}$ (resp. $\effortB{\theta_\true}$) denote effort exerted by a student with true ranking $\theta_\true$ from group $\A$ (resp. group $\B$).
The welfare gap $\welfgap(\theta_\true)$ is non-negative for every $\theta_\true$, and strictly positive for $\theta_\true \geq \thresB{c}$. In contrast, the one-level pure randomization policy has $\welfgap(\theta_\true) \equiv 0$.


%\begin{enumerate}
%    \item ``Low''. For $0 \le \theta_\true < \thresA{c}$,
%    \begin{enumerate}
%        \item $\swelf^\A(\theta_\true) = \swelf^\B(\theta_\true) = 0$ (both have zero welfare)
%        \item $\E[\outcome \mid \A, \theta_\true ] = \E[\outcome \mid \B, \theta_\true ] = 0$ (both have zero probability of admission)
%        \item $\effortA{\theta_\true} = \effortB{\theta_\true} =0$ (both exert zero effort)
%    \end{enumerate}
%    \item ``Middle''. For $\thresA{c} \le \theta_\true < \thresB{c}$
%    \begin{enumerate}
%        \item $\swelf^\A(\theta_\true) > \swelf^\B(\theta_\true) = 0$ (group $\A$ student has higher welfare)
%        \item $\E[\outcome \mid \A, \theta_\true ] > \E[\outcome \mid \B, \theta_\true ] = 0$ (group $\A$ student has higher probability of admission)
%        \item $\effortA{\theta_\true} > \effortB{\theta_\true} =0$ (group $\A$ student exerts more effort)
%    \end{enumerate}
%    \item ``High". For $\thresB{c} \le \theta_\true \le 1 $,
%        \begin{enumerate}
%        \item $\swelf^\A(\theta_\true) > \swelf^\B(\theta_\true)$ (group $\A$ student has higher welfare)
%        \item $\E[\outcome \mid \A, \theta_\true ] =\E[\outcome \mid \B, \theta_\true ] >0$ (the admissions probabilities are the same)
%        \item $\effortA{\theta_\true} < \effortB{\theta_\true} =0$ (group $\A$ student exerts less effort)
%    \end{enumerate}
%\end{enumerate}
\end{restatable}

[Next verbatim source excerpt]

\parbold{Applicants}
There is a unit mass of \textit{applicants}, indexed by an observed index $\omega \in [0,1]$ distributed uniformly.\footnote{The index $\omega$ should be interpreted as each applicant's ``name,'' uncorrelated with skill, used solely for tie-breaking.}  Each applicant has a latent (unobserved) skill level represented by some measurable function of $\omega$.
We assume that {the distribution of the skills} has no atoms and that the CDF of this distribution is strictly increasing.
Using the CDF to map the skill of an applicant to a rank in~$[0,1]$, each applicant gets an (unobserved) rank $\theta_\pre =\theta_\pre(\omega)$ which by our assumption on the CDF is again uniformly distributed in $[0, 1]$ (the higher the rank the better).  With this setup, the skill of an applicant with rank $\theta_\pre$ can be written as $f(\theta_\pre)$ where $f$ is a strictly increasing, continuous function.  {It will be notationally convenient to label applicants by their rank $\theta_\pre$, though the reader should note that in contrast to $\omega$, $\theta_\pre(\omega)$ is assumed to be unobservable.}

%3. Each applicant exerts some effort e, at cost \cost(e).


Each applicant chooses an \textit{effort} level $e \ge 0$, the result of which is an observed, post-effort \textit{score}, {$v=v(e,\theta_\pre) = {g(e)}\cdot{f(\theta_\pre)}$}. \lledit{In other words, we assume the post-effort score is the product of two components, associated with the effort and the pre-effort skill respectively.}
{We assume} that the effort transfer function $g: [0, \infty) \mapsto [0, \infty)$ is a continuous, concave, strictly increasing function, representing that marginal effort improves one's score but has diminishing returns.
{%Since applicants can be labelled by their skill level $f(\theta_\pre)$, or equivalently, by their rank $\theta_{\pre}$,
{T}he strategies of the applicants can {then} be described by a function $\theta_\pre \mapsto e(\theta_\pre)$.
%\cb{$\omega\mapsto e(\omega)$.}%, leading to a distribution over post-effort scores $v=v(\theta_\pre))$ via the uniform distribution of $\theta_\pre$ in $[0,1]$.
}
  Each {applicant} is then ranked according to their score $v$%(where we break ties uniformly at random, if they appear)
, resulting in a \textit{post-effort rank} $\theta_\post$.
Note that the ranking $\theta_\post$ is slightly less trivial than a ranking of the skills, since the scores might have ties, which have to be resolved.

{\textbf{Tie Breaking.}
Given a choice of strategies $\theta_\pre \mapsto e(\theta_\pre)$, let $F$ be the CDF of the scores $v(e(\theta_\pre),\theta_\pre)$.
Since atoms for the distribution of $v$ would lead to ties for ranks defined as $F(v)$, we will use the labels of the applicants to break ties with the help of a (publicly announced) tie-breaking function $\Gamma(\omega)$, defined, e.g., via a collision free hash of the applicant names. Here  we require  that $\Gamma$  is a measurable function from $[0,1]$ to $[0,1]$ that maps different applicant labels to different values.
We then use $\Gamma$ to resolve the atoms of $F$, leading to a ranking function $v\mapsto \gamma(\omega, v)$ which is equal to $F(v)$ except when $v$ is an atom of the score distribution, {in which case it takes values in the ``gap interval'' $[F_-(v),  F(v)]$, where $F_-(v)$ is the left limit of $F$ at $v$.}  We  construct $\gamma$ in such a way that it
gives the uniform distribution for $\theta_\post(\omega)$ if we set  $\theta_\post(\omega)=\gamma(\omega,v(e(\theta_\pre(\omega)),\theta_\pre(\omega)))$.\footnote{See Remark~\ref{rem:gamma} for details on this construction.}

[Next verbatim source excerpt]

\parbold{Ranking designer (school)} A single \textit{school} is admitting applicants, based on their ranking. In particular, the school can choose a ranking reward function $\lambda : [0, 1] \mapsto [0, 1]$, such that an applicant with post-effort rank $\theta_\post$ is admitted with probability $\lambda(\theta_\post)$. We assume that $\lambda$ is non-decreasing and that the school has a constraint on the overall probability, such that in expectation it admits a number of applicants equal to {a} capacity constraint  $\rho \in (0,1)$, i.e., \lledit{$\E_{\theta_\post}[\lambda(\theta_\post)] = \rho $}.\footnote{\lledit{We use $\E_{\theta_\post}[\cdot]$ to denote an integral over $\theta_\post$ (and $\E[\cdot]$ to denote an integral over $\omega$)} with respect to the Lebesgue measure; informally, this can be thought of as averaging over the applicant population.} We may also refer to $\lambda$, informally, as the admission policy.



For simplicity, we further assume that $\lambda$ is a step-function with $K$ distinct levels $\ell_0 < \dots < \ell_{K-1}$, and $K-1$ cut-points parameterized by $c_1< \dots < c_{K-1}$ (with $c_0 = 0$, $c_K = 1$). %$\ell_0, \dots, \ell_{K-1}$.
In other words, we have $\lambda(\theta) = \ell_k$, for all $\theta \in \psi_k \triangleq [c_{k}, c_{k+1})$. Thus, applicants in the same post-rank interval $\theta \in \psi_k$ receive the same reward.

[Next verbatim source excerpt]

\parbold{Individual applicant welfare and equilibrium} Given the designer's function $\lambda$ and the effort levels of other applicants, each applicant chooses effort $e$ to maximize their individual welfare, % $W$,
\begin{equation*}
	\iwelf(e, \lambda(\theta_\post)) = \lambda(\theta_\post) - \cost(e),
\end{equation*}
where the effort cost function $\cost$ is non-negative, continuous, and strictly convex {on $[0,\infty)$, with $e_0:=\argmin_e \cost(e)$ and $p(e_0) = 0$.}
%\lnote{also want to assume that $\cost$ attains its minimum at $0$?}
%Let $e_0:=\argmin_e \cost(e)$ and $p(e_0) = 0$.

Applicants are assumed not to personally benefit from increasing their score $v$, except through the corresponding increase in their rank and reward. While the definition of $\cost$ does not preclude the applicant receiving any intrinsic benefit from exerting effort, we assume that the \emph{net} benefit from effort is non-positive.

[Next verbatim source excerpt]

\begin{definition} [Equilibrium]\label{def:equi}
{Given a}
%Let the
tie-breaking function $\Gamma(\omega)$
%be a pre-announced ordering over index $\omega$, to break ties in post-effort value. Then, let $\gamma(\omega, v)$ be the 1-1 ranking function defined as the CDF of post-effort values $v$ except where there are ties in $v$, in which case the ordering $\Gamma$ is used to produce unique ranks such that the range of $\gamma$ is uniform on $[0, 1]$.\footnote{More formally, note that each discontinuity in the CDF of post-effort values $v$ has a height equal to the measure of the applicants who are tied at that value, and so $\gamma$ for those applicants can be filled in using the CDF of $\Gamma(\omega)$ restricted to the applicants $\omega$ who are tied. Thus, the range of $\gamma$ is uniform on $[0,1]$.
%
{and}
%	Then, given
	a ranking probability function $\lambda$, an \textit{equilibrium} is a set of effort levels and post-effort ranking rewards for each applicant, $\{e(\theta_\pre),  \lambda(\theta_\post(\theta_\pre))\}$ such that, for all
	$\omega$, %$\theta_\pre$,
	\begin{align*}
		e(\theta_\pre(\omega)) &\in \argmax_{e} \iwelf\left(e, \lambda\left(\gamma\left(\omega, v\left(e,\theta_\pre(\omega)\right)\right)\right)\right) %\triangleq \argmax_{e} \left[ \lambda(\gamma(v(e,\theta_\pre)) - \cost(e) \right].
		\\
		\theta_\post(\theta_\pre(\omega)) &= \gamma(\omega, v(e(\theta_\pre(\omega)),\theta_\pre(\omega))),
	\end{align*}
	{where $\gamma(\omega,v)$ is the ranking induced by the CDF of the scores resulting from effort levels $\{e(\theta_\pre)\}$ and the tie-breaking function $\Gamma$.\footnote{See Remark~\ref{rem:gamma-effort} on $\gamma$ and the set of efforts.%Formally, $\gamma$ depends on the set of efforts. Given a fixed set and $\gamma$, for each applicant $\omega$ the first condition considers the counter-factual ranking of $\omega$ with different post-effort values but using the same ranking function. As defined, $\gamma$ yields a uniform distribution of ranks with such measure $0$ changes. In Appendix \Cref{lem:deviationsmeasure0}, we further prove that two effort sets equal up to sets of measure $0$ induce the same ranking function $\gamma$, and so the condition is consistent.%, and so we set it at the $\gamma$ induced by the effort levels $\{e(\theta_\pre)\}$.%, and so it does not vary.
	}
	}
	%In particular, we denote $\theta_\post(\theta_\pre) := \gamma(v(e(\theta_\pre),\theta_\pre))$.
	%\lledit{and for any $\theta_\pre, \theta'_\pre$, we have $\lambda(\theta_\post(\theta_\pre)) > \lambda(\theta_\post(\theta'_\pre)) \implies v(e(\theta_\pre),\theta_\pre) > v(e(\theta_\pre'),\theta_\pre')$.}
\end{definition}

[Next verbatim source excerpt]

\begin{definition}[Two-level policy]\label{assump:2-level-fixed-cap} In our baseline two-level function parameterized by cut-off $c \in (0,1-\rho]$, each applicant with post-effort rank  $\theta_\post \geq c$ is admitted with probability $\ell_1 =\frac{\rho}{1-c} \in (0,1]$. Others are rejected, $\ell_0 = 0$.
\end{definition}

Note that standard non-randomized admissions policies are equivalent to case where $c = 1 - \rho$: the highest scoring applicants up to the capacity constraint are accepted with probability 1 and all others are rejected. We call this case \textbf{non-randomized admissions}. The other extreme is the one-level policy, \textbf{pure randomization}, where each applicant is admitted with probability ${\rho}$, i.e., $\ell_0 =\rho$. Decreasing the cut-off $c$ can be viewed as increasing the level of randomization in the admissions policy.


{We now reason about how various welfare and fairness metrics vary with $\lambda$ (in the two-level policy class, just $c$). To simplify the presentation, we assume that $f, g$ and $\cost$ are differentiable and that baseline effort $e_0$ is~$0$.}

[Next verbatim source excerpt]

\paragraph{Model and equilibria characterization} We now denote each applicant's latent skill rank as $\theta_\true \in [0,1]$. In addition to the latent skill, each applicant has an (unobserved) environmental factor $\psi \in \Psi$ that represents how favorable their environment is for attaining a higher score. Because a favorable environment results in a higher rate of return for effort, we model the environment as a multiplicative factor in the score (see~e.g., \citet{calsamiglia09decent}). Formally, the post-effort score is a function of the latent skill, the environmental factor, and the effort level:
\begin{equation*}\label{eq:env-score}
v = \psi \cdot  g(e)\cdot f(\theta_\true),
\end{equation*}
where $g, f$ are as defined in Section~\ref{sec:model}. %We also require $f$ to be strictly monotone.

%We model structural inequalities as heterogeneous returns to effort. In contrast to results in Section~\ref{sec:equillemmas}, the ranks of applicants are no longer preserved when there are heterogeneous returns to effort.

 We assume there are two groups of applicants, $\A$ and $\B$, and the distribution of skill is the same in both groups. Group $\A$ has a more favorable environment factor, that is,  $\Psi = \{\psi_\A, \psi_\B\}$ and $\psi_\A > \psi_\B$. Thus we will also refer to $\B$ as the ``disadvantaged group". To simplify our presentation, we assume each group is half of the total applicant population, though the results in this section generalize.% to the setting where groups are not equal-sized.
%
We defer all proofs in this section to Appendix~\ref{app:environment}.

We begin by characterizing the equilibrium ranking under the designer's policy $\lambda$. The equilibrium effort levels and post-effort ranks are as defined in Definition~\ref{def:equi}, except they are now group-dependent, that is, we have $e(\theta_\true, \psi)$ and $\theta_\post(\theta_\true, \psi)$. Because of the differences in $\psi$, the post-effort ranking $\theta_\post$ is now group-dependent, and in general is not equal to $\theta_\true$. To apply Proposition~\ref{lem:rankpreservedindex} as before, we construct an ``environment-scaled pre-effort rank'' (denoted $\theta_\pre$).

[Next verbatim source excerpt]

\begin{restatable}[Equilibrium under group differences]{proposition}{propEquiEnv}\label{prop:equi-env}
Define $\theta_\pre$ as:
%	\begin{equation*}
	\(\theta_\pre := f^{-1}_\mix(f(\theta_\true)\cdot \psi)\),
%	\end{equation*}
		where $f^{-1}_\mix$	is the CDF for the environment-scaled skill, $f(\theta_\true)\cdot \psi$:
	\begin{equation*}
	f^{-1}_\mix(x) := \frac{1}{2}f^{-1}(x/\psi_\A)+ \frac{1}{2}f^{-1}(x/\psi_\B).
	\end{equation*}


[Next verbatim source excerpt]

\begin{definition}[Welfare gap]
	Let $\swelf^\G(\theta_\true)$ denote post-effort welfare of an applicant with latent skill ranking $\theta_\true$ from group $\G \in \{\A, \B\}$, i.e.,
	\begin{equation}
	\swelf^\G(\theta_\true):= \lambda(\theta_\post(\theta_\true, \psi_\G)) - \cost(e(\theta_\true, \psi_\G)).
	\end{equation}
	%\lnote{Using function notation for $\theta_\post$ and $e$? or some other way to express the welfare as a function of $\theta_\true$.}
	We define the \emph{%pointwise
		welfare gap} as %a function of rank $\theta_\true$ is
	\begin{equation*}
		\welfgap(\theta_\true):= \swelf^\A(\theta_\true) - \swelf^\B(\theta_\true).
	\end{equation*}
%	We define the \emph{subpopulation welfare gap} as %a function of baseline rank $\overline{\theta}_\true$ is
%	\begin{equation*}
%		\popwelfgap(\overline{\theta}_\true):=\E[\welfgap(\theta_\true) \mid \theta_\true > \overline{\theta}_\true].
%	\end{equation*}
\end{definition}

The welfare gap captures differences in admission probabilities \textit{and} in the effort needed to achieve such probabilities. Our next notion, \textit{access}, captures whether a decision policy \emph{includes} the disadvantaged group in the admitted class, regardless of effort. %This is captured by the notion of \emph{access}.

\begin{definition}[Access]
 \emph{Access} is the overall probability of admission of the disadvantaged group.
	\begin{equation*}
		\access :=  \E_{\theta_\true}[\lambda(\theta_\post(\theta_\true, \psi_\B))].
	\end{equation*}
% \nikhil{Defn is a bit confusing. What is the expectation over, and what does that conditioning mean?}
\end{definition}

[Next verbatim source excerpt]

	Denote the group-specific rank threshold for group $\G$ as
	\begin{equation*}
	\thresG{c} := f^{-1}\left(\frac{f_\mix(c)}{\psi_\G}\right).
	\end{equation*}
\end{ReviewVerbatim}

% report-context-presentation-sha256: 16265e6401473cf3c85b47c3a2f48907c1f14d36cd38cc80fc86055d51df0113
\paragraph{Settled source reading}
\begin{sloppypar}
Equilibrium means that almost every applicant best responds against every feasible deviation;\allowbreak{} uniqueness is up to a null set.\allowbreak{}
\end{sloppypar}
\paragraph{Settled source reading}
\begin{sloppypar}
Applicants know the fixed priority used at score-\allowbreak{}boundary ties.\allowbreak{} This clarifies the stated model, not a corrected admission policy or a new economic assumption.\allowbreak{}
\end{sloppypar}
\paragraph{Settled source reading}
\begin{sloppypar}
The main-\allowbreak{}text group comparisons are exact under clarified weak/\allowbreak{}strict interpretation.\allowbreak{} Strictness requires positive disadvantaged effort;\allowbreak{} classical derivative existence is qualified separately.\allowbreak{}
\end{sloppypar}

\paragraph{Expanded PaperInterface specification}
\begin{ReviewVerbatim}
∀ (cost production skill : ℝ → ℝ) (effort score tie : Bool × ℝ → ℝ) (rho c psiA psiB : ℝ),
  LBG22StrategicRanking.RankingPrimitives cost production skill 0 →
    rho ∈ Set.Ioo 0 1 →
      c ∈ Set.Ioc 0 (1 - rho) →
        0 < psiB →
          psiB < psiA →
            DifferentiableOn ℝ cost (Set.Ioi 0) →
              DifferentiableOn ℝ production (Set.Ioi 0) →
                DifferentiableOn ℝ skill (Set.Ioo 0 1) →
                  Measurable tie →
                    Set.InjOn tie {x : Bool × ℝ | x.2 ∈ Set.Ioc 0 1} →
                      LBG22StrategicRanking.EnvironmentEquilibrium cost production skill psiA psiB tie 1
                          (LBG22StrategicRanking.sourceTwoLevelCutoff c)
                          (LBG22StrategicRanking.sourceTwoLevelReward rho c) effort score →
                        (∀ (group : Bool),
                            ∀ᵐ (t : ℝ) ∂LBG22StrategicRanking.unitRankMeasure,
                              LBG22StrategicRanking.sourceTwoLevelReward rho c
                                  ↑(LBG22StrategicRanking.finiteLowerRankBand
                                      (fun (i : Fin 2) => LBG22StrategicRanking.sourceTwoLevelCutoff c ↑i)
                                      (LBG22StrategicRanking.tieBrokenRank LBG22StrategicRanking.sourceTwoGroupMeasure
                                        score tie (group, t))) =
                                if LBG22StrategicRanking.sourceGroupThreshold skill psiA psiB group c < t then
                                  rho / (1 - c)
                                else 0) ∧
                          (∀ᵐ (t : ℝ) ∂LBG22StrategicRanking.unitRankMeasure,
                              0 ≤
                                  LBG22StrategicRanking.sourceActualGroupWelfareGap cost effort score tie
                                    (fun (i : Fin 2) => LBG22StrategicRanking.sourceTwoLevelCutoff c ↑i)
                                    (LBG22StrategicRanking.sourceTwoLevelReward rho c) t ∧
                                (LBG22StrategicRanking.sourceGroupThreshold skill psiA psiB false c < t →
                                  (0 <
                                      LBG22StrategicRanking.sourceActualGroupWelfareGap cost effort score tie
                                        (fun (i : Fin 2) => LBG22StrategicRanking.sourceTwoLevelCutoff c ↑i)
                                        (LBG22StrategicRanking.sourceTwoLevelReward rho c) t ↔
                                    0 < effort (false, t)))) ∧
                            ∀ (pureEffort pureScore : Bool × ℝ → ℝ),
                              LBG22StrategicRanking.EnvironmentEquilibrium cost production skill psiA psiB tie 0
                                  (fun (k : ℕ) => if k = 0 then 0 else 1) (fun (x : ℕ) => rho) pureEffort pureScore →
                                ∀ᵐ (t : ℝ) ∂LBG22StrategicRanking.unitRankMeasure,
                                  LBG22StrategicRanking.sourceActualGroupWelfareGap cost pureEffort pureScore tie
                                      (fun (x : Fin 1) => 0) (fun (x : ℕ) => rho) t =
                                    0
\end{ReviewVerbatim}

\paragraph{Retained governing declarations}
\begin{itemize}\raggedright
\item \hyperlink{paper-prerequisite-b8e06f39e94d0049}{LBG22\allowbreak{}Strategic\allowbreak{}Ranking.\allowbreak{}Environment\allowbreak{}Equilibrium}
\item \hyperlink{paper-prerequisite-9bbcc12f4a1736bf}{LBG22\allowbreak{}Strategic\allowbreak{}Ranking.\allowbreak{}Ranking\allowbreak{}Primitives}
\item \hyperlink{paper-prerequisite-a42f440ddc254fa7}{LBG22\allowbreak{}Strategic\allowbreak{}Ranking.\allowbreak{}finite\allowbreak{}Lower\allowbreak{}Rank\allowbreak{}Band}
\item \hyperlink{paper-prerequisite-d3e90f3908818ca5}{LBG22\allowbreak{}Strategic\allowbreak{}Ranking.\allowbreak{}source\allowbreak{}Actual\allowbreak{}Group\allowbreak{}Welfare\allowbreak{}Gap}
\item \hyperlink{paper-prerequisite-b2f8283910c81320}{LBG22\allowbreak{}Strategic\allowbreak{}Ranking.\allowbreak{}source\allowbreak{}Group\allowbreak{}Threshold}
\item \hyperlink{paper-prerequisite-5b3be494cde30a37}{LBG22\allowbreak{}Strategic\allowbreak{}Ranking.\allowbreak{}source\allowbreak{}Two\allowbreak{}Group\allowbreak{}Measure}
\item \hyperlink{paper-prerequisite-f88aac58507e44ca}{LBG22\allowbreak{}Strategic\allowbreak{}Ranking.\allowbreak{}source\allowbreak{}Two\allowbreak{}Level\allowbreak{}Cutoff}
\item \hyperlink{paper-prerequisite-30d97cd43eb79017}{LBG22\allowbreak{}Strategic\allowbreak{}Ranking.\allowbreak{}source\allowbreak{}Two\allowbreak{}Level\allowbreak{}Reward}
\item \hyperlink{paper-prerequisite-b7c71e53177b76c6}{LBG22\allowbreak{}Strategic\allowbreak{}Ranking.\allowbreak{}tie\allowbreak{}Broken\allowbreak{}Rank}
\item \hyperlink{governing-declaration-678cf48d549cb620}{LBG22\allowbreak{}Strategic\allowbreak{}Ranking.\allowbreak{}unit\allowbreak{}Rank\allowbreak{}Measure}
\end{itemize}
\paragraph{Lean-elaborated claim roles}\begin{ReviewVerbatim}
1. parameter: ℝ → ℝ
2. parameter: ℝ → ℝ
3. parameter: ℝ → ℝ
4. parameter: Bool × ℝ → ℝ
5. parameter: Bool × ℝ → ℝ
6. parameter: Bool × ℝ → ℝ
7. parameter: ℝ
8. parameter: ℝ
9. parameter: ℝ
10. parameter: ℝ
11. assumption: LBG22StrategicRanking.RankingPrimitives cost production skill 0
12. assumption: rho ∈ Set.Ioo 0 1
13. assumption: c ∈ Set.Ioc 0 (1 - rho)
14. assumption: 0 < psiB
15. assumption: psiB < psiA
16. assumption: DifferentiableOn ℝ cost (Set.Ioi 0)
17. assumption: DifferentiableOn ℝ production (Set.Ioi 0)
18. assumption: DifferentiableOn ℝ skill (Set.Ioo 0 1)
19. assumption: Measurable tie
20. assumption: Set.InjOn tie {x : Bool × ℝ | x.2 ∈ Set.Ioc 0 1}
21. assumption: LBG22StrategicRanking.EnvironmentEquilibrium cost production skill psiA psiB tie 1
  (LBG22StrategicRanking.sourceTwoLevelCutoff c) (LBG22StrategicRanking.sourceTwoLevelReward rho c) effort score
22. conclusion: (∀ (group : Bool),
    ∀ᵐ (t : ℝ) ∂LBG22StrategicRanking.unitRankMeasure,
      LBG22StrategicRanking.sourceTwoLevelReward rho c
          ↑(LBG22StrategicRanking.finiteLowerRankBand
              (fun (i : Fin 2) => LBG22StrategicRanking.sourceTwoLevelCutoff c ↑i)
              (LBG22StrategicRanking.tieBrokenRank LBG22StrategicRanking.sourceTwoGroupMeasure score tie (group, t))) =
        if LBG22StrategicRanking.sourceGroupThreshold skill psiA psiB group c < t then rho / (1 - c) else 0) ∧
  (∀ᵐ (t : ℝ) ∂LBG22StrategicRanking.unitRankMeasure,
      0 ≤
          LBG22StrategicRanking.sourceActualGroupWelfareGap cost effort score tie
            (fun (i : Fin 2) => LBG22StrategicRanking.sourceTwoLevelCutoff c ↑i)
            (LBG22StrategicRanking.sourceTwoLevelReward rho c) t ∧
        (LBG22StrategicRanking.sourceGroupThreshold skill psiA psiB false c < t →
          (0 <
              LBG22StrategicRanking.sourceActualGroupWelfareGap cost effort score tie
                (fun (i : Fin 2) => LBG22StrategicRanking.sourceTwoLevelCutoff c ↑i)
                (LBG22StrategicRanking.sourceTwoLevelReward rho c) t ↔
            0 < effort (false, t)))) ∧
    ∀ (pureEffort pureScore : Bool × ℝ → ℝ),
      LBG22StrategicRanking.EnvironmentEquilibrium cost production skill psiA psiB tie 0
          (fun (k : ℕ) => if k = 0 then 0 else 1) (fun (x : ℕ) => rho) pureEffort pureScore →
        ∀ᵐ (t : ℝ) ∂LBG22StrategicRanking.unitRankMeasure,
          LBG22StrategicRanking.sourceActualGroupWelfareGap cost pureEffort pureScore tie (fun (x : Fin 1) => 0)
              (fun (x : ℕ) => rho) t =
            0
\end{ReviewVerbatim}

\paragraph{Lean proof endpoint (built separately)}\begin{ReviewVerbatim}
LBG22StrategicRanking.ProofInterface.environmentWelfare
\end{ReviewVerbatim}

\paragraph{Recorded source-to-Spec screening}
\textbf{Verdict:} \texttt{matches}
\\
{\raggedright
\textbf{Reason:} The target reproduces group-\allowbreak{}specific two-\allowbreak{}level admission thresholds, nonnegative pointwise welfare gap, strict positivity exactly when disadvantaged effort is positive in the common-\allowbreak{}admission region, and identically zero gap under pure randomization.\allowbreak{} The weak/\allowbreak{}strict qualification is precisely the approved LBG22-\allowbreak{}WEAK-\allowbreak{}STRICT-\allowbreak{}01 reading, and the displayed welfare, threshold, mixture, equilibrium, and priority dependencies preserve the source formulas and domains.\allowbreak{}
\par}
\reviewmatch{Matches source input}
\noindent\textbf{Reviewer annotation}\par
\reviewerbox
\clearpage
\hypertarget{source-claim-7af287b5fa9fa319}{}
\section*{Review row 9}
\textbf{Source locator:} {\footnotesize\raggedright cited publication, lines 321-\allowbreak{}326\par}
\paragraph{Verbatim source input}
\begin{ReviewVerbatim}
\begin{restatable}[Welfare gap increases with $c$]{proposition}{propWelfareGap}\label{prop:deriv-welfare-gap}
	Consider the setting in Definition~\ref{assump:2-level-fixed-cap}, with the school's chosen admissions policy $c=\bar{c}$, where $\bar{c} \le  (f_\mix)^{-1}(f(1)\cdot \psi_\B)$. Then, for any $\theta_\true \geq \thresB{\bar{c}}$, that is, $\theta_\true$ is in the ``High'' region, we have
%	\begin{equation*}
	\(\left.\frac{\partial \welfgap(\theta_\true)}{\partial c}\right\vert_{c=\bar{c}}> 0\).
%	\end{equation*}
\end{restatable}

[Next verbatim source excerpt]

\parbold{Applicants}
There is a unit mass of \textit{applicants}, indexed by an observed index $\omega \in [0,1]$ distributed uniformly.\footnote{The index $\omega$ should be interpreted as each applicant's ``name,'' uncorrelated with skill, used solely for tie-breaking.}  Each applicant has a latent (unobserved) skill level represented by some measurable function of $\omega$.
We assume that {the distribution of the skills} has no atoms and that the CDF of this distribution is strictly increasing.
Using the CDF to map the skill of an applicant to a rank in~$[0,1]$, each applicant gets an (unobserved) rank $\theta_\pre =\theta_\pre(\omega)$ which by our assumption on the CDF is again uniformly distributed in $[0, 1]$ (the higher the rank the better).  With this setup, the skill of an applicant with rank $\theta_\pre$ can be written as $f(\theta_\pre)$ where $f$ is a strictly increasing, continuous function.  {It will be notationally convenient to label applicants by their rank $\theta_\pre$, though the reader should note that in contrast to $\omega$, $\theta_\pre(\omega)$ is assumed to be unobservable.}

%3. Each applicant exerts some effort e, at cost \cost(e).


Each applicant chooses an \textit{effort} level $e \ge 0$, the result of which is an observed, post-effort \textit{score}, {$v=v(e,\theta_\pre) = {g(e)}\cdot{f(\theta_\pre)}$}. \lledit{In other words, we assume the post-effort score is the product of two components, associated with the effort and the pre-effort skill respectively.}
{We assume} that the effort transfer function $g: [0, \infty) \mapsto [0, \infty)$ is a continuous, concave, strictly increasing function, representing that marginal effort improves one's score but has diminishing returns.
{%Since applicants can be labelled by their skill level $f(\theta_\pre)$, or equivalently, by their rank $\theta_{\pre}$,
{T}he strategies of the applicants can {then} be described by a function $\theta_\pre \mapsto e(\theta_\pre)$.
%\cb{$\omega\mapsto e(\omega)$.}%, leading to a distribution over post-effort scores $v=v(\theta_\pre))$ via the uniform distribution of $\theta_\pre$ in $[0,1]$.
}
  Each {applicant} is then ranked according to their score $v$%(where we break ties uniformly at random, if they appear)
, resulting in a \textit{post-effort rank} $\theta_\post$.
Note that the ranking $\theta_\post$ is slightly less trivial than a ranking of the skills, since the scores might have ties, which have to be resolved.

{\textbf{Tie Breaking.}
Given a choice of strategies $\theta_\pre \mapsto e(\theta_\pre)$, let $F$ be the CDF of the scores $v(e(\theta_\pre),\theta_\pre)$.
Since atoms for the distribution of $v$ would lead to ties for ranks defined as $F(v)$, we will use the labels of the applicants to break ties with the help of a (publicly announced) tie-breaking function $\Gamma(\omega)$, defined, e.g., via a collision free hash of the applicant names. Here  we require  that $\Gamma$  is a measurable function from $[0,1]$ to $[0,1]$ that maps different applicant labels to different values.
We then use $\Gamma$ to resolve the atoms of $F$, leading to a ranking function $v\mapsto \gamma(\omega, v)$ which is equal to $F(v)$ except when $v$ is an atom of the score distribution, {in which case it takes values in the ``gap interval'' $[F_-(v),  F(v)]$, where $F_-(v)$ is the left limit of $F$ at $v$.}  We  construct $\gamma$ in such a way that it
gives the uniform distribution for $\theta_\post(\omega)$ if we set  $\theta_\post(\omega)=\gamma(\omega,v(e(\theta_\pre(\omega)),\theta_\pre(\omega)))$.\footnote{See Remark~\ref{rem:gamma} for details on this construction.}

[Next verbatim source excerpt]

\parbold{Ranking designer (school)} A single \textit{school} is admitting applicants, based on their ranking. In particular, the school can choose a ranking reward function $\lambda : [0, 1] \mapsto [0, 1]$, such that an applicant with post-effort rank $\theta_\post$ is admitted with probability $\lambda(\theta_\post)$. We assume that $\lambda$ is non-decreasing and that the school has a constraint on the overall probability, such that in expectation it admits a number of applicants equal to {a} capacity constraint  $\rho \in (0,1)$, i.e., \lledit{$\E_{\theta_\post}[\lambda(\theta_\post)] = \rho $}.\footnote{\lledit{We use $\E_{\theta_\post}[\cdot]$ to denote an integral over $\theta_\post$ (and $\E[\cdot]$ to denote an integral over $\omega$)} with respect to the Lebesgue measure; informally, this can be thought of as averaging over the applicant population.} We may also refer to $\lambda$, informally, as the admission policy.



For simplicity, we further assume that $\lambda$ is a step-function with $K$ distinct levels $\ell_0 < \dots < \ell_{K-1}$, and $K-1$ cut-points parameterized by $c_1< \dots < c_{K-1}$ (with $c_0 = 0$, $c_K = 1$). %$\ell_0, \dots, \ell_{K-1}$.
In other words, we have $\lambda(\theta) = \ell_k$, for all $\theta \in \psi_k \triangleq [c_{k}, c_{k+1})$. Thus, applicants in the same post-rank interval $\theta \in \psi_k$ receive the same reward.

[Next verbatim source excerpt]

\parbold{Individual applicant welfare and equilibrium} Given the designer's function $\lambda$ and the effort levels of other applicants, each applicant chooses effort $e$ to maximize their individual welfare, % $W$,
\begin{equation*}
	\iwelf(e, \lambda(\theta_\post)) = \lambda(\theta_\post) - \cost(e),
\end{equation*}
where the effort cost function $\cost$ is non-negative, continuous, and strictly convex {on $[0,\infty)$, with $e_0:=\argmin_e \cost(e)$ and $p(e_0) = 0$.}
%\lnote{also want to assume that $\cost$ attains its minimum at $0$?}
%Let $e_0:=\argmin_e \cost(e)$ and $p(e_0) = 0$.

Applicants are assumed not to personally benefit from increasing their score $v$, except through the corresponding increase in their rank and reward. While the definition of $\cost$ does not preclude the applicant receiving any intrinsic benefit from exerting effort, we assume that the \emph{net} benefit from effort is non-positive.

[Next verbatim source excerpt]

\begin{definition} [Equilibrium]\label{def:equi}
{Given a}
%Let the
tie-breaking function $\Gamma(\omega)$
%be a pre-announced ordering over index $\omega$, to break ties in post-effort value. Then, let $\gamma(\omega, v)$ be the 1-1 ranking function defined as the CDF of post-effort values $v$ except where there are ties in $v$, in which case the ordering $\Gamma$ is used to produce unique ranks such that the range of $\gamma$ is uniform on $[0, 1]$.\footnote{More formally, note that each discontinuity in the CDF of post-effort values $v$ has a height equal to the measure of the applicants who are tied at that value, and so $\gamma$ for those applicants can be filled in using the CDF of $\Gamma(\omega)$ restricted to the applicants $\omega$ who are tied. Thus, the range of $\gamma$ is uniform on $[0,1]$.
%
{and}
%	Then, given
	a ranking probability function $\lambda$, an \textit{equilibrium} is a set of effort levels and post-effort ranking rewards for each applicant, $\{e(\theta_\pre),  \lambda(\theta_\post(\theta_\pre))\}$ such that, for all
	$\omega$, %$\theta_\pre$,
	\begin{align*}
		e(\theta_\pre(\omega)) &\in \argmax_{e} \iwelf\left(e, \lambda\left(\gamma\left(\omega, v\left(e,\theta_\pre(\omega)\right)\right)\right)\right) %\triangleq \argmax_{e} \left[ \lambda(\gamma(v(e,\theta_\pre)) - \cost(e) \right].
		\\
		\theta_\post(\theta_\pre(\omega)) &= \gamma(\omega, v(e(\theta_\pre(\omega)),\theta_\pre(\omega))),
	\end{align*}
	{where $\gamma(\omega,v)$ is the ranking induced by the CDF of the scores resulting from effort levels $\{e(\theta_\pre)\}$ and the tie-breaking function $\Gamma$.\footnote{See Remark~\ref{rem:gamma-effort} on $\gamma$ and the set of efforts.%Formally, $\gamma$ depends on the set of efforts. Given a fixed set and $\gamma$, for each applicant $\omega$ the first condition considers the counter-factual ranking of $\omega$ with different post-effort values but using the same ranking function. As defined, $\gamma$ yields a uniform distribution of ranks with such measure $0$ changes. In Appendix \Cref{lem:deviationsmeasure0}, we further prove that two effort sets equal up to sets of measure $0$ induce the same ranking function $\gamma$, and so the condition is consistent.%, and so we set it at the $\gamma$ induced by the effort levels $\{e(\theta_\pre)\}$.%, and so it does not vary.
	}
	}
	%In particular, we denote $\theta_\post(\theta_\pre) := \gamma(v(e(\theta_\pre),\theta_\pre))$.
	%\lledit{and for any $\theta_\pre, \theta'_\pre$, we have $\lambda(\theta_\post(\theta_\pre)) > \lambda(\theta_\post(\theta'_\pre)) \implies v(e(\theta_\pre),\theta_\pre) > v(e(\theta_\pre'),\theta_\pre')$.}
\end{definition}

[Next verbatim source excerpt]

\begin{definition}[Two-level policy]\label{assump:2-level-fixed-cap} In our baseline two-level function parameterized by cut-off $c \in (0,1-\rho]$, each applicant with post-effort rank  $\theta_\post \geq c$ is admitted with probability $\ell_1 =\frac{\rho}{1-c} \in (0,1]$. Others are rejected, $\ell_0 = 0$.
\end{definition}

Note that standard non-randomized admissions policies are equivalent to case where $c = 1 - \rho$: the highest scoring applicants up to the capacity constraint are accepted with probability 1 and all others are rejected. We call this case \textbf{non-randomized admissions}. The other extreme is the one-level policy, \textbf{pure randomization}, where each applicant is admitted with probability ${\rho}$, i.e., $\ell_0 =\rho$. Decreasing the cut-off $c$ can be viewed as increasing the level of randomization in the admissions policy.


{We now reason about how various welfare and fairness metrics vary with $\lambda$ (in the two-level policy class, just $c$). To simplify the presentation, we assume that $f, g$ and $\cost$ are differentiable and that baseline effort $e_0$ is~$0$.}

[Next verbatim source excerpt]

\paragraph{Model and equilibria characterization} We now denote each applicant's latent skill rank as $\theta_\true \in [0,1]$. In addition to the latent skill, each applicant has an (unobserved) environmental factor $\psi \in \Psi$ that represents how favorable their environment is for attaining a higher score. Because a favorable environment results in a higher rate of return for effort, we model the environment as a multiplicative factor in the score (see~e.g., \citet{calsamiglia09decent}). Formally, the post-effort score is a function of the latent skill, the environmental factor, and the effort level:
\begin{equation*}\label{eq:env-score}
v = \psi \cdot  g(e)\cdot f(\theta_\true),
\end{equation*}
where $g, f$ are as defined in Section~\ref{sec:model}. %We also require $f$ to be strictly monotone.

%We model structural inequalities as heterogeneous returns to effort. In contrast to results in Section~\ref{sec:equillemmas}, the ranks of applicants are no longer preserved when there are heterogeneous returns to effort.

 We assume there are two groups of applicants, $\A$ and $\B$, and the distribution of skill is the same in both groups. Group $\A$ has a more favorable environment factor, that is,  $\Psi = \{\psi_\A, \psi_\B\}$ and $\psi_\A > \psi_\B$. Thus we will also refer to $\B$ as the ``disadvantaged group". To simplify our presentation, we assume each group is half of the total applicant population, though the results in this section generalize.% to the setting where groups are not equal-sized.
%
We defer all proofs in this section to Appendix~\ref{app:environment}.

We begin by characterizing the equilibrium ranking under the designer's policy $\lambda$. The equilibrium effort levels and post-effort ranks are as defined in Definition~\ref{def:equi}, except they are now group-dependent, that is, we have $e(\theta_\true, \psi)$ and $\theta_\post(\theta_\true, \psi)$. Because of the differences in $\psi$, the post-effort ranking $\theta_\post$ is now group-dependent, and in general is not equal to $\theta_\true$. To apply Proposition~\ref{lem:rankpreservedindex} as before, we construct an ``environment-scaled pre-effort rank'' (denoted $\theta_\pre$).

[Next verbatim source excerpt]

\begin{restatable}[Equilibrium under group differences]{proposition}{propEquiEnv}\label{prop:equi-env}
Define $\theta_\pre$ as:
%	\begin{equation*}
	\(\theta_\pre := f^{-1}_\mix(f(\theta_\true)\cdot \psi)\),
%	\end{equation*}
		where $f^{-1}_\mix$	is the CDF for the environment-scaled skill, $f(\theta_\true)\cdot \psi$:
	\begin{equation*}
	f^{-1}_\mix(x) := \frac{1}{2}f^{-1}(x/\psi_\A)+ \frac{1}{2}f^{-1}(x/\psi_\B).
	\end{equation*}


[Next verbatim source excerpt]

\begin{definition}[Welfare gap]
	Let $\swelf^\G(\theta_\true)$ denote post-effort welfare of an applicant with latent skill ranking $\theta_\true$ from group $\G \in \{\A, \B\}$, i.e.,
	\begin{equation}
	\swelf^\G(\theta_\true):= \lambda(\theta_\post(\theta_\true, \psi_\G)) - \cost(e(\theta_\true, \psi_\G)).
	\end{equation}
	%\lnote{Using function notation for $\theta_\post$ and $e$? or some other way to express the welfare as a function of $\theta_\true$.}
	We define the \emph{%pointwise
		welfare gap} as %a function of rank $\theta_\true$ is
	\begin{equation*}
		\welfgap(\theta_\true):= \swelf^\A(\theta_\true) - \swelf^\B(\theta_\true).
	\end{equation*}
%	We define the \emph{subpopulation welfare gap} as %a function of baseline rank $\overline{\theta}_\true$ is
%	\begin{equation*}
%		\popwelfgap(\overline{\theta}_\true):=\E[\welfgap(\theta_\true) \mid \theta_\true > \overline{\theta}_\true].
%	\end{equation*}
\end{definition}

The welfare gap captures differences in admission probabilities \textit{and} in the effort needed to achieve such probabilities. Our next notion, \textit{access}, captures whether a decision policy \emph{includes} the disadvantaged group in the admitted class, regardless of effort. %This is captured by the notion of \emph{access}.

\begin{definition}[Access]
 \emph{Access} is the overall probability of admission of the disadvantaged group.
	\begin{equation*}
		\access :=  \E_{\theta_\true}[\lambda(\theta_\post(\theta_\true, \psi_\B))].
	\end{equation*}
% \nikhil{Defn is a bit confusing. What is the expectation over, and what does that conditioning mean?}
\end{definition}

[Next verbatim source excerpt]

	Denote the group-specific rank threshold for group $\G$ as
	\begin{equation*}
	\thresG{c} := f^{-1}\left(\frac{f_\mix(c)}{\psi_\G}\right).
	\end{equation*}
\end{ReviewVerbatim}

% report-context-presentation-sha256: 16265e6401473cf3c85b47c3a2f48907c1f14d36cd38cc80fc86055d51df0113
\paragraph{Settled source reading}
\begin{sloppypar}
Equilibrium means that almost every applicant best responds against every feasible deviation;\allowbreak{} uniqueness is up to a null set.\allowbreak{}
\end{sloppypar}
\paragraph{Settled source reading}
\begin{sloppypar}
Applicants know the fixed priority used at score-\allowbreak{}boundary ties.\allowbreak{} This clarifies the stated model, not a corrected admission policy or a new economic assumption.\allowbreak{}
\end{sloppypar}
\paragraph{Settled source reading}
\begin{sloppypar}
The main-\allowbreak{}text group comparisons are exact under clarified weak/\allowbreak{}strict interpretation.\allowbreak{} Strictness requires positive disadvantaged effort;\allowbreak{} classical derivative existence is qualified separately.\allowbreak{}
\end{sloppypar}
\paragraph{Formalized review target}
The archival source above is not asserted equivalent to this formalized target.
\begin{ReviewVerbatim}
Keep the source group primitives, differentiability and concavity conditions, capacity, population, and tie order fixed. In any family of actual equilibria, for each 0<c<=d<=1-rho and almost every latent quantile admitted in both groups at d, Delta(d,t)>=Delta(c,t), strictly if c<d and disadvantaged effort at c is positive. At each feasible cutoff c and almost every commonly admitted latent quantile t, every derivative of Delta along the feasible policy set is nonnegative and is positive when disadvantaged effort is positive. This derivative-sign clause assumes that the specified derivative exists and uses the feasible one-sided derivative at the deterministic endpoint; it does not assert everywhere classical differentiability. These are finite comparisons and where-defined derivative consequences, with no new economic shape assumption.
\end{ReviewVerbatim}

\textbf{Archival source anchor:} {\footnotesize\raggedright cited publication, lines 321-\allowbreak{}326\par}
\paragraph{Expanded PaperInterface specification}
\begin{ReviewVerbatim}
∀ (cost production skill : ℝ → ℝ) (effort score : ℝ → Bool × ℝ → ℝ) (tie : Bool × ℝ → ℝ) (rho psiA psiB : ℝ),
  LBG22StrategicRanking.RankingPrimitives cost production skill 0 →
    rho ∈ Set.Ioo 0 1 →
      0 < psiB →
        psiB < psiA →
          DifferentiableOn ℝ cost (Set.Ioi 0) →
            DifferentiableOn ℝ production (Set.Ioi 0) →
              DifferentiableOn ℝ skill (Set.Ioo 0 1) →
                Measurable tie →
                  Set.InjOn tie {x : Bool × ℝ | x.2 ∈ Set.Ioc 0 1} →
                    (∀ c ∈ Set.Ioc 0 (1 - rho),
                        LBG22StrategicRanking.EnvironmentEquilibrium cost production skill psiA psiB tie 1
                          (LBG22StrategicRanking.sourceTwoLevelCutoff c)
                          (LBG22StrategicRanking.sourceTwoLevelReward rho c) (effort c) (score c)) →
                      (∀ c ∈ Set.Ioc 0 (1 - rho),
                          ∀ d ∈ Set.Ioc 0 (1 - rho),
                            c ≤ d →
                              ∀ᵐ (t : ℝ) ∂LBG22StrategicRanking.unitRankMeasure,
                                LBG22StrategicRanking.sourceGroupThreshold skill psiA psiB false d < t →
                                  LBG22StrategicRanking.sourceActualGroupWelfareGap cost (effort c) (score c) tie
                                        (fun (i : Fin 2) => LBG22StrategicRanking.sourceTwoLevelCutoff c ↑i)
                                        (LBG22StrategicRanking.sourceTwoLevelReward rho c) t ≤
                                      LBG22StrategicRanking.sourceActualGroupWelfareGap cost (effort d) (score d) tie
                                        (fun (i : Fin 2) => LBG22StrategicRanking.sourceTwoLevelCutoff d ↑i)
                                        (LBG22StrategicRanking.sourceTwoLevelReward rho d) t ∧
                                    (c < d →
                                      0 < effort c (false, t) →
                                        LBG22StrategicRanking.sourceActualGroupWelfareGap cost (effort c) (score c) tie
                                            (fun (i : Fin 2) => LBG22StrategicRanking.sourceTwoLevelCutoff c ↑i)
                                            (LBG22StrategicRanking.sourceTwoLevelReward rho c) t <
                                          LBG22StrategicRanking.sourceActualGroupWelfareGap cost (effort d) (score d)
                                            tie (fun (i : Fin 2) => LBG22StrategicRanking.sourceTwoLevelCutoff d ↑i)
                                            (LBG22StrategicRanking.sourceTwoLevelReward rho d) t)) ∧
                        ∀ c ∈ Set.Ioc 0 (1 - rho),
                          ∀ᵐ (t : ℝ) ∂LBG22StrategicRanking.unitRankMeasure,
                            LBG22StrategicRanking.sourceGroupThreshold skill psiA psiB false c < t →
                              ∀ (v : ℝ),
                                HasDerivWithinAt
                                    (fun (d : ℝ) =>
                                      LBG22StrategicRanking.sourceActualGroupWelfareGap cost (effort d) (score d) tie
                                        (fun (i : Fin 2) => LBG22StrategicRanking.sourceTwoLevelCutoff d ↑i)
                                        (LBG22StrategicRanking.sourceTwoLevelReward rho d) t)
                                    v (Set.Ioc 0 (1 - rho)) c →
                                  0 ≤ v ∧ (0 < effort c (false, t) → 0 < v)
\end{ReviewVerbatim}

\paragraph{Retained governing declarations}
\begin{itemize}\raggedright
\item \hyperlink{paper-prerequisite-b8e06f39e94d0049}{LBG22\allowbreak{}Strategic\allowbreak{}Ranking.\allowbreak{}Environment\allowbreak{}Equilibrium}
\item \hyperlink{paper-prerequisite-9bbcc12f4a1736bf}{LBG22\allowbreak{}Strategic\allowbreak{}Ranking.\allowbreak{}Ranking\allowbreak{}Primitives}
\item \hyperlink{paper-prerequisite-d3e90f3908818ca5}{LBG22\allowbreak{}Strategic\allowbreak{}Ranking.\allowbreak{}source\allowbreak{}Actual\allowbreak{}Group\allowbreak{}Welfare\allowbreak{}Gap}
\item \hyperlink{paper-prerequisite-b2f8283910c81320}{LBG22\allowbreak{}Strategic\allowbreak{}Ranking.\allowbreak{}source\allowbreak{}Group\allowbreak{}Threshold}
\item \hyperlink{paper-prerequisite-f88aac58507e44ca}{LBG22\allowbreak{}Strategic\allowbreak{}Ranking.\allowbreak{}source\allowbreak{}Two\allowbreak{}Level\allowbreak{}Cutoff}
\item \hyperlink{paper-prerequisite-30d97cd43eb79017}{LBG22\allowbreak{}Strategic\allowbreak{}Ranking.\allowbreak{}source\allowbreak{}Two\allowbreak{}Level\allowbreak{}Reward}
\item \hyperlink{governing-declaration-678cf48d549cb620}{LBG22\allowbreak{}Strategic\allowbreak{}Ranking.\allowbreak{}unit\allowbreak{}Rank\allowbreak{}Measure}
\end{itemize}
\paragraph{Lean-elaborated claim roles}\begin{ReviewVerbatim}
1. parameter: ℝ → ℝ
2. parameter: ℝ → ℝ
3. parameter: ℝ → ℝ
4. parameter: ℝ → Bool × ℝ → ℝ
5. parameter: ℝ → Bool × ℝ → ℝ
6. parameter: Bool × ℝ → ℝ
7. parameter: ℝ
8. parameter: ℝ
9. parameter: ℝ
10. assumption: LBG22StrategicRanking.RankingPrimitives cost production skill 0
11. assumption: rho ∈ Set.Ioo 0 1
12. assumption: 0 < psiB
13. assumption: psiB < psiA
14. assumption: DifferentiableOn ℝ cost (Set.Ioi 0)
15. assumption: DifferentiableOn ℝ production (Set.Ioi 0)
16. assumption: DifferentiableOn ℝ skill (Set.Ioo 0 1)
17. assumption: Measurable tie
18. assumption: Set.InjOn tie {x : Bool × ℝ | x.2 ∈ Set.Ioc 0 1}
19. assumption: ∀ c ∈ Set.Ioc 0 (1 - rho),
  LBG22StrategicRanking.EnvironmentEquilibrium cost production skill psiA psiB tie 1
    (LBG22StrategicRanking.sourceTwoLevelCutoff c) (LBG22StrategicRanking.sourceTwoLevelReward rho c) (effort c)
    (score c)
20. conclusion: (∀ c ∈ Set.Ioc 0 (1 - rho),
    ∀ d ∈ Set.Ioc 0 (1 - rho),
      c ≤ d →
        ∀ᵐ (t : ℝ) ∂LBG22StrategicRanking.unitRankMeasure,
          LBG22StrategicRanking.sourceGroupThreshold skill psiA psiB false d < t →
            LBG22StrategicRanking.sourceActualGroupWelfareGap cost (effort c) (score c) tie
                  (fun (i : Fin 2) => LBG22StrategicRanking.sourceTwoLevelCutoff c ↑i)
                  (LBG22StrategicRanking.sourceTwoLevelReward rho c) t ≤
                LBG22StrategicRanking.sourceActualGroupWelfareGap cost (effort d) (score d) tie
                  (fun (i : Fin 2) => LBG22StrategicRanking.sourceTwoLevelCutoff d ↑i)
                  (LBG22StrategicRanking.sourceTwoLevelReward rho d) t ∧
              (c < d →
                0 < effort c (false, t) →
                  LBG22StrategicRanking.sourceActualGroupWelfareGap cost (effort c) (score c) tie
                      (fun (i : Fin 2) => LBG22StrategicRanking.sourceTwoLevelCutoff c ↑i)
                      (LBG22StrategicRanking.sourceTwoLevelReward rho c) t <
                    LBG22StrategicRanking.sourceActualGroupWelfareGap cost (effort d) (score d) tie
                      (fun (i : Fin 2) => LBG22StrategicRanking.sourceTwoLevelCutoff d ↑i)
                      (LBG22StrategicRanking.sourceTwoLevelReward rho d) t)) ∧
  ∀ c ∈ Set.Ioc 0 (1 - rho),
    ∀ᵐ (t : ℝ) ∂LBG22StrategicRanking.unitRankMeasure,
      LBG22StrategicRanking.sourceGroupThreshold skill psiA psiB false c < t →
        ∀ (v : ℝ),
          HasDerivWithinAt
              (fun (d : ℝ) =>
                LBG22StrategicRanking.sourceActualGroupWelfareGap cost (effort d) (score d) tie
                  (fun (i : Fin 2) => LBG22StrategicRanking.sourceTwoLevelCutoff d ↑i)
                  (LBG22StrategicRanking.sourceTwoLevelReward rho d) t)
              v (Set.Ioc 0 (1 - rho)) c →
            0 ≤ v ∧ (0 < effort c (false, t) → 0 < v)
\end{ReviewVerbatim}

\paragraph{Lean proof endpoint (built separately)}\begin{ReviewVerbatim}
LBG22StrategicRanking.ProofInterface.welfareGapDerivative
\end{ReviewVerbatim}

\paragraph{Recorded source-to-Spec screening}
\textbf{Verdict:} \texttt{matches\_approved\_corrected\_target}
\\
{\raggedright
\textbf{Reason:} The selected route disclaims archival equivalence and binds defect LBG22-\allowbreak{}C3 to the complete approved correction.\allowbreak{} The target states finite weak welfare-\allowbreak{}gap monotonicity for c<=d in the common-\allowbreak{}admission region, strictness when c advantaged policy changes and disadvantaged effort is positive, and, only conditional on a displayed HasDerivWithinAt witness, a nonnegative feasible derivative that is positive with positive disadvantaged effort.\allowbreak{} This matches the approved where-\allowbreak{}defined, one-\allowbreak{}sided-\allowbreak{}at-\allowbreak{}endpoint derivative reading and its displayed group-\allowbreak{}welfare, threshold, equilibrium, and priority prerequisites.\allowbreak{}
\par}
\reviewmatch{Matches formalized target}
\noindent\textbf{Reviewer annotation}\par
\reviewerbox
\clearpage
\hypertarget{source-claim-dcc0514ddd510442}{}
\section*{Review row 10}
\textbf{Source locator:} {\footnotesize\raggedright cited publication, lines 335-\allowbreak{}339\par}
\paragraph{Verbatim source input}
\begin{ReviewVerbatim}
\begin{restatable}[Access decreases with $c$]{proposition}{propAccess}
	Pure randomization has higher $\access$ than any two-level policy.
	Moreover, if $f^{-1}$ is convex, $\access$ for  two-level policies is non-increasing in $c$. %\nikhil{Why "non-strictly"?} \nikhil{All K with $K \leq 2$ reads funny}
	\label{prop:access}
\end{restatable}

[Next verbatim source excerpt]

\parbold{Applicants}
There is a unit mass of \textit{applicants}, indexed by an observed index $\omega \in [0,1]$ distributed uniformly.\footnote{The index $\omega$ should be interpreted as each applicant's ``name,'' uncorrelated with skill, used solely for tie-breaking.}  Each applicant has a latent (unobserved) skill level represented by some measurable function of $\omega$.
We assume that {the distribution of the skills} has no atoms and that the CDF of this distribution is strictly increasing.
Using the CDF to map the skill of an applicant to a rank in~$[0,1]$, each applicant gets an (unobserved) rank $\theta_\pre =\theta_\pre(\omega)$ which by our assumption on the CDF is again uniformly distributed in $[0, 1]$ (the higher the rank the better).  With this setup, the skill of an applicant with rank $\theta_\pre$ can be written as $f(\theta_\pre)$ where $f$ is a strictly increasing, continuous function.  {It will be notationally convenient to label applicants by their rank $\theta_\pre$, though the reader should note that in contrast to $\omega$, $\theta_\pre(\omega)$ is assumed to be unobservable.}

%3. Each applicant exerts some effort e, at cost \cost(e).


Each applicant chooses an \textit{effort} level $e \ge 0$, the result of which is an observed, post-effort \textit{score}, {$v=v(e,\theta_\pre) = {g(e)}\cdot{f(\theta_\pre)}$}. \lledit{In other words, we assume the post-effort score is the product of two components, associated with the effort and the pre-effort skill respectively.}
{We assume} that the effort transfer function $g: [0, \infty) \mapsto [0, \infty)$ is a continuous, concave, strictly increasing function, representing that marginal effort improves one's score but has diminishing returns.
{%Since applicants can be labelled by their skill level $f(\theta_\pre)$, or equivalently, by their rank $\theta_{\pre}$,
{T}he strategies of the applicants can {then} be described by a function $\theta_\pre \mapsto e(\theta_\pre)$.
%\cb{$\omega\mapsto e(\omega)$.}%, leading to a distribution over post-effort scores $v=v(\theta_\pre))$ via the uniform distribution of $\theta_\pre$ in $[0,1]$.
}
  Each {applicant} is then ranked according to their score $v$%(where we break ties uniformly at random, if they appear)
, resulting in a \textit{post-effort rank} $\theta_\post$.
Note that the ranking $\theta_\post$ is slightly less trivial than a ranking of the skills, since the scores might have ties, which have to be resolved.

{\textbf{Tie Breaking.}
Given a choice of strategies $\theta_\pre \mapsto e(\theta_\pre)$, let $F$ be the CDF of the scores $v(e(\theta_\pre),\theta_\pre)$.
Since atoms for the distribution of $v$ would lead to ties for ranks defined as $F(v)$, we will use the labels of the applicants to break ties with the help of a (publicly announced) tie-breaking function $\Gamma(\omega)$, defined, e.g., via a collision free hash of the applicant names. Here  we require  that $\Gamma$  is a measurable function from $[0,1]$ to $[0,1]$ that maps different applicant labels to different values.
We then use $\Gamma$ to resolve the atoms of $F$, leading to a ranking function $v\mapsto \gamma(\omega, v)$ which is equal to $F(v)$ except when $v$ is an atom of the score distribution, {in which case it takes values in the ``gap interval'' $[F_-(v),  F(v)]$, where $F_-(v)$ is the left limit of $F$ at $v$.}  We  construct $\gamma$ in such a way that it
gives the uniform distribution for $\theta_\post(\omega)$ if we set  $\theta_\post(\omega)=\gamma(\omega,v(e(\theta_\pre(\omega)),\theta_\pre(\omega)))$.\footnote{See Remark~\ref{rem:gamma} for details on this construction.}

[Next verbatim source excerpt]

\parbold{Ranking designer (school)} A single \textit{school} is admitting applicants, based on their ranking. In particular, the school can choose a ranking reward function $\lambda : [0, 1] \mapsto [0, 1]$, such that an applicant with post-effort rank $\theta_\post$ is admitted with probability $\lambda(\theta_\post)$. We assume that $\lambda$ is non-decreasing and that the school has a constraint on the overall probability, such that in expectation it admits a number of applicants equal to {a} capacity constraint  $\rho \in (0,1)$, i.e., \lledit{$\E_{\theta_\post}[\lambda(\theta_\post)] = \rho $}.\footnote{\lledit{We use $\E_{\theta_\post}[\cdot]$ to denote an integral over $\theta_\post$ (and $\E[\cdot]$ to denote an integral over $\omega$)} with respect to the Lebesgue measure; informally, this can be thought of as averaging over the applicant population.} We may also refer to $\lambda$, informally, as the admission policy.



For simplicity, we further assume that $\lambda$ is a step-function with $K$ distinct levels $\ell_0 < \dots < \ell_{K-1}$, and $K-1$ cut-points parameterized by $c_1< \dots < c_{K-1}$ (with $c_0 = 0$, $c_K = 1$). %$\ell_0, \dots, \ell_{K-1}$.
In other words, we have $\lambda(\theta) = \ell_k$, for all $\theta \in \psi_k \triangleq [c_{k}, c_{k+1})$. Thus, applicants in the same post-rank interval $\theta \in \psi_k$ receive the same reward.

[Next verbatim source excerpt]

\parbold{Individual applicant welfare and equilibrium} Given the designer's function $\lambda$ and the effort levels of other applicants, each applicant chooses effort $e$ to maximize their individual welfare, % $W$,
\begin{equation*}
	\iwelf(e, \lambda(\theta_\post)) = \lambda(\theta_\post) - \cost(e),
\end{equation*}
where the effort cost function $\cost$ is non-negative, continuous, and strictly convex {on $[0,\infty)$, with $e_0:=\argmin_e \cost(e)$ and $p(e_0) = 0$.}
%\lnote{also want to assume that $\cost$ attains its minimum at $0$?}
%Let $e_0:=\argmin_e \cost(e)$ and $p(e_0) = 0$.

Applicants are assumed not to personally benefit from increasing their score $v$, except through the corresponding increase in their rank and reward. While the definition of $\cost$ does not preclude the applicant receiving any intrinsic benefit from exerting effort, we assume that the \emph{net} benefit from effort is non-positive.

[Next verbatim source excerpt]

\begin{definition} [Equilibrium]\label{def:equi}
{Given a}
%Let the
tie-breaking function $\Gamma(\omega)$
%be a pre-announced ordering over index $\omega$, to break ties in post-effort value. Then, let $\gamma(\omega, v)$ be the 1-1 ranking function defined as the CDF of post-effort values $v$ except where there are ties in $v$, in which case the ordering $\Gamma$ is used to produce unique ranks such that the range of $\gamma$ is uniform on $[0, 1]$.\footnote{More formally, note that each discontinuity in the CDF of post-effort values $v$ has a height equal to the measure of the applicants who are tied at that value, and so $\gamma$ for those applicants can be filled in using the CDF of $\Gamma(\omega)$ restricted to the applicants $\omega$ who are tied. Thus, the range of $\gamma$ is uniform on $[0,1]$.
%
{and}
%	Then, given
	a ranking probability function $\lambda$, an \textit{equilibrium} is a set of effort levels and post-effort ranking rewards for each applicant, $\{e(\theta_\pre),  \lambda(\theta_\post(\theta_\pre))\}$ such that, for all
	$\omega$, %$\theta_\pre$,
	\begin{align*}
		e(\theta_\pre(\omega)) &\in \argmax_{e} \iwelf\left(e, \lambda\left(\gamma\left(\omega, v\left(e,\theta_\pre(\omega)\right)\right)\right)\right) %\triangleq \argmax_{e} \left[ \lambda(\gamma(v(e,\theta_\pre)) - \cost(e) \right].
		\\
		\theta_\post(\theta_\pre(\omega)) &= \gamma(\omega, v(e(\theta_\pre(\omega)),\theta_\pre(\omega))),
	\end{align*}
	{where $\gamma(\omega,v)$ is the ranking induced by the CDF of the scores resulting from effort levels $\{e(\theta_\pre)\}$ and the tie-breaking function $\Gamma$.\footnote{See Remark~\ref{rem:gamma-effort} on $\gamma$ and the set of efforts.%Formally, $\gamma$ depends on the set of efforts. Given a fixed set and $\gamma$, for each applicant $\omega$ the first condition considers the counter-factual ranking of $\omega$ with different post-effort values but using the same ranking function. As defined, $\gamma$ yields a uniform distribution of ranks with such measure $0$ changes. In Appendix \Cref{lem:deviationsmeasure0}, we further prove that two effort sets equal up to sets of measure $0$ induce the same ranking function $\gamma$, and so the condition is consistent.%, and so we set it at the $\gamma$ induced by the effort levels $\{e(\theta_\pre)\}$.%, and so it does not vary.
	}
	}
	%In particular, we denote $\theta_\post(\theta_\pre) := \gamma(v(e(\theta_\pre),\theta_\pre))$.
	%\lledit{and for any $\theta_\pre, \theta'_\pre$, we have $\lambda(\theta_\post(\theta_\pre)) > \lambda(\theta_\post(\theta'_\pre)) \implies v(e(\theta_\pre),\theta_\pre) > v(e(\theta_\pre'),\theta_\pre')$.}
\end{definition}

[Next verbatim source excerpt]

\begin{definition}[Two-level policy]\label{assump:2-level-fixed-cap} In our baseline two-level function parameterized by cut-off $c \in (0,1-\rho]$, each applicant with post-effort rank  $\theta_\post \geq c$ is admitted with probability $\ell_1 =\frac{\rho}{1-c} \in (0,1]$. Others are rejected, $\ell_0 = 0$.
\end{definition}

Note that standard non-randomized admissions policies are equivalent to case where $c = 1 - \rho$: the highest scoring applicants up to the capacity constraint are accepted with probability 1 and all others are rejected. We call this case \textbf{non-randomized admissions}. The other extreme is the one-level policy, \textbf{pure randomization}, where each applicant is admitted with probability ${\rho}$, i.e., $\ell_0 =\rho$. Decreasing the cut-off $c$ can be viewed as increasing the level of randomization in the admissions policy.


{We now reason about how various welfare and fairness metrics vary with $\lambda$ (in the two-level policy class, just $c$). To simplify the presentation, we assume that $f, g$ and $\cost$ are differentiable and that baseline effort $e_0$ is~$0$.}

[Next verbatim source excerpt]

\paragraph{Model and equilibria characterization} We now denote each applicant's latent skill rank as $\theta_\true \in [0,1]$. In addition to the latent skill, each applicant has an (unobserved) environmental factor $\psi \in \Psi$ that represents how favorable their environment is for attaining a higher score. Because a favorable environment results in a higher rate of return for effort, we model the environment as a multiplicative factor in the score (see~e.g., \citet{calsamiglia09decent}). Formally, the post-effort score is a function of the latent skill, the environmental factor, and the effort level:
\begin{equation*}\label{eq:env-score}
v = \psi \cdot  g(e)\cdot f(\theta_\true),
\end{equation*}
where $g, f$ are as defined in Section~\ref{sec:model}. %We also require $f$ to be strictly monotone.

%We model structural inequalities as heterogeneous returns to effort. In contrast to results in Section~\ref{sec:equillemmas}, the ranks of applicants are no longer preserved when there are heterogeneous returns to effort.

 We assume there are two groups of applicants, $\A$ and $\B$, and the distribution of skill is the same in both groups. Group $\A$ has a more favorable environment factor, that is,  $\Psi = \{\psi_\A, \psi_\B\}$ and $\psi_\A > \psi_\B$. Thus we will also refer to $\B$ as the ``disadvantaged group". To simplify our presentation, we assume each group is half of the total applicant population, though the results in this section generalize.% to the setting where groups are not equal-sized.
%
We defer all proofs in this section to Appendix~\ref{app:environment}.

We begin by characterizing the equilibrium ranking under the designer's policy $\lambda$. The equilibrium effort levels and post-effort ranks are as defined in Definition~\ref{def:equi}, except they are now group-dependent, that is, we have $e(\theta_\true, \psi)$ and $\theta_\post(\theta_\true, \psi)$. Because of the differences in $\psi$, the post-effort ranking $\theta_\post$ is now group-dependent, and in general is not equal to $\theta_\true$. To apply Proposition~\ref{lem:rankpreservedindex} as before, we construct an ``environment-scaled pre-effort rank'' (denoted $\theta_\pre$).

[Next verbatim source excerpt]

\begin{restatable}[Equilibrium under group differences]{proposition}{propEquiEnv}\label{prop:equi-env}
Define $\theta_\pre$ as:
%	\begin{equation*}
	\(\theta_\pre := f^{-1}_\mix(f(\theta_\true)\cdot \psi)\),
%	\end{equation*}
		where $f^{-1}_\mix$	is the CDF for the environment-scaled skill, $f(\theta_\true)\cdot \psi$:
	\begin{equation*}
	f^{-1}_\mix(x) := \frac{1}{2}f^{-1}(x/\psi_\A)+ \frac{1}{2}f^{-1}(x/\psi_\B).
	\end{equation*}


[Next verbatim source excerpt]

\begin{definition}[Welfare gap]
	Let $\swelf^\G(\theta_\true)$ denote post-effort welfare of an applicant with latent skill ranking $\theta_\true$ from group $\G \in \{\A, \B\}$, i.e.,
	\begin{equation}
	\swelf^\G(\theta_\true):= \lambda(\theta_\post(\theta_\true, \psi_\G)) - \cost(e(\theta_\true, \psi_\G)).
	\end{equation}
	%\lnote{Using function notation for $\theta_\post$ and $e$? or some other way to express the welfare as a function of $\theta_\true$.}
	We define the \emph{%pointwise
		welfare gap} as %a function of rank $\theta_\true$ is
	\begin{equation*}
		\welfgap(\theta_\true):= \swelf^\A(\theta_\true) - \swelf^\B(\theta_\true).
	\end{equation*}
%	We define the \emph{subpopulation welfare gap} as %a function of baseline rank $\overline{\theta}_\true$ is
%	\begin{equation*}
%		\popwelfgap(\overline{\theta}_\true):=\E[\welfgap(\theta_\true) \mid \theta_\true > \overline{\theta}_\true].
%	\end{equation*}
\end{definition}

The welfare gap captures differences in admission probabilities \textit{and} in the effort needed to achieve such probabilities. Our next notion, \textit{access}, captures whether a decision policy \emph{includes} the disadvantaged group in the admitted class, regardless of effort. %This is captured by the notion of \emph{access}.

\begin{definition}[Access]
 \emph{Access} is the overall probability of admission of the disadvantaged group.
	\begin{equation*}
		\access :=  \E_{\theta_\true}[\lambda(\theta_\post(\theta_\true, \psi_\B))].
	\end{equation*}
% \nikhil{Defn is a bit confusing. What is the expectation over, and what does that conditioning mean?}
\end{definition}

[Next verbatim source excerpt]

	Denote the group-specific rank threshold for group $\G$ as
	\begin{equation*}
	\thresG{c} := f^{-1}\left(\frac{f_\mix(c)}{\psi_\G}\right).
	\end{equation*}
\end{ReviewVerbatim}

% report-context-presentation-sha256: 90f9116e292dd71fddfd6b445ca0b060a481d026855bdb9725f2833c2cf62a31
\paragraph{Settled source reading}
\begin{sloppypar}
Equilibrium means that almost every applicant best responds against every feasible deviation;\allowbreak{} uniqueness is up to a null set.\allowbreak{}
\end{sloppypar}
\paragraph{Settled source reading}
\begin{sloppypar}
Applicants know the fixed priority used at score-\allowbreak{}boundary ties.\allowbreak{} This clarifies the stated model, not a corrected admission policy or a new economic assumption.\allowbreak{}
\end{sloppypar}

\paragraph{Expanded PaperInterface specification}
\begin{ReviewVerbatim}
∀ (cost production skill : ℝ → ℝ) (effort score : ℝ → Bool × ℝ → ℝ) (tie : Bool × ℝ → ℝ) (baseline rho psiA psiB : ℝ),
  LBG22StrategicRanking.RankingPrimitives cost production skill baseline →
    rho ∈ Set.Ioo 0 1 →
      0 < psiB →
        psiB < psiA →
          Measurable tie →
            Set.InjOn tie {x : Bool × ℝ | x.2 ∈ Set.Ioc 0 1} →
              (∀ c ∈ Set.Ioc 0 (1 - rho),
                  LBG22StrategicRanking.EnvironmentEquilibrium cost production skill psiA psiB tie 1
                    (LBG22StrategicRanking.sourceTwoLevelCutoff c) (LBG22StrategicRanking.sourceTwoLevelReward rho c)
                    (effort c) (score c)) →
                have pureAccess : ℝ :=
                  LBG22StrategicRanking.sourceActualDisadvantagedAccess
                    (fun (x : Bool × ℝ) =>
                      production baseline * LBG22StrategicRanking.sourceEnvironmentSkill skill psiA psiB x)
                    tie (fun (x : Fin 1) => 0) fun (x : ℕ) => rho;
                have access : ℝ → ℝ := fun (c : ℝ) =>
                  LBG22StrategicRanking.sourceActualDisadvantagedAccess (score c) tie
                    (fun (i : Fin 2) => LBG22StrategicRanking.sourceTwoLevelCutoff c ↑i)
                    (LBG22StrategicRanking.sourceTwoLevelReward rho c);
                pureAccess = rho ∧
                  (∀ c ∈ Set.Ioc 0 (1 - rho), access c < pureAccess) ∧
                    (ConvexOn ℝ (Set.Icc (skill 0) (skill 1)) (LBG22StrategicRanking.sourceSkillCDF skill) →
                      AntitoneOn access (Set.Ioc 0 (1 - rho)))
\end{ReviewVerbatim}

\paragraph{Retained governing declarations}
\begin{itemize}\raggedright
\item \hyperlink{paper-prerequisite-b8e06f39e94d0049}{LBG22\allowbreak{}Strategic\allowbreak{}Ranking.\allowbreak{}Environment\allowbreak{}Equilibrium}
\item \hyperlink{paper-prerequisite-9bbcc12f4a1736bf}{LBG22\allowbreak{}Strategic\allowbreak{}Ranking.\allowbreak{}Ranking\allowbreak{}Primitives}
\item \hyperlink{paper-prerequisite-6f1ee9b200124096}{LBG22\allowbreak{}Strategic\allowbreak{}Ranking.\allowbreak{}source\allowbreak{}Actual\allowbreak{}Disadvantaged\allowbreak{}Access}
\item \hyperlink{paper-prerequisite-73413973c7594479}{LBG22\allowbreak{}Strategic\allowbreak{}Ranking.\allowbreak{}source\allowbreak{}Environment\allowbreak{}Skill}
\item \hyperlink{governing-declaration-15b79809af308d40}{LBG22\allowbreak{}Strategic\allowbreak{}Ranking.\allowbreak{}source\allowbreak{}Skill\allowbreak{}CDF}
\item \hyperlink{paper-prerequisite-f88aac58507e44ca}{LBG22\allowbreak{}Strategic\allowbreak{}Ranking.\allowbreak{}source\allowbreak{}Two\allowbreak{}Level\allowbreak{}Cutoff}
\item \hyperlink{paper-prerequisite-30d97cd43eb79017}{LBG22\allowbreak{}Strategic\allowbreak{}Ranking.\allowbreak{}source\allowbreak{}Two\allowbreak{}Level\allowbreak{}Reward}
\end{itemize}
\paragraph{Lean-elaborated claim roles}\begin{ReviewVerbatim}
1. parameter: ℝ → ℝ
2. parameter: ℝ → ℝ
3. parameter: ℝ → ℝ
4. parameter: ℝ → Bool × ℝ → ℝ
5. parameter: ℝ → Bool × ℝ → ℝ
6. parameter: Bool × ℝ → ℝ
7. parameter: ℝ
8. parameter: ℝ
9. parameter: ℝ
10. parameter: ℝ
11. assumption: LBG22StrategicRanking.RankingPrimitives cost production skill baseline
12. assumption: rho ∈ Set.Ioo 0 1
13. assumption: 0 < psiB
14. assumption: psiB < psiA
15. assumption: Measurable tie
16. assumption: Set.InjOn tie {x : Bool × ℝ | x.2 ∈ Set.Ioc 0 1}
17. assumption: ∀ c ∈ Set.Ioc 0 (1 - rho),
  LBG22StrategicRanking.EnvironmentEquilibrium cost production skill psiA psiB tie 1
    (LBG22StrategicRanking.sourceTwoLevelCutoff c) (LBG22StrategicRanking.sourceTwoLevelReward rho c) (effort c)
    (score c)
18. conclusion: (LBG22StrategicRanking.sourceActualDisadvantagedAccess
      (fun (x : Bool × ℝ) => production baseline * LBG22StrategicRanking.sourceEnvironmentSkill skill psiA psiB x) tie
      (fun (x : Fin 1) => 0) fun (x : ℕ) => rho) =
    rho ∧
  (∀ c ∈ Set.Ioc 0 (1 - rho),
      (fun (c : ℝ) =>
            LBG22StrategicRanking.sourceActualDisadvantagedAccess (score c) tie
              (fun (i : Fin 2) => LBG22StrategicRanking.sourceTwoLevelCutoff c ↑i)
              (LBG22StrategicRanking.sourceTwoLevelReward rho c))
          c <
        LBG22StrategicRanking.sourceActualDisadvantagedAccess
          (fun (x : Bool × ℝ) => production baseline * LBG22StrategicRanking.sourceEnvironmentSkill skill psiA psiB x)
          tie (fun (x : Fin 1) => 0) fun (x : ℕ) => rho) ∧
    (ConvexOn ℝ (Set.Icc (skill 0) (skill 1)) (LBG22StrategicRanking.sourceSkillCDF skill) →
      AntitoneOn
        (fun (c : ℝ) =>
          LBG22StrategicRanking.sourceActualDisadvantagedAccess (score c) tie
            (fun (i : Fin 2) => LBG22StrategicRanking.sourceTwoLevelCutoff c ↑i)
            (LBG22StrategicRanking.sourceTwoLevelReward rho c))
        (Set.Ioc 0 (1 - rho)))
\end{ReviewVerbatim}

\paragraph{Lean proof endpoint (built separately)}\begin{ReviewVerbatim}
LBG22StrategicRanking.ProofInterface.access
\end{ReviewVerbatim}

\paragraph{Recorded source-to-Spec screening}
\textbf{Verdict:} \texttt{matches}
\\
{\raggedright
\textbf{Reason:} The expanded target matches Proposition 4.\allowbreak{}4:\allowbreak{} pure randomization gives disadvantaged access rho, every admissible two-\allowbreak{}level cutoff gives strictly lower access, and convexity of the displayed skill CDF on the skill range implies antitonicity in the cutoff.\allowbreak{} The displayed group mixture, threshold/\allowbreak{}access integral, fixed-\allowbreak{}priority rank, policy, and a.\allowbreak{}e.\allowbreak{} all-\allowbreak{}deviation equilibrium prerequisites implement the source domains under the approved continuum and known-\allowbreak{}priority conventions.\allowbreak{}
\par}
\reviewmatch{Matches source input}
\noindent\textbf{Reviewer annotation}\par
\reviewerbox
\clearpage
\hypertarget{source-claim-6cc5784dc863c5dc}{}
\section*{Review row 11}
\textbf{Source locator:} {\footnotesize\raggedright cited publication, lines 1698-\allowbreak{}1703\par}
\paragraph{Verbatim source input}
\begin{ReviewVerbatim}
\begin{corollary}[Effort comparative statics]
	\label{corr:effectcomparativestatics}
	{Assume that $g$ and $\cost$ are differentiable.}
	{If} $\ell_k$ increases {and $\ell_j$ decreases for some $j>k$} (fixing all other parameters), then {$e_{i}(\theta)$ for all $i < k$ are unaffected,
		$e_k$ is weakly increasing, and efforts $e_{k+1}\dots$, $e_j$ are weakly decreasing.}
\end{corollary}

[Next verbatim source excerpt]

\parbold{Applicants}
There is a unit mass of \textit{applicants}, indexed by an observed index $\omega \in [0,1]$ distributed uniformly.\footnote{The index $\omega$ should be interpreted as each applicant's ``name,'' uncorrelated with skill, used solely for tie-breaking.}  Each applicant has a latent (unobserved) skill level represented by some measurable function of $\omega$.
We assume that {the distribution of the skills} has no atoms and that the CDF of this distribution is strictly increasing.
Using the CDF to map the skill of an applicant to a rank in~$[0,1]$, each applicant gets an (unobserved) rank $\theta_\pre =\theta_\pre(\omega)$ which by our assumption on the CDF is again uniformly distributed in $[0, 1]$ (the higher the rank the better).  With this setup, the skill of an applicant with rank $\theta_\pre$ can be written as $f(\theta_\pre)$ where $f$ is a strictly increasing, continuous function.  {It will be notationally convenient to label applicants by their rank $\theta_\pre$, though the reader should note that in contrast to $\omega$, $\theta_\pre(\omega)$ is assumed to be unobservable.}

%3. Each applicant exerts some effort e, at cost \cost(e).


Each applicant chooses an \textit{effort} level $e \ge 0$, the result of which is an observed, post-effort \textit{score}, {$v=v(e,\theta_\pre) = {g(e)}\cdot{f(\theta_\pre)}$}. \lledit{In other words, we assume the post-effort score is the product of two components, associated with the effort and the pre-effort skill respectively.}
{We assume} that the effort transfer function $g: [0, \infty) \mapsto [0, \infty)$ is a continuous, concave, strictly increasing function, representing that marginal effort improves one's score but has diminishing returns.
{%Since applicants can be labelled by their skill level $f(\theta_\pre)$, or equivalently, by their rank $\theta_{\pre}$,
{T}he strategies of the applicants can {then} be described by a function $\theta_\pre \mapsto e(\theta_\pre)$.
%\cb{$\omega\mapsto e(\omega)$.}%, leading to a distribution over post-effort scores $v=v(\theta_\pre))$ via the uniform distribution of $\theta_\pre$ in $[0,1]$.
}
  Each {applicant} is then ranked according to their score $v$%(where we break ties uniformly at random, if they appear)
, resulting in a \textit{post-effort rank} $\theta_\post$.
Note that the ranking $\theta_\post$ is slightly less trivial than a ranking of the skills, since the scores might have ties, which have to be resolved.

{\textbf{Tie Breaking.}
Given a choice of strategies $\theta_\pre \mapsto e(\theta_\pre)$, let $F$ be the CDF of the scores $v(e(\theta_\pre),\theta_\pre)$.
Since atoms for the distribution of $v$ would lead to ties for ranks defined as $F(v)$, we will use the labels of the applicants to break ties with the help of a (publicly announced) tie-breaking function $\Gamma(\omega)$, defined, e.g., via a collision free hash of the applicant names. Here  we require  that $\Gamma$  is a measurable function from $[0,1]$ to $[0,1]$ that maps different applicant labels to different values.
We then use $\Gamma$ to resolve the atoms of $F$, leading to a ranking function $v\mapsto \gamma(\omega, v)$ which is equal to $F(v)$ except when $v$ is an atom of the score distribution, {in which case it takes values in the ``gap interval'' $[F_-(v),  F(v)]$, where $F_-(v)$ is the left limit of $F$ at $v$.}  We  construct $\gamma$ in such a way that it
gives the uniform distribution for $\theta_\post(\omega)$ if we set  $\theta_\post(\omega)=\gamma(\omega,v(e(\theta_\pre(\omega)),\theta_\pre(\omega)))$.\footnote{See Remark~\ref{rem:gamma} for details on this construction.}

[Next verbatim source excerpt]

\parbold{Ranking designer (school)} A single \textit{school} is admitting applicants, based on their ranking. In particular, the school can choose a ranking reward function $\lambda : [0, 1] \mapsto [0, 1]$, such that an applicant with post-effort rank $\theta_\post$ is admitted with probability $\lambda(\theta_\post)$. We assume that $\lambda$ is non-decreasing and that the school has a constraint on the overall probability, such that in expectation it admits a number of applicants equal to {a} capacity constraint  $\rho \in (0,1)$, i.e., \lledit{$\E_{\theta_\post}[\lambda(\theta_\post)] = \rho $}.\footnote{\lledit{We use $\E_{\theta_\post}[\cdot]$ to denote an integral over $\theta_\post$ (and $\E[\cdot]$ to denote an integral over $\omega$)} with respect to the Lebesgue measure; informally, this can be thought of as averaging over the applicant population.} We may also refer to $\lambda$, informally, as the admission policy.



For simplicity, we further assume that $\lambda$ is a step-function with $K$ distinct levels $\ell_0 < \dots < \ell_{K-1}$, and $K-1$ cut-points parameterized by $c_1< \dots < c_{K-1}$ (with $c_0 = 0$, $c_K = 1$). %$\ell_0, \dots, \ell_{K-1}$.
In other words, we have $\lambda(\theta) = \ell_k$, for all $\theta \in \psi_k \triangleq [c_{k}, c_{k+1})$. Thus, applicants in the same post-rank interval $\theta \in \psi_k$ receive the same reward.

[Next verbatim source excerpt]

\parbold{Individual applicant welfare and equilibrium} Given the designer's function $\lambda$ and the effort levels of other applicants, each applicant chooses effort $e$ to maximize their individual welfare, % $W$,
\begin{equation*}
	\iwelf(e, \lambda(\theta_\post)) = \lambda(\theta_\post) - \cost(e),
\end{equation*}
where the effort cost function $\cost$ is non-negative, continuous, and strictly convex {on $[0,\infty)$, with $e_0:=\argmin_e \cost(e)$ and $p(e_0) = 0$.}
%\lnote{also want to assume that $\cost$ attains its minimum at $0$?}
%Let $e_0:=\argmin_e \cost(e)$ and $p(e_0) = 0$.

Applicants are assumed not to personally benefit from increasing their score $v$, except through the corresponding increase in their rank and reward. While the definition of $\cost$ does not preclude the applicant receiving any intrinsic benefit from exerting effort, we assume that the \emph{net} benefit from effort is non-positive.

[Next verbatim source excerpt]

\begin{definition} [Equilibrium]\label{def:equi}
{Given a}
%Let the
tie-breaking function $\Gamma(\omega)$
%be a pre-announced ordering over index $\omega$, to break ties in post-effort value. Then, let $\gamma(\omega, v)$ be the 1-1 ranking function defined as the CDF of post-effort values $v$ except where there are ties in $v$, in which case the ordering $\Gamma$ is used to produce unique ranks such that the range of $\gamma$ is uniform on $[0, 1]$.\footnote{More formally, note that each discontinuity in the CDF of post-effort values $v$ has a height equal to the measure of the applicants who are tied at that value, and so $\gamma$ for those applicants can be filled in using the CDF of $\Gamma(\omega)$ restricted to the applicants $\omega$ who are tied. Thus, the range of $\gamma$ is uniform on $[0,1]$.
%
{and}
%	Then, given
	a ranking probability function $\lambda$, an \textit{equilibrium} is a set of effort levels and post-effort ranking rewards for each applicant, $\{e(\theta_\pre),  \lambda(\theta_\post(\theta_\pre))\}$ such that, for all
	$\omega$, %$\theta_\pre$,
	\begin{align*}
		e(\theta_\pre(\omega)) &\in \argmax_{e} \iwelf\left(e, \lambda\left(\gamma\left(\omega, v\left(e,\theta_\pre(\omega)\right)\right)\right)\right) %\triangleq \argmax_{e} \left[ \lambda(\gamma(v(e,\theta_\pre)) - \cost(e) \right].
		\\
		\theta_\post(\theta_\pre(\omega)) &= \gamma(\omega, v(e(\theta_\pre(\omega)),\theta_\pre(\omega))),
	\end{align*}
	{where $\gamma(\omega,v)$ is the ranking induced by the CDF of the scores resulting from effort levels $\{e(\theta_\pre)\}$ and the tie-breaking function $\Gamma$.\footnote{See Remark~\ref{rem:gamma-effort} on $\gamma$ and the set of efforts.%Formally, $\gamma$ depends on the set of efforts. Given a fixed set and $\gamma$, for each applicant $\omega$ the first condition considers the counter-factual ranking of $\omega$ with different post-effort values but using the same ranking function. As defined, $\gamma$ yields a uniform distribution of ranks with such measure $0$ changes. In Appendix \Cref{lem:deviationsmeasure0}, we further prove that two effort sets equal up to sets of measure $0$ induce the same ranking function $\gamma$, and so the condition is consistent.%, and so we set it at the $\gamma$ induced by the effort levels $\{e(\theta_\pre)\}$.%, and so it does not vary.
	}
	}
	%In particular, we denote $\theta_\post(\theta_\pre) := \gamma(v(e(\theta_\pre),\theta_\pre))$.
	%\lledit{and for any $\theta_\pre, \theta'_\pre$, we have $\lambda(\theta_\post(\theta_\pre)) > \lambda(\theta_\post(\theta'_\pre)) \implies v(e(\theta_\pre),\theta_\pre) > v(e(\theta_\pre'),\theta_\pre')$.}
\end{definition}
\end{ReviewVerbatim}

% report-context-presentation-sha256: 90f9116e292dd71fddfd6b445ca0b060a481d026855bdb9725f2833c2cf62a31
\paragraph{Settled source reading}
\begin{sloppypar}
Equilibrium means that almost every applicant best responds against every feasible deviation;\allowbreak{} uniqueness is up to a null set.\allowbreak{}
\end{sloppypar}
\paragraph{Settled source reading}
\begin{sloppypar}
Applicants know the fixed priority used at score-\allowbreak{}boundary ties.\allowbreak{} This clarifies the stated model, not a corrected admission policy or a new economic assumption.\allowbreak{}
\end{sloppypar}

\paragraph{Expanded PaperInterface specification}
\begin{ReviewVerbatim}
∀ (cost production skill tie effort newEffort score newScore : ℝ → ℝ) (baseline rho : ℝ) (n k j : ℕ)
  (cutoff reward newReward : ℕ → ℝ),
  LBG22StrategicRanking.RankingPrimitives cost production skill baseline →
    LBG22StrategicRanking.FiniteAdmissionPolicy n cutoff reward rho →
      LBG22StrategicRanking.FiniteAdmissionPolicy n cutoff newReward rho →
        DifferentiableOn ℝ cost (Set.Ioi 0) →
          DifferentiableOn ℝ production (Set.Ioi 0) →
            k < j →
              j ≤ n →
                reward k < newReward k →
                  newReward j < reward j →
                    (∀ i ≤ n, i ≠ k → i ≠ j → newReward i = reward i) →
                      Measurable tie →
                        Set.InjOn tie (Set.Ioc 0 1) →
                          LBG22StrategicRanking.RankingEquilibrium cost production skill tie n cutoff reward effort
                              score →
                            LBG22StrategicRanking.RankingEquilibrium cost production skill tie n cutoff newReward
                                newEffort newScore →
                              ∀ᵐ (t : ℝ) ∂LBG22StrategicRanking.unitRankMeasure,
                                have i : ℕ :=
                                  ↑(LBG22StrategicRanking.finiteLowerRankBand (fun (l : Fin (n + 1)) => cutoff ↑l) t);
                                (i < k → newEffort t = effort t) ∧
                                  (i = k → effort t ≤ newEffort t) ∧ (k < i → i ≤ j → newEffort t ≤ effort t)
\end{ReviewVerbatim}

\paragraph{Retained governing declarations}
\begin{itemize}\raggedright
\item \hyperlink{paper-prerequisite-d056dd7c8060653c}{LBG22\allowbreak{}Strategic\allowbreak{}Ranking.\allowbreak{}Finite\allowbreak{}Admission\allowbreak{}Policy}
\item \hyperlink{paper-prerequisite-2a1b692f4e3136c6}{LBG22\allowbreak{}Strategic\allowbreak{}Ranking.\allowbreak{}Ranking\allowbreak{}Equilibrium}
\item \hyperlink{paper-prerequisite-9bbcc12f4a1736bf}{LBG22\allowbreak{}Strategic\allowbreak{}Ranking.\allowbreak{}Ranking\allowbreak{}Primitives}
\item \hyperlink{paper-prerequisite-a42f440ddc254fa7}{LBG22\allowbreak{}Strategic\allowbreak{}Ranking.\allowbreak{}finite\allowbreak{}Lower\allowbreak{}Rank\allowbreak{}Band}
\item \hyperlink{governing-declaration-678cf48d549cb620}{LBG22\allowbreak{}Strategic\allowbreak{}Ranking.\allowbreak{}unit\allowbreak{}Rank\allowbreak{}Measure}
\end{itemize}
\paragraph{Lean-elaborated claim roles}\begin{ReviewVerbatim}
1. parameter: ℝ → ℝ
2. parameter: ℝ → ℝ
3. parameter: ℝ → ℝ
4. parameter: ℝ → ℝ
5. parameter: ℝ → ℝ
6. parameter: ℝ → ℝ
7. parameter: ℝ → ℝ
8. parameter: ℝ → ℝ
9. parameter: ℝ
10. parameter: ℝ
11. parameter: ℕ
12. parameter: ℕ
13. parameter: ℕ
14. parameter: ℕ → ℝ
15. parameter: ℕ → ℝ
16. parameter: ℕ → ℝ
17. assumption: LBG22StrategicRanking.RankingPrimitives cost production skill baseline
18. assumption: LBG22StrategicRanking.FiniteAdmissionPolicy n cutoff reward rho
19. assumption: LBG22StrategicRanking.FiniteAdmissionPolicy n cutoff newReward rho
20. assumption: DifferentiableOn ℝ cost (Set.Ioi 0)
21. assumption: DifferentiableOn ℝ production (Set.Ioi 0)
22. assumption: k < j
23. assumption: j ≤ n
24. assumption: reward k < newReward k
25. assumption: newReward j < reward j
26. assumption: ∀ i ≤ n, i ≠ k → i ≠ j → newReward i = reward i
27. assumption: Measurable tie
28. assumption: Set.InjOn tie (Set.Ioc 0 1)
29. assumption: LBG22StrategicRanking.RankingEquilibrium cost production skill tie n cutoff reward effort score
30. assumption: LBG22StrategicRanking.RankingEquilibrium cost production skill tie n cutoff newReward newEffort newScore
31. conclusion: ∀ᵐ (t : ℝ) ∂LBG22StrategicRanking.unitRankMeasure,
  have i : ℕ := ↑(LBG22StrategicRanking.finiteLowerRankBand (fun (l : Fin (n + 1)) => cutoff ↑l) t);
  (i < k → newEffort t = effort t) ∧ (i = k → effort t ≤ newEffort t) ∧ (k < i → i ≤ j → newEffort t ≤ effort t)
\end{ReviewVerbatim}

\paragraph{Lean proof endpoint (built separately)}\begin{ReviewVerbatim}
LBG22StrategicRanking.ProofInterface.effortComparativeStatics
\end{ReviewVerbatim}

\paragraph{Recorded source-to-Spec screening}
\textbf{Verdict:} \texttt{matches}
\\
{\raggedright
\textbf{Reason:} For policies differing only by an increase at k and decrease at j>k, the target states almost everywhere that lower bands i<k are unchanged, band k effort weakly increases, and bands k<i<=j weakly decrease.\allowbreak{} Its differentiability, feasible-\allowbreak{}policy, equilibrium, rank-\allowbreak{}band, and fixed-\allowbreak{}priority premises match Corollary A.\allowbreak{}1 and the approved continuum quantifier order.\allowbreak{}
\par}
\reviewmatch{Matches source input}
\noindent\textbf{Reviewer annotation}\par
\reviewerbox
\clearpage
\hypertarget{source-claim-ed834a4cf517d7fc}{}
\section*{Review row 12}
\textbf{Source locator:} {\footnotesize\raggedright cited publication, lines 1768-\allowbreak{}1780\par}
\paragraph{Verbatim source input}
\begin{ReviewVerbatim}
\begin{restatable}[Multi-dimensional rank preservation for linear $g$]{proposition}{propLinearG}
	Suppose $g$ is a
	linear function such that $g(e) = h x$, $h > 0$. %Also assume that $\f{i}$'s are strictly increasing.
	Suppose the school picks some $\alpha$ and $\lambda$. Define the combined pre-effort index as:
	\begin{align*}
		v_\pre^\alpha := \max_i \alpha_i\f{i}(\theta_\pre^i).
	\end{align*}
	Then in every equilibrium, for any two applicants with combined pre-effort indices $v_\pre^\alpha, \overline{v_\pre^\alpha}$ and  combined post-effort ranks $\theta_\post^\alpha, \overline{\theta_\post^\alpha}$ we have
	\begin{equation*}
	\lambda(\theta_\post^\alpha) >  \lambda(\overline{\theta_\post^\alpha}) \iff v_\pre^\alpha >  \overline{v_\pre^\alpha}.
	\end{equation*}
	\label{prop:linearg}
\end{restatable}

[Next verbatim source excerpt]

\parbold{Applicants}
There is a unit mass of \textit{applicants}, indexed by an observed index $\omega \in [0,1]$ distributed uniformly.\footnote{The index $\omega$ should be interpreted as each applicant's ``name,'' uncorrelated with skill, used solely for tie-breaking.}  Each applicant has a latent (unobserved) skill level represented by some measurable function of $\omega$.
We assume that {the distribution of the skills} has no atoms and that the CDF of this distribution is strictly increasing.
Using the CDF to map the skill of an applicant to a rank in~$[0,1]$, each applicant gets an (unobserved) rank $\theta_\pre =\theta_\pre(\omega)$ which by our assumption on the CDF is again uniformly distributed in $[0, 1]$ (the higher the rank the better).  With this setup, the skill of an applicant with rank $\theta_\pre$ can be written as $f(\theta_\pre)$ where $f$ is a strictly increasing, continuous function.  {It will be notationally convenient to label applicants by their rank $\theta_\pre$, though the reader should note that in contrast to $\omega$, $\theta_\pre(\omega)$ is assumed to be unobservable.}

%3. Each applicant exerts some effort e, at cost \cost(e).


Each applicant chooses an \textit{effort} level $e \ge 0$, the result of which is an observed, post-effort \textit{score}, {$v=v(e,\theta_\pre) = {g(e)}\cdot{f(\theta_\pre)}$}. \lledit{In other words, we assume the post-effort score is the product of two components, associated with the effort and the pre-effort skill respectively.}
{We assume} that the effort transfer function $g: [0, \infty) \mapsto [0, \infty)$ is a continuous, concave, strictly increasing function, representing that marginal effort improves one's score but has diminishing returns.
{%Since applicants can be labelled by their skill level $f(\theta_\pre)$, or equivalently, by their rank $\theta_{\pre}$,
{T}he strategies of the applicants can {then} be described by a function $\theta_\pre \mapsto e(\theta_\pre)$.
%\cb{$\omega\mapsto e(\omega)$.}%, leading to a distribution over post-effort scores $v=v(\theta_\pre))$ via the uniform distribution of $\theta_\pre$ in $[0,1]$.
}
  Each {applicant} is then ranked according to their score $v$%(where we break ties uniformly at random, if they appear)
, resulting in a \textit{post-effort rank} $\theta_\post$.
Note that the ranking $\theta_\post$ is slightly less trivial than a ranking of the skills, since the scores might have ties, which have to be resolved.

{\textbf{Tie Breaking.}
Given a choice of strategies $\theta_\pre \mapsto e(\theta_\pre)$, let $F$ be the CDF of the scores $v(e(\theta_\pre),\theta_\pre)$.
Since atoms for the distribution of $v$ would lead to ties for ranks defined as $F(v)$, we will use the labels of the applicants to break ties with the help of a (publicly announced) tie-breaking function $\Gamma(\omega)$, defined, e.g., via a collision free hash of the applicant names. Here  we require  that $\Gamma$  is a measurable function from $[0,1]$ to $[0,1]$ that maps different applicant labels to different values.
We then use $\Gamma$ to resolve the atoms of $F$, leading to a ranking function $v\mapsto \gamma(\omega, v)$ which is equal to $F(v)$ except when $v$ is an atom of the score distribution, {in which case it takes values in the ``gap interval'' $[F_-(v),  F(v)]$, where $F_-(v)$ is the left limit of $F$ at $v$.}  We  construct $\gamma$ in such a way that it
gives the uniform distribution for $\theta_\post(\omega)$ if we set  $\theta_\post(\omega)=\gamma(\omega,v(e(\theta_\pre(\omega)),\theta_\pre(\omega)))$.\footnote{See Remark~\ref{rem:gamma} for details on this construction.}

[Next verbatim source excerpt]

\parbold{Ranking designer (school)} A single \textit{school} is admitting applicants, based on their ranking. In particular, the school can choose a ranking reward function $\lambda : [0, 1] \mapsto [0, 1]$, such that an applicant with post-effort rank $\theta_\post$ is admitted with probability $\lambda(\theta_\post)$. We assume that $\lambda$ is non-decreasing and that the school has a constraint on the overall probability, such that in expectation it admits a number of applicants equal to {a} capacity constraint  $\rho \in (0,1)$, i.e., \lledit{$\E_{\theta_\post}[\lambda(\theta_\post)] = \rho $}.\footnote{\lledit{We use $\E_{\theta_\post}[\cdot]$ to denote an integral over $\theta_\post$ (and $\E[\cdot]$ to denote an integral over $\omega$)} with respect to the Lebesgue measure; informally, this can be thought of as averaging over the applicant population.} We may also refer to $\lambda$, informally, as the admission policy.



For simplicity, we further assume that $\lambda$ is a step-function with $K$ distinct levels $\ell_0 < \dots < \ell_{K-1}$, and $K-1$ cut-points parameterized by $c_1< \dots < c_{K-1}$ (with $c_0 = 0$, $c_K = 1$). %$\ell_0, \dots, \ell_{K-1}$.
In other words, we have $\lambda(\theta) = \ell_k$, for all $\theta \in \psi_k \triangleq [c_{k}, c_{k+1})$. Thus, applicants in the same post-rank interval $\theta \in \psi_k$ receive the same reward.

[Next verbatim source excerpt]

\parbold{Individual applicant welfare and equilibrium} Given the designer's function $\lambda$ and the effort levels of other applicants, each applicant chooses effort $e$ to maximize their individual welfare, % $W$,
\begin{equation*}
	\iwelf(e, \lambda(\theta_\post)) = \lambda(\theta_\post) - \cost(e),
\end{equation*}
where the effort cost function $\cost$ is non-negative, continuous, and strictly convex {on $[0,\infty)$, with $e_0:=\argmin_e \cost(e)$ and $p(e_0) = 0$.}
%\lnote{also want to assume that $\cost$ attains its minimum at $0$?}
%Let $e_0:=\argmin_e \cost(e)$ and $p(e_0) = 0$.

Applicants are assumed not to personally benefit from increasing their score $v$, except through the corresponding increase in their rank and reward. While the definition of $\cost$ does not preclude the applicant receiving any intrinsic benefit from exerting effort, we assume that the \emph{net} benefit from effort is non-positive.

[Next verbatim source excerpt]

\begin{definition} [Equilibrium]\label{def:equi}
{Given a}
%Let the
tie-breaking function $\Gamma(\omega)$
%be a pre-announced ordering over index $\omega$, to break ties in post-effort value. Then, let $\gamma(\omega, v)$ be the 1-1 ranking function defined as the CDF of post-effort values $v$ except where there are ties in $v$, in which case the ordering $\Gamma$ is used to produce unique ranks such that the range of $\gamma$ is uniform on $[0, 1]$.\footnote{More formally, note that each discontinuity in the CDF of post-effort values $v$ has a height equal to the measure of the applicants who are tied at that value, and so $\gamma$ for those applicants can be filled in using the CDF of $\Gamma(\omega)$ restricted to the applicants $\omega$ who are tied. Thus, the range of $\gamma$ is uniform on $[0,1]$.
%
{and}
%	Then, given
	a ranking probability function $\lambda$, an \textit{equilibrium} is a set of effort levels and post-effort ranking rewards for each applicant, $\{e(\theta_\pre),  \lambda(\theta_\post(\theta_\pre))\}$ such that, for all
	$\omega$, %$\theta_\pre$,
	\begin{align*}
		e(\theta_\pre(\omega)) &\in \argmax_{e} \iwelf\left(e, \lambda\left(\gamma\left(\omega, v\left(e,\theta_\pre(\omega)\right)\right)\right)\right) %\triangleq \argmax_{e} \left[ \lambda(\gamma(v(e,\theta_\pre)) - \cost(e) \right].
		\\
		\theta_\post(\theta_\pre(\omega)) &= \gamma(\omega, v(e(\theta_\pre(\omega)),\theta_\pre(\omega))),
	\end{align*}
	{where $\gamma(\omega,v)$ is the ranking induced by the CDF of the scores resulting from effort levels $\{e(\theta_\pre)\}$ and the tie-breaking function $\Gamma$.\footnote{See Remark~\ref{rem:gamma-effort} on $\gamma$ and the set of efforts.%Formally, $\gamma$ depends on the set of efforts. Given a fixed set and $\gamma$, for each applicant $\omega$ the first condition considers the counter-factual ranking of $\omega$ with different post-effort values but using the same ranking function. As defined, $\gamma$ yields a uniform distribution of ranks with such measure $0$ changes. In Appendix \Cref{lem:deviationsmeasure0}, we further prove that two effort sets equal up to sets of measure $0$ induce the same ranking function $\gamma$, and so the condition is consistent.%, and so we set it at the $\gamma$ induced by the effort levels $\{e(\theta_\pre)\}$.%, and so it does not vary.
	}
	}
	%In particular, we denote $\theta_\post(\theta_\pre) := \gamma(v(e(\theta_\pre),\theta_\pre))$.
	%\lledit{and for any $\theta_\pre, \theta'_\pre$, we have $\lambda(\theta_\post(\theta_\pre)) > \lambda(\theta_\post(\theta'_\pre)) \implies v(e(\theta_\pre),\theta_\pre) > v(e(\theta_\pre'),\theta_\pre')$.}
\end{definition}

[Next verbatim source excerpt]

\parbold{Model.} The model is similar to our base model \Cref{sec:model}; each applicant now has $m$ latent skill levels with respective ranks $\theta_\pre^i \in [0,1]$ for $i\in \{1, \cdots, m\}$. Each rank is drawn independently from the uniform distribution over $[0,1]$, and the skill $i$ of applicant with rank $\theta_\pre^i$ is $\f{i}(\theta_\pre^i)$. Each applicant now chooses effort levels $\{ \e{i} \}$, at cost $p^m( \e{1}, \cdots, \e{m})$, resulting in post-effort scores $\{\score{i}\}$,
\begin{equation*}
\score{i} = g(\e{i})\cdot \f{i}(\theta_\pre^i),
\end{equation*}
where $g$ is concave, increasing, as before, and $\f{i}$ is a continuous, strictly increasing quantile function (${\f{i}}^{-1}$ is the CDF function of the scores on dimension $i$). As before, the school observes the post-effort scores for each applicant, now for each dimension $i$, and designs a non-decreasing function $\lambda : [0,1] \to [0,1]$ denoting the admissions probability $\lambda(\theta_\post)$ for applicant with post-effort rank $\theta_\post$.

%Suppose the school observes $v_i$.

How does the school construct post-effort rank $\theta_\post$? In general, each applicant's \textit{combined} post-effort score may be any function of the scores on each dimension, $\{\score{i}\}$. Here, we assume the following linear score function. The school announces weights $\alpha = (\alpha_1, \cdots, \alpha_m) \in \Delta_{m-1}$ (denoting the $m$-dimensional simplex); the weights $\alpha$ represent the relative emphasis placed on each skill for the admissions decision. Then, each applicant is ranked according to their combined score, $\combscore := \sum_{i=1}^m \alpha_i\score{i}$, resulting in their combined post-effort rank $\theta_\post^\alpha$, and is admitted with probability $\lambda(\theta_\post^\alpha)$.\footnote{The linear combination of skill is similar to the ``linear mechanism'' in \citet{kleinberg18investeffort}, which studied how reward design incentivizes strategic agents to exert effort on different skill dimensions. Compared to their work, where efforts are connected to skills via an effort graph, we consider a simplified setting where each effort maps to one skill, and study how reward design affects equilibrium rankings, in presence of competition.}

%\lnote{\textbf{weighted score or weighted rank?} alternatively, the school combines the post-effort ranks for each skill. i.e. an applicant with \emph{weighted} post-effort rank $\theta_\post^\alpha := \sum_{i=1}^m \alpha_i\theta_\post^i$ is admitted with probability $\lambda(\theta_\post^\alpha)$. This case seems interesting as well. Not clear in which case it's possible to reduce to the single-dimensional case?}

Putting things together, the applicant's individual welfare is:
\begin{equation*}
\iwelf(\{ \e{i} \}_{i=1}^m, \lambda(\theta_\post^\alpha)) = \lambda(\theta_\post^\alpha) - p^m( \e{1}, \cdots, \e{m})
\end{equation*}

[Next verbatim source excerpt]

We show this result by characterizing the equilibrium of the following simplified setting, with further assumptions on the effort cost function $\cost$ and the effort transfer function $g$; the cost function $\cost$ is assumed to be a function of the sum of the efforts exerted:
\[p(e^1, \cdots, e^m) \triangleq p\left(\sum_{i=1}^m \e{i}\right),\] where $p$ is convex and increasing.
This assumption says that the effort exerted for any skill is entirely exchangeable, for example, two hours spent on studying math is as costly as two hours spent on studying chemistry.
%
Effort transfer function $g$ is assumed to be linear, that is, there are constant returns to effort.
\end{ReviewVerbatim}

% report-context-presentation-sha256: 16265e6401473cf3c85b47c3a2f48907c1f14d36cd38cc80fc86055d51df0113
\paragraph{Settled source reading}
\begin{sloppypar}
Equilibrium means that almost every applicant best responds against every feasible deviation;\allowbreak{} uniqueness is up to a null set.\allowbreak{}
\end{sloppypar}
\paragraph{Settled source reading}
\begin{sloppypar}
Applicants know the fixed priority used at score-\allowbreak{}boundary ties.\allowbreak{} This clarifies the stated model, not a corrected admission policy or a new economic assumption.\allowbreak{}
\end{sloppypar}
\paragraph{Settled source reading}
\begin{sloppypar}
The main-\allowbreak{}text group comparisons are exact under clarified weak/\allowbreak{}strict interpretation.\allowbreak{} Strictness requires positive disadvantaged effort;\allowbreak{} classical derivative existence is qualified separately.\allowbreak{}
\end{sloppypar}
\paragraph{Formalized review target}
The archival source above is not asserted equivalent to this formalized target.
\begin{ReviewVerbatim}
Under the source independent continuous skill ranks, positive coordinate weights, exchangeable strictly convex total-effort cost, linear production, fixed known priority order, and any finite-level policy, every equilibrium over all nonnegative effort vectors preserves reward bands almost everywhere: the reward at actual post-effort rank equals the reward at the population CDF rank of X=max_i alpha_i f_i(t_i). Almost every applicant has a coordinate attaining that maximum and allocates all total effort to it. There is no hard budget and no assertion that distinct indices must receive strictly different rewards within one constant reward band.
\end{ReviewVerbatim}

\textbf{Archival source anchor:} {\footnotesize\raggedright cited publication, lines 1760-\allowbreak{}1780\par}
\paragraph{Expanded PaperInterface specification}
\begin{ReviewVerbatim}
∀ (m n : ℕ) (cost : ℝ → ℝ) (h rho : ℝ) (weight : Fin (m + 1) → ℝ) (skill : Fin (m + 1) → ℝ → ℝ) (cutoff reward : ℕ → ℝ)
  (effort : (Fin (m + 1) → ℝ) → Fin (m + 1) → ℝ) (score tie : (Fin (m + 1) → ℝ) → ℝ),
  LBG22StrategicRanking.MultidimensionalPrimitives m cost h weight skill →
    LBG22StrategicRanking.FiniteAdmissionPolicy n cutoff reward rho →
      Measurable tie →
        Set.InjOn tie {x : Fin (m + 1) → ℝ | ∀ (i : Fin (m + 1)), x i ∈ Set.Ioc 0 1} →
          LBG22StrategicRanking.MultidimensionalEquilibrium m n cost h weight skill cutoff reward effort score tie →
            ∀ᵐ (x : Fin (m + 1) → ℝ) ∂LBG22StrategicRanking.sourceMultidimensionalPopulation (Fin (m + 1)),
              reward
                    ↑(LBG22StrategicRanking.finiteLowerRankBand (fun (i : Fin (n + 1)) => cutoff ↑i)
                        (LBG22StrategicRanking.tieBrokenRank
                          (LBG22StrategicRanking.sourceMultidimensionalPopulation (Fin (m + 1))) score tie x)) =
                  reward
                    ↑(LBG22StrategicRanking.finiteLowerRankBand (fun (i : Fin (n + 1)) => cutoff ↑i)
                        (LBG22StrategicRanking.sourceMultidimensionalPreRank weight skill x)) ∧
                ∃ (best : Fin (m + 1)),
                  LBG22StrategicRanking.sourceMultidimensionalCoefficient weight skill x best =
                      LBG22StrategicRanking.sourceCombinedSkill
                        (LBG22StrategicRanking.sourceMultidimensionalCoefficient weight skill x) ∧
                    ∀ (i : Fin (m + 1)), effort x i = if i = best then ∑ j : Fin (m + 1), effort x j else 0
\end{ReviewVerbatim}

\paragraph{Retained governing declarations}
\begin{itemize}\raggedright
\item \hyperlink{paper-prerequisite-d056dd7c8060653c}{LBG22\allowbreak{}Strategic\allowbreak{}Ranking.\allowbreak{}Finite\allowbreak{}Admission\allowbreak{}Policy}
\item \hyperlink{paper-prerequisite-b201854e4e4dd891}{LBG22\allowbreak{}Strategic\allowbreak{}Ranking.\allowbreak{}Multidimensional\allowbreak{}Equilibrium}
\item \hyperlink{paper-prerequisite-8e74a377842bef95}{LBG22\allowbreak{}Strategic\allowbreak{}Ranking.\allowbreak{}Multidimensional\allowbreak{}Primitives}
\item \hyperlink{paper-prerequisite-a42f440ddc254fa7}{LBG22\allowbreak{}Strategic\allowbreak{}Ranking.\allowbreak{}finite\allowbreak{}Lower\allowbreak{}Rank\allowbreak{}Band}
\item \hyperlink{governing-declaration-e17c2a81c5471276}{LBG22\allowbreak{}Strategic\allowbreak{}Ranking.\allowbreak{}source\allowbreak{}Combined\allowbreak{}Skill}
\item \hyperlink{paper-prerequisite-e57e55071c75ce91}{LBG22\allowbreak{}Strategic\allowbreak{}Ranking.\allowbreak{}source\allowbreak{}Multidimensional\allowbreak{}Coefficient}
\item \hyperlink{paper-prerequisite-16d135799a0f0491}{LBG22\allowbreak{}Strategic\allowbreak{}Ranking.\allowbreak{}source\allowbreak{}Multidimensional\allowbreak{}Population}
\item \hyperlink{governing-declaration-9631d67a19f5df1f}{LBG22\allowbreak{}Strategic\allowbreak{}Ranking.\allowbreak{}source\allowbreak{}Multidimensional\allowbreak{}Pre\allowbreak{}Rank}
\item \hyperlink{paper-prerequisite-b7c71e53177b76c6}{LBG22\allowbreak{}Strategic\allowbreak{}Ranking.\allowbreak{}tie\allowbreak{}Broken\allowbreak{}Rank}
\end{itemize}
\paragraph{Lean-elaborated claim roles}\begin{ReviewVerbatim}
1. parameter: ℕ
2. parameter: ℕ
3. parameter: ℝ → ℝ
4. parameter: ℝ
5. parameter: ℝ
6. parameter: Fin (m + 1) → ℝ
7. parameter: Fin (m + 1) → ℝ → ℝ
8. parameter: ℕ → ℝ
9. parameter: ℕ → ℝ
10. parameter: (Fin (m + 1) → ℝ) → Fin (m + 1) → ℝ
11. parameter: (Fin (m + 1) → ℝ) → ℝ
12. parameter: (Fin (m + 1) → ℝ) → ℝ
13. assumption: LBG22StrategicRanking.MultidimensionalPrimitives m cost h weight skill
14. assumption: LBG22StrategicRanking.FiniteAdmissionPolicy n cutoff reward rho
15. assumption: Measurable tie
16. assumption: Set.InjOn tie {x : Fin (m + 1) → ℝ | ∀ (i : Fin (m + 1)), x i ∈ Set.Ioc 0 1}
17. assumption: LBG22StrategicRanking.MultidimensionalEquilibrium m n cost h weight skill cutoff reward effort score tie
18. conclusion: ∀ᵐ (x : Fin (m + 1) → ℝ) ∂LBG22StrategicRanking.sourceMultidimensionalPopulation (Fin (m + 1)),
  reward
        ↑(LBG22StrategicRanking.finiteLowerRankBand (fun (i : Fin (n + 1)) => cutoff ↑i)
            (LBG22StrategicRanking.tieBrokenRank (LBG22StrategicRanking.sourceMultidimensionalPopulation (Fin (m + 1)))
              score tie x)) =
      reward
        ↑(LBG22StrategicRanking.finiteLowerRankBand (fun (i : Fin (n + 1)) => cutoff ↑i)
            (LBG22StrategicRanking.sourceMultidimensionalPreRank weight skill x)) ∧
    ∃ (best : Fin (m + 1)),
      LBG22StrategicRanking.sourceMultidimensionalCoefficient weight skill x best =
          LBG22StrategicRanking.sourceCombinedSkill
            (LBG22StrategicRanking.sourceMultidimensionalCoefficient weight skill x) ∧
        ∀ (i : Fin (m + 1)), effort x i = if i = best then ∑ j : Fin (m + 1), effort x j else 0
\end{ReviewVerbatim}

\paragraph{Lean proof endpoint (built separately)}\begin{ReviewVerbatim}
LBG22StrategicRanking.ProofInterface.multidimensionalRank
\end{ReviewVerbatim}

\paragraph{Recorded source-to-Spec screening}
\textbf{Verdict:} \texttt{matches\_approved\_corrected\_target}
\\
{\raggedright
\textbf{Reason:} The selected route expressly disclaims equivalence to the archival strict-\allowbreak{}index statement and binds defect LBG22-\allowbreak{}C4.\allowbreak{} The expanded target matches the complete correction:\allowbreak{} under the independent multidimensional population, simplex weights, exchangeable strictly convex total-\allowbreak{}effort cost, linear production, finite policy, and fixed priority, actual reward equals the reward at the CDF rank of the maximum weighted skill coefficient almost everywhere, and all total effort is placed on a coordinate attaining that maximum.\allowbreak{} Nonnegative rather than strictly positive weights is a harmless broader premise surface that includes every approved positive-\allowbreak{}weight instance.\allowbreak{}
\par}
\reviewmatch{Matches formalized target}
\noindent\textbf{Reviewer annotation}\par
\reviewerbox
\clearpage
\hypertarget{source-claim-ef5c828d0d65c93b}{}
\section*{Review row 13}
\textbf{Source locator:} {\footnotesize\raggedright cited publication, lines 1826-\allowbreak{}1828\par}
\paragraph{Verbatim source input}
\begin{ReviewVerbatim}
\begin{restatable}[The school's weighted private utility is maximized by some randomization]{proposition}{propWeightedUtil}\label{prop:2levels-fixedcap}
	Consider the class of two-level policies. For any $c \in (0,1-\rho)$, there exist some $\beta \in (0,1)$ such that the school's utility $\Util\pri_\alpha$ is maximized at that value of $c$.
\end{restatable}

[Next verbatim source excerpt]

\parbold{Applicants}
There is a unit mass of \textit{applicants}, indexed by an observed index $\omega \in [0,1]$ distributed uniformly.\footnote{The index $\omega$ should be interpreted as each applicant's ``name,'' uncorrelated with skill, used solely for tie-breaking.}  Each applicant has a latent (unobserved) skill level represented by some measurable function of $\omega$.
We assume that {the distribution of the skills} has no atoms and that the CDF of this distribution is strictly increasing.
Using the CDF to map the skill of an applicant to a rank in~$[0,1]$, each applicant gets an (unobserved) rank $\theta_\pre =\theta_\pre(\omega)$ which by our assumption on the CDF is again uniformly distributed in $[0, 1]$ (the higher the rank the better).  With this setup, the skill of an applicant with rank $\theta_\pre$ can be written as $f(\theta_\pre)$ where $f$ is a strictly increasing, continuous function.  {It will be notationally convenient to label applicants by their rank $\theta_\pre$, though the reader should note that in contrast to $\omega$, $\theta_\pre(\omega)$ is assumed to be unobservable.}

%3. Each applicant exerts some effort e, at cost \cost(e).


Each applicant chooses an \textit{effort} level $e \ge 0$, the result of which is an observed, post-effort \textit{score}, {$v=v(e,\theta_\pre) = {g(e)}\cdot{f(\theta_\pre)}$}. \lledit{In other words, we assume the post-effort score is the product of two components, associated with the effort and the pre-effort skill respectively.}
{We assume} that the effort transfer function $g: [0, \infty) \mapsto [0, \infty)$ is a continuous, concave, strictly increasing function, representing that marginal effort improves one's score but has diminishing returns.
{%Since applicants can be labelled by their skill level $f(\theta_\pre)$, or equivalently, by their rank $\theta_{\pre}$,
{T}he strategies of the applicants can {then} be described by a function $\theta_\pre \mapsto e(\theta_\pre)$.
%\cb{$\omega\mapsto e(\omega)$.}%, leading to a distribution over post-effort scores $v=v(\theta_\pre))$ via the uniform distribution of $\theta_\pre$ in $[0,1]$.
}
  Each {applicant} is then ranked according to their score $v$%(where we break ties uniformly at random, if they appear)
, resulting in a \textit{post-effort rank} $\theta_\post$.
Note that the ranking $\theta_\post$ is slightly less trivial than a ranking of the skills, since the scores might have ties, which have to be resolved.

{\textbf{Tie Breaking.}
Given a choice of strategies $\theta_\pre \mapsto e(\theta_\pre)$, let $F$ be the CDF of the scores $v(e(\theta_\pre),\theta_\pre)$.
Since atoms for the distribution of $v$ would lead to ties for ranks defined as $F(v)$, we will use the labels of the applicants to break ties with the help of a (publicly announced) tie-breaking function $\Gamma(\omega)$, defined, e.g., via a collision free hash of the applicant names. Here  we require  that $\Gamma$  is a measurable function from $[0,1]$ to $[0,1]$ that maps different applicant labels to different values.
We then use $\Gamma$ to resolve the atoms of $F$, leading to a ranking function $v\mapsto \gamma(\omega, v)$ which is equal to $F(v)$ except when $v$ is an atom of the score distribution, {in which case it takes values in the ``gap interval'' $[F_-(v),  F(v)]$, where $F_-(v)$ is the left limit of $F$ at $v$.}  We  construct $\gamma$ in such a way that it
gives the uniform distribution for $\theta_\post(\omega)$ if we set  $\theta_\post(\omega)=\gamma(\omega,v(e(\theta_\pre(\omega)),\theta_\pre(\omega)))$.\footnote{See Remark~\ref{rem:gamma} for details on this construction.}

[Next verbatim source excerpt]

\parbold{Ranking designer (school)} A single \textit{school} is admitting applicants, based on their ranking. In particular, the school can choose a ranking reward function $\lambda : [0, 1] \mapsto [0, 1]$, such that an applicant with post-effort rank $\theta_\post$ is admitted with probability $\lambda(\theta_\post)$. We assume that $\lambda$ is non-decreasing and that the school has a constraint on the overall probability, such that in expectation it admits a number of applicants equal to {a} capacity constraint  $\rho \in (0,1)$, i.e., \lledit{$\E_{\theta_\post}[\lambda(\theta_\post)] = \rho $}.\footnote{\lledit{We use $\E_{\theta_\post}[\cdot]$ to denote an integral over $\theta_\post$ (and $\E[\cdot]$ to denote an integral over $\omega$)} with respect to the Lebesgue measure; informally, this can be thought of as averaging over the applicant population.} We may also refer to $\lambda$, informally, as the admission policy.



For simplicity, we further assume that $\lambda$ is a step-function with $K$ distinct levels $\ell_0 < \dots < \ell_{K-1}$, and $K-1$ cut-points parameterized by $c_1< \dots < c_{K-1}$ (with $c_0 = 0$, $c_K = 1$). %$\ell_0, \dots, \ell_{K-1}$.
In other words, we have $\lambda(\theta) = \ell_k$, for all $\theta \in \psi_k \triangleq [c_{k}, c_{k+1})$. Thus, applicants in the same post-rank interval $\theta \in \psi_k$ receive the same reward.

[Next verbatim source excerpt]

\parbold{Individual applicant welfare and equilibrium} Given the designer's function $\lambda$ and the effort levels of other applicants, each applicant chooses effort $e$ to maximize their individual welfare, % $W$,
\begin{equation*}
	\iwelf(e, \lambda(\theta_\post)) = \lambda(\theta_\post) - \cost(e),
\end{equation*}
where the effort cost function $\cost$ is non-negative, continuous, and strictly convex {on $[0,\infty)$, with $e_0:=\argmin_e \cost(e)$ and $p(e_0) = 0$.}
%\lnote{also want to assume that $\cost$ attains its minimum at $0$?}
%Let $e_0:=\argmin_e \cost(e)$ and $p(e_0) = 0$.

Applicants are assumed not to personally benefit from increasing their score $v$, except through the corresponding increase in their rank and reward. While the definition of $\cost$ does not preclude the applicant receiving any intrinsic benefit from exerting effort, we assume that the \emph{net} benefit from effort is non-positive.

[Next verbatim source excerpt]

\begin{definition} [Equilibrium]\label{def:equi}
{Given a}
%Let the
tie-breaking function $\Gamma(\omega)$
%be a pre-announced ordering over index $\omega$, to break ties in post-effort value. Then, let $\gamma(\omega, v)$ be the 1-1 ranking function defined as the CDF of post-effort values $v$ except where there are ties in $v$, in which case the ordering $\Gamma$ is used to produce unique ranks such that the range of $\gamma$ is uniform on $[0, 1]$.\footnote{More formally, note that each discontinuity in the CDF of post-effort values $v$ has a height equal to the measure of the applicants who are tied at that value, and so $\gamma$ for those applicants can be filled in using the CDF of $\Gamma(\omega)$ restricted to the applicants $\omega$ who are tied. Thus, the range of $\gamma$ is uniform on $[0,1]$.
%
{and}
%	Then, given
	a ranking probability function $\lambda$, an \textit{equilibrium} is a set of effort levels and post-effort ranking rewards for each applicant, $\{e(\theta_\pre),  \lambda(\theta_\post(\theta_\pre))\}$ such that, for all
	$\omega$, %$\theta_\pre$,
	\begin{align*}
		e(\theta_\pre(\omega)) &\in \argmax_{e} \iwelf\left(e, \lambda\left(\gamma\left(\omega, v\left(e,\theta_\pre(\omega)\right)\right)\right)\right) %\triangleq \argmax_{e} \left[ \lambda(\gamma(v(e,\theta_\pre)) - \cost(e) \right].
		\\
		\theta_\post(\theta_\pre(\omega)) &= \gamma(\omega, v(e(\theta_\pre(\omega)),\theta_\pre(\omega))),
	\end{align*}
	{where $\gamma(\omega,v)$ is the ranking induced by the CDF of the scores resulting from effort levels $\{e(\theta_\pre)\}$ and the tie-breaking function $\Gamma$.\footnote{See Remark~\ref{rem:gamma-effort} on $\gamma$ and the set of efforts.%Formally, $\gamma$ depends on the set of efforts. Given a fixed set and $\gamma$, for each applicant $\omega$ the first condition considers the counter-factual ranking of $\omega$ with different post-effort values but using the same ranking function. As defined, $\gamma$ yields a uniform distribution of ranks with such measure $0$ changes. In Appendix \Cref{lem:deviationsmeasure0}, we further prove that two effort sets equal up to sets of measure $0$ induce the same ranking function $\gamma$, and so the condition is consistent.%, and so we set it at the $\gamma$ induced by the effort levels $\{e(\theta_\pre)\}$.%, and so it does not vary.
	}
	}
	%In particular, we denote $\theta_\post(\theta_\pre) := \gamma(v(e(\theta_\pre),\theta_\pre))$.
	%\lledit{and for any $\theta_\pre, \theta'_\pre$, we have $\lambda(\theta_\post(\theta_\pre)) > \lambda(\theta_\post(\theta'_\pre)) \implies v(e(\theta_\pre),\theta_\pre) > v(e(\theta_\pre'),\theta_\pre')$.}
\end{definition}

[Next verbatim source excerpt]

\parbold{Model.} The model is similar to our base model \Cref{sec:model}; each applicant now has $m$ latent skill levels with respective ranks $\theta_\pre^i \in [0,1]$ for $i\in \{1, \cdots, m\}$. Each rank is drawn independently from the uniform distribution over $[0,1]$, and the skill $i$ of applicant with rank $\theta_\pre^i$ is $\f{i}(\theta_\pre^i)$. Each applicant now chooses effort levels $\{ \e{i} \}$, at cost $p^m( \e{1}, \cdots, \e{m})$, resulting in post-effort scores $\{\score{i}\}$,
\begin{equation*}
\score{i} = g(\e{i})\cdot \f{i}(\theta_\pre^i),
\end{equation*}
where $g$ is concave, increasing, as before, and $\f{i}$ is a continuous, strictly increasing quantile function (${\f{i}}^{-1}$ is the CDF function of the scores on dimension $i$). As before, the school observes the post-effort scores for each applicant, now for each dimension $i$, and designs a non-decreasing function $\lambda : [0,1] \to [0,1]$ denoting the admissions probability $\lambda(\theta_\post)$ for applicant with post-effort rank $\theta_\post$.

%Suppose the school observes $v_i$.

How does the school construct post-effort rank $\theta_\post$? In general, each applicant's \textit{combined} post-effort score may be any function of the scores on each dimension, $\{\score{i}\}$. Here, we assume the following linear score function. The school announces weights $\alpha = (\alpha_1, \cdots, \alpha_m) \in \Delta_{m-1}$ (denoting the $m$-dimensional simplex); the weights $\alpha$ represent the relative emphasis placed on each skill for the admissions decision. Then, each applicant is ranked according to their combined score, $\combscore := \sum_{i=1}^m \alpha_i\score{i}$, resulting in their combined post-effort rank $\theta_\post^\alpha$, and is admitted with probability $\lambda(\theta_\post^\alpha)$.\footnote{The linear combination of skill is similar to the ``linear mechanism'' in \citet{kleinberg18investeffort}, which studied how reward design incentivizes strategic agents to exert effort on different skill dimensions. Compared to their work, where efforts are connected to skills via an effort graph, we consider a simplified setting where each effort maps to one skill, and study how reward design affects equilibrium rankings, in presence of competition.}

%\lnote{\textbf{weighted score or weighted rank?} alternatively, the school combines the post-effort ranks for each skill. i.e. an applicant with \emph{weighted} post-effort rank $\theta_\post^\alpha := \sum_{i=1}^m \alpha_i\theta_\post^i$ is admitted with probability $\lambda(\theta_\post^\alpha)$. This case seems interesting as well. Not clear in which case it's possible to reduce to the single-dimensional case?}

Putting things together, the applicant's individual welfare is:
\begin{equation*}
\iwelf(\{ \e{i} \}_{i=1}^m, \lambda(\theta_\post^\alpha)) = \lambda(\theta_\post^\alpha) - p^m( \e{1}, \cdots, \e{m})
\end{equation*}

\subsection{Equilibrium under multi-dimensional competition}\label{sec:multi:linear}


[Next verbatim source excerpt]

In our simplified setting, there are two skills $M$ and $U$: $M$ has a measurable score $v^M$ and $U$ has an unmeasurable score $v^U$. For example, $M$ could be scholastic achievement as measured by SAT scores, and $U$ could be ``creativity", a personal quality that is valued by the school but is not directly measurable.
%
Since the school cannot observe $v^U$, its admission policy is based on $\theta_\post^M$ only, that is, $\alpha_U = 0$ and $\alpha_M = 1$.

We further assume that each applicant has a fixed effort budget of $B > 0$, and is intrinsically motivated to exert effort in the unmeasurable skill $U$; in fact, they will always exert effort $e^U = B - e^M$. Formally this corresponds to the effort cost function \[\cost^m(e_M, e_U) = \cost(e^M) - (\max(0,B-(e^M+e^U)))^2\] where $\cost$ is convex and increasing. We can write the applicant's individual welfare as:
\begin{equation*}
W(e^M, e^U, \lambda(\theta_\post^M)) = \lambda(\theta_\post^M) - \cost^m(e_M, e_U).
\end{equation*}


The school's private utility is now weighted by $\beta \in (0,1)$, which quantifies the relative value the school places on the measurable skill over the unmeasurable skill.
\begin{equation*}
\Util\pri_\beta = \E[\beta\cdot v^M + (1-\beta)\cdot v^U\mid Z=1].
\end{equation*}
The smaller that $\beta$ is, the more the school places value on the unmeasurable skill $v^U$.

[Next verbatim source excerpt]

\begin{definition}[Two-level policy]\label{assump:2-level-fixed-cap} In our baseline two-level function parameterized by cut-off $c \in (0,1-\rho]$, each applicant with post-effort rank  $\theta_\post \geq c$ is admitted with probability $\ell_1 =\frac{\rho}{1-c} \in (0,1]$. Others are rejected, $\ell_0 = 0$.
\end{definition}

Note that standard non-randomized admissions policies are equivalent to case where $c = 1 - \rho$: the highest scoring applicants up to the capacity constraint are accepted with probability 1 and all others are rejected. We call this case \textbf{non-randomized admissions}. The other extreme is the one-level policy, \textbf{pure randomization}, where each applicant is admitted with probability ${\rho}$, i.e., $\ell_0 =\rho$. Decreasing the cut-off $c$ can be viewed as increasing the level of randomization in the admissions policy.
\end{ReviewVerbatim}

% report-context-presentation-sha256: ce5679112671da8414975164981e5740bcad9608aa7765c30f1069706683e7b1
\paragraph{Settled source reading}
\begin{sloppypar}
Equilibrium means that almost every applicant best responds against every feasible deviation;\allowbreak{} uniqueness is up to a null set.\allowbreak{}
\end{sloppypar}
\paragraph{Settled source reading}
\begin{sloppypar}
The multitask model has nonnegative effort on both tasks summing to the stated budget.\allowbreak{} This is distinct from the unrestricted total effort in the multidimensional ranking model.\allowbreak{}
\end{sloppypar}
\paragraph{Settled source reading}
\begin{sloppypar}
Applicants know the fixed priority used at score-\allowbreak{}boundary ties.\allowbreak{} This clarifies the stated model, not a corrected admission policy or a new economic assumption.\allowbreak{}
\end{sloppypar}
\paragraph{Additional assumptions}
\begin{sloppypar}
Only for Proposition B.2: measurable skill is affine a+b*t with a>=0 and b>0; g(0)=0; the hard budget B is at least E=p\textasciicircum{}\{-1\}(1).; Only for Proposition B.2: g squared is concave on (0,B) and x/(p(g\textasciicircum{}\{-1\}(x))) is convex on [g(p\textasciicircum{}\{-1\}(rho)),g(E)].; Only for Proposition B.2: cost is twice differentiable on (0,E), production is twice differentiable on (0,B), and its second derivative is continuous there. Source continuity and shape requirements remain. The population retains independent regular unmeasurable skill with positive finite mean. No frontier or optimizer property is assumed.
\end{sloppypar}
\paragraph{Formalized review target}
The archival source above is not asserted equivalent to this formalized target.
\begin{ReviewVerbatim}
Under the hard-budget feasible set eM>=0, eU>=0, eM+eU=B, source scalar cost/production primitives with baseline zero, capacity 0<rho<1, affine measurable skill fM(t)=a+b*t with a>=0 and b>0, and independent regular nonnegative strictly increasing continuous unmeasurable skill, let E>0 satisfy p(E)=1. Suppose B>=E, g(0)=0, g squared is concave on (0,B), and x/C(x) is convex on [g(p^{-1}(rho)),g(E)], where C=p composed with the inverse of g. Require p twice differentiable on (0,E), g twice differentiable on (0,B), and g'' continuous there, retaining the source continuity at endpoints. For each c in (0,1-rho), there exists beta in (0,1), chosen before equilibrium selection, such that for every equilibrium family with these fixed primitives, independent population and known tie order, beta*M(d)+(1-beta)*U(d)<=beta*M(c)+(1-beta)*U(c) for all d in (0,1-rho]. M and U are actual conditional admitted-population means, including the independent unmeasurable skill mean. These are local sufficient primitive restrictions, not an assumption of a supporting line or optimality, and do not modify unrelated main-text results.
\end{ReviewVerbatim}

\textbf{Archival source anchor:} {\footnotesize\raggedright cited publication, lines 1826-\allowbreak{}1828\par}
\paragraph{Expanded PaperInterface specification}
\begin{ReviewVerbatim}
∀ (cost production skillU : ℝ → ℝ) (B E rho lower width : ℝ) (tie : ℝ × ℝ → ℝ),
  LBG22StrategicRanking.RankingPrimitives cost production (fun (t : ℝ) => lower + width * t) 0 →
    rho ∈ Set.Ioo 0 1 →
      0 ≤ lower →
        0 < width →
          0 < E →
            cost E = 1 →
              E ≤ B →
                ContinuousOn skillU (Set.Icc 0 1) →
                  StrictMonoOn skillU (Set.Icc 0 1) →
                    0 ≤ skillU 0 →
                      production 0 = 0 →
                        (ConcaveOn ℝ (Set.Ioo 0 B) fun (e : ℝ) => production e ^ 2) →
                          (ConvexOn ℝ
                              (Set.Icc (production (LBG22StrategicRanking.effortIntervalInverse cost E rho))
                                (production E))
                              fun (x : ℝ) => x / cost (LBG22StrategicRanking.effortIntervalInverse production E x)) →
                            DifferentiableOn ℝ cost (Set.Ioo 0 E) →
                              DifferentiableOn ℝ (deriv cost) (Set.Ioo 0 E) →
                                DifferentiableOn ℝ production (Set.Ioo 0 B) →
                                  DifferentiableOn ℝ (deriv production) (Set.Ioo 0 B) →
                                    ContinuousOn (deriv (deriv production)) (Set.Ioo 0 B) →
                                      Measurable tie →
                                        Set.InjOn tie {x : ℝ × ℝ | x.1 ∈ Set.Ioc 0 1 ∧ x.2 ∈ Set.Ioc 0 1} →
                                          ∀ c ∈ Set.Ioo 0 (1 - rho),
                                            ∃ beta ∈ Set.Ioo 0 1,
                                              ∀ (effortM effortU score : ℝ → ℝ × ℝ → ℝ),
                                                (∀ d ∈ Set.Ioc 0 (1 - rho),
                                                    LBG22StrategicRanking.MultitaskEquilibrium cost production
                                                      (fun (t : ℝ) => lower + width * t) B rho d (effortM d) (effortU d)
                                                      (score d) tie) →
                                                  ∀ d ∈ Set.Ioc 0 (1 - rho),
                                                    LBG22StrategicRanking.sourceActualMultitaskSchoolUtility
                                                        LBG22StrategicRanking.unitRankMeasure production
                                                        (fun (t : ℝ) => lower + width * t) skillU (effortM d)
                                                        (effortU d) (score d) tie rho beta d ≤
                                                      LBG22StrategicRanking.sourceActualMultitaskSchoolUtility
                                                        LBG22StrategicRanking.unitRankMeasure production
                                                        (fun (t : ℝ) => lower + width * t) skillU (effortM c)
                                                        (effortU c) (score c) tie rho beta c
\end{ReviewVerbatim}

\paragraph{Retained governing declarations}
\begin{itemize}\raggedright
\item \hyperlink{paper-prerequisite-b504d72e1e0de867}{LBG22\allowbreak{}Strategic\allowbreak{}Ranking.\allowbreak{}Multitask\allowbreak{}Equilibrium}
\item \hyperlink{paper-prerequisite-9bbcc12f4a1736bf}{LBG22\allowbreak{}Strategic\allowbreak{}Ranking.\allowbreak{}Ranking\allowbreak{}Primitives}
\item \hyperlink{governing-declaration-c5b2c07abbff7224}{LBG22\allowbreak{}Strategic\allowbreak{}Ranking.\allowbreak{}effort\allowbreak{}Interval\allowbreak{}Inverse}
\item \hyperlink{paper-prerequisite-c194f1bd51c4a4cc}{LBG22\allowbreak{}Strategic\allowbreak{}Ranking.\allowbreak{}source\allowbreak{}Actual\allowbreak{}Multitask\allowbreak{}School\allowbreak{}Utility}
\item \hyperlink{governing-declaration-678cf48d549cb620}{LBG22\allowbreak{}Strategic\allowbreak{}Ranking.\allowbreak{}unit\allowbreak{}Rank\allowbreak{}Measure}
\end{itemize}
\paragraph{Lean-elaborated claim roles}\begin{ReviewVerbatim}
1. parameter: ℝ → ℝ
2. parameter: ℝ → ℝ
3. parameter: ℝ → ℝ
4. parameter: ℝ
5. parameter: ℝ
6. parameter: ℝ
7. parameter: ℝ
8. parameter: ℝ
9. parameter: ℝ × ℝ → ℝ
10. assumption: LBG22StrategicRanking.RankingPrimitives cost production (fun (t : ℝ) => lower + width * t) 0
11. assumption: rho ∈ Set.Ioo 0 1
12. assumption: 0 ≤ lower
13. assumption: 0 < width
14. assumption: 0 < E
15. assumption: cost E = 1
16. assumption: E ≤ B
17. assumption: ContinuousOn skillU (Set.Icc 0 1)
18. assumption: StrictMonoOn skillU (Set.Icc 0 1)
19. assumption: 0 ≤ skillU 0
20. assumption: production 0 = 0
21. assumption: ConcaveOn ℝ (Set.Ioo 0 B) fun (e : ℝ) => production e ^ 2
22. assumption: ConvexOn ℝ (Set.Icc (production (LBG22StrategicRanking.effortIntervalInverse cost E rho)) (production E)) fun (x : ℝ) =>
  x / cost (LBG22StrategicRanking.effortIntervalInverse production E x)
23. assumption: DifferentiableOn ℝ cost (Set.Ioo 0 E)
24. assumption: DifferentiableOn ℝ (deriv cost) (Set.Ioo 0 E)
25. assumption: DifferentiableOn ℝ production (Set.Ioo 0 B)
26. assumption: DifferentiableOn ℝ (deriv production) (Set.Ioo 0 B)
27. assumption: ContinuousOn (deriv (deriv production)) (Set.Ioo 0 B)
28. assumption: Measurable tie
29. assumption: Set.InjOn tie {x : ℝ × ℝ | x.1 ∈ Set.Ioc 0 1 ∧ x.2 ∈ Set.Ioc 0 1}
30. parameter: ℝ
31. assumption: c ∈ Set.Ioo 0 (1 - rho)
32. conclusion: ∃ beta ∈ Set.Ioo 0 1,
  ∀ (effortM effortU score : ℝ → ℝ × ℝ → ℝ),
    (∀ d ∈ Set.Ioc 0 (1 - rho),
        LBG22StrategicRanking.MultitaskEquilibrium cost production (fun (t : ℝ) => lower + width * t) B rho d
          (effortM d) (effortU d) (score d) tie) →
      ∀ d ∈ Set.Ioc 0 (1 - rho),
        LBG22StrategicRanking.sourceActualMultitaskSchoolUtility LBG22StrategicRanking.unitRankMeasure production
            (fun (t : ℝ) => lower + width * t) skillU (effortM d) (effortU d) (score d) tie rho beta d ≤
          LBG22StrategicRanking.sourceActualMultitaskSchoolUtility LBG22StrategicRanking.unitRankMeasure production
            (fun (t : ℝ) => lower + width * t) skillU (effortM c) (effortU c) (score c) tie rho beta c
\end{ReviewVerbatim}

\paragraph{Lean proof endpoint (built separately)}\begin{ReviewVerbatim}
LBG22StrategicRanking.ProofInterface.weightedPrivateUtility
\end{ReviewVerbatim}

\paragraph{Recorded source-to-Spec screening}
\textbf{Verdict:} \texttt{matches\_approved\_corrected\_target}
\\
{\raggedright
\textbf{Reason:} The selected route disclaims archival equivalence and binds defects LBG22-\allowbreak{}C5 and LBG22-\allowbreak{}P8 to the complete approved primitive repair.\allowbreak{} The target includes affine measurable skill with nonnegative intercept and positive slope, E>0 with cost E=1, B>=E, production 0=0, concavity of production squared, convexity of x/\allowbreak{}C(x) on the approved interval, the required twice-\allowbreak{}differentiability/\allowbreak{}continuous-\allowbreak{}second-\allowbreak{}derivative conditions, and an independent regular unmeasurable skill.\allowbreak{} For each interior c it chooses beta in (0,1) before universally quantifying over equilibrium families and compares the actual conditional admitted-\allowbreak{}population weighted utility against every feasible d, exactly as approved.\allowbreak{}
\par}
\reviewmatch{Matches formalized target}
\noindent\textbf{Reviewer annotation}\par
\reviewerbox
\clearpage
\hypertarget{source-claim-53bb19f6829b5622}{}
\section*{Review row 14}
\textbf{Source locator:} {\footnotesize\raggedright cited publication, lines 1844-\allowbreak{}1852\par}
\paragraph{Verbatim source input}
\begin{ReviewVerbatim}
\begin{lemma}
In any equilibrium, tie-breaking %does not matter:
{is not necessary}: ties in post-effort scores lead to ties in post-effort rewards. For any $\Gamma$, in any equilibrium the distribution of post effort scores $v$ is such that, for all $\omega, \omega'$
\[
v\left(e(\theta_\pre(\omega)),\theta_\pre(\omega)\right) = v\left(e(\theta_\pre(\omega')),\theta_\pre(\omega')\right) \triangleq v \implies
\lambda\left(\gamma\left(\omega, v \right)\right) = \lambda\left(\gamma\left(\omega', v \right)\right).
\]
\label{lem:tienotmatter}
\end{lemma}

[Next verbatim source excerpt]

\parbold{Applicants}
There is a unit mass of \textit{applicants}, indexed by an observed index $\omega \in [0,1]$ distributed uniformly.\footnote{The index $\omega$ should be interpreted as each applicant's ``name,'' uncorrelated with skill, used solely for tie-breaking.}  Each applicant has a latent (unobserved) skill level represented by some measurable function of $\omega$.
We assume that {the distribution of the skills} has no atoms and that the CDF of this distribution is strictly increasing.
Using the CDF to map the skill of an applicant to a rank in~$[0,1]$, each applicant gets an (unobserved) rank $\theta_\pre =\theta_\pre(\omega)$ which by our assumption on the CDF is again uniformly distributed in $[0, 1]$ (the higher the rank the better).  With this setup, the skill of an applicant with rank $\theta_\pre$ can be written as $f(\theta_\pre)$ where $f$ is a strictly increasing, continuous function.  {It will be notationally convenient to label applicants by their rank $\theta_\pre$, though the reader should note that in contrast to $\omega$, $\theta_\pre(\omega)$ is assumed to be unobservable.}

%3. Each applicant exerts some effort e, at cost \cost(e).


Each applicant chooses an \textit{effort} level $e \ge 0$, the result of which is an observed, post-effort \textit{score}, {$v=v(e,\theta_\pre) = {g(e)}\cdot{f(\theta_\pre)}$}. \lledit{In other words, we assume the post-effort score is the product of two components, associated with the effort and the pre-effort skill respectively.}
{We assume} that the effort transfer function $g: [0, \infty) \mapsto [0, \infty)$ is a continuous, concave, strictly increasing function, representing that marginal effort improves one's score but has diminishing returns.
{%Since applicants can be labelled by their skill level $f(\theta_\pre)$, or equivalently, by their rank $\theta_{\pre}$,
{T}he strategies of the applicants can {then} be described by a function $\theta_\pre \mapsto e(\theta_\pre)$.
%\cb{$\omega\mapsto e(\omega)$.}%, leading to a distribution over post-effort scores $v=v(\theta_\pre))$ via the uniform distribution of $\theta_\pre$ in $[0,1]$.
}
  Each {applicant} is then ranked according to their score $v$%(where we break ties uniformly at random, if they appear)
, resulting in a \textit{post-effort rank} $\theta_\post$.
Note that the ranking $\theta_\post$ is slightly less trivial than a ranking of the skills, since the scores might have ties, which have to be resolved.

{\textbf{Tie Breaking.}
Given a choice of strategies $\theta_\pre \mapsto e(\theta_\pre)$, let $F$ be the CDF of the scores $v(e(\theta_\pre),\theta_\pre)$.
Since atoms for the distribution of $v$ would lead to ties for ranks defined as $F(v)$, we will use the labels of the applicants to break ties with the help of a (publicly announced) tie-breaking function $\Gamma(\omega)$, defined, e.g., via a collision free hash of the applicant names. Here  we require  that $\Gamma$  is a measurable function from $[0,1]$ to $[0,1]$ that maps different applicant labels to different values.
We then use $\Gamma$ to resolve the atoms of $F$, leading to a ranking function $v\mapsto \gamma(\omega, v)$ which is equal to $F(v)$ except when $v$ is an atom of the score distribution, {in which case it takes values in the ``gap interval'' $[F_-(v),  F(v)]$, where $F_-(v)$ is the left limit of $F$ at $v$.}  We  construct $\gamma$ in such a way that it
gives the uniform distribution for $\theta_\post(\omega)$ if we set  $\theta_\post(\omega)=\gamma(\omega,v(e(\theta_\pre(\omega)),\theta_\pre(\omega)))$.\footnote{See Remark~\ref{rem:gamma} for details on this construction.}

[Next verbatim source excerpt]

\parbold{Ranking designer (school)} A single \textit{school} is admitting applicants, based on their ranking. In particular, the school can choose a ranking reward function $\lambda : [0, 1] \mapsto [0, 1]$, such that an applicant with post-effort rank $\theta_\post$ is admitted with probability $\lambda(\theta_\post)$. We assume that $\lambda$ is non-decreasing and that the school has a constraint on the overall probability, such that in expectation it admits a number of applicants equal to {a} capacity constraint  $\rho \in (0,1)$, i.e., \lledit{$\E_{\theta_\post}[\lambda(\theta_\post)] = \rho $}.\footnote{\lledit{We use $\E_{\theta_\post}[\cdot]$ to denote an integral over $\theta_\post$ (and $\E[\cdot]$ to denote an integral over $\omega$)} with respect to the Lebesgue measure; informally, this can be thought of as averaging over the applicant population.} We may also refer to $\lambda$, informally, as the admission policy.



For simplicity, we further assume that $\lambda$ is a step-function with $K$ distinct levels $\ell_0 < \dots < \ell_{K-1}$, and $K-1$ cut-points parameterized by $c_1< \dots < c_{K-1}$ (with $c_0 = 0$, $c_K = 1$). %$\ell_0, \dots, \ell_{K-1}$.
In other words, we have $\lambda(\theta) = \ell_k$, for all $\theta \in \psi_k \triangleq [c_{k}, c_{k+1})$. Thus, applicants in the same post-rank interval $\theta \in \psi_k$ receive the same reward.

[Next verbatim source excerpt]

\parbold{Individual applicant welfare and equilibrium} Given the designer's function $\lambda$ and the effort levels of other applicants, each applicant chooses effort $e$ to maximize their individual welfare, % $W$,
\begin{equation*}
	\iwelf(e, \lambda(\theta_\post)) = \lambda(\theta_\post) - \cost(e),
\end{equation*}
where the effort cost function $\cost$ is non-negative, continuous, and strictly convex {on $[0,\infty)$, with $e_0:=\argmin_e \cost(e)$ and $p(e_0) = 0$.}
%\lnote{also want to assume that $\cost$ attains its minimum at $0$?}
%Let $e_0:=\argmin_e \cost(e)$ and $p(e_0) = 0$.

Applicants are assumed not to personally benefit from increasing their score $v$, except through the corresponding increase in their rank and reward. While the definition of $\cost$ does not preclude the applicant receiving any intrinsic benefit from exerting effort, we assume that the \emph{net} benefit from effort is non-positive.

[Next verbatim source excerpt]

\begin{definition} [Equilibrium]\label{def:equi}
{Given a}
%Let the
tie-breaking function $\Gamma(\omega)$
%be a pre-announced ordering over index $\omega$, to break ties in post-effort value. Then, let $\gamma(\omega, v)$ be the 1-1 ranking function defined as the CDF of post-effort values $v$ except where there are ties in $v$, in which case the ordering $\Gamma$ is used to produce unique ranks such that the range of $\gamma$ is uniform on $[0, 1]$.\footnote{More formally, note that each discontinuity in the CDF of post-effort values $v$ has a height equal to the measure of the applicants who are tied at that value, and so $\gamma$ for those applicants can be filled in using the CDF of $\Gamma(\omega)$ restricted to the applicants $\omega$ who are tied. Thus, the range of $\gamma$ is uniform on $[0,1]$.
%
{and}
%	Then, given
	a ranking probability function $\lambda$, an \textit{equilibrium} is a set of effort levels and post-effort ranking rewards for each applicant, $\{e(\theta_\pre),  \lambda(\theta_\post(\theta_\pre))\}$ such that, for all
	$\omega$, %$\theta_\pre$,
	\begin{align*}
		e(\theta_\pre(\omega)) &\in \argmax_{e} \iwelf\left(e, \lambda\left(\gamma\left(\omega, v\left(e,\theta_\pre(\omega)\right)\right)\right)\right) %\triangleq \argmax_{e} \left[ \lambda(\gamma(v(e,\theta_\pre)) - \cost(e) \right].
		\\
		\theta_\post(\theta_\pre(\omega)) &= \gamma(\omega, v(e(\theta_\pre(\omega)),\theta_\pre(\omega))),
	\end{align*}
	{where $\gamma(\omega,v)$ is the ranking induced by the CDF of the scores resulting from effort levels $\{e(\theta_\pre)\}$ and the tie-breaking function $\Gamma$.\footnote{See Remark~\ref{rem:gamma-effort} on $\gamma$ and the set of efforts.%Formally, $\gamma$ depends on the set of efforts. Given a fixed set and $\gamma$, for each applicant $\omega$ the first condition considers the counter-factual ranking of $\omega$ with different post-effort values but using the same ranking function. As defined, $\gamma$ yields a uniform distribution of ranks with such measure $0$ changes. In Appendix \Cref{lem:deviationsmeasure0}, we further prove that two effort sets equal up to sets of measure $0$ induce the same ranking function $\gamma$, and so the condition is consistent.%, and so we set it at the $\gamma$ induced by the effort levels $\{e(\theta_\pre)\}$.%, and so it does not vary.
	}
	}
	%In particular, we denote $\theta_\post(\theta_\pre) := \gamma(v(e(\theta_\pre),\theta_\pre))$.
	%\lledit{and for any $\theta_\pre, \theta'_\pre$, we have $\lambda(\theta_\post(\theta_\pre)) > \lambda(\theta_\post(\theta'_\pre)) \implies v(e(\theta_\pre),\theta_\pre) > v(e(\theta_\pre'),\theta_\pre')$.}
\end{definition}
\end{ReviewVerbatim}

% report-context-presentation-sha256: 90f9116e292dd71fddfd6b445ca0b060a481d026855bdb9725f2833c2cf62a31
\paragraph{Settled source reading}
\begin{sloppypar}
Equilibrium means that almost every applicant best responds against every feasible deviation;\allowbreak{} uniqueness is up to a null set.\allowbreak{}
\end{sloppypar}
\paragraph{Settled source reading}
\begin{sloppypar}
Applicants know the fixed priority used at score-\allowbreak{}boundary ties.\allowbreak{} This clarifies the stated model, not a corrected admission policy or a new economic assumption.\allowbreak{}
\end{sloppypar}

\paragraph{Expanded PaperInterface specification}
\begin{ReviewVerbatim}
∀ (cost production skill tie effort score : ℝ → ℝ) (baseline rho : ℝ) (n : ℕ) (cutoff reward : ℕ → ℝ),
  LBG22StrategicRanking.RankingPrimitives cost production skill baseline →
    LBG22StrategicRanking.FiniteAdmissionPolicy n cutoff reward rho →
      LBG22StrategicRanking.RankingEquilibrium cost production skill tie n cutoff reward effort score →
        ∃ (good : Set ℝ),
          (∀ᵐ (t : ℝ) ∂LBG22StrategicRanking.unitRankMeasure, t ∈ good) ∧
            ∀ x ∈ good,
              ∀ y ∈ good,
                score x = score y →
                  reward
                      ↑(LBG22StrategicRanking.finiteLowerRankBand (fun (i : Fin (n + 1)) => cutoff ↑i)
                          (LBG22StrategicRanking.tieBrokenRank LBG22StrategicRanking.unitRankMeasure score tie x)) =
                    reward
                      ↑(LBG22StrategicRanking.finiteLowerRankBand (fun (i : Fin (n + 1)) => cutoff ↑i)
                          (LBG22StrategicRanking.tieBrokenRank LBG22StrategicRanking.unitRankMeasure score tie y))
\end{ReviewVerbatim}

\paragraph{Retained governing declarations}
\begin{itemize}\raggedright
\item \hyperlink{paper-prerequisite-d056dd7c8060653c}{LBG22\allowbreak{}Strategic\allowbreak{}Ranking.\allowbreak{}Finite\allowbreak{}Admission\allowbreak{}Policy}
\item \hyperlink{paper-prerequisite-2a1b692f4e3136c6}{LBG22\allowbreak{}Strategic\allowbreak{}Ranking.\allowbreak{}Ranking\allowbreak{}Equilibrium}
\item \hyperlink{paper-prerequisite-9bbcc12f4a1736bf}{LBG22\allowbreak{}Strategic\allowbreak{}Ranking.\allowbreak{}Ranking\allowbreak{}Primitives}
\item \hyperlink{paper-prerequisite-a42f440ddc254fa7}{LBG22\allowbreak{}Strategic\allowbreak{}Ranking.\allowbreak{}finite\allowbreak{}Lower\allowbreak{}Rank\allowbreak{}Band}
\item \hyperlink{paper-prerequisite-b7c71e53177b76c6}{LBG22\allowbreak{}Strategic\allowbreak{}Ranking.\allowbreak{}tie\allowbreak{}Broken\allowbreak{}Rank}
\item \hyperlink{governing-declaration-678cf48d549cb620}{LBG22\allowbreak{}Strategic\allowbreak{}Ranking.\allowbreak{}unit\allowbreak{}Rank\allowbreak{}Measure}
\end{itemize}
\paragraph{Lean-elaborated claim roles}\begin{ReviewVerbatim}
1. parameter: ℝ → ℝ
2. parameter: ℝ → ℝ
3. parameter: ℝ → ℝ
4. parameter: ℝ → ℝ
5. parameter: ℝ → ℝ
6. parameter: ℝ → ℝ
7. parameter: ℝ
8. parameter: ℝ
9. parameter: ℕ
10. parameter: ℕ → ℝ
11. parameter: ℕ → ℝ
12. assumption: LBG22StrategicRanking.RankingPrimitives cost production skill baseline
13. assumption: LBG22StrategicRanking.FiniteAdmissionPolicy n cutoff reward rho
14. assumption: LBG22StrategicRanking.RankingEquilibrium cost production skill tie n cutoff reward effort score
15. conclusion: ∃ (good : Set ℝ),
  (∀ᵐ (t : ℝ) ∂LBG22StrategicRanking.unitRankMeasure, t ∈ good) ∧
    ∀ x ∈ good,
      ∀ y ∈ good,
        score x = score y →
          reward
              ↑(LBG22StrategicRanking.finiteLowerRankBand (fun (i : Fin (n + 1)) => cutoff ↑i)
                  (LBG22StrategicRanking.tieBrokenRank LBG22StrategicRanking.unitRankMeasure score tie x)) =
            reward
              ↑(LBG22StrategicRanking.finiteLowerRankBand (fun (i : Fin (n + 1)) => cutoff ↑i)
                  (LBG22StrategicRanking.tieBrokenRank LBG22StrategicRanking.unitRankMeasure score tie y))
\end{ReviewVerbatim}

\paragraph{Lean proof endpoint (built separately)}\begin{ReviewVerbatim}
LBG22StrategicRanking.ProofInterface.equalScoreRewards
\end{ReviewVerbatim}

\paragraph{Recorded source-to-Spec screening}
\textbf{Verdict:} \texttt{matches}
\\
{\raggedright
\textbf{Reason:} On one full-\allowbreak{}measure set, the target quantifies over every pair of applicants and concludes that equal actual scores receive equal finite-\allowbreak{}policy rewards.\allowbreak{} This is the source lemma with the approved a.\allowbreak{}e.\allowbreak{} equilibrium convention;\allowbreak{} the displayed rank, band, policy, and fixed-\allowbreak{}priority prerequisites do not change the conclusion.\allowbreak{}
\par}
\reviewmatch{Matches source input}
\noindent\textbf{Reviewer annotation}\par
\reviewerbox
\clearpage
\hypertarget{source-claim-1693f763578ff49e}{}
\section*{Review row 15}
\textbf{Source locator:} {\footnotesize\raggedright cited publication, lines 1869-\allowbreak{}1883\par}
\paragraph{Verbatim source input}
\begin{ReviewVerbatim}
\begin{lemma}[Deviations of measure 0]
\label{lem:deviationsmeasure0}
Fix $\Gamma$ and $\lambda$. Consider strategy set $\{e(\theta_\pre(\omega)\}$, and corresponding CDF $F$ of post-effort scores. As defined, $\gamma(\omega, v)$ is uniquely determined by $F$ and $\Gamma$.
Now, suppose a measure $0$ set $\{\omega\}$ deviates, leading to strategy set $\{\tilde e(\theta_\pre(\omega)\}$, post-effort value distribution $\tilde F$, and ranking function $\tilde\gamma$.

Then, $\tilde F=F$, and for all $\omega$ and $v$,
$$\lambda(\tilde\gamma(\omega, v)) =  \lambda(\gamma(\omega, v)).$$
% i.e., for each $v$, rewards under the new ranking function $\tilde\gamma$ are identical to those under the old ranking function $\gamma$.

% If $\{e(\theta_\pre(\omega)\}$ induces an equilibrium, then we further have for all $\omega, \omega'$ and $v$,

% $$\lambda(\tilde\gamma(\omega, v)) =  \lambda(\gamma(\omega, v)).$$

% Then, $F=F'$, and for all $\omega$, $$\lambda(\gamma'(\omega, v(e'(\theta_\pre(\omega)), \theta_\pre(\omega)))) =  \lambda(\gamma(\omega, v(e'(\theta_\pre(\omega)), \theta_\pre(\omega)))),$$ i.e., rewards under the new ranking function $\gamma'$ and new post-effort values $v(e'(\theta_\pre(\omega)), \theta_\pre(\omega))$ are identical to those under the old ranking function $\gamma$ and new effort levels $e'$.
\end{lemma}

[Next verbatim source excerpt]

\parbold{Applicants}
There is a unit mass of \textit{applicants}, indexed by an observed index $\omega \in [0,1]$ distributed uniformly.\footnote{The index $\omega$ should be interpreted as each applicant's ``name,'' uncorrelated with skill, used solely for tie-breaking.}  Each applicant has a latent (unobserved) skill level represented by some measurable function of $\omega$.
We assume that {the distribution of the skills} has no atoms and that the CDF of this distribution is strictly increasing.
Using the CDF to map the skill of an applicant to a rank in~$[0,1]$, each applicant gets an (unobserved) rank $\theta_\pre =\theta_\pre(\omega)$ which by our assumption on the CDF is again uniformly distributed in $[0, 1]$ (the higher the rank the better).  With this setup, the skill of an applicant with rank $\theta_\pre$ can be written as $f(\theta_\pre)$ where $f$ is a strictly increasing, continuous function.  {It will be notationally convenient to label applicants by their rank $\theta_\pre$, though the reader should note that in contrast to $\omega$, $\theta_\pre(\omega)$ is assumed to be unobservable.}

%3. Each applicant exerts some effort e, at cost \cost(e).


Each applicant chooses an \textit{effort} level $e \ge 0$, the result of which is an observed, post-effort \textit{score}, {$v=v(e,\theta_\pre) = {g(e)}\cdot{f(\theta_\pre)}$}. \lledit{In other words, we assume the post-effort score is the product of two components, associated with the effort and the pre-effort skill respectively.}
{We assume} that the effort transfer function $g: [0, \infty) \mapsto [0, \infty)$ is a continuous, concave, strictly increasing function, representing that marginal effort improves one's score but has diminishing returns.
{%Since applicants can be labelled by their skill level $f(\theta_\pre)$, or equivalently, by their rank $\theta_{\pre}$,
{T}he strategies of the applicants can {then} be described by a function $\theta_\pre \mapsto e(\theta_\pre)$.
%\cb{$\omega\mapsto e(\omega)$.}%, leading to a distribution over post-effort scores $v=v(\theta_\pre))$ via the uniform distribution of $\theta_\pre$ in $[0,1]$.
}
  Each {applicant} is then ranked according to their score $v$%(where we break ties uniformly at random, if they appear)
, resulting in a \textit{post-effort rank} $\theta_\post$.
Note that the ranking $\theta_\post$ is slightly less trivial than a ranking of the skills, since the scores might have ties, which have to be resolved.

{\textbf{Tie Breaking.}
Given a choice of strategies $\theta_\pre \mapsto e(\theta_\pre)$, let $F$ be the CDF of the scores $v(e(\theta_\pre),\theta_\pre)$.
Since atoms for the distribution of $v$ would lead to ties for ranks defined as $F(v)$, we will use the labels of the applicants to break ties with the help of a (publicly announced) tie-breaking function $\Gamma(\omega)$, defined, e.g., via a collision free hash of the applicant names. Here  we require  that $\Gamma$  is a measurable function from $[0,1]$ to $[0,1]$ that maps different applicant labels to different values.
We then use $\Gamma$ to resolve the atoms of $F$, leading to a ranking function $v\mapsto \gamma(\omega, v)$ which is equal to $F(v)$ except when $v$ is an atom of the score distribution, {in which case it takes values in the ``gap interval'' $[F_-(v),  F(v)]$, where $F_-(v)$ is the left limit of $F$ at $v$.}  We  construct $\gamma$ in such a way that it
gives the uniform distribution for $\theta_\post(\omega)$ if we set  $\theta_\post(\omega)=\gamma(\omega,v(e(\theta_\pre(\omega)),\theta_\pre(\omega)))$.\footnote{See Remark~\ref{rem:gamma} for details on this construction.}

[Next verbatim source excerpt]

\parbold{Ranking designer (school)} A single \textit{school} is admitting applicants, based on their ranking. In particular, the school can choose a ranking reward function $\lambda : [0, 1] \mapsto [0, 1]$, such that an applicant with post-effort rank $\theta_\post$ is admitted with probability $\lambda(\theta_\post)$. We assume that $\lambda$ is non-decreasing and that the school has a constraint on the overall probability, such that in expectation it admits a number of applicants equal to {a} capacity constraint  $\rho \in (0,1)$, i.e., \lledit{$\E_{\theta_\post}[\lambda(\theta_\post)] = \rho $}.\footnote{\lledit{We use $\E_{\theta_\post}[\cdot]$ to denote an integral over $\theta_\post$ (and $\E[\cdot]$ to denote an integral over $\omega$)} with respect to the Lebesgue measure; informally, this can be thought of as averaging over the applicant population.} We may also refer to $\lambda$, informally, as the admission policy.



For simplicity, we further assume that $\lambda$ is a step-function with $K$ distinct levels $\ell_0 < \dots < \ell_{K-1}$, and $K-1$ cut-points parameterized by $c_1< \dots < c_{K-1}$ (with $c_0 = 0$, $c_K = 1$). %$\ell_0, \dots, \ell_{K-1}$.
In other words, we have $\lambda(\theta) = \ell_k$, for all $\theta \in \psi_k \triangleq [c_{k}, c_{k+1})$. Thus, applicants in the same post-rank interval $\theta \in \psi_k$ receive the same reward.

[Next verbatim source excerpt]

\parbold{Individual applicant welfare and equilibrium} Given the designer's function $\lambda$ and the effort levels of other applicants, each applicant chooses effort $e$ to maximize their individual welfare, % $W$,
\begin{equation*}
	\iwelf(e, \lambda(\theta_\post)) = \lambda(\theta_\post) - \cost(e),
\end{equation*}
where the effort cost function $\cost$ is non-negative, continuous, and strictly convex {on $[0,\infty)$, with $e_0:=\argmin_e \cost(e)$ and $p(e_0) = 0$.}
%\lnote{also want to assume that $\cost$ attains its minimum at $0$?}
%Let $e_0:=\argmin_e \cost(e)$ and $p(e_0) = 0$.

Applicants are assumed not to personally benefit from increasing their score $v$, except through the corresponding increase in their rank and reward. While the definition of $\cost$ does not preclude the applicant receiving any intrinsic benefit from exerting effort, we assume that the \emph{net} benefit from effort is non-positive.

[Next verbatim source excerpt]

\begin{definition} [Equilibrium]\label{def:equi}
{Given a}
%Let the
tie-breaking function $\Gamma(\omega)$
%be a pre-announced ordering over index $\omega$, to break ties in post-effort value. Then, let $\gamma(\omega, v)$ be the 1-1 ranking function defined as the CDF of post-effort values $v$ except where there are ties in $v$, in which case the ordering $\Gamma$ is used to produce unique ranks such that the range of $\gamma$ is uniform on $[0, 1]$.\footnote{More formally, note that each discontinuity in the CDF of post-effort values $v$ has a height equal to the measure of the applicants who are tied at that value, and so $\gamma$ for those applicants can be filled in using the CDF of $\Gamma(\omega)$ restricted to the applicants $\omega$ who are tied. Thus, the range of $\gamma$ is uniform on $[0,1]$.
%
{and}
%	Then, given
	a ranking probability function $\lambda$, an \textit{equilibrium} is a set of effort levels and post-effort ranking rewards for each applicant, $\{e(\theta_\pre),  \lambda(\theta_\post(\theta_\pre))\}$ such that, for all
	$\omega$, %$\theta_\pre$,
	\begin{align*}
		e(\theta_\pre(\omega)) &\in \argmax_{e} \iwelf\left(e, \lambda\left(\gamma\left(\omega, v\left(e,\theta_\pre(\omega)\right)\right)\right)\right) %\triangleq \argmax_{e} \left[ \lambda(\gamma(v(e,\theta_\pre)) - \cost(e) \right].
		\\
		\theta_\post(\theta_\pre(\omega)) &= \gamma(\omega, v(e(\theta_\pre(\omega)),\theta_\pre(\omega))),
	\end{align*}
	{where $\gamma(\omega,v)$ is the ranking induced by the CDF of the scores resulting from effort levels $\{e(\theta_\pre)\}$ and the tie-breaking function $\Gamma$.\footnote{See Remark~\ref{rem:gamma-effort} on $\gamma$ and the set of efforts.%Formally, $\gamma$ depends on the set of efforts. Given a fixed set and $\gamma$, for each applicant $\omega$ the first condition considers the counter-factual ranking of $\omega$ with different post-effort values but using the same ranking function. As defined, $\gamma$ yields a uniform distribution of ranks with such measure $0$ changes. In Appendix \Cref{lem:deviationsmeasure0}, we further prove that two effort sets equal up to sets of measure $0$ induce the same ranking function $\gamma$, and so the condition is consistent.%, and so we set it at the $\gamma$ induced by the effort levels $\{e(\theta_\pre)\}$.%, and so it does not vary.
	}
	}
	%In particular, we denote $\theta_\post(\theta_\pre) := \gamma(v(e(\theta_\pre),\theta_\pre))$.
	%\lledit{and for any $\theta_\pre, \theta'_\pre$, we have $\lambda(\theta_\post(\theta_\pre)) > \lambda(\theta_\post(\theta'_\pre)) \implies v(e(\theta_\pre),\theta_\pre) > v(e(\theta_\pre'),\theta_\pre')$.}
\end{definition}
\end{ReviewVerbatim}

% report-context-presentation-sha256: 827a19afe0274b1faef13e344ddbcfe40d9841f2ff34b92ac9d53c0f83e29f56
\paragraph{Settled source reading}
\begin{sloppypar}
Applicants know the fixed priority used at score-\allowbreak{}boundary ties.\allowbreak{} This clarifies the stated model, not a corrected admission policy or a new economic assumption.\allowbreak{}
\end{sloppypar}

\paragraph{Expanded PaperInterface specification}
\begin{ReviewVerbatim}
∀ (production skill tie effort deviatedEffort admission : ℝ → ℝ),
  effort =ᵐ[LBG22StrategicRanking.unitRankMeasure] deviatedEffort →
    have score : ℝ → ℝ := fun (t : ℝ) => production (effort t) * skill t;
    have deviatedScore : ℝ → ℝ := fun (t : ℝ) => production (deviatedEffort t) * skill t;
    MeasureTheory.Measure.map score LBG22StrategicRanking.unitRankMeasure =
        MeasureTheory.Measure.map deviatedScore LBG22StrategicRanking.unitRankMeasure ∧
      ∀ (t v : ℝ),
        admission
            (LBG22StrategicRanking.counterfactualTieBrokenRank LBG22StrategicRanking.unitRankMeasure score tie t v) =
          admission
            (LBG22StrategicRanking.counterfactualTieBrokenRank LBG22StrategicRanking.unitRankMeasure deviatedScore tie t
              v)
\end{ReviewVerbatim}

\paragraph{Retained governing declarations}
\begin{itemize}\raggedright
\item \hyperlink{paper-prerequisite-6ed45e46d01eadeb}{LBG22\allowbreak{}Strategic\allowbreak{}Ranking.\allowbreak{}counterfactual\allowbreak{}Tie\allowbreak{}Broken\allowbreak{}Rank}
\item \hyperlink{governing-declaration-678cf48d549cb620}{LBG22\allowbreak{}Strategic\allowbreak{}Ranking.\allowbreak{}unit\allowbreak{}Rank\allowbreak{}Measure}
\end{itemize}
\paragraph{Lean-elaborated claim roles}\begin{ReviewVerbatim}
1. parameter: ℝ → ℝ
2. parameter: ℝ → ℝ
3. parameter: ℝ → ℝ
4. parameter: ℝ → ℝ
5. parameter: ℝ → ℝ
6. parameter: ℝ → ℝ
7. assumption: effort =ᵐ[LBG22StrategicRanking.unitRankMeasure] deviatedEffort
8. conclusion: MeasureTheory.Measure.map (fun (t : ℝ) => production (effort t) * skill t) LBG22StrategicRanking.unitRankMeasure =
    MeasureTheory.Measure.map (fun (t : ℝ) => production (deviatedEffort t) * skill t)
      LBG22StrategicRanking.unitRankMeasure ∧
  ∀ (t v : ℝ),
    admission
        (LBG22StrategicRanking.counterfactualTieBrokenRank LBG22StrategicRanking.unitRankMeasure
          (fun (t : ℝ) => production (effort t) * skill t) tie t v) =
      admission
        (LBG22StrategicRanking.counterfactualTieBrokenRank LBG22StrategicRanking.unitRankMeasure
          (fun (t : ℝ) => production (deviatedEffort t) * skill t) tie t v)
\end{ReviewVerbatim}

\paragraph{Lean proof endpoint (built separately)}\begin{ReviewVerbatim}
LBG22StrategicRanking.ProofInterface.nullDeviations
\end{ReviewVerbatim}

\paragraph{Recorded source-to-Spec screening}
\textbf{Verdict:} \texttt{matches}
\\
{\raggedright
\textbf{Reason:} A.\allowbreak{}e.\allowbreak{}-\allowbreak{}equal effort profiles induce equal pushforward score measures, and the target then gives identical fixed-\allowbreak{}priority counterfactual rank rewards for every applicant and proposed score.\allowbreak{} This is exactly Lemma C.\allowbreak{}2 under the displayed fixed-\allowbreak{}priority convention.\allowbreak{}
\par}
\reviewmatch{Matches source input}
\noindent\textbf{Reviewer annotation}\par
\reviewerbox
\clearpage
\hypertarget{source-claim-12dd3b600e58e716}{}
\section*{Review row 16}
\textbf{Source locator:} {\footnotesize\raggedright cited publication, lines 2005-\allowbreak{}2016\par}
\paragraph{Verbatim source input}
\begin{ReviewVerbatim}
\begin{lemma}\label{lem:order_p}
Suppose there are applicants $\omega, \omega'$ with $\theta_\pre(\omega') <\theta_\pre(\omega) $ and effort levels, $e<\{\tilde e', \tilde e \}<e'$ such that $\lambda(\gamma(\omega, e)) = \lambda(\gamma(\omega', \tilde e'))$, and $\lambda(\gamma(\omega', e')) = \lambda(\gamma(\omega, \tilde e))$. Then the following inequalities
\begin{align}
    	\lambda(\gamma(\omega, e)) - \cost(e) &{\geq }
		%>
		\lambda(\gamma(\omega, \tilde e)) - \cost(\tilde e) \label{eq:OPT1} \\
		\lambda(\gamma(\omega', e')) - \cost(e')
		%>
		&{\geq}\lambda(\gamma(\omega', \tilde e')) - \cost(\tilde e') \label{eq:OPT2}
\end{align}
imply $\cost(e') - \cost(\tilde e') \leq \cost(\tilde e) - \cost(e). $
\end{lemma}

[Next verbatim source excerpt]

\parbold{Applicants}
There is a unit mass of \textit{applicants}, indexed by an observed index $\omega \in [0,1]$ distributed uniformly.\footnote{The index $\omega$ should be interpreted as each applicant's ``name,'' uncorrelated with skill, used solely for tie-breaking.}  Each applicant has a latent (unobserved) skill level represented by some measurable function of $\omega$.
We assume that {the distribution of the skills} has no atoms and that the CDF of this distribution is strictly increasing.
Using the CDF to map the skill of an applicant to a rank in~$[0,1]$, each applicant gets an (unobserved) rank $\theta_\pre =\theta_\pre(\omega)$ which by our assumption on the CDF is again uniformly distributed in $[0, 1]$ (the higher the rank the better).  With this setup, the skill of an applicant with rank $\theta_\pre$ can be written as $f(\theta_\pre)$ where $f$ is a strictly increasing, continuous function.  {It will be notationally convenient to label applicants by their rank $\theta_\pre$, though the reader should note that in contrast to $\omega$, $\theta_\pre(\omega)$ is assumed to be unobservable.}

%3. Each applicant exerts some effort e, at cost \cost(e).


Each applicant chooses an \textit{effort} level $e \ge 0$, the result of which is an observed, post-effort \textit{score}, {$v=v(e,\theta_\pre) = {g(e)}\cdot{f(\theta_\pre)}$}. \lledit{In other words, we assume the post-effort score is the product of two components, associated with the effort and the pre-effort skill respectively.}
{We assume} that the effort transfer function $g: [0, \infty) \mapsto [0, \infty)$ is a continuous, concave, strictly increasing function, representing that marginal effort improves one's score but has diminishing returns.
{%Since applicants can be labelled by their skill level $f(\theta_\pre)$, or equivalently, by their rank $\theta_{\pre}$,
{T}he strategies of the applicants can {then} be described by a function $\theta_\pre \mapsto e(\theta_\pre)$.
%\cb{$\omega\mapsto e(\omega)$.}%, leading to a distribution over post-effort scores $v=v(\theta_\pre))$ via the uniform distribution of $\theta_\pre$ in $[0,1]$.
}
  Each {applicant} is then ranked according to their score $v$%(where we break ties uniformly at random, if they appear)
, resulting in a \textit{post-effort rank} $\theta_\post$.
Note that the ranking $\theta_\post$ is slightly less trivial than a ranking of the skills, since the scores might have ties, which have to be resolved.

{\textbf{Tie Breaking.}
Given a choice of strategies $\theta_\pre \mapsto e(\theta_\pre)$, let $F$ be the CDF of the scores $v(e(\theta_\pre),\theta_\pre)$.
Since atoms for the distribution of $v$ would lead to ties for ranks defined as $F(v)$, we will use the labels of the applicants to break ties with the help of a (publicly announced) tie-breaking function $\Gamma(\omega)$, defined, e.g., via a collision free hash of the applicant names. Here  we require  that $\Gamma$  is a measurable function from $[0,1]$ to $[0,1]$ that maps different applicant labels to different values.
We then use $\Gamma$ to resolve the atoms of $F$, leading to a ranking function $v\mapsto \gamma(\omega, v)$ which is equal to $F(v)$ except when $v$ is an atom of the score distribution, {in which case it takes values in the ``gap interval'' $[F_-(v),  F(v)]$, where $F_-(v)$ is the left limit of $F$ at $v$.}  We  construct $\gamma$ in such a way that it
gives the uniform distribution for $\theta_\post(\omega)$ if we set  $\theta_\post(\omega)=\gamma(\omega,v(e(\theta_\pre(\omega)),\theta_\pre(\omega)))$.\footnote{See Remark~\ref{rem:gamma} for details on this construction.}

[Next verbatim source excerpt]

\parbold{Ranking designer (school)} A single \textit{school} is admitting applicants, based on their ranking. In particular, the school can choose a ranking reward function $\lambda : [0, 1] \mapsto [0, 1]$, such that an applicant with post-effort rank $\theta_\post$ is admitted with probability $\lambda(\theta_\post)$. We assume that $\lambda$ is non-decreasing and that the school has a constraint on the overall probability, such that in expectation it admits a number of applicants equal to {a} capacity constraint  $\rho \in (0,1)$, i.e., \lledit{$\E_{\theta_\post}[\lambda(\theta_\post)] = \rho $}.\footnote{\lledit{We use $\E_{\theta_\post}[\cdot]$ to denote an integral over $\theta_\post$ (and $\E[\cdot]$ to denote an integral over $\omega$)} with respect to the Lebesgue measure; informally, this can be thought of as averaging over the applicant population.} We may also refer to $\lambda$, informally, as the admission policy.



For simplicity, we further assume that $\lambda$ is a step-function with $K$ distinct levels $\ell_0 < \dots < \ell_{K-1}$, and $K-1$ cut-points parameterized by $c_1< \dots < c_{K-1}$ (with $c_0 = 0$, $c_K = 1$). %$\ell_0, \dots, \ell_{K-1}$.
In other words, we have $\lambda(\theta) = \ell_k$, for all $\theta \in \psi_k \triangleq [c_{k}, c_{k+1})$. Thus, applicants in the same post-rank interval $\theta \in \psi_k$ receive the same reward.

[Next verbatim source excerpt]

\parbold{Individual applicant welfare and equilibrium} Given the designer's function $\lambda$ and the effort levels of other applicants, each applicant chooses effort $e$ to maximize their individual welfare, % $W$,
\begin{equation*}
	\iwelf(e, \lambda(\theta_\post)) = \lambda(\theta_\post) - \cost(e),
\end{equation*}
where the effort cost function $\cost$ is non-negative, continuous, and strictly convex {on $[0,\infty)$, with $e_0:=\argmin_e \cost(e)$ and $p(e_0) = 0$.}
%\lnote{also want to assume that $\cost$ attains its minimum at $0$?}
%Let $e_0:=\argmin_e \cost(e)$ and $p(e_0) = 0$.

Applicants are assumed not to personally benefit from increasing their score $v$, except through the corresponding increase in their rank and reward. While the definition of $\cost$ does not preclude the applicant receiving any intrinsic benefit from exerting effort, we assume that the \emph{net} benefit from effort is non-positive.

[Next verbatim source excerpt]

\begin{definition} [Equilibrium]\label{def:equi}
{Given a}
%Let the
tie-breaking function $\Gamma(\omega)$
%be a pre-announced ordering over index $\omega$, to break ties in post-effort value. Then, let $\gamma(\omega, v)$ be the 1-1 ranking function defined as the CDF of post-effort values $v$ except where there are ties in $v$, in which case the ordering $\Gamma$ is used to produce unique ranks such that the range of $\gamma$ is uniform on $[0, 1]$.\footnote{More formally, note that each discontinuity in the CDF of post-effort values $v$ has a height equal to the measure of the applicants who are tied at that value, and so $\gamma$ for those applicants can be filled in using the CDF of $\Gamma(\omega)$ restricted to the applicants $\omega$ who are tied. Thus, the range of $\gamma$ is uniform on $[0,1]$.
%
{and}
%	Then, given
	a ranking probability function $\lambda$, an \textit{equilibrium} is a set of effort levels and post-effort ranking rewards for each applicant, $\{e(\theta_\pre),  \lambda(\theta_\post(\theta_\pre))\}$ such that, for all
	$\omega$, %$\theta_\pre$,
	\begin{align*}
		e(\theta_\pre(\omega)) &\in \argmax_{e} \iwelf\left(e, \lambda\left(\gamma\left(\omega, v\left(e,\theta_\pre(\omega)\right)\right)\right)\right) %\triangleq \argmax_{e} \left[ \lambda(\gamma(v(e,\theta_\pre)) - \cost(e) \right].
		\\
		\theta_\post(\theta_\pre(\omega)) &= \gamma(\omega, v(e(\theta_\pre(\omega)),\theta_\pre(\omega))),
	\end{align*}
	{where $\gamma(\omega,v)$ is the ranking induced by the CDF of the scores resulting from effort levels $\{e(\theta_\pre)\}$ and the tie-breaking function $\Gamma$.\footnote{See Remark~\ref{rem:gamma-effort} on $\gamma$ and the set of efforts.%Formally, $\gamma$ depends on the set of efforts. Given a fixed set and $\gamma$, for each applicant $\omega$ the first condition considers the counter-factual ranking of $\omega$ with different post-effort values but using the same ranking function. As defined, $\gamma$ yields a uniform distribution of ranks with such measure $0$ changes. In Appendix \Cref{lem:deviationsmeasure0}, we further prove that two effort sets equal up to sets of measure $0$ induce the same ranking function $\gamma$, and so the condition is consistent.%, and so we set it at the $\gamma$ induced by the effort levels $\{e(\theta_\pre)\}$.%, and so it does not vary.
	}
	}
	%In particular, we denote $\theta_\post(\theta_\pre) := \gamma(v(e(\theta_\pre),\theta_\pre))$.
	%\lledit{and for any $\theta_\pre, \theta'_\pre$, we have $\lambda(\theta_\post(\theta_\pre)) > \lambda(\theta_\post(\theta'_\pre)) \implies v(e(\theta_\pre),\theta_\pre) > v(e(\theta_\pre'),\theta_\pre')$.}
\end{definition}
\end{ReviewVerbatim}

% report-context-presentation-sha256: 90f9116e292dd71fddfd6b445ca0b060a481d026855bdb9725f2833c2cf62a31
\paragraph{Settled source reading}
\begin{sloppypar}
Equilibrium means that almost every applicant best responds against every feasible deviation;\allowbreak{} uniqueness is up to a null set.\allowbreak{}
\end{sloppypar}
\paragraph{Settled source reading}
\begin{sloppypar}
Applicants know the fixed priority used at score-\allowbreak{}boundary ties.\allowbreak{} This clarifies the stated model, not a corrected admission policy or a new economic assumption.\allowbreak{}
\end{sloppypar}

\paragraph{Expanded PaperInterface specification}
\begin{ReviewVerbatim}
∀ (cost production skill score tie : ℝ → ℝ) (baseline rho : ℝ) (n : ℕ) (cutoff reward : ℕ → ℝ)
  (x y e eToHigh eHighToLow eHigh : ℝ),
  LBG22StrategicRanking.RankingPrimitives cost production skill baseline →
    LBG22StrategicRanking.FiniteAdmissionPolicy n cutoff reward rho →
      x ∈ Set.Icc 0 1 →
        y ∈ Set.Icc 0 1 →
          x < y →
            0 ≤ e →
              e < eToHigh →
                e < eHighToLow →
                  eToHigh < eHigh →
                    eHighToLow < eHigh →
                      have R : ℝ → ℝ → ℝ := fun (t d : ℝ) =>
                        reward
                          ↑(LBG22StrategicRanking.finiteLowerRankBand (fun (i : Fin (n + 1)) => cutoff ↑i)
                              (LBG22StrategicRanking.counterfactualTieBrokenRank LBG22StrategicRanking.unitRankMeasure
                                score tie t (production d * skill t)));
                      R y e = R x eHighToLow →
                        R x eHigh = R y eToHigh →
                          R y e - cost e ≥ R y eToHigh - cost eToHigh →
                            R x eHigh - cost eHigh ≥ R x eHighToLow - cost eHighToLow →
                              cost eHigh - cost eHighToLow ≤ cost eToHigh - cost e
\end{ReviewVerbatim}

\paragraph{Retained governing declarations}
\begin{itemize}\raggedright
\item \hyperlink{paper-prerequisite-d056dd7c8060653c}{LBG22\allowbreak{}Strategic\allowbreak{}Ranking.\allowbreak{}Finite\allowbreak{}Admission\allowbreak{}Policy}
\item \hyperlink{paper-prerequisite-9bbcc12f4a1736bf}{LBG22\allowbreak{}Strategic\allowbreak{}Ranking.\allowbreak{}Ranking\allowbreak{}Primitives}
\item \hyperlink{paper-prerequisite-6ed45e46d01eadeb}{LBG22\allowbreak{}Strategic\allowbreak{}Ranking.\allowbreak{}counterfactual\allowbreak{}Tie\allowbreak{}Broken\allowbreak{}Rank}
\item \hyperlink{paper-prerequisite-a42f440ddc254fa7}{LBG22\allowbreak{}Strategic\allowbreak{}Ranking.\allowbreak{}finite\allowbreak{}Lower\allowbreak{}Rank\allowbreak{}Band}
\item \hyperlink{governing-declaration-678cf48d549cb620}{LBG22\allowbreak{}Strategic\allowbreak{}Ranking.\allowbreak{}unit\allowbreak{}Rank\allowbreak{}Measure}
\end{itemize}
\paragraph{Lean-elaborated claim roles}\begin{ReviewVerbatim}
1. parameter: ℝ → ℝ
2. parameter: ℝ → ℝ
3. parameter: ℝ → ℝ
4. parameter: ℝ → ℝ
5. parameter: ℝ → ℝ
6. parameter: ℝ
7. parameter: ℝ
8. parameter: ℕ
9. parameter: ℕ → ℝ
10. parameter: ℕ → ℝ
11. parameter: ℝ
12. parameter: ℝ
13. parameter: ℝ
14. parameter: ℝ
15. parameter: ℝ
16. parameter: ℝ
17. assumption: LBG22StrategicRanking.RankingPrimitives cost production skill baseline
18. assumption: LBG22StrategicRanking.FiniteAdmissionPolicy n cutoff reward rho
19. assumption: x ∈ Set.Icc 0 1
20. assumption: y ∈ Set.Icc 0 1
21. assumption: x < y
22. assumption: 0 ≤ e
23. assumption: e < eToHigh
24. assumption: e < eHighToLow
25. assumption: eToHigh < eHigh
26. assumption: eHighToLow < eHigh
27. conclusion: (fun (t d : ℝ) =>
        reward
          ↑(LBG22StrategicRanking.finiteLowerRankBand (fun (i : Fin (n + 1)) => cutoff ↑i)
              (LBG22StrategicRanking.counterfactualTieBrokenRank LBG22StrategicRanking.unitRankMeasure score tie t
                (production d * skill t))))
      y e =
    (fun (t d : ℝ) =>
        reward
          ↑(LBG22StrategicRanking.finiteLowerRankBand (fun (i : Fin (n + 1)) => cutoff ↑i)
              (LBG22StrategicRanking.counterfactualTieBrokenRank LBG22StrategicRanking.unitRankMeasure score tie t
                (production d * skill t))))
      x eHighToLow →
  (fun (t d : ℝ) =>
          reward
            ↑(LBG22StrategicRanking.finiteLowerRankBand (fun (i : Fin (n + 1)) => cutoff ↑i)
                (LBG22StrategicRanking.counterfactualTieBrokenRank LBG22StrategicRanking.unitRankMeasure score tie t
                  (production d * skill t))))
        x eHigh =
      (fun (t d : ℝ) =>
          reward
            ↑(LBG22StrategicRanking.finiteLowerRankBand (fun (i : Fin (n + 1)) => cutoff ↑i)
                (LBG22StrategicRanking.counterfactualTieBrokenRank LBG22StrategicRanking.unitRankMeasure score tie t
                  (production d * skill t))))
        y eToHigh →
    (fun (t d : ℝ) =>
              reward
                ↑(LBG22StrategicRanking.finiteLowerRankBand (fun (i : Fin (n + 1)) => cutoff ↑i)
                    (LBG22StrategicRanking.counterfactualTieBrokenRank LBG22StrategicRanking.unitRankMeasure score tie t
                      (production d * skill t))))
            y e -
          cost e ≥
        (fun (t d : ℝ) =>
              reward
                ↑(LBG22StrategicRanking.finiteLowerRankBand (fun (i : Fin (n + 1)) => cutoff ↑i)
                    (LBG22StrategicRanking.counterfactualTieBrokenRank LBG22StrategicRanking.unitRankMeasure score tie t
                      (production d * skill t))))
            y eToHigh -
          cost eToHigh →
      (fun (t d : ℝ) =>
                reward
                  ↑(LBG22StrategicRanking.finiteLowerRankBand (fun (i : Fin (n + 1)) => cutoff ↑i)
                      (LBG22StrategicRanking.counterfactualTieBrokenRank LBG22StrategicRanking.unitRankMeasure score tie
                        t (production d * skill t))))
              x eHigh -
            cost eHigh ≥
          (fun (t d : ℝ) =>
                reward
                  ↑(LBG22StrategicRanking.finiteLowerRankBand (fun (i : Fin (n + 1)) => cutoff ↑i)
                      (LBG22StrategicRanking.counterfactualTieBrokenRank LBG22StrategicRanking.unitRankMeasure score tie
                        t (production d * skill t))))
              x eHighToLow -
            cost eHighToLow →
        cost eHigh - cost eHighToLow ≤ cost eToHigh - cost e
\end{ReviewVerbatim}

\paragraph{Lean proof endpoint (built separately)}\begin{ReviewVerbatim}
LBG22StrategicRanking.ProofInterface.costGap
\end{ReviewVerbatim}

\paragraph{Recorded source-to-Spec screening}
\textbf{Verdict:} \texttt{matches}
\\
{\raggedright
\textbf{Reason:} The displayed reward equalities and the two best-\allowbreak{}response inequalities are the source hypotheses, with the four efforts in the stated strict order;\allowbreak{} adding the inequalities cancels the equal rewards and yields exactly cost eHigh minus cost eHighToLow at most cost eToHigh minus cost e.\allowbreak{} The supporting finite-\allowbreak{}band and counterfactual-\allowbreak{}rank displays preserve the source reward meaning under the approved fixed-\allowbreak{}priority convention.\allowbreak{}
\par}
\reviewmatch{Matches source input}
\noindent\textbf{Reviewer annotation}\par
\reviewerbox
\clearpage
\hypertarget{source-claim-e24c6aaefba531d0}{}
\section*{Review row 17}
\textbf{Source locator:} {\footnotesize\raggedright cited publication, lines 2049-\allowbreak{}2051\par}
\paragraph{Verbatim source input}
\begin{ReviewVerbatim}
\begin{lemma}\label{lem:order_g}
Suppose there are applicants $\omega, \omega'$ with $\theta_\pre(\omega') <\theta_\pre(\omega) $ and effort levels, $e<\{\tilde e', \tilde e \}<e'$ such that $v = v(e, \theta_\pre(\omega)) =v(\tilde{e}', \theta_\pre(\omega'))  $ and $v' =v(e', \theta_\pre(\omega'))=v(\tilde e, \theta_\pre(\omega))$. Then we have that $ e' - \tilde e' > \tilde e - e$.
\end{lemma}

[Next verbatim source excerpt]

\parbold{Applicants}
There is a unit mass of \textit{applicants}, indexed by an observed index $\omega \in [0,1]$ distributed uniformly.\footnote{The index $\omega$ should be interpreted as each applicant's ``name,'' uncorrelated with skill, used solely for tie-breaking.}  Each applicant has a latent (unobserved) skill level represented by some measurable function of $\omega$.
We assume that {the distribution of the skills} has no atoms and that the CDF of this distribution is strictly increasing.
Using the CDF to map the skill of an applicant to a rank in~$[0,1]$, each applicant gets an (unobserved) rank $\theta_\pre =\theta_\pre(\omega)$ which by our assumption on the CDF is again uniformly distributed in $[0, 1]$ (the higher the rank the better).  With this setup, the skill of an applicant with rank $\theta_\pre$ can be written as $f(\theta_\pre)$ where $f$ is a strictly increasing, continuous function.  {It will be notationally convenient to label applicants by their rank $\theta_\pre$, though the reader should note that in contrast to $\omega$, $\theta_\pre(\omega)$ is assumed to be unobservable.}

%3. Each applicant exerts some effort e, at cost \cost(e).


Each applicant chooses an \textit{effort} level $e \ge 0$, the result of which is an observed, post-effort \textit{score}, {$v=v(e,\theta_\pre) = {g(e)}\cdot{f(\theta_\pre)}$}. \lledit{In other words, we assume the post-effort score is the product of two components, associated with the effort and the pre-effort skill respectively.}
{We assume} that the effort transfer function $g: [0, \infty) \mapsto [0, \infty)$ is a continuous, concave, strictly increasing function, representing that marginal effort improves one's score but has diminishing returns.
{%Since applicants can be labelled by their skill level $f(\theta_\pre)$, or equivalently, by their rank $\theta_{\pre}$,
{T}he strategies of the applicants can {then} be described by a function $\theta_\pre \mapsto e(\theta_\pre)$.
%\cb{$\omega\mapsto e(\omega)$.}%, leading to a distribution over post-effort scores $v=v(\theta_\pre))$ via the uniform distribution of $\theta_\pre$ in $[0,1]$.
}
  Each {applicant} is then ranked according to their score $v$%(where we break ties uniformly at random, if they appear)
, resulting in a \textit{post-effort rank} $\theta_\post$.
Note that the ranking $\theta_\post$ is slightly less trivial than a ranking of the skills, since the scores might have ties, which have to be resolved.

{\textbf{Tie Breaking.}
Given a choice of strategies $\theta_\pre \mapsto e(\theta_\pre)$, let $F$ be the CDF of the scores $v(e(\theta_\pre),\theta_\pre)$.
Since atoms for the distribution of $v$ would lead to ties for ranks defined as $F(v)$, we will use the labels of the applicants to break ties with the help of a (publicly announced) tie-breaking function $\Gamma(\omega)$, defined, e.g., via a collision free hash of the applicant names. Here  we require  that $\Gamma$  is a measurable function from $[0,1]$ to $[0,1]$ that maps different applicant labels to different values.
We then use $\Gamma$ to resolve the atoms of $F$, leading to a ranking function $v\mapsto \gamma(\omega, v)$ which is equal to $F(v)$ except when $v$ is an atom of the score distribution, {in which case it takes values in the ``gap interval'' $[F_-(v),  F(v)]$, where $F_-(v)$ is the left limit of $F$ at $v$.}  We  construct $\gamma$ in such a way that it
gives the uniform distribution for $\theta_\post(\omega)$ if we set  $\theta_\post(\omega)=\gamma(\omega,v(e(\theta_\pre(\omega)),\theta_\pre(\omega)))$.\footnote{See Remark~\ref{rem:gamma} for details on this construction.}

[Next verbatim source excerpt]

\parbold{Ranking designer (school)} A single \textit{school} is admitting applicants, based on their ranking. In particular, the school can choose a ranking reward function $\lambda : [0, 1] \mapsto [0, 1]$, such that an applicant with post-effort rank $\theta_\post$ is admitted with probability $\lambda(\theta_\post)$. We assume that $\lambda$ is non-decreasing and that the school has a constraint on the overall probability, such that in expectation it admits a number of applicants equal to {a} capacity constraint  $\rho \in (0,1)$, i.e., \lledit{$\E_{\theta_\post}[\lambda(\theta_\post)] = \rho $}.\footnote{\lledit{We use $\E_{\theta_\post}[\cdot]$ to denote an integral over $\theta_\post$ (and $\E[\cdot]$ to denote an integral over $\omega$)} with respect to the Lebesgue measure; informally, this can be thought of as averaging over the applicant population.} We may also refer to $\lambda$, informally, as the admission policy.



For simplicity, we further assume that $\lambda$ is a step-function with $K$ distinct levels $\ell_0 < \dots < \ell_{K-1}$, and $K-1$ cut-points parameterized by $c_1< \dots < c_{K-1}$ (with $c_0 = 0$, $c_K = 1$). %$\ell_0, \dots, \ell_{K-1}$.
In other words, we have $\lambda(\theta) = \ell_k$, for all $\theta \in \psi_k \triangleq [c_{k}, c_{k+1})$. Thus, applicants in the same post-rank interval $\theta \in \psi_k$ receive the same reward.

[Next verbatim source excerpt]

\parbold{Individual applicant welfare and equilibrium} Given the designer's function $\lambda$ and the effort levels of other applicants, each applicant chooses effort $e$ to maximize their individual welfare, % $W$,
\begin{equation*}
	\iwelf(e, \lambda(\theta_\post)) = \lambda(\theta_\post) - \cost(e),
\end{equation*}
where the effort cost function $\cost$ is non-negative, continuous, and strictly convex {on $[0,\infty)$, with $e_0:=\argmin_e \cost(e)$ and $p(e_0) = 0$.}
%\lnote{also want to assume that $\cost$ attains its minimum at $0$?}
%Let $e_0:=\argmin_e \cost(e)$ and $p(e_0) = 0$.

Applicants are assumed not to personally benefit from increasing their score $v$, except through the corresponding increase in their rank and reward. While the definition of $\cost$ does not preclude the applicant receiving any intrinsic benefit from exerting effort, we assume that the \emph{net} benefit from effort is non-positive.

[Next verbatim source excerpt]

\begin{definition} [Equilibrium]\label{def:equi}
{Given a}
%Let the
tie-breaking function $\Gamma(\omega)$
%be a pre-announced ordering over index $\omega$, to break ties in post-effort value. Then, let $\gamma(\omega, v)$ be the 1-1 ranking function defined as the CDF of post-effort values $v$ except where there are ties in $v$, in which case the ordering $\Gamma$ is used to produce unique ranks such that the range of $\gamma$ is uniform on $[0, 1]$.\footnote{More formally, note that each discontinuity in the CDF of post-effort values $v$ has a height equal to the measure of the applicants who are tied at that value, and so $\gamma$ for those applicants can be filled in using the CDF of $\Gamma(\omega)$ restricted to the applicants $\omega$ who are tied. Thus, the range of $\gamma$ is uniform on $[0,1]$.
%
{and}
%	Then, given
	a ranking probability function $\lambda$, an \textit{equilibrium} is a set of effort levels and post-effort ranking rewards for each applicant, $\{e(\theta_\pre),  \lambda(\theta_\post(\theta_\pre))\}$ such that, for all
	$\omega$, %$\theta_\pre$,
	\begin{align*}
		e(\theta_\pre(\omega)) &\in \argmax_{e} \iwelf\left(e, \lambda\left(\gamma\left(\omega, v\left(e,\theta_\pre(\omega)\right)\right)\right)\right) %\triangleq \argmax_{e} \left[ \lambda(\gamma(v(e,\theta_\pre)) - \cost(e) \right].
		\\
		\theta_\post(\theta_\pre(\omega)) &= \gamma(\omega, v(e(\theta_\pre(\omega)),\theta_\pre(\omega))),
	\end{align*}
	{where $\gamma(\omega,v)$ is the ranking induced by the CDF of the scores resulting from effort levels $\{e(\theta_\pre)\}$ and the tie-breaking function $\Gamma$.\footnote{See Remark~\ref{rem:gamma-effort} on $\gamma$ and the set of efforts.%Formally, $\gamma$ depends on the set of efforts. Given a fixed set and $\gamma$, for each applicant $\omega$ the first condition considers the counter-factual ranking of $\omega$ with different post-effort values but using the same ranking function. As defined, $\gamma$ yields a uniform distribution of ranks with such measure $0$ changes. In Appendix \Cref{lem:deviationsmeasure0}, we further prove that two effort sets equal up to sets of measure $0$ induce the same ranking function $\gamma$, and so the condition is consistent.%, and so we set it at the $\gamma$ induced by the effort levels $\{e(\theta_\pre)\}$.%, and so it does not vary.
	}
	}
	%In particular, we denote $\theta_\post(\theta_\pre) := \gamma(v(e(\theta_\pre),\theta_\pre))$.
	%\lledit{and for any $\theta_\pre, \theta'_\pre$, we have $\lambda(\theta_\post(\theta_\pre)) > \lambda(\theta_\post(\theta'_\pre)) \implies v(e(\theta_\pre),\theta_\pre) > v(e(\theta_\pre'),\theta_\pre')$.}
\end{definition}
\end{ReviewVerbatim}



\paragraph{Expanded PaperInterface specification}
\begin{ReviewVerbatim}
∀ (cost production skill : ℝ → ℝ) (baseline x y e eToHigh eHighToLow eHigh : ℝ),
  LBG22StrategicRanking.RankingPrimitives cost production skill baseline →
    x ∈ Set.Icc 0 1 →
      y ∈ Set.Icc 0 1 →
        x < y →
          0 ≤ e →
            e < eToHigh →
              e < eHighToLow →
                eToHigh < eHigh →
                  eHighToLow < eHigh →
                    production e * skill y = production eHighToLow * skill x →
                      production eHigh * skill x = production eToHigh * skill y → eToHigh - e < eHigh - eHighToLow
\end{ReviewVerbatim}

\paragraph{Retained governing declarations}
\begin{itemize}\raggedright
\item \hyperlink{paper-prerequisite-9bbcc12f4a1736bf}{LBG22\allowbreak{}Strategic\allowbreak{}Ranking.\allowbreak{}Ranking\allowbreak{}Primitives}
\end{itemize}
\paragraph{Lean-elaborated claim roles}\begin{ReviewVerbatim}
1. parameter: ℝ → ℝ
2. parameter: ℝ → ℝ
3. parameter: ℝ → ℝ
4. parameter: ℝ
5. parameter: ℝ
6. parameter: ℝ
7. parameter: ℝ
8. parameter: ℝ
9. parameter: ℝ
10. parameter: ℝ
11. assumption: LBG22StrategicRanking.RankingPrimitives cost production skill baseline
12. assumption: x ∈ Set.Icc 0 1
13. assumption: y ∈ Set.Icc 0 1
14. assumption: x < y
15. assumption: 0 ≤ e
16. assumption: e < eToHigh
17. assumption: e < eHighToLow
18. assumption: eToHigh < eHigh
19. assumption: eHighToLow < eHigh
20. assumption: production e * skill y = production eHighToLow * skill x
21. assumption: production eHigh * skill x = production eToHigh * skill y
22. conclusion: eToHigh - e < eHigh - eHighToLow
\end{ReviewVerbatim}

\paragraph{Lean proof endpoint (built separately)}\begin{ReviewVerbatim}
LBG22StrategicRanking.ProofInterface.productionGap
\end{ReviewVerbatim}

\paragraph{Recorded source-to-Spec screening}
\textbf{Verdict:} \texttt{matches}
\\
{\raggedright
\textbf{Reason:} Under the source rank and effort ordering and the two equal-\allowbreak{}score equations, the target concludes eToHigh-\allowbreak{}e < eHigh-\allowbreak{}eHighToLow, which is the source strict production-\allowbreak{}gap inequality after matching its tilde-\allowbreak{}effort notation.\allowbreak{} The displayed scalar primitives supply the required strict production/\allowbreak{}skill and concavity domain.\allowbreak{}
\par}
\reviewmatch{Matches source input}
\noindent\textbf{Reviewer annotation}\par
\reviewerbox
\clearpage
\hypertarget{source-claim-4fa39a045f045b81}{}
\section*{Review row 18}
\textbf{Source locator:} {\footnotesize\raggedright cited publication, lines 1674-\allowbreak{}1676\par}
\paragraph{Verbatim source input}
\begin{ReviewVerbatim}
\begin{remark}[Ties in pre-effort skill]
	With probability one, there are no two applicants with the same skill level, and the support of the skill distribution has no gap: given two applicants with different skills, the probability of finding an applicant with skill in between these two is always non-zero.
\end{remark}

[Next verbatim source excerpt]

\parbold{Applicants}
There is a unit mass of \textit{applicants}, indexed by an observed index $\omega \in [0,1]$ distributed uniformly.\footnote{The index $\omega$ should be interpreted as each applicant's ``name,'' uncorrelated with skill, used solely for tie-breaking.}  Each applicant has a latent (unobserved) skill level represented by some measurable function of $\omega$.
We assume that {the distribution of the skills} has no atoms and that the CDF of this distribution is strictly increasing.
Using the CDF to map the skill of an applicant to a rank in~$[0,1]$, each applicant gets an (unobserved) rank $\theta_\pre =\theta_\pre(\omega)$ which by our assumption on the CDF is again uniformly distributed in $[0, 1]$ (the higher the rank the better).  With this setup, the skill of an applicant with rank $\theta_\pre$ can be written as $f(\theta_\pre)$ where $f$ is a strictly increasing, continuous function.  {It will be notationally convenient to label applicants by their rank $\theta_\pre$, though the reader should note that in contrast to $\omega$, $\theta_\pre(\omega)$ is assumed to be unobservable.}
\end{ReviewVerbatim}



\paragraph{Expanded PaperInterface specification}
\begin{ReviewVerbatim}
∀ (skill : ℝ → ℝ),
  ContinuousOn skill (Set.Icc 0 1) →
    StrictMonoOn skill (Set.Icc 0 1) →
      (∀ᵐ (t : ℝ) ∂LBG22StrategicRanking.unitRankMeasure, ∀ u ∈ Set.Icc 0 1, skill t = skill u → t = u) ∧
        ∀ s ∈ Set.Icc 0 1,
          ∀ t ∈ Set.Icc 0 1,
            s < t → 0 < LBG22StrategicRanking.unitRankMeasure {u : ℝ | skill s < skill u ∧ skill u < skill t}
\end{ReviewVerbatim}

\paragraph{Retained governing declarations}
\begin{itemize}\raggedright
\item \hyperlink{governing-declaration-678cf48d549cb620}{LBG22\allowbreak{}Strategic\allowbreak{}Ranking.\allowbreak{}unit\allowbreak{}Rank\allowbreak{}Measure}
\end{itemize}
\paragraph{Lean-elaborated claim roles}\begin{ReviewVerbatim}
1. parameter: ℝ → ℝ
2. assumption: ContinuousOn skill (Set.Icc 0 1)
3. assumption: StrictMonoOn skill (Set.Icc 0 1)
4. conclusion: (∀ᵐ (t : ℝ) ∂LBG22StrategicRanking.unitRankMeasure, ∀ u ∈ Set.Icc 0 1, skill t = skill u → t = u) ∧
  ∀ s ∈ Set.Icc 0 1,
    ∀ t ∈ Set.Icc 0 1, s < t → 0 < LBG22StrategicRanking.unitRankMeasure {u : ℝ | skill s < skill u ∧ skill u < skill t}
\end{ReviewVerbatim}

\paragraph{Lean proof endpoint (built separately)}\begin{ReviewVerbatim}
LBG22StrategicRanking.ProofInterface.preEffortSkillRegularity
\end{ReviewVerbatim}

\paragraph{Recorded source-to-Spec screening}
\textbf{Verdict:} \texttt{matches}
\\
{\raggedright
\textbf{Reason:} Continuous strict monotonicity of the skill quantile on [0,1] yields almost-\allowbreak{}sure injectivity and positive uniform measure of ranks whose skills lie strictly between any two ordered endpoint skills.\allowbreak{} These are exactly the two assertions in the supplemental no-\allowbreak{}ties/\allowbreak{}no-\allowbreak{}support-\allowbreak{}gap remark.\allowbreak{}
\par}
\reviewmatch{Matches source input}
\noindent\textbf{Reviewer annotation}\par
\reviewerbox

\end{document}
